\documentclass[12pt,oneside]{book}

% -------------------------------------------------
% PAQUETES
% -------------------------------------------------
\usepackage[spanish,es-nodecimaldot]{babel}
\usepackage[utf8]{inputenc}
\usepackage[T1]{fontenc}
\usepackage{lmodern}
\usepackage[a4paper,margin=2.5cm]{geometry}
\usepackage{microtype}
\usepackage{setspace}
\usepackage{graphicx}
\usepackage{xcolor}
\usepackage{fancyhdr}
\usepackage{titlesec}
\usepackage{hyperref}
\usepackage{bookmark}
\usepackage{tocloft}
\usepackage{tcolorbox}
\tcbuselibrary{theorems, skins, breakable}
\usepackage{amsmath,amssymb,amsthm}
\usepackage{mathtools}
\usepackage{enumitem}
\usepackage{csquotes}
\usepackage[backend=biber, style=alphabetic, sorting=nyt]{biblatex}
\addbibresource{referencias.bib}

% -------------------------------------------------
% CONFIGURACIÓN VISUAL
% -------------------------------------------------
\definecolor{azulMath}{HTML}{026937}
\definecolor{grisClaro}{HTML}{F4F6F8}
\definecolor{grisTexto}{HTML}{2E2E2E}
\definecolor{grisSec}{HTML}{6C757D}
\definecolor{colorDef}{HTML}{EBF5FB}
\definecolor{colorTeo}{HTML}{EAFAF1}
\definecolor{colorEj}{HTML}{FEF9E7}
\definecolor{colorObs}{HTML}{F2F3F4}

\hypersetup{
    colorlinks=true,
    linkcolor=azulMath,
    urlcolor=azulMath,
    citecolor=azulMath,
    pdftitle={Notas de Clase - Números Transfinitos},
    pdfauthor={Daniel Soto Villada}
}

\setstretch{1.12}
\setlength{\parskip}{0.4em}
\setlength{\parindent}{0pt}
\setcounter{secnumdepth}{3}
\setcounter{tocdepth}{2}
\setlength{\headheight}{26pt}

% Estilo de capítulos y secciones
\titleformat{\chapter}[display]
  {\bfseries\Huge\color{azulMath}}
  {\filleft\Large\chaptername\ \thechapter}
  {1ex}
  {\titlerule\vspace{1ex}\filright}
  [\vspace{1ex}\titlerule]

\titleformat{\section}
  {\bfseries\Large\color{azulMath}}
  {\thesection}{0.7em}{}

\titleformat{\subsection}
  {\bfseries\large\color{grisTexto}}
  {\thesubsection}{0.6em}{}

\titleformat{\subsubsection}
  {\bfseries\normalsize\color{grisTexto}}
  {\thesubsubsection}{0.5em}{}

% Encabezados y pies
\pagestyle{fancy}
\fancyhf{}
\fancyhead[L]{\textcolor{grisSec}{\nouppercase{\leftmark}}}
\fancyhead[R]{\textcolor{grisSec}{Daniel Soto Villada}}
\fancyfoot[C]{\textcolor{grisSec}{\thepage}}
\renewcommand{\headrulewidth}{0.4pt}
\renewcommand{\headrule}{\hbox to\headwidth{\color{azulMath}\leaders\hrule height \headrulewidth\hfill}}

% Índice más limpio
\renewcommand{\cftchapfont}{\bfseries\color{azulMath}}
\renewcommand{\cftsecfont}{\color{grisTexto}}
\renewcommand{\cftchappagefont}{\bfseries\color{azulMath}}
\renewcommand{\cftsecpagefont}{\color{grisTexto}}

% -------------------------------------------------
% ENTORNOS TCOLORBOX
% -------------------------------------------------

% Caja genérica para ideas importantes
\newtcolorbox{importante}{
  colback=grisClaro,
  colframe=azulMath,
  boxrule=0.8pt,
  arc=2mm,
  left=3mm, right=3mm, top=1.5mm, bottom=1.5mm
}

% Definición
\newtcbtheorem[number within=chapter]{definicion}{Definición}{
  enhanced,
  colback=colorDef,
  colframe=azulMath,
  coltitle=white,
  colbacktitle=azulMath,
  attach boxed title to top left={yshift=-2mm, xshift=4mm},
  boxed title style={enhanced, arc=1mm, boxrule=0pt, colframe=azulMath},
  fonttitle=\bfseries,
  breakable,
  arc=1.5mm,
  boxrule=0.8pt,
  left=4mm, right=4mm, top=5mm, bottom=2mm,
  separator sign none,
  description delimiters={(}{)}
}{def}

% Teorema (define el contador compartido)
\newtcbtheorem[number within=chapter]{teorema}{Teorema}{
  enhanced,
  colback=colorTeo,
  colframe=azulMath,
  coltitle=white,
  colbacktitle=azulMath,
  attach boxed title to top left={yshift=-2mm, xshift=4mm},
  boxed title style={enhanced, arc=1mm, boxrule=0pt, colframe=azulMath},
  fonttitle=\bfseries,
  breakable,
  arc=1.5mm,
  boxrule=0.8pt,
  left=4mm, right=4mm, top=5mm, bottom=2mm,
  separator sign none,
  description delimiters={(}{)}
}{teo}

% Lema (comparte contador con teorema)
\newtcbtheorem[use counter from=teorema]{lema}{Lema}{
  enhanced,
  colback=colorTeo,
  colframe=azulMath,
  coltitle=white,
  colbacktitle=azulMath,
  attach boxed title to top left={yshift=-2mm, xshift=4mm},
  boxed title style={enhanced, arc=1mm, boxrule=0pt, colframe=azulMath},
  fonttitle=\bfseries,
  breakable,
  arc=1.5mm,
  boxrule=0.8pt,
  left=4mm, right=4mm, top=5mm, bottom=2mm,
  separator sign none,
  description delimiters={(}{)}
}{lem}

% Corolario (comparte contador con teorema)
\newtcbtheorem[use counter from=teorema]{corolario}{Corolario}{
  enhanced,
  colback=colorTeo,
  colframe=azulMath,
  coltitle=white,
  colbacktitle=azulMath,
  attach boxed title to top left={yshift=-2mm, xshift=4mm},
  boxed title style={enhanced, arc=1mm, boxrule=0pt, colframe=azulMath},
  fonttitle=\bfseries,
  breakable,
  arc=1.5mm,
  boxrule=0.8pt,
  left=4mm, right=4mm, top=5mm, bottom=2mm,
  separator sign none,
  description delimiters={(}{)}
}{cor}

% Proposición (comparte contador con teorema)
\newtcbtheorem[use counter from=teorema]{proposicion}{Proposición}{
  enhanced,
  colback=colorTeo,
  colframe=azulMath,
  coltitle=white,
  colbacktitle=azulMath,
  attach boxed title to top left={yshift=-2mm, xshift=4mm},
  boxed title style={enhanced, arc=1mm, boxrule=0pt, colframe=azulMath},
  fonttitle=\bfseries,
  breakable,
  arc=1.5mm,
  boxrule=0.8pt,
  left=4mm, right=4mm, top=5mm, bottom=2mm,
  separator sign none,
  description delimiters={(}{)}
}{prop}

% Ejemplo
\newtcbtheorem[number within=chapter]{ejemplo}{Ejemplo}{
  enhanced,
  colback=colorEj,
  colframe=azulMath,
  coltitle=white,
  colbacktitle=azulMath,
  attach boxed title to top left={yshift=-2mm, xshift=4mm},
  boxed title style={enhanced, arc=1mm, boxrule=0pt, colframe=azulMath},
  fonttitle=\bfseries,
  breakable,
  arc=1.5mm,
  boxrule=0.8pt,
  left=4mm, right=4mm, top=5mm, bottom=2mm,
  separator sign none,
  description delimiters={(}{)}
}{ej}

% Observación
\newtcbtheorem[number within=chapter]{observacion}{Observación}{
  enhanced,
  colback=colorObs,
  colframe=grisSec,
  coltitle=white,
  colbacktitle=grisSec,
  attach boxed title to top left={yshift=-2mm, xshift=4mm},
  boxed title style={enhanced, arc=1mm, boxrule=0pt, colframe=grisSec},
  fonttitle=\bfseries,
  breakable,
  arc=1.5mm,
  boxrule=0.6pt,
  left=4mm, right=4mm, top=5mm, bottom=2mm,
  separator sign none,
  description delimiters={(}{)}
}{obs}

% -------------------------------------------------
% ENTORNO DE DEMOSTRACIÓN
% -------------------------------------------------
\renewcommand{\proofname}{\textbf{Demostración}}

% -------------------------------------------------
% MACROS MATEMÁTICOS
% -------------------------------------------------
% Regla de inferencia: \infr[nombre]{premisas}{conclusión}
\newcommand{\infr}[3][]{\dfrac{\displaystyle #2}{\displaystyle #3}\;\text{(}#1\text{)}}

% Lenguajes y estructuras
\newcommand{\Lang}{\mathcal{L}}
\newcommand{\Struct}[1]{\mathfrak{#1}}
\newcommand{\FV}{\mathrm{FV}}
\newcommand{\dom}[1]{|{#1}|}

% Consecuencia
\newcommand{\entails}{\models}
\newcommand{\proves}{\vdash}

% -------------------------------------------------
% DATOS DEL DOCUMENTO
% -------------------------------------------------
\newcommand{\NombreEstudiante}{Daniel Soto Villada}
\newcommand{\Programa}{Estudiante de astronomía y matemáticas}
\newcommand{\Universidad}{Universidad de Antioquia}
\newcommand{\Facultad}{Facultad de Ciencias Exactas y Naturales}
\newcommand{\Instituto}{Instituto de Matemáticas}
\newcommand{\TituloDocumento}{Notas de Clase\\[0.4em]Números Transfinitos}

% -------------------------------------------------
% DOCUMENTO
% -------------------------------------------------
\begin{document}
\hypersetup{pageanchor=false}

% ---------- PORTADA ----------
\begin{titlepage}
    \centering
    \vspace*{1.8cm}

    {\Large\textcolor{grisSec}{\Universidad}\par}
    \vspace{0.25cm}
    {\large\textcolor{grisSec}{\Facultad}\par}
    {\large\textcolor{grisSec}{\Instituto}\par}

    \vspace{2.2cm}
    {\Huge\bfseries\textcolor{azulMath}{\TituloDocumento}\par}
    \vspace{0.5cm}
    {\Large\textcolor{grisSec}{Fundamentos de Teoría de Conjuntos y Lógica Matemática}\par}

    \vspace{2.4cm}
    \begin{tcolorbox}[colback=grisClaro,colframe=azulMath,boxrule=0.9pt,arc=2mm,width=0.78\textwidth]
        \centering
        {\Large\bfseries \NombreEstudiante}\par
        \vspace{0.2cm}
        {\large \Programa}
    \end{tcolorbox}

    \vfill
    {\large\textcolor{grisSec}{\today}\par}
    \vspace{0.5cm}
\end{titlepage}

% ---------- PÁGINAS PRELIMINARES ----------
\frontmatter
\tableofcontents
\clearpage

% ---------- CONTENIDO ----------
\clearpage
\hypersetup{pageanchor=true}
\mainmatter
\chapter{Clase 1: Lógica de Primer Orden --- Sintaxis, Semántica y Deducción}

La lógica de primer orden —también llamada \textbf{lógica de predicados}— constituye el lenguaje universal en el que se expresan y demuestran la casi totalidad de los resultados de las matemáticas modernas. Antes de construir la teoría de conjuntos de Zermelo–Fraenkel, estudiar los números cardinales y ordinales o demostrar resultados sobre conjuntos transfinitos, es imprescindible comprender con precisión el substrato formal sobre el que todo esto descansa: un lenguaje con reglas sintácticas precisas, una semántica que asigna significado a las fórmulas, y un aparato deductivo que permite derivar consecuencias a partir de hipótesis.

Esta clase presenta dicho substrato de manera rigurosa y con los detalles suficientes para que el estudiante pueda trabajar con la lógica de primer orden con plena confianza. Los resultados culminantes —los teoremas de Completitud y Compacidad de Gödel— son fundamentales no solo para la lógica misma, sino también para entender los límites y alcances de la teoría de conjuntos y la aritmética formal \cite{godel1930, enderton2001}.

% ===========================================================
\section{El lenguaje formal de primer orden}
% ===========================================================

\subsection{Motivación: ¿por qué un lenguaje formal?}

En matemáticas es habitual escribir enunciados como:
\begin{itemize}
  \item \textit{``Para todo número natural $n$, existe un primo mayor que $n$.''}
  \item \textit{``Todo grupo abeliano finito es isomorfo a un producto de grupos cíclicos.''}
  \item \textit{``Si $A \subseteq B$ y $B \subseteq C$, entonces $A \subseteq C$.''}
\end{itemize}

Estos enunciados se expresan en lenguaje natural, lo cual introduce ambigüedades y dependencias culturales. La \textbf{lógica de primer orden} ofrece un lenguaje completamente formal: cada símbolo tiene un significado exacto, las reglas de formación de fórmulas son algorítmicas y las demostraciones son secuencias de pasos verificables mecánicamente.

Esta formalización no es un capricho filosófico: es necesaria para poder enunciar y demostrar con rigor resultados como los teoremas de incompletitud de Gödel, el teorema de Löwenheim–Skolem, y para estudiar los modelos de la teoría de conjuntos.

\subsection{Firma o signatura}

El primer paso para construir un lenguaje de primer orden es especificar los \emph{símbolos no lógicos} que lo componen. Estos símbolos codifican la estructura matemática que queremos hablar.

\begin{definicion}{Firma (o signatura)}{firma}
Una \textbf{firma} (o \textbf{signatura}) es una tupla $\Lang = (C, \mathcal{F}, \mathcal{R})$ donde:
\begin{itemize}[leftmargin=2em]
  \item $C$ es un conjunto (posiblemente vacío) de \textbf{símbolos de constante}: $c_0, c_1, \ldots$
  \item $\mathcal{F}$ es un conjunto de \textbf{símbolos de función}, cada uno con una aridad $n \geq 1$ fija: $f^{(n)}, g^{(m)}, \ldots$
  \item $\mathcal{R}$ es un conjunto de \textbf{símbolos de relación}, cada uno con una aridad $n \geq 1$ fija: $R^{(n)}, P^{(m)}, \ldots$
\end{itemize}
\end{definicion}

\begin{ejemplo}{Firmas matemáticas usuales}{}
\begin{enumerate}[leftmargin=2em, label=(\alph*)]
  \item \textbf{Lenguaje de la aritmética de Peano:}
    \[
      \Lang_{\mathrm{ar}} = (\{0,1\},\; \{s^{(1)}, +^{(2)}, \times^{(2)}\},\; \{\leq^{(2)}\})
    \]
    donde $0,1$ son constantes, $s$ es la función sucesor, $+$ y $\times$ son operaciones binarias, y $\leq$ es una relación binaria.

  \item \textbf{Lenguaje de la teoría de grupos:}
    \[
      \Lang_{\mathrm{gr}} = (\{e\},\; \{\cdot^{(2)},\, (-)^{-1\,(1)}\},\; \emptyset)
    \]
    donde $e$ es la identidad, $\cdot$ es la operación de grupo y $(-)^{-1}$ es la operación de inversión.

  \item \textbf{Lenguaje de la teoría de conjuntos ($\mathrm{ZFC}$):}
    \[
      \Lang_{\in} = (\emptyset,\; \emptyset,\; \{\in^{(2)}\})
    \]
    Solo contiene la relación binaria de pertenencia. Esta es la firma que utilizaremos a lo largo del curso.
\end{enumerate}
\end{ejemplo}

\subsection{El alfabeto de \texorpdfstring{$\mathcal{L}_1$}{L1}}

A partir de una firma $\Lang$, el alfabeto de la lógica de primer orden $\mathcal{L}_1(\Lang)$ consta de los siguientes símbolos:

\begin{enumerate}[leftmargin=2em, label=\textbf{\arabic*.}]
  \item \textbf{Variables:} un conjunto contable de símbolos $x_0, x_1, x_2, \ldots$ (o $x, y, z, w, \ldots$ para mayor legibilidad).
  \item \textbf{Símbolo de igualdad:} $=$
  \item \textbf{Conectivos lógicos:} $\neg$ (negación), $\wedge$ (conjunción), $\vee$ (disyunción), $\to$ (implicación), $\leftrightarrow$ (bicondicional).
  \item \textbf{Cuantificadores:} $\forall$ (universal), $\exists$ (existencial).
  \item \textbf{Signos de puntuación:} paréntesis $($ y $)$, coma $,$.
  \item \textbf{Símbolos de la firma:} todos los símbolos en $C$, $\mathcal{F}$ y $\mathcal{R}$.
\end{enumerate}

\begin{importante}
La \textbf{lógica de primer orden} (abreviada $\mathcal{L}_1$) recibe ese nombre porque los cuantificadores $\forall$ y $\exists$ solo pueden cuantificar sobre \emph{individuos} (elementos del dominio), no sobre propiedades, funciones ni conjuntos de individuos. La \emph{lógica de segundo orden} ($\mathcal{L}_2$) extiende $\mathcal{L}_1$ permitiendo cuantificar sobre relaciones y funciones, pero pierde propiedades cruciales como la completitud.
\end{importante}

\subsection{Términos}

Los \textbf{términos} son las expresiones que denotan \emph{objetos} o \emph{individuos} del dominio. Su definición es inductiva.

\begin{definicion}{Términos de $\mathcal{L}_1(\Lang)$}{terminos}
El conjunto $\mathrm{Term}(\Lang)$ de \textbf{términos} se define inductivamente:
\begin{enumerate}[leftmargin=2em]
  \item Toda variable $x_i$ es un término.
  \item Todo símbolo de constante $c \in C$ es un término.
  \item Si $f \in \mathcal{F}$ tiene aridad $n$ y $t_1, \ldots, t_n$ son términos, entonces $f(t_1, \ldots, t_n)$ es un término.
  \item \textbf{Nada más es un término.}
\end{enumerate}
\end{definicion}

La cláusula (4) —llamada \emph{cláusula de cierre}— garantiza que la definición sea precisa: un término es exactamente lo que puede construirse mediante las reglas (1)--(3).

\begin{ejemplo}{Términos en el lenguaje de la aritmética}{}
Usando $\Lang_{\mathrm{ar}} = (\{0,1\}, \{s,+,\times\}, \{\leq\})$:
\begin{itemize}
  \item $x$, $y$, $z$ son términos (variables).
  \item $0$, $1$ son términos (constantes).
  \item $s(0)$, $s(1)$, $s(s(0))$ son términos. Intuitivamente representan $1, 2, 2$ (o $2$ si $1$ ya es la constante para el uno).
  \item $+(x, s(0))$ es un término; en notación infija: $x + s(0)$, que representa $x + 1$.
  \item $\times(+(x,y), s(0))$ es un término; en notación infija: $(x+y) \times s(0) = (x+y)\cdot 1$.
  \item $+(x, \times(y, z))$ representa $x + (y \cdot z)$.
\end{itemize}
\end{ejemplo}

\begin{ejemplo}{Términos en el lenguaje de los grupos}{}
Usando $\Lang_{\mathrm{gr}} = (\{e\}, \{\cdot, (-)^{-1}\}, \emptyset)$:
\begin{itemize}
  \item $x$, $y$ son términos.
  \item $e$ es un término (la identidad).
  \item $x^{-1}$ es un término.
  \item $x \cdot y^{-1}$ es el término $\cdot(x, (-)^{-1}(y))$.
  \item $e \cdot (x \cdot (y \cdot z))$ es un término más complejo.
\end{itemize}
\end{ejemplo}

\subsection{Fórmulas bien formadas}

A partir de los términos, construimos las \textbf{fórmulas}, que son las expresiones que tienen valor de verdad.

\begin{definicion}{Fórmulas atómicas}{atomicas}
Las \textbf{fórmulas atómicas} de $\mathcal{L}_1(\Lang)$ son:
\begin{enumerate}[leftmargin=2em]
  \item Si $t_1, t_2$ son términos, entonces $t_1 = t_2$ es una fórmula atómica.
  \item Si $R \in \mathcal{R}$ tiene aridad $n$ y $t_1, \ldots, t_n$ son términos, entonces $R(t_1, \ldots, t_n)$ es una fórmula atómica.
\end{enumerate}
\end{definicion}

\begin{definicion}{Fórmulas bien formadas (fbf)}{fbf}
El conjunto $\mathrm{Form}(\Lang)$ de \textbf{fórmulas bien formadas} se define inductivamente:
\begin{enumerate}[leftmargin=2em]
  \item Toda fórmula atómica es una fbf.
  \item Si $\varphi$ es una fbf, entonces $\neg \varphi$ es una fbf.
  \item Si $\varphi$ y $\psi$ son fbf, entonces $(\varphi \wedge \psi)$, $(\varphi \vee \psi)$, $(\varphi \to \psi)$ y $(\varphi \leftrightarrow \psi)$ son fbf.
  \item Si $\varphi$ es una fbf y $x$ es una variable, entonces $\forall x\, \varphi$ y $\exists x\, \varphi$ son fbf.
  \item \textbf{Nada más es una fbf.}
\end{enumerate}
\end{definicion}

\begin{ejemplo}{Fórmulas bien formadas en $\Lang_{\mathrm{ar}}$}{}
\begin{itemize}
  \item $0 = 0$ es una fbf atómica (y es verdadera).
  \item $x \leq y$ es una fbf atómica.
  \item $\neg(x = 0)$ es una fbf (``$x$ no es cero'').
  \item $\forall x\, \forall y\, (x + y = y + x)$ es una fbf (conmutatividad de la suma).
  \item $\exists x\, (x + x = s(0))$ es una fbf (``existe $x$ tal que $2x = 1$'').
  \item $\forall x\, \exists y\, (x < y)$ es una fbf (``no hay máximo'').
\end{itemize}
\end{ejemplo}

\subsection{Variables libres y ligadas}

En la fórmula $\forall x\, \exists y\, (x < y)$, las variables $x$ e $y$ están \emph{ligadas} por los cuantificadores. En $x + y = z$, las variables $x, y, z$ son \emph{libres}: no están bajo el alcance de ningún cuantificador y, por tanto, actúan como parámetros.

\begin{definicion}{Variables libres}{libres}
El conjunto $\FV(\varphi)$ de \textbf{variables libres} de una fórmula $\varphi$ se define inductivamente:
\begin{align*}
  \FV(t_1 = t_2) &= \FV(t_1) \cup \FV(t_2) \\
  \FV(R(t_1,\ldots,t_n)) &= \FV(t_1) \cup \cdots \cup \FV(t_n) \\
  \FV(\neg\varphi) &= \FV(\varphi) \\
  \FV(\varphi \circ \psi) &= \FV(\varphi) \cup \FV(\psi) \quad (\circ \in \{\wedge, \vee, \to, \leftrightarrow\})\\
  \FV(\forall x\, \varphi) &= \FV(\varphi) \setminus \{x\} \\
  \FV(\exists x\, \varphi) &= \FV(\varphi) \setminus \{x\}
\end{align*}
donde $\FV(x) = \{x\}$ para una variable $x$ y $\FV(c) = \emptyset$ para una constante $c$.
\end{definicion}

\begin{definicion}{Sentencia (fórmula cerrada)}{sentencia}
Una fórmula $\varphi$ es una \textbf{sentencia} (o \textbf{fórmula cerrada}) si $\FV(\varphi) = \emptyset$. El conjunto de sentencias de $\mathcal{L}_1(\Lang)$ se denota $\mathrm{Sent}(\Lang)$.
\end{definicion}

\begin{ejemplo}{Variables libres y ligadas}{}
\begin{enumerate}[leftmargin=2em, label=(\alph*)]
  \item En $\varphi := \forall x\, (x + y = z)$: $x$ está ligada; $y$ y $z$ son libres. Luego $\FV(\varphi) = \{y, z\}$. No es una sentencia.
  \item En $\psi := \forall x\, \forall y\, (x + y = y + x)$: todas las variables están ligadas. $\FV(\psi) = \emptyset$. Es una sentencia.
  \item En $\chi := \exists x\, (x < y) \wedge (y < z)$: $x$ está ligada; $y$ aparece tanto dentro del alcance de $\exists x$ (libre) como en $y < z$ (libre). $\FV(\chi) = \{y, z\}$.
\end{enumerate}
\end{ejemplo}

\subsection{Sustitución}

La \textbf{sustitución} $\varphi[t/x]$ denota la fórmula obtenida al reemplazar todas las ocurrencias libres de $x$ en $\varphi$ por el término $t$. Para que esta operación sea semánticamente correcta, debe hacerse con cuidado para evitar \emph{capturas de variable}.

\begin{definicion}{Sustitución libre}{sustitucion}
Un término $t$ es \textbf{libre para $x$ en $\varphi$} si ninguna variable libre de $t$ queda capturada por un cuantificador al sustituir. Formalmente, para cada ocurrencia libre de $x$ en $\varphi$ dentro del alcance de $\forall y\, (\ldots)$ o $\exists y\, (\ldots)$, se requiere que $y \notin \FV(t)$.

La \textbf{sustitución} $\varphi[t/x]$ se define inductivamente de manera obvia en los casos atómicos y conectivos, y para los cuantificadores:
\begin{itemize}
  \item Si $\varphi = \forall y\, \psi$ y $y \neq x$ y $t$ es libre para $x$ en $\varphi$:
  \[
    (\forall y\, \psi)[t/x] \;=\; \forall y\, (\psi[t/x]).
  \]
  \item Si $\varphi = \forall x\, \psi$: $(\forall x\, \psi)[t/x] = \forall x\, \psi$ (la variable $x$ no ocurre libre, no hay nada que sustituir).
\end{itemize}
\end{definicion}

\begin{ejemplo}{Sustitución correcta e incorrecta}{}
Sea $\varphi := \exists y\, (x < y)$ y $t := y + 1$.

\textbf{Sustitución problemática:} $\varphi[y+1/x] = \exists y\, (y+1 < y)$.
Aquí la variable $y$ de $t = y+1$ queda capturada por el cuantificador $\exists y$. El resultado no tiene el significado deseado: en lugar de decir ``$y+1$ tiene un sucesor mayor'', dice ``existe $y$ tal que $y+1 < y$'', que es falso.

\textbf{Solución (renombramiento):} Antes de sustituir, renombramos la variable ligada: $\varphi \equiv \exists z\,(x < z)$. Luego $\varphi[y+1/x] = \exists z\,(y+1 < z)$, que dice correctamente que $y+1$ tiene algo mayor. $t = y+1$ es libre para $x$ en $\exists z\,(x<z)$.
\end{ejemplo}

% ===========================================================
\section{Semántica: interpretaciones y modelos}
% ===========================================================

\subsection{Estructuras de primer orden}

Para dotar de \emph{significado} a las fórmulas de $\mathcal{L}_1(\Lang)$, se recurre al concepto de \textbf{estructura} (o \textbf{interpretación}).

\begin{definicion}{Estructura de primer orden}{estructura}
Una \textbf{$\Lang$-estructura} (o \textbf{$\Lang$-interpretación}) es una tupla
\[
  \Struct{A} = \bigl(A,\; (c^{\Struct{A}})_{c \in C},\; (f^{\Struct{A}})_{f \in \mathcal{F}},\; (R^{\Struct{A}})_{R \in \mathcal{R}}\bigr)
\]
donde:
\begin{itemize}[leftmargin=2em]
  \item $A \neq \emptyset$ es el \textbf{dominio} o \textbf{universo} de la estructura (también escrito $|\Struct{A}|$).
  \item Para cada $c \in C$: $c^{\Struct{A}} \in A$ (una constante designada del dominio).
  \item Para cada $f \in \mathcal{F}$ con aridad $n$: $f^{\Struct{A}} : A^n \to A$ (una función $n$-aria sobre $A$).
  \item Para cada $R \in \mathcal{R}$ con aridad $n$: $R^{\Struct{A}} \subseteq A^n$ (una relación $n$-aria sobre $A$).
\end{itemize}
\end{definicion}

\begin{ejemplo}{Estructuras para $\Lang_{\mathrm{ar}}$}{}
\begin{enumerate}[leftmargin=2em, label=(\alph*)]
  \item \textbf{Modelo estándar de la aritmética:}
  \[
    \mathbb{N} = \bigl(\mathbb{N},\; 0^{\mathbb{N}}=0,\; 1^{\mathbb{N}}=1,\; s^{\mathbb{N}}=\mathrm{sucesor},\; +^{\mathbb{N}}=\mathrm{suma},\; \times^{\mathbb{N}}=\mathrm{producto},\; \leq^{\mathbb{N}}={\leq}\bigr).
  \]
  \item \textbf{Modelo no estándar (anticipación):} Existen estructuras $\Struct{A}$ para $\Lang_{\mathrm{ar}}$ en las que el dominio $A$ contiene elementos ``infinitamente grandes'' (mayores que $n$ para todo $n$ estándar). Estas surgen como consecuencia del Teorema de Compacidad.
  \item \textbf{Interpretación en $\mathbb{Z}$:} Se puede interpretar $\Lang_{\mathrm{ar}}$ en $\mathbb{Z}$ con las operaciones usuales. Nótese que en $\mathbb{Z}$ la fórmula $\exists x\,(x + x = s(0))$ es verdadera (tiene solución: $x = 1/2$... no, en $\mathbb{Z}$ no. Pero la fórmula $\forall x\, \exists y\, (x = y + y)$ falla en $\mathbb{N}$ pero es relevante estudiarla.)
\end{enumerate}
\end{ejemplo}

\subsection{Asignación de variables e interpretación de términos}

Para evaluar una fórmula con variables libres necesitamos asignar elementos del dominio a las variables.

\begin{definicion}{Asignación de variables}{asignacion}
Dada una $\Lang$-estructura $\Struct{A}$, una \textbf{asignación de variables} es una función $\sigma : \mathrm{Var} \to A$.

Si $a \in A$ y $x$ es una variable, escribimos $\sigma[x \mapsto a]$ para la asignación que coincide con $\sigma$ en todas las variables excepto en $x$, donde toma el valor $a$.
\end{definicion}

\begin{definicion}{Interpretación de términos}{interpretacion-terminos}
Dada $(\Struct{A}, \sigma)$, la \textbf{interpretación} $\overline{\sigma}(t) \in A$ de un término $t$ se define inductivamente:
\begin{align*}
  \overline{\sigma}(x) &= \sigma(x) \quad \text{(variable)} \\
  \overline{\sigma}(c) &= c^{\Struct{A}} \quad \text{(constante)} \\
  \overline{\sigma}(f(t_1,\ldots,t_n)) &= f^{\Struct{A}}(\overline{\sigma}(t_1),\ldots,\overline{\sigma}(t_n))
\end{align*}
\end{definicion}

\subsection{La relación de satisfacción}

El corazón de la semántica de primer orden es la \textbf{relación de satisfacción} $\Struct{A} \models \varphi[\sigma]$, que se lee ``la estructura $\Struct{A}$ satisface $\varphi$ bajo la asignación $\sigma$''.

\begin{definicion}{Satisfacción ($\models$)}{satisfaccion}
Dada una $\Lang$-estructura $\Struct{A}$ y una asignación $\sigma$, definimos $\Struct{A} \models \varphi[\sigma]$ inductivamente sobre la estructura de $\varphi$:

\begin{enumerate}[leftmargin=2em, label=\textbf{(\roman*)}]
  \item $\Struct{A} \models (t_1 = t_2)[\sigma]$ \quad sii \quad $\overline{\sigma}(t_1) = \overline{\sigma}(t_2)$ en $A$.
  \item $\Struct{A} \models R(t_1,\ldots,t_n)[\sigma]$ \quad sii \quad $(\overline{\sigma}(t_1),\ldots,\overline{\sigma}(t_n)) \in R^{\Struct{A}}$.
  \item $\Struct{A} \models \neg\varphi[\sigma]$ \quad sii \quad $\Struct{A} \not\models \varphi[\sigma]$.
  \item $\Struct{A} \models (\varphi \wedge \psi)[\sigma]$ \quad sii \quad $\Struct{A} \models \varphi[\sigma]$ \textbf{ y } $\Struct{A} \models \psi[\sigma]$.
  \item $\Struct{A} \models (\varphi \vee \psi)[\sigma]$ \quad sii \quad $\Struct{A} \models \varphi[\sigma]$ \textbf{ o } $\Struct{A} \models \psi[\sigma]$.
  \item $\Struct{A} \models (\varphi \to \psi)[\sigma]$ \quad sii \quad $\Struct{A} \not\models \varphi[\sigma]$ \textbf{ o } $\Struct{A} \models \psi[\sigma]$.
  \item $\Struct{A} \models (\varphi \leftrightarrow \psi)[\sigma]$ \quad sii \quad $\Struct{A} \models \varphi[\sigma]$ $\Leftrightarrow$ $\Struct{A} \models \psi[\sigma]$.
  \item $\Struct{A} \models \forall x\, \varphi[\sigma]$ \quad sii \quad para todo $a \in A$: $\Struct{A} \models \varphi[\sigma[x \mapsto a]]$.
  \item $\Struct{A} \models \exists x\, \varphi[\sigma]$ \quad sii \quad existe $a \in A$ tal que $\Struct{A} \models \varphi[\sigma[x \mapsto a]]$.
\end{enumerate}
\end{definicion}

\begin{importante}
La regla para $\forall$ en el punto (viii) expresa que $\forall x\, \varphi$ es verdadera cuando $\varphi$ es verdadera \emph{para cada elección posible} del valor de $x$ en el dominio. La regla para $\exists$ en (ix) expresa que basta con que $\varphi$ sea verdadera \emph{para algún} valor de $x$. Estas son las reglas semánticas definitivas de los cuantificadores.
\end{importante}

Para sentencias (fórmulas sin variables libres), la asignación $\sigma$ no afecta el resultado, por lo que escribimos simplemente $\Struct{A} \models \varphi$.

\begin{definicion}{Modelo de un conjunto de sentencias}{modelo}
Sea $\Gamma$ un conjunto de sentencias. Decimos que $\Struct{A}$ es un \textbf{modelo} de $\Gamma$, y escribimos $\Struct{A} \models \Gamma$, si $\Struct{A} \models \varphi$ para toda $\varphi \in \Gamma$.
\end{definicion}

\subsection{Verificación semántica: ejemplos desarrollados}

\begin{ejemplo}{Verificación de verdad en $\mathbb{N}$}{}
Tomemos $\Lang_{\mathrm{ar}}$ y la estructura $\mathbb{N}$ con interpretación estándar. Verificamos semánticamente:

\medskip
\textbf{(a)} $\varphi_1 := \forall x\, \forall y\, (x + y = y + x)$.

Para verificar $\mathbb{N} \models \varphi_1$: tomamos cualquier $m, n \in \mathbb{N}$ (cualquier asignación $\sigma[x \mapsto m, y \mapsto n]$) y comprobamos $m + n = n + m$ en $\mathbb{N}$. Esto es la conmutatividad de la suma natural, que es un hecho aritmético conocido. Luego $\mathbb{N} \models \varphi_1$. $\checkmark$

\medskip
\textbf{(b)} $\varphi_2 := \forall x\, \exists y\, (x < y)$ (``no hay máximo'').

Para verificar $\mathbb{N} \models \varphi_2$: dado cualquier $m \in \mathbb{N}$, tomemos $n = m + 1 \in \mathbb{N}$. Entonces $m < m+1$, luego $\mathbb{N} \models (x < y)[\sigma[x\mapsto m, y \mapsto m+1]]$. Como esto vale para todo $m$, se tiene $\mathbb{N} \models \varphi_2$. $\checkmark$

\medskip
\textbf{(c)} $\varphi_3 := \exists x\, (x + x = s(0))$ (``existe $x$ tal que $2x = 1$'').

En $\mathbb{N}$, no existe ningún número natural $m$ con $m + m = 1$. Por tanto $\mathbb{N} \not\models \varphi_3$.

Sin embargo, en $\mathbb{Q}$ (con la interpretación obvia) sí existe tal $x$ (a saber, $x = 1/2$). Esto ilustra que la verdad depende de la estructura elegida.
\end{ejemplo}

\begin{ejemplo}{Construcción de un modelo}{}
Sea $\Gamma = \{\exists x\, \exists y\, (x \neq y),\; \forall x\, (x = x)\}$.

Buscamos una estructura $\Struct{A}$ con $\Struct{A} \models \Gamma$.

Tomemos $A = \{0, 1\}$ (dos elementos) con la igualdad estándar. Entonces:
\begin{itemize}
  \item $\Struct{A} \models \exists x\, \exists y\, (x \neq y)$: tomamos $x \mapsto 0, y \mapsto 1$; claramente $0 \neq 1$. $\checkmark$
  \item $\Struct{A} \models \forall x\, (x = x)$: para todo $a \in A$, $a = a$. $\checkmark$
\end{itemize}
Por tanto $\Struct{A} \models \Gamma$. Cualquier conjunto con al menos dos elementos sirve como modelo de $\Gamma$.
\end{ejemplo}

% ===========================================================
\section{Consecuencia semántica y validez}
% ===========================================================

\subsection{Fórmulas válidas, satisfacibles e insatisfacibles}

\begin{definicion}{Validez y satisfacibilidad}{validez}
Sea $\varphi$ una sentencia de $\mathcal{L}_1$.
\begin{itemize}[leftmargin=2em]
  \item $\varphi$ es \textbf{lógicamente válida} (o \textbf{tautología de primer orden}) si $\Struct{A} \models \varphi$ para toda $\Lang$-estructura $\Struct{A}$. Notación: $\models \varphi$.
  \item $\varphi$ es \textbf{satisfacible} si existe alguna $\Lang$-estructura $\Struct{A}$ con $\Struct{A} \models \varphi$.
  \item $\varphi$ es \textbf{insatisfacible} (o \textbf{contradicción}) si no existe ninguna estructura que la satisfaga.
\end{itemize}
\end{definicion}

\begin{ejemplo}{Tautologías de primer orden}{}
\begin{itemize}
  \item $\models \forall x\, (x = x)$ \quad (reflexividad de la igualdad).
  \item $\models \forall x\, (P(x) \to P(x))$ \quad para cualquier símbolo de relación $P$.
  \item $\models \forall x\, P(x) \to \exists x\, P(x)$ \quad (lo universal implica lo existencial).
  \item $\models \neg(\varphi \wedge \neg\varphi)$ \quad (ley de no contradicción).
\end{itemize}
\end{ejemplo}

\subsection{Consecuencia semántica}

\begin{definicion}{Consecuencia semántica ($\models$)}{consecuencia-semantica}
Sea $\Gamma$ un conjunto de sentencias y $\varphi$ una sentencia. Decimos que $\varphi$ es \textbf{consecuencia semántica} de $\Gamma$, y escribimos $\Gamma \models \varphi$, si para toda $\Lang$-estructura $\Struct{A}$:
\[
  \Struct{A} \models \Gamma \implies \Struct{A} \models \varphi.
\]
En palabras: todo modelo de $\Gamma$ es también modelo de $\varphi$.
\end{definicion}

\begin{observacion}{Casos especiales}{}
\begin{itemize}[leftmargin=2em]
  \item $\emptyset \models \varphi$ equivale a $\models \varphi$ (tautología).
  \item $\Gamma \models \varphi$ para $\varphi \in \Gamma$ siempre (reflexividad).
\end{itemize}
\end{observacion}

\begin{proposicion}{Propiedades de la consecuencia semántica}{}
Sean $\Gamma, \Delta$ conjuntos de sentencias y $\varphi, \psi$ sentencias.
\begin{enumerate}[leftmargin=2em, label=(\roman*)]
  \item \textbf{Reflexividad:} Si $\varphi \in \Gamma$, entonces $\Gamma \models \varphi$.
  \item \textbf{Monotonía:} Si $\Gamma \models \varphi$ y $\Gamma \subseteq \Delta$, entonces $\Delta \models \varphi$.
  \item \textbf{Corte:} Si $\Gamma \models \psi$ y $\Gamma \cup \{\psi\} \models \varphi$, entonces $\Gamma \models \varphi$.
  \item \textbf{Equivalencia con insatisfacibilidad:} $\Gamma \models \varphi$ si y solo si $\Gamma \cup \{\neg\varphi\}$ es insatisfacible.
\end{enumerate}
\end{proposicion}

\begin{proof}
Todas las propiedades se derivan directamente de la definición. Verificamos (iv):

($\Rightarrow$) Si $\Gamma \models \varphi$, todo modelo de $\Gamma$ satisface $\varphi$, luego ningún modelo de $\Gamma$ satisface $\neg\varphi$. Por tanto $\Gamma \cup \{\neg\varphi\}$ es insatisfacible.

($\Leftarrow$) Si $\Gamma \cup \{\neg\varphi\}$ es insatisfacible, no existe modelo de $\Gamma$ que satisfaga $\neg\varphi$; es decir, todo modelo de $\Gamma$ satisface $\varphi$, o sea $\Gamma \models \varphi$.
\end{proof}

\begin{ejemplo}{Ejemplos de consecuencia semántica}{}
\textbf{(a)} Verificar: $\{\forall x\, P(x)\} \models P(c)$ para cualquier término cerrado $c$.

Sea $\Struct{A}$ cualquier estructura con $\Struct{A} \models \forall x\, P(x)$. Esto significa que $a \in P^{\Struct{A}}$ para todo $a \in A$. En particular, $c^{\Struct{A}} \in P^{\Struct{A}}$, luego $\Struct{A} \models P(c)$. $\checkmark$

\medskip
\textbf{(b)} Verificar: $\{\forall x\,(P(x) \to Q(x)),\, \forall x\, P(x)\} \models \forall x\, Q(x)$.

Sea $\Struct{A}$ modelo de ambas hipótesis. Tomemos cualquier $a \in A$. De $\forall x\, P(x)$ obtenemos $a \in P^{\Struct{A}}$. De $\forall x\,(P(x)\to Q(x))$ obtenemos $P^{\Struct{A}}(a) \Rightarrow Q^{\Struct{A}}(a)$. Como la premisa se cumple, $a \in Q^{\Struct{A}}$. Como $a$ fue arbitrario, $\Struct{A} \models \forall x\, Q(x)$. $\checkmark$

\medskip
\textbf{(c)} Mostrar que $\models \forall x\, P(x) \to \exists x\, P(x)$.

Sea $\Struct{A}$ cualquier estructura y supongamos $\Struct{A} \models \forall x\, P(x)$. Como $A \neq \emptyset$ (por definición de estructura), existe $a \in A$ con $a \in P^{\Struct{A}}$. Luego $\Struct{A} \models \exists x\, P(x)$. $\checkmark$
\end{ejemplo}

% ===========================================================
\section{Sistemas de deducción natural}
% ===========================================================

\subsection{Introducción y motivación}

Los sistemas de deducción natural fueron introducidos por Gerhard Gentzen en 1935 \cite{vandalen2013}. La idea central es imitar el estilo de razonamiento matemático informal: uno asume hipótesis temporalmente, realiza pasos intermedios y eventualmente concluye. Cada paso está justificado por una \emph{regla de inferencia}.

En el sistema NK (deducción natural clásica), una \textbf{derivación} de $\varphi$ a partir de hipótesis $\Gamma$ es un árbol de fórmulas en el que:
\begin{itemize}
  \item Las hojas son hipótesis (elementos de $\Gamma$) o supuestos temporales (que deben ser descargados).
  \item Cada nodo interior se obtiene de sus hijos mediante una regla de inferencia.
  \item La raíz es la fórmula $\varphi$ que se demuestra.
\end{itemize}

En la práctica, las derivaciones se presentan en forma lineal numerada, lo cual es más legible.

\subsection{Reglas para la conjunción}

\[
  \infr[\wedge\mathrm{I}]{\varphi \qquad \psi}{\varphi \wedge \psi}
  \qquad\qquad
  \infr[\wedge\mathrm{E}_1]{\varphi \wedge \psi}{\varphi}
  \qquad\qquad
  \infr[\wedge\mathrm{E}_2]{\varphi \wedge \psi}{\psi}
\]

\textbf{Lectura:} $\wedge$I dice que si tenemos demostradas $\varphi$ y $\psi$ por separado, podemos concluir su conjunción. $\wedge$E$_1$ (resp. $\wedge$E$_2$) dice que de una conjunción podemos extraer el primer (resp. segundo) conjunto.

\subsection{Reglas para la disyunción}

\[
  \infr[\vee\mathrm{I}_1]{\varphi}{\varphi \vee \psi}
  \qquad\qquad
  \infr[\vee\mathrm{I}_2]{\psi}{\varphi \vee \psi}
  \qquad\qquad
  \infr[\vee\mathrm{E}]{\varphi \vee \psi \quad [\varphi] \vdash \chi \quad [\psi] \vdash \chi}{\chi}
\]

La regla $\vee$E es la eliminación por casos: si tenemos $\varphi \vee \psi$ y en cada caso (asumiendo $\varphi$, o asumiendo $\psi$) llegamos a $\chi$, entonces podemos concluir $\chi$ incondicionalmente. Los corchetes $[\,\cdot\,]$ indican que esa hipótesis es \emph{descargada} (eliminada) al aplicar la regla.

\subsection{Reglas para la implicación}

\[
  \infr[\to\mathrm{I}]{[\varphi] \vdash \psi}{\varphi \to \psi}
  \qquad\qquad\qquad
  \infr[\to\mathrm{E}]{\varphi \to \psi \qquad \varphi}{\psi}
\]

La regla $\to$I (introducción de implicación) formaliza la \emph{prueba condicional}: para demostrar $\varphi \to \psi$, asumimos temporalmente $\varphi$ y derivamos $\psi$; al terminar, descargamos la hipótesis $\varphi$.

La regla $\to$E es el célebre \emph{modus ponens}: de $\varphi \to \psi$ y $\varphi$, se concluye $\psi$.

\subsection{Reglas para la negación y contradicción}

Sea $\bot$ el símbolo de la \textbf{contradicción} (``lo falso'').

\[
  \infr[\neg\mathrm{I}]{[\varphi] \vdash \bot}{\neg\varphi}
  \qquad\qquad
  \infr[\bot\mathrm{E}]{\bot}{\psi}
  \qquad\qquad
  \infr[\mathrm{DNE}]{[\neg\varphi] \vdash \bot}{\varphi}
\]

$\neg$I (introducción de negación): si asumiendo $\varphi$ llegamos a una contradicción $\bot$, concluimos $\neg\varphi$ y descargamos $\varphi$.

$\bot$E (ex falso quodlibet / explosión): de $\bot$ puede concluirse cualquier cosa.

$\mathrm{DNE}$ (eliminación de doble negación): si asumiendo $\neg\varphi$ llegamos a $\bot$, concluimos $\varphi$. Esta regla corresponde a la \emph{reducción al absurdo clásica} y diferencia la lógica clásica de la intuicionista.

La forma en que se obtiene $\bot$ es:
\[
  \infr[\neg\mathrm{E}]{\varphi \qquad \neg\varphi}{\bot}
\]
De una fórmula y su negación se obtiene $\bot$.

\subsection{Reglas para los cuantificadores}

Estas reglas son la característica distintiva de la lógica de primer orden. Se usa la notación $\varphi[t/x]$ para indicar la sustitución (libre) de $t$ por $x$ en $\varphi$.

\[
  \infr[\forall\mathrm{E}]{\forall x\, \varphi}{\varphi[t/x]}
  \qquad\qquad\qquad
  \infr[\forall\mathrm{I}]{\varphi[a/x] \quad (a \text{ fresca})}{\forall x\,\varphi}
\]

$\forall$E (eliminación universal / instanciación): de $\forall x\, \varphi$ podemos instanciar cualquier término $t$ libre para $x$.

$\forall$I (introducción universal / generalización): si hemos demostrado $\varphi[a/x]$ donde $a$ es un \emph{parámetro fresco} (no ocurre libre en ninguna hipótesis no descargada de la derivación), entonces concluimos $\forall x\, \varphi$. La condición de frescura es crucial: garantiza que $a$ sea realmente arbitrario.

\[
  \infr[\exists\mathrm{I}]{\varphi[t/x]}{\exists x\, \varphi}
  \qquad\qquad\qquad
  \infr[\exists\mathrm{E}]{\exists x\,\varphi \quad [\varphi[a/x]] \vdash \psi \quad (a \text{ fresca en } \psi)}{\psi}
\]

$\exists$I (introducción existencial): si tenemos $\varphi[t/x]$, podemos concluir que existe algo con la propiedad $\varphi$.

$\exists$E (eliminación existencial): si sabemos que $\exists x\,\varphi$ y que asumiendo $\varphi[a/x]$ (con $a$ fresca) podemos derivar $\psi$, entonces concluimos $\psi$. La condición de frescura garantiza que no usamos nada específico sobre $a$.

\subsection{Reglas de igualdad}

\[
  \infr[{=}\mathrm{I}]{\quad}{t = t}
  \qquad\qquad\qquad
  \infr[{=}\mathrm{E}]{t_1 = t_2 \qquad \varphi[t_1/x]}{\varphi[t_2/x]}
\]

${=}$I dice que la igualdad es reflexiva: $t = t$ es siempre demostrable sin hipótesis.

${=}$E (sustitución por iguales): si $t_1 = t_2$ y $\varphi$ es verdadera con $t_1$, entonces también con $t_2$.

\subsection{Derivaciones formales: ejemplos desarrollados}

\begin{ejemplo}{Modus Barbara: $\forall x\,(P(x)\to Q(x)),\, \forall x\, P(x) \vdash \forall x\, Q(x)$}{}
Este es el silogismo bárbaro de la lógica aristotélica, formalizado en deducción natural.

\begin{align*}
(1)\quad &\forall x\,(P(x) \to Q(x)) && \text{hipótesis} \\
(2)\quad &\forall x\, P(x) && \text{hipótesis} \\
(3)\quad &P(a) \to Q(a) && \forall\mathrm{E} \text{ a (1), instanciando } a \text{ fresca} \\
(4)\quad &P(a) && \forall\mathrm{E} \text{ a (2), instanciando } a \\
(5)\quad &Q(a) && \to\mathrm{E} \text{ a (3) y (4)} \\
(6)\quad &\forall x\, Q(x) && \forall\mathrm{I} \text{ a (5), } a \text{ fresca}
\end{align*}

Observación: $a$ es una variable o parámetro fresco que no ocurre en las hipótesis (1) y (2), de modo que la generalización en el paso (6) es legítima.
\end{ejemplo}

\begin{ejemplo}{Dualidad cuantificadores: $\neg\exists x\, P(x) \vdash \forall x\, \neg P(x)$}{}
\begin{align*}
(1)\quad &\neg\exists x\, P(x) && \text{hipótesis} \\
(2)\quad &[P(a)] && \text{supuesto (para } \neg\mathrm{I} \text{), } a \text{ fresca} \\
(3)\quad &\exists x\, P(x) && \exists\mathrm{I} \text{ a (2)} \\
(4)\quad &\bot && \neg\mathrm{E} \text{ a (1) y (3)} \\
(5)\quad &\neg P(a) && \neg\mathrm{I} \text{ a (2)--(4), descargando } [P(a)] \\
(6)\quad &\forall x\, \neg P(x) && \forall\mathrm{I} \text{ a (5), } a \text{ fresca respecto a (1)}
\end{align*}
\end{ejemplo}

\begin{ejemplo}{Transitividad de la igualdad: $x=y,\, y=z \vdash x=z$}{}
\begin{align*}
(1)\quad & x = y && \text{hipótesis} \\
(2)\quad & y = z && \text{hipótesis} \\
(3)\quad & x = x && {=}\mathrm{I} \\
(4)\quad & x = z && {=}\mathrm{E} \text{ a (1) y (3), sustituyendo } x \text{ por } y \text{ en } x = y \text{ ... }
\end{align*}

Más precisamente: de $x = y$ (línea 1) y $x = z$ [donde usamos (2) $y=z$ para reescribir]:

Usando ${=}$E con $y = z$ (línea 2) y la fórmula $x = y$ (línea 1), donde sustituimos $y$ por $z$:
\begin{align*}
(4)\quad & x = z && {=}\mathrm{E} \text{ a (2) y (1): de } y=z \text{ y } x=y, \text{ concluir } x=z
\end{align*}
Aquí aplicamos ${=}$E con $t_1 = y$, $t_2 = z$ y $\varphi(w) = (x = w)$: de $y = z$ y $\varphi[y/w] = (x=y)$, obtenemos $\varphi[z/w] = (x = z)$.
\end{ejemplo}

\begin{ejemplo}{Importación-exportación: $\varphi \to (\psi \to \chi) \vdash (\varphi \wedge \psi) \to \chi$}{}
Esta derivación muestra que la conjunción en la premisa puede ``importarse'' al antecedente de la implicación.

\begin{align*}
(1)\quad &\varphi \to (\psi \to \chi) && \text{hipótesis} \\
(2)\quad &[\varphi \wedge \psi] && \text{supuesto (para }\to\mathrm{I}\text{)} \\
(3)\quad &\varphi && \wedge\mathrm{E}_1 \text{ a (2)} \\
(4)\quad &\psi && \wedge\mathrm{E}_2 \text{ a (2)} \\
(5)\quad &\psi \to \chi && \to\mathrm{E} \text{ a (1) y (3)} \\
(6)\quad &\chi && \to\mathrm{E} \text{ a (5) y (4)} \\
(7)\quad &(\varphi \wedge \psi) \to \chi && \to\mathrm{I} \text{ a (2)--(6), descargando }[\varphi\wedge\psi]
\end{align*}
\end{ejemplo}

% ===========================================================
\section{Consecuencia sintáctica}
% ===========================================================

\subsection{Definición y notación}

\begin{definicion}{Consecuencia sintáctica ($\vdash$)}{consecuencia-sintactica}
Sea $\Gamma$ un conjunto de sentencias y $\varphi$ una sentencia. Decimos que $\varphi$ es \textbf{consecuencia sintáctica} de $\Gamma$ en el sistema NK, y escribimos $\Gamma \vdash \varphi$, si existe una derivación en NK de $\varphi$ a partir de hipótesis en $\Gamma$.

Una sentencia $\varphi$ es un \textbf{teorema} si $\emptyset \vdash \varphi$, escrito simplemente $\vdash \varphi$.
\end{definicion}

La relación $\vdash$ es puramente \emph{sintáctica}: no requiere hablar de ningún modelo. Es una relación de derivabilidad formal mediante reglas simbólicas.

\begin{proposicion}{Propiedades de $\vdash$}{}
\begin{enumerate}[leftmargin=2em, label=(\roman*)]
  \item \textbf{Reflexividad:} $\{\varphi\} \vdash \varphi$.
  \item \textbf{Monotonía:} Si $\Gamma \vdash \varphi$ y $\Gamma \subseteq \Delta$, entonces $\Delta \vdash \varphi$.
  \item \textbf{Corte:} Si $\Gamma \vdash \psi$ y $\Gamma \cup \{\psi\} \vdash \varphi$, entonces $\Gamma \vdash \varphi$.
  \item \textbf{Teorema de la deducción:} $\Gamma \cup \{\varphi\} \vdash \psi$ si y solo si $\Gamma \vdash \varphi \to \psi$.
\end{enumerate}
\end{proposicion}

\subsection{Consistencia}

\begin{definicion}{Consistencia}{consistencia}
Un conjunto de sentencias $\Gamma$ es \textbf{consistente} si $\Gamma \not\vdash \bot$, es decir, si no existe derivación de $\bot$ (la contradicción) a partir de $\Gamma$.

$\Gamma$ es \textbf{inconsistente} si $\Gamma \vdash \bot$, lo que implica $\Gamma \vdash \varphi$ para toda fórmula $\varphi$ (por la regla $\bot$E).
\end{definicion}

% ===========================================================
\section{Los teoremas fundamentales: Corrección, Completitud y Compacidad}
% ===========================================================

\subsection{Teorema de Corrección (Soundness)}

El teorema de corrección establece que el sistema de deducción natural NK no demuestra nada falso: toda demostración formal refleja una verdad semántica.

\begin{teorema}{Corrección de NK}{correccion}
Para todo conjunto $\Gamma$ de sentencias y toda sentencia $\varphi$:
\[
  \Gamma \vdash \varphi \implies \Gamma \models \varphi.
\]
\end{teorema}

\begin{proof}[Esquema de demostración]
Se procede por inducción sobre la estructura de la derivación. El caso base corresponde a hipótesis: si $\varphi \in \Gamma$, todo modelo de $\Gamma$ satisface $\varphi$, luego $\Gamma \models \varphi$.

Para el paso inductivo, se verifica que cada regla de inferencia preserva la consecuencia semántica. Ilustramos con dos casos:

\textbf{Modus ponens ($\to$E):} Si $\Gamma \models \varphi \to \psi$ y $\Gamma \models \varphi$, entonces en todo modelo $\Struct{A}$ de $\Gamma$ se cumplen tanto $\varphi \to \psi$ como $\varphi$. Por la semántica de $\to$, también se cumple $\psi$. Luego $\Gamma \models \psi$.

\textbf{Introducción universal ($\forall$I):} Si $\Gamma \models \varphi[a/x]$ con $a$ no libre en $\Gamma$, entonces en todo modelo $\Struct{A}$ de $\Gamma$, la fórmula $\varphi[a/x]$ es verdadera. Como $a$ no ocurre libre en $\Gamma$, su interpretación puede ser cualquier elemento de $A$. Luego $\varphi[b/x]$ es verdadera para todo $b \in A$, i.e., $\Struct{A} \models \forall x\, \varphi$.

Los demás casos son análogos.
\end{proof}

\begin{importante}
El teorema de corrección garantiza que NK es una herramienta \emph{confiable}: si se logra una demostración formal de $\varphi$, entonces $\varphi$ es semánticamente verdadera en todo modelo de las hipótesis. No puede haber ``demostraciones incorrectas'' en un sistema correcto.
\end{importante}

\subsection{Teorema de Completitud de Gödel}

La completitud es la propiedad recíproca a la corrección: todo lo que es semánticamente consecuencia puede demostrarse formalmente. Este resultado fue establecido por Kurt Gödel en su tesis doctoral de 1930 \cite{godel1930, enderton2001}.

\begin{teorema}{Completitud de Gödel (1930)}{completitud}
Para todo conjunto $\Gamma$ de sentencias y toda sentencia $\varphi$:
\[
  \Gamma \models \varphi \implies \Gamma \vdash \varphi.
\]
En particular, $\models \varphi \iff \vdash \varphi$.

Equivalentemente: $\Gamma$ es consistente (i.e., $\Gamma \not\vdash \bot$) si y solo si $\Gamma$ tiene un modelo.
\end{teorema}

La equivalencia entre las dos formulaciones del teorema se deduce del Lema~\ref{lem:equiv-completitud}.

\begin{lema}{Equivalencia de formulaciones}{equiv-completitud}
Las siguientes afirmaciones son equivalentes para todo $\Gamma$ y $\varphi$:
\begin{enumerate}[label=(\roman*), leftmargin=2em]
  \item $\Gamma \models \varphi \implies \Gamma \vdash \varphi$.
  \item Todo $\Gamma$ consistente tiene un modelo.
\end{enumerate}
\end{lema}

\begin{proof}
$(i) \Rightarrow (ii)$: Supongamos que $\Gamma$ es consistente. Si $\Gamma$ no tuviera modelo, entonces $\Gamma \models \varphi$ para toda $\varphi$ (vacuamente), y en particular $\Gamma \models \bot$. Por (i), $\Gamma \vdash \bot$, contradiciendo la consistencia.

$(ii) \Rightarrow (i)$: Supongamos $\Gamma \models \varphi$. Entonces $\Gamma \cup \{\neg\varphi\}$ es insatisfacible. Por (ii) aplicado en su contrarrecíproco, $\Gamma \cup \{\neg\varphi\}$ es inconsistente, i.e., $\Gamma \cup \{\neg\varphi\} \vdash \bot$. Por la regla DNE, $\Gamma \vdash \varphi$.
\end{proof}

\subsubsection{Idea de la demostración: la construcción de Henkin}

La demostración de la completitud utiliza el método de \textbf{construcción de modelos canónicos}, desarrollado por Leon Henkin (1949) y que es una de las técnicas más elegantes de la lógica matemática. La presentamos como esquema:

\medskip
\textbf{Paso 1 — Lema de Lindenbaum.} Sea $\Gamma$ un conjunto consistente de sentencias en un lenguaje $\Lang$. Podemos extender $\Gamma$ a un conjunto $\Gamma^*$ de sentencias que sea \emph{maximalmente consistente}: todo $\Gamma^*$ con $\Gamma \subseteq \Gamma^*$ que sea consistente y tal que, para toda sentencia $\varphi$, exactamente una de $\varphi$ o $\neg\varphi$ pertenece a $\Gamma^*$.

La construcción de $\Gamma^*$ usa el axioma de elección o la enumeración de todas las sentencias: se añaden una a una, cuidando de no crear inconsistencia.

\textbf{Paso 2 — Extensión de Henkin.} Ampliamos el lenguaje $\Lang$ a $\Lang^+$ añadiendo una cantidad contable de nuevas constantes $\{c_{\varphi}\}$, una por cada sentencia existencial $\varphi = \exists x\, \psi$. Luego extendemos $\Gamma^*$ añadiendo los \emph{testigos de Henkin}:
\[
  \exists x\, \psi(x) \to \psi(c_{\exists x\,\psi(x)}).
\]
Esto garantiza que para toda existencial en $\Gamma^*$, hay una constante que ``testimonia'' dicha existencia. El conjunto extendido sigue siendo consistente y maximalmente consistente en $\Lang^+$.

\textbf{Paso 3 — Modelo canónico o de términos.} Definimos la estructura $\Struct{M}$:
\begin{itemize}
  \item \textbf{Dominio:} $|\Struct{M}|$ = clases de equivalencia de términos cerrados de $\Lang^+$ bajo la relación $t_1 \sim t_2 \iff (t_1 = t_2) \in \Gamma^*$.
  \item \textbf{Funciones y relaciones:} interpretadas de manera natural a través de los términos.
\end{itemize}

\textbf{Paso 4 — Lema de verdad.} El resultado central es:
\[
  \Struct{M} \models \varphi \iff \varphi \in \Gamma^*.
\]
Este lema se demuestra por inducción sobre la complejidad de $\varphi$. El caso de los cuantificadores usa esencialmente los testigos de Henkin del Paso 2.

\textbf{Conclusión.} Como $\Gamma \subseteq \Gamma^*$, el modelo $\Struct{M}$ satisface toda sentencia de $\Gamma$. Luego $\Struct{M} \models \Gamma$, y $\Gamma$ tiene un modelo.

\begin{importante}
La completitud de Gödel tiene una consecuencia filosófica profunda: la distinción entre ``semánticamente verdadero en todos los modelos'' y ``formalmente demostrable'' desaparece en la lógica de primer orden. Esto no es trivial: Gödel mismo demostró después (en sus teoremas de incompletitud de 1931) que en teorías suficientemente expresivas como la aritmética de Peano, hay sentencias verdaderas en el modelo estándar que no son demostrables en el sistema formal. Estos resultados no se contradicen: la completitud garantiza que todo lo válido en \emph{todos} los modelos es demostrable, pero no todo lo verdadero en \emph{un} modelo particular.
\end{importante}

\subsection{Teorema de Compacidad}

El Teorema de Compacidad es una consecuencia inmediata de la Completitud y tiene aplicaciones sorprendentes en toda la matemática.

\begin{teorema}{Compacidad de la lógica de primer orden}{compacidad}
Sea $\Gamma$ un conjunto (posiblemente infinito) de sentencias. Entonces:
\[
  \Gamma \text{ tiene un modelo} \iff \text{todo subconjunto finito de } \Gamma \text{ tiene un modelo.}
\]
Equivalentemente: $\Gamma \models \varphi$ si y solo si existe un subconjunto finito $\Delta \subseteq \Gamma$ tal que $\Delta \models \varphi$.
\end{teorema}

\begin{proof}
$(\Rightarrow)$ Si $\Struct{A} \models \Gamma$, entonces en particular $\Struct{A} \models \Delta$ para todo $\Delta \subseteq \Gamma$. Trivial.

$(\Leftarrow)$ Supongamos que $\Gamma$ no tiene modelo. Por el Teorema de Completitud, $\Gamma$ es inconsistente: $\Gamma \vdash \bot$. Toda derivación en NK es una estructura finita y solo usa finitely many hipótesis; por tanto, existe un subconjunto finito $\Delta \subseteq \Gamma$ con $\Delta \vdash \bot$. Por el Teorema de Corrección, $\Delta$ es insatisfacible: $\Delta$ no tiene modelo. Luego no todo subconjunto finito de $\Gamma$ tiene modelo.
\end{proof}

\subsubsection{Aplicación 1: Modelos no estándar de la aritmética}

Sea $\mathrm{PA}$ el conjunto de axiomas de Peano para $\Lang_{\mathrm{ar}}$, y definamos la sentencia $\varphi_n := s(\cdots s(0)\cdots) < x$ (``$x$ es mayor que el numeral $n$'', con $n$ aplicaciones del sucesor). Consideremos:
\[
  \Gamma = \mathrm{PA} \cup \{\varphi_n : n \in \mathbb{N}\}.
\]

\textbf{Todo subconjunto finito de $\Gamma$ tiene modelo:} dado $\Delta$ finito, contiene $\varphi_0, \ldots, \varphi_k$ para algún $k$. El modelo estándar $\mathbb{N}$ con $x$ interpretado como $k+1$ satisface todos los $\varphi_j$ para $j \leq k$.

\textbf{Por Compacidad}, $\Gamma$ tiene un modelo $\Struct{A}$. En $\Struct{A}$ existe un elemento $a = x^{\Struct{A}}$ que es mayor que el numeral $n$ para todo $n \in \mathbb{N}$. Este elemento $a$ es ``infinitamente grande'' y no corresponde a ningún número natural estándar. Así, $\Struct{A}$ es un \textbf{modelo no estándar de la aritmética}.

\begin{importante}
Este resultado tiene consecuencias filosóficas profundas: ningún conjunto de axiomas de primer orden puede \emph{caracterizar} completamente los números naturales. Siempre existe un modelo no estándar con elementos ``infinitos''. Esto está relacionado con el Teorema de Löwenheim–Skolem, que estudiaremos más adelante.
\end{importante}

\subsubsection{Aplicación 2: Coloración de grafos infinitos}

\begin{corolario}{Coloración de grafos}{coloracion}
Si todo subgrafo finito de un grafo $G$ es $k$-colorable (para algún $k$ fijo), entonces $G$ es $k$-colorable.
\end{corolario}

\begin{proof}[Esquema]
Se codifica el problema de coloración como un conjunto $\Gamma_G$ de sentencias de primer orden: para cada vértice $v$ se introduce una constante $c_v$, para cada color $1,\ldots,k$ un predicado unario $C_i$, y se añaden sentencias que expresan:
\begin{enumerate}
  \item Cada vértice tiene al menos un color: $\bigvee_{i=1}^k C_i(c_v)$.
  \item Vértices adyacentes tienen colores distintos: $\neg(C_i(c_v) \wedge C_i(c_w))$ para cada arista $\{v,w\}$ y cada $i$.
\end{enumerate}

Todo subconjunto finito de $\Gamma_G$ involucra finitely many vértices, por lo que corresponde a un subgrafo finito, que por hipótesis es $k$-colorable. Por Compacidad, $\Gamma_G$ tiene un modelo, que proporciona una $k$-coloración de $G$.
\end{proof}

\subsubsection{Aplicación 3: Teoría de modelos}

La Compacidad es la herramienta central de la teoría de modelos \cite{marker2002}. Entre sus usos:
\begin{itemize}
  \item \textbf{Construcción de extensiones elementales}: toda estructura infinita tiene una extensión elemental de cardinalidad mayor (Teorema de Löwenheim–Skolem hacia arriba).
  \item \textbf{Transferencia de propiedades}: si una propiedad es expresable en $\mathcal{L}_1$ y falla en estructuras arbitrariamente grandes, entonces falla en alguna estructura infinita.
  \item \textbf{Teorema de los cuatro colores}: la versión infinita de la conjetura de los cuatro colores (todo grafo planar es 4-colorable) se reduce al caso finito via Compacidad.
\end{itemize}

% ===========================================================
\section{Síntesis y perspectiva}
% ===========================================================

A lo largo de esta clase hemos construido el andamiaje lógico sobre el que se apoya todo el curso. Las nociones fundamentales son:

\begin{enumerate}[leftmargin=2em, label=\textbf{\arabic*.}]
  \item \textbf{Sintaxis:} el lenguaje $\mathcal{L}_1(\Lang)$ con sus términos, fórmulas, variables libres y sustitución.
  \item \textbf{Semántica:} las $\Lang$-estructuras como modelos, la relación de satisfacción $\models$ que conecta símbolos formales con objetos matemáticos.
  \item \textbf{Consecuencia semántica ($\models$):} verdad en todos los modelos de las hipótesis.
  \item \textbf{Deducción natural NK:} reglas de inferencia que formalizan el razonamiento matemático.
  \item \textbf{Consecuencia sintáctica ($\vdash$):} derivabilidad formal.
  \item \textbf{Corrección + Completitud} (Gödel, 1930): $\Gamma \models \varphi \iff \Gamma \vdash \varphi$.
  \item \textbf{Compacidad:} satisfacibilidad es un asunto local (finito).
\end{enumerate}

Estos resultados garantizan que podamos trabajar con la teoría de conjuntos $\mathrm{ZFC}$ con plena confianza: la firma $\Lang_\in = \{\in\}$, los axiomas de $\mathrm{ZFC}$ como un conjunto $\Gamma$, y las demostraciones matemáticas como derivaciones en NK (o en sistemas equivalentes). El Teorema de Completitud asegura que toda consecuencia semántica de $\mathrm{ZFC}$ es alcanzable por medios formales.

% ===========================================================
\section{Problemas y tareas}
% ===========================================================

Los siguientes problemas consolidan y amplían los conceptos de la clase. Se recomienda resolver cada uno con el nivel de formalidad adecuado: los problemas de derivación deben presentarse con pasos numerados y justificaciones explícitas.

\medskip

\begin{enumerate}[leftmargin=2em, label=\textbf{Problema \arabic*.}, itemsep=1em]

  \item \textbf{[Términos y fórmulas]}
  Usando el lenguaje de los grupos $\Lang_{\mathrm{gr}} = (\{e\}, \{\cdot, (-)^{-1}\}, \emptyset)$:
  \begin{enumerate}[label=(\alph*)]
    \item Escribe cinco términos distintos que involucren variables $x$ e $y$.
    \item Formula en $\mathcal{L}_1(\Lang_{\mathrm{gr}})$ los axiomas de grupo: asociatividad, identidad por la izquierda e inverso por la izquierda.
    \item Determina $\FV(\varphi)$ para cada uno de los axiomas escritos.
  \end{enumerate}

  \item \textbf{[Variables libres y sustitución]}
  Considera la fórmula $\varphi := \forall y\, (R(x,y) \to \exists z\, R(y,z))$ en el lenguaje con única relación binaria $R$.
  \begin{enumerate}[label=(\alph*)]
    \item Calcula $\FV(\varphi)$.
    \item Calcula $\varphi[c/x]$ donde $c$ es una constante.
    \item Explica qué problemas surgen al intentar calcular $\varphi[y/x]$ y cómo resolverlos.
    \item Calcula $\varphi[f(y)/x]$ donde $f$ es un símbolo de función unaria, renombrando variables si es necesario.
  \end{enumerate}

  \item \textbf{[Satisfacción semántica]}
  Sea $\Lang = \{+^{(2)}, 0^{(0)}, s^{(1)}\}$ y considera las estructuras $\mathbb{N}$ (naturales estándar) y $\mathbb{Z}$ (enteros) con sus interpretaciones naturales.
  \begin{enumerate}[label=(\alph*)]
    \item Verifica semánticamente en $\mathbb{N}$ la fórmula: $\forall x\,(x + s(0) \neq x)$.
    \item Encuentra una fórmula $\varphi$ que sea verdadera en $\mathbb{Z}$ pero falsa en $\mathbb{N}$.
    \item Encuentra una fórmula $\psi$ verdadera en $\mathbb{N}$ pero falsa en la estructura $(\{0\}, 0^A = 0, s^A(0) = 0, +^A = \text{trivial})$.
  \end{enumerate}

  \item \textbf{[Consecuencia semántica]}
  Determina, justificando, si las siguientes relaciones de consecuencia semántica son válidas. Si son válidas, da un argumento semántico. Si no, exhibe un contraejemplo (una estructura que satisfaga las hipótesis pero no la conclusión).
  \begin{enumerate}[label=(\alph*)]
    \item $\{\exists x\, P(x),\; \forall x\,(P(x) \to Q(x))\} \models \exists x\, Q(x)$.
    \item $\{\forall x\, P(x)\} \models \forall x\, \forall y\, (P(x) \wedge P(y))$.
    \item $\{\exists x\, P(x),\; \exists x\, Q(x)\} \models \exists x\,(P(x) \wedge Q(x))$.
    \item $\emptyset \models \forall x\, \exists y\, (x = y)$.
  \end{enumerate}

  \item \textbf{[Derivaciones en NK -- Nivel básico]}
  Construye derivaciones formales en NK para:
  \begin{enumerate}[label=(\alph*)]
    \item $\{\varphi \to \psi,\; \psi \to \chi\} \vdash \varphi \to \chi$ \hfill (silogismo hipotético).
    \item $\{\varphi \wedge \psi\} \vdash \psi \wedge \varphi$ \hfill (conmutatividad de la conjunción).
    \item $\{\varphi \to \psi,\; \neg\psi\} \vdash \neg\varphi$ \hfill (modus tollens).
    \item $\vdash \varphi \to (\psi \to \varphi)$ \hfill (afirmación del consecuente).
  \end{enumerate}

  \item \textbf{[Derivaciones en NK -- Cuantificadores]}
  Construye derivaciones formales en NK para:
  \begin{enumerate}[label=(\alph*)]
    \item $\{\forall x\, \forall y\, R(x,y)\} \vdash \forall y\, \forall x\, R(x,y)$.
    \item $\{\exists x\, \neg P(x)\} \vdash \neg \forall x\, P(x)$.
    \item $\{\neg \forall x\, P(x)\} \vdash \exists x\, \neg P(x)$ \hfill (\textit{Sugerencia: usa DNE}).
    \item $\{\forall x\,(P(x) \wedge Q(x))\} \vdash \forall x\, P(x) \wedge \forall x\, Q(x)$.
  \end{enumerate}

  \item \textbf{[Igualdad]}
  \begin{enumerate}[label=(\alph*)]
    \item Demuestra formalmente que la igualdad es simétrica: $x = y \vdash y = x$.
    \item Demuestra que la igualdad es transitiva: $x = y,\; y = z \vdash x = z$ (usa ${=}$E con la fórmula apropiada).
    \item Si $f$ es un símbolo de función unaria, prueba: $x = y \vdash f(x) = f(y)$.
  \end{enumerate}

  \item \textbf{[Corrección e incompletitud]}
  \begin{enumerate}[label=(\alph*)]
    \item Enuncia con precisión el Teorema de Corrección y explica por qué su demostración va por inducción sobre las derivaciones.
    \item Explica (informalmente) por qué el Teorema de Completitud de Gödel (1930) \emph{no} contradice los Teoremas de Incompletitud de Gödel (1931). ¿Qué demuestra cada uno?
    \item Sea $T$ una teoría consistente en $\mathcal{L}_1$. Argumenta usando el Teorema de Completitud por qué $T$ tiene un modelo.
  \end{enumerate}

  \item \textbf{[Compacidad -- Aplicaciones]}
  \begin{enumerate}[label=(\alph*)]
    \item Sea $\mathrm{DLO}$ la teoría del orden linear denso sin extremos (con axiomas: orden total, densidad y no extremos). Usando Compacidad, demuestra que cualquier teoría que extienda $\mathrm{DLO}$ con una cantidad finita de constantes que satisfacen un orden específico tiene un modelo.
    \item Enuncia y demuestra (usando Compacidad) el siguiente resultado: si una propiedad $P$ sobre grafos es preservada bajo subgrafos finitos y todo grafo finito satisface $P$, entonces todo grafo infinito satisface $P$. ¿Esta afirmación es cierta en general? Da un contraejemplo o una prueba.
    \item (Difícil) Sea $\Gamma = \mathrm{PA} \cup \{\overline{n} < c : n \in \mathbb{N}\}$ donde $\overline{n}$ es el numeral para $n$ y $c$ es una nueva constante. Usando Compacidad, demuestra que $\Gamma$ tiene un modelo. ¿Qué dice esto sobre la teoría de los números naturales?
  \end{enumerate}

  \item \textbf{[Proyecto: La construcción de Henkin]}
  Investiga y escribe una presentación detallada (2--4 páginas) de la demostración completa del Teorema de Completitud de Gödel siguiendo la construcción de Henkin. Tu presentación debe incluir:
  \begin{enumerate}[label=(\alph*)]
    \item La definición precisa de conjunto maximalmente consistente.
    \item El Lema de Lindenbaum con su demostración.
    \item La extensión de Henkin: cómo y por qué se agregan los testigos existenciales.
    \item La construcción del modelo canónico: la relación de equivalencia $\sim$ en los términos y la definición de las funciones y relaciones.
    \item El Lema de Verdad (por inducción sobre fórmulas).
    \item La conclusión: todo conjunto consistente tiene un modelo.
  \end{enumerate}
  Referencias sugeridas: \cite{enderton2001}, \cite{ebbinghaus1994}, \cite{vandalen2013}.

\end{enumerate}

% ===========================================================
\chapter{Clase 2: Axiomática de Zermelo–Fraenkel con el Axioma de Elección (ZFC)}
% ===========================================================

La teoría de conjuntos de \textbf{Zermelo–Fraenkel con el Axioma de Elección}, conocida universalmente como \textbf{ZFC}, es el sistema axiomático que sirve de fundamento estándar para la totalidad de las matemáticas contemporáneas. A partir de un único símbolo primitivo —la relación de pertenencia $\in$— y de nueve axiomas más el Axioma de Elección, es posible construir los números naturales, reales y complejos, toda la teoría de la medida, el álgebra abstracta, el análisis funcional, la geometría diferencial, y —lo que más nos interesa en este curso— la jerarquía transfinita de números ordinales y cardinales.

En esta clase presentamos los axiomas de ZF de manera formal y comentada, añadimos el Axioma de Elección, y estudiamos sus equivalencias más importantes: el \textbf{Lema de Zorn} y el \textbf{Teorema del Buen Orden}. Discutiremos también por qué ZFC constituye el sistema preferido y cuáles son sus limitaciones inherentes \cite{kunen2011, jech2003, halmos1960}.

% ===========================================================
\section{El lenguaje de la teoría de conjuntos}
% ===========================================================

Toda la teoría ZFC se expresa en el lenguaje de primer orden con una única relación no lógica:
\[
  \mathcal{L}_{\in} = (\emptyset,\, \emptyset,\, \{\in^{(2)}\}).
\]
Los únicos objetos que existen son los \emph{conjuntos}. En ZFC no hay \emph{ur-elementos} (objetos que no sean conjuntos); todo lo que existe es un conjunto, y los conjuntos se relacionan únicamente mediante la pertenencia $\in$.

La relación de \textbf{inclusión} se introduce como abreviatura:
\[
  A \subseteq B \;\stackrel{\mathrm{def}}{\iff}\; \forall x\,(x \in A \to x \in B).
\]

\begin{observacion}{Notación y convenios}{conv}
A lo largo de esta clase usaremos las siguientes abreviaturas estándar:
\begin{itemize}[leftmargin=2em]
  \item $A \subsetneq B$ (inclusión estricta): $A \subseteq B \wedge A \neq B$.
  \item $\exists! x\,\varphi(x)$ (existencia y unicidad): $\exists x\,(\varphi(x) \wedge \forall y\,(\varphi(y) \to y = x))$.
  \item $\{x \mid \varphi(x)\}$ denota la clase de todos los conjuntos que satisfacen $\varphi$. Esta notación es \emph{informal} dentro de ZFC (puede no ser un conjunto), como veremos al estudiar el Axioma de Separación.
\end{itemize}
\end{observacion}

% ===========================================================
\section{Los nueve axiomas de Zermelo–Fraenkel}
% ===========================================================

Presentamos cada axioma con su nombre, su fórmula en $\mathcal{L}_{\in}$, una explicación intuitiva y ejemplos desarrollados.

% -----------------------------------------------------------
\subsection{Axioma 1: Extensionalidad}
% -----------------------------------------------------------

\begin{definicion}{Axioma de Extensionalidad}{ext}
\[
  \forall A\,\forall B\,\bigl(\forall x\,(x \in A \leftrightarrow x \in B) \to A = B\bigr).
\]
\end{definicion}

\textbf{Significado.} Dos conjuntos son iguales si y solo si tienen los mismos elementos. Este axioma captura la idea de que un conjunto queda completamente determinado por su extensión (sus miembros), no por la forma en que se lo describe.

\begin{ejemplo}{Extensionalidad en acción}{}
Consideremos las siguientes descripciones:
\begin{itemize}[leftmargin=2em]
  \item $A = \{x \in \mathbb{N} \mid x^2 - 3x + 2 = 0\} = \{1, 2\}$.
  \item $B = \{x \in \mathbb{N} \mid x < 3 \wedge x \geq 1\} = \{1, 2\}$.
\end{itemize}
Aunque $A$ y $B$ se describen mediante condiciones distintas, el Axioma de Extensionalidad garantiza que $A = B$, puesto que $\forall x\,(x \in A \leftrightarrow x \in B)$.

De manera más abstracta: $\{1, 2, 1\} = \{1, 2\}$, ya que la repetición de elementos no altera la extensión.
\end{ejemplo}

\begin{observacion}{Consecuencia fundamental}{unicidad}
El Axioma de Extensionalidad implica que, si un axioma garantiza la existencia de un conjunto con cierta propiedad, dicho conjunto es \emph{único}. Esto justifica el uso del artículo determinado («el conjunto vacío», «el par $\{a,b\}$», etc.) y del símbolo $\exists!$ en muchos resultados.
\end{observacion}

% -----------------------------------------------------------
\subsection{Axioma 2: Vacío (o Existencia)}
% -----------------------------------------------------------

\begin{definicion}{Axioma del Conjunto Vacío}{vacio}
\[
  \exists A\,\forall x\,(x \notin A).
\]
\end{definicion}

\textbf{Significado.} Existe al menos un conjunto que no tiene elementos. Combinado con el Axioma de Extensionalidad, este conjunto es único y se denota $\emptyset$.

\begin{ejemplo}{Unicidad del conjunto vacío}{}
Supongamos que $A$ y $B$ son dos conjuntos vacíos. Entonces para todo $x$: $x \notin A$ y $x \notin B$, de modo que trivialmente $x \in A \leftrightarrow x \in B$ (ambas partes son falsas). Por el Axioma de Extensionalidad, $A = B$.
\end{ejemplo}

\begin{observacion}{Necesidad del axioma}{nec-vacio}
En algunos sistemas axiomáticos la existencia del conjunto vacío se deduce a partir del Axioma de Separación (aplicado a cualquier conjunto existente). Sin embargo, para que Separación tenga algo sobre qué operar, se necesita al menos un conjunto; de ahí que muchos autores incluyan el Axioma del Vacío de manera explícita como garantía de que el universo conjuntista no es vacío.
\end{observacion}

% -----------------------------------------------------------
\subsection{Axioma 3: Pares}
% -----------------------------------------------------------

\begin{definicion}{Axioma de Pares}{par}
\[
  \forall a\,\forall b\,\exists P\,\forall x\,(x \in P \leftrightarrow x = a \vee x = b).
\]
\end{definicion}

\textbf{Significado.} Para cualesquiera dos conjuntos $a$ y $b$, existe el conjunto $\{a, b\}$ cuyos únicos miembros son $a$ y $b$.

\begin{ejemplo}{Construcción de singletons y pares}{}
\begin{enumerate}[leftmargin=2em, label=(\alph*)]
  \item Tomando $a = b = \emptyset$, el axioma garantiza la existencia de $\{\emptyset, \emptyset\} = \{\emptyset\}$. Así obtenemos el \textbf{singleton} del vacío.

  \item A partir de $\emptyset$ y $\{\emptyset\}$, el axioma produce $\{\emptyset, \{\emptyset\}\}$, que en la construcción de von Neumann representa el número natural $2$.

  \item La construcción iterada de pares nos permite construir todos los naturales: $0 = \emptyset$, $1 = \{0\} = \{\emptyset\}$, $2 = \{0,1\} = \{\emptyset, \{\emptyset\}\}$, y en general $n+1 = n \cup \{n\}$ (donde la unión requiere el Axioma de Unión).
\end{enumerate}
\end{ejemplo}

% -----------------------------------------------------------
\subsection{Axioma 4: Unión}
% -----------------------------------------------------------

\begin{definicion}{Axioma de Unión}{union}
\[
  \forall \mathcal{F}\,\exists U\,\forall x\bigl(x \in U \leftrightarrow \exists A\,(A \in \mathcal{F} \wedge x \in A)\bigr).
\]
\end{definicion}

\textbf{Significado.} Para toda familia de conjuntos $\mathcal{F}$, existe la unión $\bigcup \mathcal{F}$, cuyos elementos son exactamente los elementos de los miembros de $\mathcal{F}$.

\begin{ejemplo}{Unión de una familia}{}
\begin{enumerate}[leftmargin=2em, label=(\alph*)]
  \item Sea $\mathcal{F} = \{\{1,2\},\{2,3\},\{3,4\}\}$.
    Entonces $\bigcup \mathcal{F} = \{1,2,3,4\}$.

  \item Sea $\mathcal{F} = \{a, b\}$. Entonces $\bigcup\{a,b\} = a \cup b$ en el sentido usual.

  \item La unión binaria $a \cup b$ se define como $\bigcup\{a,b\}$, que existe por los Axiomas de Pares y Unión.

  \item La sucesión de von Neumann: $n+1 = n \cup \{n\} = \bigcup\{n, \{n\}\}$. El Axioma de Unión garantiza que este sucesor existe.
\end{enumerate}
\end{ejemplo}

% -----------------------------------------------------------
\subsection{Axioma 5: Potencia (Conjunto de Partes)}
% -----------------------------------------------------------

\begin{definicion}{Axioma de Potencia}{potencia}
\[
  \forall A\,\exists P\,\forall x\,(x \in P \leftrightarrow x \subseteq A).
\]
\end{definicion}

\textbf{Significado.} Para todo conjunto $A$, existe el conjunto $\mathcal{P}(A)$ —llamado \textbf{conjunto potencia} o \textbf{conjunto de partes} de $A$— que contiene exactamente todos los subconjuntos de $A$.

\begin{ejemplo}{Cómputo de conjuntos potencia}{}
\begin{enumerate}[leftmargin=2em, label=(\alph*)]
  \item $\mathcal{P}(\emptyset) = \{\emptyset\}$. El único subconjunto del vacío es el propio vacío.

  \item $\mathcal{P}(\{\emptyset\}) = \{\emptyset,\, \{\emptyset\}\}$. Tiene $2^1 = 2$ elementos.

  \item $\mathcal{P}(\{0,1,2\}) = \{\emptyset,\{0\},\{1\},\{2\},\{0,1\},\{0,2\},\{1,2\},\{0,1,2\}\}$. Tiene $2^3 = 8$ elementos.

  \item Para un conjunto finito con $n$ elementos, $|\mathcal{P}(A)| = 2^n$. Para conjuntos infinitos, el teorema de Cantor establece que $|\mathcal{P}(A)| > |A|$ en cardinalidad (ver Clase 3).
\end{enumerate}
\end{ejemplo}

\begin{observacion}{Importancia del Axioma de Potencia}{imp-pot}
El Axioma de Potencia es responsable de la existencia de los números reales (como subconjuntos de $\mathbb{N}$ o secuencias de bits), de espacios funcionales $B^A$, y de toda la jerarquía de cardinalidades infinitas más allá de $\aleph_0$. Sin él, el universo conjuntista sería demasiado pequeño para albergar el análisis real.
\end{observacion}

% -----------------------------------------------------------
\subsection{Axioma 6: Infinito}
% -----------------------------------------------------------

\begin{definicion}{Axioma del Infinito}{infinito}
\[
  \exists I\,\bigl(\emptyset \in I \;\wedge\; \forall x\,(x \in I \to x \cup \{x\} \in I)\bigr).
\]
\end{definicion}

\textbf{Significado.} Existe al menos un conjunto \emph{inductivo}: un conjunto que contiene $\emptyset$ y es cerrado bajo la operación sucesor $x \mapsto x \cup \{x\}$.

\begin{ejemplo}{El conjunto de von Neumann de los naturales}{}
El Axioma del Infinito garantiza la existencia de un conjunto inductivo $I$. Mediante el Axioma de Separación (ver más adelante) se puede extraer el \textbf{menor} conjunto inductivo:
\[
  \omega = \{x \in I \mid x \text{ pertenece a todo conjunto inductivo incluido en } I\}.
\]
Este conjunto $\omega$ contiene exactamente:
\[
  \omega = \{0, 1, 2, 3, \ldots\} = \{\emptyset,\, \{\emptyset\},\, \{\emptyset,\{\emptyset\}\},\, \ldots\}
\]
y sirve como modelo estándar de los números naturales dentro de ZFC. La construcción explicita los primeros elementos:
\begin{align*}
  0 &= \emptyset, \\
  1 &= \{0\} = \{\emptyset\}, \\
  2 &= \{0,1\} = \{\emptyset, \{\emptyset\}\}, \\
  3 &= \{0,1,2\} = \{\emptyset, \{\emptyset\}, \{\emptyset,\{\emptyset\}\}\}.
\end{align*}
\end{ejemplo}

% -----------------------------------------------------------
\subsection{Axioma 7: Esquema de Reemplazo}
% -----------------------------------------------------------

\begin{definicion}{Esquema de Axiomas de Reemplazo}{reemplazo}
Para toda fórmula $\varphi(x, y, p_1, \ldots, p_k)$ de $\mathcal{L}_{\in}$ que sea \emph{funcional} en los parámetros dados (es decir, $\forall x\,\exists! y\,\varphi(x,y,\vec{p})$), el siguiente enunciado es un axioma:
\[
  \forall \vec{p}\,\forall A\,\bigl(\forall x \in A\,\exists! y\,\varphi(x,y,\vec{p})\bigr)
  \to \exists B\,\forall y\,\bigl(y \in B \leftrightarrow \exists x \in A\,\varphi(x,y,\vec{p})\bigr).
\]
\end{definicion}

\textbf{Significado.} Si $\varphi$ define una función $F: A \to V$ sobre un conjunto $A$ (donde $V$ es el universo conjuntista), entonces la imagen $F[A] = \{F(x) \mid x \in A\}$ es también un conjunto.

\begin{ejemplo}{Reemplazo para construir conjuntos de ordinales}{}
Sea $A = \omega$ y $\varphi(n, y)$ la fórmula ``$y$ es el $n$-ésimo número primo''. Esta es una función bien definida $F: \omega \to \omega$. El Esquema de Reemplazo garantiza que el conjunto de todos los números primos:
\[
  \{F(0), F(1), F(2), \ldots\} = \{2, 3, 5, 7, 11, \ldots\}
\]
es un conjunto.

Otro ejemplo: sea $\varphi(n, y)$ la fórmula ``$y = \mathcal{P}^n(\omega)$'' (la $n$-ésima iteración del conjunto potencia de $\omega$). Entonces el Esquema de Reemplazo garantiza la existencia de:
\[
  \{\omega,\, \mathcal{P}(\omega),\, \mathcal{P}(\mathcal{P}(\omega)),\, \ldots\},
\]
que es fundamental para construir los ordinales de la jerarquía acumulativa $V_\alpha$.
\end{ejemplo}

\begin{observacion}{Reemplazo versus Separación}{rem-sep}
El Esquema de Separación (que presentamos a continuación) es un caso especial del Esquema de Reemplazo: separar los elementos de $A$ que satisfacen $\psi$ equivale a aplicar la función identidad sobre el subconjunto $\{x \in A \mid \psi(x)\}$. En muchas presentaciones modernas solo se incluye Reemplazo, y Separación se deduce de él.
\end{observacion}

% -----------------------------------------------------------
\subsection{Axioma 8: Regularidad (o Fundación)}
% -----------------------------------------------------------

\begin{definicion}{Axioma de Regularidad}{regularidad}
\[
  \forall A\,\bigl(A \neq \emptyset \to \exists x \in A\,(x \cap A = \emptyset)\bigr).
\]
\end{definicion}

\textbf{Significado.} Todo conjunto no vacío tiene un \textbf{elemento $\in$-minimal}: un miembro $x$ que no comparte ningún elemento con $A$. Esto equivale a decir que la relación $\in$ sobre cualquier conjunto es \emph{bien fundada}.

\begin{ejemplo}{Consecuencias del Axioma de Regularidad}{}
\begin{enumerate}[leftmargin=2em, label=(\alph*)]
  \item \textbf{No existen conjuntos que se contengan a sí mismos.}
    Supongamos $A \in A$. Sea $M = \{A\}$. Entonces $M \neq \emptyset$ y su único elemento es $A$. El Axioma de Regularidad exige $A \cap M = \emptyset$. Pero $A \in A$ y $A \in M$, así que $A \in A \cap M$, contradicción.

  \item \textbf{No existe cadena $\in$-descendente infinita} $\cdots \in A_2 \in A_1 \in A_0$.
    En efecto, el conjunto $\{A_i \mid i \in \omega\}$ (si existe como conjunto, lo que garantiza Reemplazo) no tendría elemento $\in$-minimal, contradiciendo el axioma.

  \item \textbf{La jerarquía cumulativa.} El Axioma de Regularidad implica que todo conjunto pertenece a algún estrato $V_\alpha$ de la jerarquía de von Neumann:
    \[
      V_0 = \emptyset, \quad V_{\alpha+1} = \mathcal{P}(V_\alpha), \quad V_\lambda = \bigcup_{\alpha < \lambda} V_\alpha \;(\lambda \text{ límite}).
    \]
    El universo de ZFC es $V = \bigcup_{\alpha \in \mathrm{Ord}} V_\alpha$.
\end{enumerate}
\end{ejemplo}

% -----------------------------------------------------------
\subsection{Axioma 9: Esquema de Separación (o Comprensión Restringida)}
% -----------------------------------------------------------

\begin{definicion}{Esquema de Axiomas de Separación}{separacion}
Para toda fórmula $\varphi(x, p_1, \ldots, p_k)$ de $\mathcal{L}_{\in}$:
\[
  \forall \vec{p}\,\forall A\,\exists B\,\forall x\,\bigl(x \in B \leftrightarrow x \in A \wedge \varphi(x, \vec{p})\bigr).
\]
\end{definicion}

\textbf{Significado.} A partir de cualquier conjunto $A$ y cualquier propiedad $\varphi$, podemos separar los elementos de $A$ que satisfacen $\varphi$. El conjunto resultante se denota $\{x \in A \mid \varphi(x)\}$ y es un subconjunto de $A$.

\begin{ejemplo}{Separación: usos fundamentales}{}
\begin{enumerate}[leftmargin=2em, label=(\alph*)]
  \item \textbf{Intersección binaria.} La intersección $A \cap B = \{x \in A \mid x \in B\}$ existe por Separación (con $\varphi(x) = x \in B$ y parámetro $B$).

  \item \textbf{Complemento relativo.} $A \setminus B = \{x \in A \mid x \notin B\}$ existe por Separación.

  \item \textbf{Conjunto de los números pares.} El conjunto $\{n \in \omega \mid \exists k \in \omega\,(n = k+k)\}$ existe por Separación sobre $\omega$.

  \item \textbf{Derivación del Teorema de Russell.} La fórmula «no restringida» $R = \{x \mid x \notin x\}$ lleva a la paradoja de Russell. Con el Esquema de Separación, solo podemos formar $R_A = \{x \in A \mid x \notin x\}$ para un conjunto $A$ dado. Se puede demostrar que $R_A \notin A$ (pero $R_A$ sí existe como subconjunto de $A$), evitando así la paradoja.
\end{enumerate}
\end{ejemplo}

\begin{importante}
La paradoja de Russell surge al intentar definir $\{x \mid x \notin x\}$ \emph{sin restricción}. La solución de Zermelo fue exigir que la comprensión sea siempre \emph{relativa a un conjunto preexistente}: no se puede construir un conjunto desde la nada usando solo una propiedad; hay que partir de un conjunto ya existente y seleccionar sus elementos que satisfacen dicha propiedad.
\end{importante}

% ===========================================================
\section{Tabla resumen de los axiomas de ZF}
% ===========================================================

\begin{center}
\renewcommand{\arraystretch}{1.4}
\begin{tabular}{|c|l|p{7cm}|}
\hline
\textbf{N.º} & \textbf{Nombre} & \textbf{Garantía informal} \\
\hline
1 & Extensionalidad & Igualdad determinada por los miembros \\
2 & Vacío & Existe $\emptyset$ \\
3 & Pares & Existen $\{a, b\}$ \\
4 & Unión & Existe $\bigcup \mathcal{F}$ \\
5 & Potencia & Existe $\mathcal{P}(A)$ \\
6 & Infinito & Existe $\omega$ \\
7 & Reemplazo (esquema) & La imagen funcional de un conjunto es un conjunto \\
8 & Regularidad & No hay descenso $\in$-infinito \\
9 & Separación (esquema) & Subconjuntos por propiedad son conjuntos \\
\hline
\end{tabular}
\end{center}

% ===========================================================
\section{El Axioma de Elección}
% ===========================================================

\subsection{Formulación del Axioma de Elección}

\begin{definicion}{Axioma de Elección (AC)}{AC}
Para toda familia $\mathcal{F}$ de conjuntos no vacíos y mutuamente disjuntos, existe un conjunto $C$ que contiene exactamente un elemento de cada miembro de $\mathcal{F}$:
\[
  \forall \mathcal{F}\,\Bigl(\bigl(\forall A \in \mathcal{F}\, A \neq \emptyset\bigr) \wedge \bigl(\forall A,B \in \mathcal{F}\,(A \neq B \to A \cap B = \emptyset)\bigr)
  \to \exists C\,\forall A \in \mathcal{F}\,\exists! x\,(x \in C \cap A)\Bigr).
\]
\end{definicion}

Una formulación equivalente —quizás más transparente— es mediante \textbf{funciones de elección}:

\begin{definicion}{Función de elección}{fun-elec}
Sea $\mathcal{F}$ una familia de conjuntos no vacíos. Una \textbf{función de elección} para $\mathcal{F}$ es una función $f: \mathcal{F} \to \bigcup \mathcal{F}$ tal que $f(A) \in A$ para todo $A \in \mathcal{F}$.

\medskip
El \textbf{Axioma de Elección} (AC) afirma que toda familia de conjuntos no vacíos admite una función de elección.
\end{definicion}

\begin{ejemplo}{Elección finita vs. elección infinita}{}
\begin{enumerate}[leftmargin=2em, label=(\alph*)]
  \item \textbf{Caso finito.} Sea $\mathcal{F} = \{\{1,2\},\{3,5\},\{7\}\}$. Podemos elegir explícitamente $f(\{1,2\}) = 1$, $f(\{3,5\}) = 3$, $f(\{7\}) = 7$. No se necesita el AC; basta con una especificación finita.

  \item \textbf{Pares de calcetines (Russell).} Supongamos que tenemos una familia infinita de pares de calcetines $\{P_n\}_{n \in \omega}$, donde cada $P_n = \{c_n^L, c_n^R\}$ y los calcetines del par $n$ son \emph{indistinguibles} entre sí. No hay ninguna regla que nos permita elegir uno de cada par de manera uniforme. El AC postula que, a pesar de ello, existe una función de elección.

  \item \textbf{Subconjuntos de $\mathbb{R}$.} La construcción de un conjunto de Vitale (un conjunto no medible en $\mathbb{R}$) requiere esencialmente el AC: se elige un representante de cada clase de equivalencia de $\mathbb{R}/\mathbb{Q}$.
\end{enumerate}
\end{ejemplo}

\subsection{Independencia del AC}

Un resultado profundo de la teoría de conjuntos es que el Axioma de Elección es \textbf{independiente} de los axiomas ZF:
\begin{itemize}[leftmargin=2em]
  \item \textbf{Gödel (1938):} Si ZF es consistente, entonces ZFC también lo es. (El AC es consistente con ZF.) Gödel lo demostró construyendo el \textit{universo constructible} $L$.
  \item \textbf{Cohen (1963):} Si ZF es consistente, entonces ZF $+ \neg$AC también lo es. (La negación del AC es consistente con ZF.) Cohen lo demostró mediante la técnica de \textit{forcing}.
\end{itemize}

\begin{observacion}{Consecuencias de la independencia}{ind-AC}
La independencia del AC significa que, dentro del sistema ZF, no es posible ni demostrar ni refutar el Axioma de Elección. Esto nos libera de «probar» el AC; simplemente lo adoptamos como un axioma adicional, dando lugar a ZFC. Las matemáticas que requieren el AC (como la teoría de medida, el álgebra lineal sobre espacios de dimensión infinita, o la teoría de cardinales arbitrarios) dependen de esta elección axiomática.
\end{observacion}

% ===========================================================
\section{El Lema de Zorn}
% ===========================================================

\subsection{Órdenes parciales y cadenas}

\begin{definicion}{Orden parcial}{ord-parcial}
Un \textbf{orden parcial} es un par $(P, \leq)$ donde $P$ es un conjunto y $\leq$ es una relación binaria sobre $P$ que satisface:
\begin{enumerate}[leftmargin=2em]
  \item \textbf{Reflexividad:} $\forall x \in P\,(x \leq x)$.
  \item \textbf{Antisimetría:} $\forall x, y \in P\,(x \leq y \wedge y \leq x \to x = y)$.
  \item \textbf{Transitividad:} $\forall x, y, z \in P\,(x \leq y \wedge y \leq z \to x \leq z)$.
\end{enumerate}
\end{definicion}

\begin{definicion}{Cadena y cota superior}{cadena}
Sea $(P, \leq)$ un orden parcial.
\begin{itemize}[leftmargin=2em]
  \item Una \textbf{cadena} (o subconjunto totalmente ordenado) de $P$ es un subconjunto $C \subseteq P$ tal que $\forall x, y \in C\,(x \leq y \vee y \leq x)$.
  \item Una \textbf{cota superior} de un subconjunto $S \subseteq P$ es un elemento $u \in P$ tal que $\forall x \in S\,(x \leq u)$.
  \item Un \textbf{elemento maximal} de $P$ es un elemento $m \in P$ tal que $\forall x \in P\,(m \leq x \to x = m)$.
\end{itemize}
\end{definicion}

\subsection{Enunciado y demostración del Lema de Zorn}

\begin{lema}{Lema de Zorn}{zorn}
Si $(P, \leq)$ es un orden parcial no vacío en el que toda cadena tiene una cota superior en $P$, entonces $P$ tiene al menos un elemento maximal.
\end{lema}

\begin{proof}
Demostraremos que el Lema de Zorn se deduce del Axioma de Elección. La demostración se realiza por el absurdo.

Supongamos que $(P, \leq)$ satisface las hipótesis pero que $P$ no tiene ningún elemento maximal. Entonces para cada cadena $C \subseteq P$ existe una cota superior $u_C \in P$ con $u_C \notin C$ (pues de lo contrario $u_C$ sería un elemento maximal, contradiciendo nuestra suposición, ya que una cota superior $u$ de $C$ con $u \in C$ satisfaría que nada en $C$ supera estrictamente a $u$, pero entonces $u$ es maximal en $P$ salvo que exista $v > u$ fuera de $C$; el argumento preciso requiere la siguiente construcción transfinita).

Sea $\mathcal{C}$ la colección de todas las cadenas de $P$. Por el AC, existe una función de elección $f: \mathcal{C} \to P$ tal que $f(C) \in P$ es una cota estricta de $C$ (es decir, $f(C) > c$ para todo $c \in C$; esta elección es posible porque $P$ no tiene maximales).

Usando esta función $f$ y el Principio de Recursión Transfinita (que se demuestra en ZF), construimos una secuencia transfinita estrictamente creciente $\{a_\alpha\}_{\alpha \in \mathrm{Ord}}$ en $P$: tomamos $a_0 \in P$ cualquiera, $a_{\alpha+1} = f(\{a_\beta \mid \beta \leq \alpha\})$, y para $\lambda$ límite $a_\lambda = f(\{a_\beta \mid \beta < \lambda\})$.

Esta secuencia es estrictamente creciente en un orden parcial, lo que implica que todos los $a_\alpha$ son distintos, formando una clase propia. Pero $P$ es un conjunto (no una clase propia), contradicción. Por lo tanto, la suposición de que $P$ no tiene maximales es falsa. $\square$
\end{proof}

\begin{ejemplo}{Aplicaciones del Lema de Zorn}{}
\begin{enumerate}[leftmargin=2em, label=(\alph*)]
  \item \textbf{Bases de espacios vectoriales.} Sea $V$ un espacio vectorial sobre un cuerpo $K$. Sea $P$ el conjunto de todos los conjuntos linealmente independientes de $V$, ordenado por inclusión. Toda cadena tiene cota superior (la unión), y $P \neq \emptyset$ (contiene conjuntos de un solo vector no nulo). Por el Lema de Zorn, $P$ tiene un elemento maximal: un conjunto linealmente independiente maximal, es decir, una \textbf{base de Hamel} de $V$.

  \item \textbf{Ideales maximales en anillos.} Sea $R$ un anillo con unidad y $P$ el conjunto de todos los ideales propios de $R$, ordenado por inclusión. Si $R \neq \{0\}$, entonces $P \neq \emptyset$ y toda cadena de ideales propios tiene como cota su unión (que sigue siendo un ideal propio). Por el Lema de Zorn, todo anillo con unidad tiene al menos un ideal maximal.

  \item \textbf{Extensión de homomorfismos.} En álgebra homológica y análisis funcional, el Lema de Zorn se usa para extender funcionales lineales acotados (Lema de Hahn–Banach en su versión algebraica) y homomorfismos de grupos, entre otros resultados.
\end{enumerate}
\end{ejemplo}

% ===========================================================
\section{El Teorema del Buen Orden}
% ===========================================================

\subsection{Órdenes totales y buenos órdenes}

\begin{definicion}{Orden total}{ord-total}
Un \textbf{orden total} (o \textbf{lineal}) sobre $P$ es un orden parcial $(P, \leq)$ en el que $\forall x, y \in P\,(x \leq y \vee y \leq x)$, es decir, todo par de elementos es comparable.
\end{definicion}

\begin{definicion}{Buen orden}{buen-orden}
Un \textbf{buen orden} sobre $P$ es un orden total $(P, \leq)$ tal que todo subconjunto no vacío de $P$ tiene un \textbf{mínimo} (elemento menor o igual que todos los demás):
\[
  \forall S \subseteq P\,\bigl(S \neq \emptyset \to \exists m \in S\,\forall s \in S\,(m \leq s)\bigr).
\]
El par $(P, \leq)$ se llama entonces un \textbf{conjunto bien ordenado}.
\end{definicion}

\begin{ejemplo}{Ejemplos de buenos órdenes}{}
\begin{enumerate}[leftmargin=2em, label=(\alph*)]
  \item $(\mathbb{N}, \leq)$ con el orden usual es un buen orden: todo subconjunto no vacío de $\mathbb{N}$ tiene un mínimo.

  \item $(\mathbb{Z}, \leq)$ con el orden usual \emph{no} es un buen orden: el subconjunto $\mathbb{Z}$ mismo no tiene mínimo.

  \item $(\mathbb{R}_{>0}, \leq)$ con el orden usual \emph{no} es un buen orden: el intervalo $(0, 1)$ no tiene mínimo.

  \item El conjunto $\omega + \omega = \{0, 1, 2, \ldots, \omega, \omega+1, \omega+2, \ldots\}$ con el orden ordinal es un buen orden (ver Clase 3).

  \item Todo subconjunto de un conjunto bien ordenado, con el orden inducido, es bien ordenado.
\end{enumerate}
\end{ejemplo}

\subsection{Enunciado y demostración del Teorema del Buen Orden}

\begin{teorema}{Teorema del Buen Orden (Zermelo, 1904)}{buen-orden}
Todo conjunto puede ser bien ordenado: para todo conjunto $A$, existe una relación $\leq$ sobre $A$ que es un buen orden.
\end{teorema}

\begin{proof}
Utilizamos el Axioma de Elección. Sea $A$ un conjunto. El AC proporciona una función de elección $f: \mathcal{P}(A) \setminus \{\emptyset\} \to A$ tal que $f(S) \in S$ para todo $S \neq \emptyset$.

Definimos una secuencia transfinita por recursión sobre los ordinales:
\begin{itemize}[leftmargin=2em]
  \item $a_0 = f(A)$.
  \item $a_{\alpha+1} = f(A \setminus \{a_\beta \mid \beta \leq \alpha\})$, si $A \setminus \{a_\beta \mid \beta \leq \alpha\} \neq \emptyset$.
  \item Para $\lambda$ límite: $a_\lambda = f(A \setminus \{a_\beta \mid \beta < \lambda\})$, si este conjunto es no vacío.
  \item El proceso se detiene cuando $A \setminus \{a_\beta \mid \beta < \alpha\} = \emptyset$.
\end{itemize}

Puesto que $A$ es un conjunto (no una clase propia), existe un ordinal $\theta$ en el que el proceso termina. Entonces $\{a_\beta \mid \beta < \theta\}$ es una enumeración biyectiva de todos los elementos de $A$, y el orden $a_\alpha \leq a_\beta \iff \alpha \leq \beta$ define un buen orden sobre $A$. $\square$
\end{proof}

\begin{ejemplo}{Buen orden sobre $\mathbb{R}$}{}
El Teorema del Buen Orden garantiza que $\mathbb{R}$ puede ser bien ordenado. Sin embargo, dicho buen orden \emph{no} puede ser construido explícitamente en ZFC (independientemente de si se cree o no en el AC). Ninguna fórmula concreta de $\mathcal{L}_{\in}$ define un buen orden de $\mathbb{R}$ en todos los modelos de ZFC. Esta es una de las manifestaciones de que el AC —y los objetos que produce— en general no son constructivos.
\end{ejemplo}

% ===========================================================
\section{Equivalencias clásicas del Axioma de Elección}
% ===========================================================

Uno de los resultados más notables de la teoría de conjuntos del siglo XX es que muchas proposiciones que a primera vista parecen independientes entre sí son de hecho equivalentes en ZF. El siguiente teorema recoge las más importantes.

\begin{teorema}{Equivalencias del Axioma de Elección}{equiv-AC}
En el sistema ZF, las siguientes afirmaciones son equivalentes:
\begin{enumerate}[leftmargin=2em, label=\textbf{(\roman*)}]
  \item \textbf{Axioma de Elección (AC):} Toda familia de conjuntos no vacíos admite una función de elección.
  \item \textbf{Lema de Zorn:} Todo orden parcial no vacío en el que toda cadena tiene cota superior posee un elemento maximal.
  \item \textbf{Teorema del Buen Orden:} Todo conjunto puede ser bien ordenado.
  \item \textbf{Principio de la base de Hamel:} Todo espacio vectorial (sobre cualquier cuerpo) tiene una base.
  \item \textbf{Teorema de Tychonoff:} El producto cartesiano de una familia arbitraria de espacios topológicos compactos es compacto.
  \item \textbf{Principio de la cardinalidad:} Para cualesquiera dos conjuntos $A$ y $B$, se cumple $|A| \leq |B|$ o $|B| \leq |A|$.
  \item \textbf{Principio de la unión numerable:} La unión de una colección numerable de conjuntos numerables es numerable.
\end{enumerate}
\end{teorema}

\begin{proof}[Esquema de la demostración]
Indicamos las implicaciones principales; las demostraciones completas se encuentran en \cite{jech2003, kunen2011}.

\medskip
\textbf{(i) $\Rightarrow$ (ii) [AC $\Rightarrow$ Zorn]:} Ya demostrado en la sección anterior (Lema de Zorn).

\medskip
\textbf{(ii) $\Rightarrow$ (iii) [Zorn $\Rightarrow$ Buen Orden]:} Sea $A$ un conjunto. Sea $P$ el conjunto de todos los pares $(B, \leq_B)$ donde $B \subseteq A$ y $\leq_B$ es un buen orden de $B$. Se ordena $P$ por: $(B, \leq_B) \preceq (C, \leq_C)$ si $B \subseteq C$ y $\leq_C$ extiende $\leq_B$ (y $B$ es un segmento inicial de $C$). Toda cadena en $P$ tiene cota superior (la unión). Por el Lema de Zorn, existe un elemento maximal $(M, \leq_M)$. Se muestra que $M = A$ (de lo contrario, se podría extender a un conjunto estrictamente mayor, contradiciendo la maximalidad).

\medskip
\textbf{(iii) $\Rightarrow$ (i) [Buen Orden $\Rightarrow$ AC]:} Sea $\mathcal{F}$ una familia de conjuntos no vacíos. Buen Ordenemos $\bigcup \mathcal{F}$ con algún buen orden $\leq$. Para cada $A \in \mathcal{F}$, definimos $f(A) = \min_{\leq}(A)$ (el mínimo de $A$ respecto a $\leq$). Esta $f$ es una función de elección para $\mathcal{F}$.
\end{proof}

\begin{ejemplo}{Equivalencias en la práctica}{}
\begin{enumerate}[leftmargin=2em, label=(\alph*)]
  \item \textbf{Sin el AC, hay espacios vectoriales sin base.} En ZF + $\neg$AC (por ejemplo, en el modelo de Solovay donde todo subconjunto de $\mathbb{R}$ es Lebesgue medible), existen espacios vectoriales sobre $\mathbb{Q}$ que no admiten base. La afirmación «todo espacio vectorial tiene base» equivale al AC.

  \item \textbf{Tychonoff y la compacidad del cubo de Hilbert.} El producto $[0,1]^{\mathbb{R}}$ (con la topología producto) es compacto si y solo si aceptamos el AC. Sin el AC, este resultado puede fallar.

  \item \textbf{La cardinalidad requiere el AC.} En ZF sin el AC, es posible tener dos conjuntos $A$ y $B$ tales que $|A| \not\leq |B|$ y $|B| \not\leq |A|$ simultáneamente; es decir, los cardinales pueden ser incomparables. Con el AC, todos los cardinales son comparables (son ordinales iniciales).
\end{enumerate}
\end{ejemplo}

% ===========================================================
\section{Formas débiles del Axioma de Elección}
% ===========================================================

En la literatura se estudian varias versiones debilitadas del AC, estrictamente más débiles que él pero suficientes para muchos propósitos:

\begin{definicion}{Axioma de Elección Numerable (CC)}{CC}
Toda familia \emph{numerable} de conjuntos no vacíos admite una función de elección.
\end{definicion}

\begin{definicion}{Axioma de Elección Dependiente (DC)}{DC}
Si $R$ es una relación binaria sobre un conjunto $A$ no vacío tal que para todo $x \in A$ existe $y \in A$ con $x \mathrel{R} y$, entonces existe una secuencia $\langle a_n \rangle_{n \in \omega}$ en $A$ con $a_n \mathrel{R} a_{n+1}$ para todo $n$.
\end{definicion}

La jerarquía de implicaciones (en ZF) es:
\[
  \mathrm{AC} \Rightarrow \mathrm{DC} \Rightarrow \mathrm{CC}.
\]
Y ninguna implicación inversa es demostrable en ZF. Las formas débiles son suficientes para gran parte del análisis real estándar: la completitud de los reales, la existencia de bases ortonormales en espacios de Hilbert separables, y los teoremas de categoría de Baire, entre otros.

% ===========================================================
\section{ZFC como fundamento estándar de las matemáticas}
% ===========================================================

\subsection{¿Por qué ZFC?}

ZFC se ha establecido como el sistema axiomático de referencia por varias razones:

\begin{enumerate}[leftmargin=2em, label=\textbf{\arabic*.}]
  \item \textbf{Suficiencia.} ZFC es suficientemente expresivo para formalizar prácticamente toda la matemática estándar: los números naturales, enteros, racionales, reales y complejos; el análisis real y complejo; la topología general; el álgebra abstracta; la teoría de medida; y la geometría diferencial.

  \item \textbf{Consistencia relativa.} Si ZFC fuera inconsistente, lo sabríamos: cualquier contradicción derivable en ZFC revelaría una falla en nuestra comprensión matemática. El consenso actual de la comunidad matemática es que ZFC es consistente (aunque, por el Segundo Teorema de Incompletitud de Gödel, ZFC no puede demostrar su propia consistencia).

  \item \textbf{Economía de primitivos.} ZFC se expresa con un único símbolo primitivo ($\in$) y un conjunto manejable de axiomas, todos en lógica de primer orden.

  \item \textbf{Compatibilidad con el AC.} El Axioma de Elección, aunque filosóficamente debatido, es necesario para resultados cotidianos como «todo espacio vectorial tiene base», y su inclusión en ZFC lo convierte en el sistema adecuado para matemáticas sin restricciones.
\end{enumerate}

\subsection{Limitaciones de ZFC}

A pesar de su predominio, ZFC tiene limitaciones importantes:

\begin{itemize}[leftmargin=2em]
  \item \textbf{Incompletitud (Gödel, 1931).} Por el Primer Teorema de Incompletitud, si ZFC es consistente, existen enunciados de $\mathcal{L}_{\in}$ que no son ni demostrables ni refutables en ZFC. El más célebre es la \textbf{Hipótesis del Continuo} (CH): $2^{\aleph_0} = \aleph_1$.

  \item \textbf{Necesidad de grandes cardinales.} Muchos enunciados de importancia matemática (como la existencia de cardinales de medida o el axioma de determinancy) requieren axiomas adicionales más allá de ZFC.

  \item \textbf{Alternativas.} Existen otros sistemas fundacionales: la teoría de tipos de Russell–Whitehead, la Teoría de Conjuntos de Von Neumann–Bernays–Gödel (NBG, que permite hablar de clases propias), la Teoría de Categorías (con sus propios fundamentos) o la Homotopy Type Theory (HoTT).
\end{itemize}

\begin{observacion}{El papel de ZFC en este curso}{rol-ZFC}
En el estudio de los números transfinitos, ZFC desempeña un papel central: los números ordinales y cardinales se construyen como conjuntos especiales en el universo de ZFC, y el Axioma de Elección es indispensable para demostrar que los cardinales son comparables, que toda cardinalidad infinita es un $\aleph$, y que el Teorema de Hartogs tiene sus consecuencias esperadas. Los próximos capítulos desarrollarán en detalle esta construcción.
\end{observacion}

% ===========================================================
\section{Problemas y tareas para el estudiante}
% ===========================================================

Los siguientes ejercicios están diseñados para afianzar la comprensión de los axiomas de ZFC y sus consecuencias. Se recomienda justificar cada respuesta con rigor, citando los axiomas utilizados.

\begin{enumerate}[leftmargin=2em, label=\textbf{P\arabic*.}, itemsep=0.8em]

  \item \textbf{(Extensionalidad y unicidad.)}
  \begin{enumerate}[label=(\alph*)]
    \item Demuestra que $\{1, 2, 3\} = \{3, 1, 2\} = \{1, 1, 2, 3\}$ usando únicamente el Axioma de Extensionalidad.
    \item Enuncia y demuestra que el conjunto vacío es único en ZF.
    \item ¿Es cierto que $\emptyset = \{\emptyset\}$? Justifica tu respuesta.
  \end{enumerate}

  \item \textbf{(Pares y unión.)}
  \begin{enumerate}[label=(\alph*)]
    \item Construye el conjunto $\{a, b, c\}$ (para conjuntos $a, b, c$ dados) a partir de los Axiomas de Pares y Unión.
    \item Demuestra que la intersección $A \cap B$ existe para cualesquiera conjuntos $A$ y $B$, usando los Axiomas de Separación (y Pares si es necesario).
    \item Muestra que $A \cup B = \bigcup\{A, B\}$ y deduce que la unión binaria existe usando los Axiomas de Pares y Unión.
  \end{enumerate}

  \item \textbf{(Potencia e infinito.)}
  \begin{enumerate}[label=(\alph*)]
    \item Calcula $\mathcal{P}(\mathcal{P}(\emptyset))$ y $\mathcal{P}(\mathcal{P}(\mathcal{P}(\emptyset)))$.
    \item Describe los primeros cinco elementos del conjunto $\omega$ en la construcción de von Neumann y verifica que cada uno satisface la condición del Axioma del Infinito.
    \item ¿Por qué la existencia de $\omega$ no se sigue solo de los axiomas de Pares, Unión y Potencia?
  \end{enumerate}

  \item \textbf{(Regularidad.)}
  \begin{enumerate}[label=(\alph*)]
    \item Demuestra que no existe ningún conjunto $A$ tal que $A \in A$, usando el Axioma de Regularidad.
    \item Demuestra que no existen conjuntos $A$ y $B$ con $A \in B$ y $B \in A$ simultáneamente.
    \item Más generalmente, demuestra que no existe ninguna cadena $\in$-cíclica: $A_1 \in A_2 \in \cdots \in A_n \in A_1$.
  \end{enumerate}

  \item \textbf{(Paradoja de Russell y Separación.)}
  \begin{enumerate}[label=(\alph*)]
    \item Sea $A$ un conjunto cualquiera. Define $R_A = \{x \in A \mid x \notin x\}$. Demuestra que $R_A \notin A$.
    \item Concluye que no existe el «conjunto de todos los conjuntos» en ZFC.
    \item ¿Qué sucedería si permitiéramos el «Principio de Comprensión Irrestricta» $\{x \mid \varphi(x)\}$ sin restricción? Explica la paradoja que surgiría.
  \end{enumerate}

  \item \textbf{(Reemplazo.)}
  \begin{enumerate}[label=(\alph*)]
    \item Sea $f: \omega \to \omega$ la función $f(n) = n^2$. Usando el Esquema de Reemplazo, argumenta que el conjunto $\{0, 1, 4, 9, 16, \ldots\}$ existe en ZFC.
    \item Demuestra que el Esquema de Separación es consecuencia del Esquema de Reemplazo (considerando la función identidad restringida al subconjunto que satisface la propiedad).
  \end{enumerate}

  \item \textbf{(Axioma de Elección.)}
  \begin{enumerate}[label=(\alph*)]
    \item Escribe explícitamente una función de elección para la familia $\mathcal{F} = \{\{2k, 2k+1\} \mid k \in \omega\}$ (pares de naturales consecutivos). ¿Es necesario el AC en este caso?
    \item Explica por qué la siguiente afirmación requiere el AC: «El producto cartesiano de una familia infinita de conjuntos no vacíos es no vacío.»
    \item Sea $V$ un espacio vectorial sobre $\mathbb{Q}$ con $\dim_{\mathbb{Q}}(V) = \aleph_0$. ¿Garantiza ZFC la existencia de una base? ¿Y ZF sin el AC?
  \end{enumerate}

  \item \textbf{(Lema de Zorn.)}
  \begin{enumerate}[label=(\alph*)]
    \item Sea $R$ un anillo conmutativo unitario con $R \neq \{0\}$. Usando el Lema de Zorn, demuestra que $R$ tiene al menos un ideal maximal.
    \item Sea $G$ un grupo abeliano. Demuestra, usando el Lema de Zorn, que todo subgrupo de $G$ está contenido en un subgrupo maximal propio (si existe alguno).
    \item Enuncia el Teorema de Hahn–Banach y explica en qué paso de su demostración se usa el Lema de Zorn.
  \end{enumerate}

  \item \textbf{(Teorema del Buen Orden y cardinalidad.)}
  \begin{enumerate}[label=(\alph*)]
    \item Demuestra que el Teorema del Buen Orden implica el Axioma de Elección en ZF (es decir, TBO $\Rightarrow$ AC).
    \item El Teorema del Buen Orden garantiza que $\mathbb{R}$ puede ser bien ordenado. Sin embargo, tal bien orden no puede ser construido explícitamente. Investiga y resume brevemente el argumento que muestra que ninguna fórmula concreta de ZFC define un buen orden de $\mathbb{R}$ en todos sus modelos.
    \item Demuestra que todo conjunto bien ordenado es isomorfo a un único ordinal. (Puede usar resultados de la Clase 3.)
  \end{enumerate}

  \item \textbf{(Formas débiles del AC y análisis.)}
  \begin{enumerate}[label=(\alph*)]
    \item El Axioma de Elección Dependiente (DC) implica que $\mathbb{R}$ es Arquimediano (todo real positivo es superado por algún natural). Explica la conexión con DC.
    \item Investiga y enuncia el \emph{Principio de Elección Numerable} (CC) y da un ejemplo de un teorema del análisis real que requiere CC pero no el AC completo.
    \item Demuestra (usando únicamente ZF + DC) que toda sucesión de Cauchy de reales converge.
  \end{enumerate}

  \item \textbf{(Proyecto: independencia del AC.)}
  Investiga y redacta un resumen de no más de dos páginas sobre la demostración de la independencia del Axioma de Elección respecto a ZF, cubriendo:
  \begin{itemize}[leftmargin=2em]
    \item La construcción del universo constructible $L$ de Gödel y por qué $L \models \mathrm{AC}$.
    \item La idea intuitiva del método de \emph{forcing} de Cohen y por qué permite construir un modelo de ZF $+ \neg$AC.
    \item Las consecuencias filosóficas de esta independencia para el fundamento de las matemáticas.
  \end{itemize}

\end{enumerate}

% ===========================================================
\section*{Referencias bibliográficas de esta clase}
% ===========================================================

Los temas tratados en esta clase se desarrollan con mayor detalle en las siguientes fuentes, que constituyen la bibliografía principal del curso:

\begin{itemize}[leftmargin=2em]
  \item \textbf{Fundamentos y axiomática de ZFC:} \cite{jech2003} es la referencia enciclopédica estándar; \cite{kunen2011} ofrece una presentación rigurosa con énfasis en independencia y forcing; \cite{halmos1960} es un texto introductorio clásico, conciso y elegante.

  \item \textbf{Axioma de Elección y equivalencias:} \cite{herrlich2006} es una monografía dedicada íntegramente al Axioma de Elección y sus consecuencias; \cite{moore1982} ofrece una historia detallada del AC desde Cantor hasta Cohen.

  \item \textbf{Lema de Zorn y aplicaciones algebraicas:} \cite{lang2002} contiene aplicaciones del Lema de Zorn en álgebra; \cite{rudin1991} ilustra su uso en análisis funcional (Teorema de Hahn–Banach).

  \item \textbf{Independencia del AC:} \cite{cohen1966} es el trabajo original de Cohen sobre forcing; \cite{kanamori2003} cubre la historia y los grandes cardinales en el contexto de ZFC.
\end{itemize}

% ===========================================================
\chapter{Clase 3: Relaciones, Órdenes y Conjuntos Bien Ordenados}
% ===========================================================

Las estructuras de orden constituyen el andamiaje conceptual sobre el que se construye toda la teoría transfinita. En esta clase estudiamos, con pleno rigor, las relaciones binarias y sus propiedades fundamentales, los distintos grados de estructura de orden —parcial, total y bien-orden—, los segmentos iniciales y los isomorfismos entre conjuntos ordenados, y culminamos con una introducción al poderoso \textbf{Principio de Inducción Transfinita}. Este último es la herramienta principal que nos permitirá, en la clase siguiente, definir los números ordinales y operar con ellos.

Todo el desarrollo se realiza dentro del sistema axiomático ZFC cuya axiomática se estudió en la Clase~2. En particular, los conceptos de «relación» y «función» son casos especiales de conjuntos, y sus propiedades se demuestran a partir de los axiomas \cite{jech2003, kunen2011, enderton1977}.

% ===========================================================
\section{Relaciones binarias}
% ===========================================================

\subsection{Definición y ejemplos}

\begin{definicion}{Relación binaria}{relbin}
Sean $A$ y $B$ conjuntos. Una \textbf{relación binaria} de $A$ en $B$ es un subconjunto
\[
  R \subseteq A \times B.
\]
Se escribe $a \mathrel{R} b$ (o también $R(a,b)$) en lugar de $(a,b) \in R$. Cuando $A = B$ se dice que $R$ es una relación binaria \emph{sobre} $A$ (o \emph{en} $A$).
\end{definicion}

\begin{observacion}{Las relaciones son conjuntos}{relconj}
En ZFC, el producto cartesiano $A \times B$ se define como
\[
  A \times B \;=\; \bigl\{z \mid \exists a \in A\,\exists b \in B\,(z = (a,b))\bigr\},
\]
donde el par ordenado se codifica como $(a,b) = \{\{a\}, \{a,b\}\}$ (par de Kuratowski). Por el Axioma de Potencia y el Esquema de Separación, $A \times B$ y toda relación $R \subseteq A \times B$ son conjuntos.
\end{observacion}

\begin{ejemplo}{Relaciones binarias elementales}{}
\begin{enumerate}[leftmargin=2em, label=(\alph*)]
  \item La relación de \textbf{orden usual} en $\mathbb{N}$:
  \[
    {\leq} \;=\; \{(m,n) \in \mathbb{N}\times\mathbb{N} \mid m \leq n\}.
  \]
  Por ejemplo, $(2,5) \in {\leq}$ ya que $2 \leq 5$; en cambio $(7,3) \notin {\leq}$.

  \item La \textbf{divisibilidad} en $\mathbb{Z}^+$:
  \[
    a \mid b \iff \exists k \in \mathbb{Z}^+\,(b = ka).
  \]
  Así $(3,12) \in {\mid}$ pero $(5,7) \notin {\mid}$.

  \item La \textbf{inclusión} en $\mathcal{P}(X)$:
  \[
    A \subseteq B \iff \forall x\,(x \in A \to x \in B).
  \]
  Esta es una de las relaciones de orden más importantes en teoría de conjuntos.

  \item La relación de \textbf{equicardinalidad}: $A \sim B$ si existe una biyección $f: A \to B$. Veremos en clases posteriores que no siempre $A \sim B$ implica $A = B$.
\end{enumerate}
\end{ejemplo}

\subsection{Propiedades fundamentales de las relaciones}

Sea $R$ una relación binaria sobre un conjunto $A$. Definimos las siguientes propiedades.

\begin{definicion}{Propiedades de las relaciones}{proprel}
\begin{enumerate}[leftmargin=2em, label=\textbf{(\arabic*)}]
  \item \textbf{Reflexividad:} $\forall a \in A,\; a \mathrel{R} a$.
  \item \textbf{Irreflexividad:} $\forall a \in A,\; \neg(a \mathrel{R} a)$.
  \item \textbf{Simetría:} $\forall a,b \in A,\; a \mathrel{R} b \to b \mathrel{R} a$.
  \item \textbf{Antisimetría:} $\forall a,b \in A,\; (a \mathrel{R} b \wedge b \mathrel{R} a) \to a = b$.
  \item \textbf{Asimetría:} $\forall a,b \in A,\; a \mathrel{R} b \to \neg(b \mathrel{R} a)$.
  \item \textbf{Transitividad:} $\forall a,b,c \in A,\; (a \mathrel{R} b \wedge b \mathrel{R} c) \to a \mathrel{R} c$.
  \item \textbf{Totalidad (o conexidad):} $\forall a,b \in A,\; a \mathrel{R} b \vee b \mathrel{R} a \vee a = b$.
  \item \textbf{Tricotomía:} $\forall a,b \in A,\; a \mathrel{R} b \vee a = b \vee b \mathrel{R} a$, y estas tres posibilidades son mutuamente excluyentes.
  \item \textbf{Buena fundamentación:} Todo subconjunto no vacío $S \subseteq A$ tiene un elemento $R$-minimal, es decir, existe $m \in S$ tal que $\neg\exists x \in S\,(x \mathrel{R} m)$.
\end{enumerate}
\end{definicion}

\begin{observacion}{Asimetría implica irreflexividad}{asimirref}
Si $R$ es asimétrica, entonces es irreflexiva. En efecto, si existiera $a \in A$ con $a \mathrel{R} a$, la asimetría daría $\neg(a \mathrel{R} a)$, contradicción.
\end{observacion}

\begin{ejemplo}{Comparación de propiedades para relaciones clásicas}{}
En la tabla siguiente, $\checkmark$ indica que la propiedad se cumple y $\times$ que no.
\begin{center}
\renewcommand{\arraystretch}{1.3}
\begin{tabular}{lcccccc}
\hline
\textbf{Relación} & \textbf{Refl.} & \textbf{Simét.} & \textbf{Antisimét.} & \textbf{Transit.} & \textbf{Total} & \textbf{B.F.} \\
\hline
$\leq$ en $\mathbb{N}$              & $\checkmark$ & $\times$ & $\checkmark$ & $\checkmark$ & $\checkmark$ & $\checkmark$ \\
$<$ en $\mathbb{N}$                 & $\times$     & $\times$ & $\checkmark$ & $\checkmark$ & $\checkmark$ & $\checkmark$ \\
$\leq$ en $\mathbb{R}$              & $\checkmark$ & $\times$ & $\checkmark$ & $\checkmark$ & $\checkmark$ & $\times$ \\
$|$ en $\mathbb{Z}^+$               & $\checkmark$ & $\times$ & $\checkmark$ & $\checkmark$ & $\times$     & $\checkmark$ \\
$\subseteq$ en $\mathcal{P}(X)$     & $\checkmark$ & $\times$ & $\checkmark$ & $\checkmark$ & $\times$     & $\times$ \\
$\sim$ (equicardinalidad)           & $\checkmark$ & $\checkmark$ & $\times$  & $\checkmark$ & $\times$     & $\times$ \\
\hline
\end{tabular}
\end{center}
\textbf{Nota:} Las filas con «buena fundamentación» $\times$ admiten subconjuntos no vacíos sin elemento minimal (por ejemplo, $(0,1) \subseteq \mathbb{R}$ con el orden usual). La relación $\sim$ en $\mathcal{P}(X)$ cuando $X$ es infinito es un ejemplo de relación de equivalencia que no es de orden.
\end{ejemplo}

\subsection{Relaciones de equivalencia y particiones}

Por contraste con las relaciones de orden, recordamos que una relación de equivalencia satisface reflexividad, simetría y transitividad, y genera una partición del conjunto base.

\begin{definicion}{Relación de equivalencia}{releq}
Una relación $R$ sobre $A$ es una \textbf{relación de equivalencia} si es reflexiva, simétrica y transitiva. La \textbf{clase de equivalencia} de $a \in A$ es $[a]_R = \{b \in A \mid a \mathrel{R} b\}$.
\end{definicion}

\begin{proposicion}{Equivalencia y partición}{eqpart}
Si $R$ es una relación de equivalencia sobre $A$, entonces las clases de equivalencia forman una \textbf{partición} de $A$: son subconjuntos no vacíos, mutuamente disjuntos, cuya unión es $A$.
\end{proposicion}

\begin{proof}
\textbf{No vacías:} Por reflexividad, $a \in [a]_R$ para todo $a$. \quad \textbf{Disjuntas:} Si $[a]_R \cap [b]_R \neq \emptyset$, sea $c$ en la intersección. Entonces $c \mathrel{R} a$ y $c \mathrel{R} b$; por simetría y transitividad, $a \mathrel{R} b$, de donde $[a]_R = [b]_R$. \quad \textbf{Cubren $A$:} Claramente $\bigcup_{a \in A} [a]_R = A$.
\end{proof}

% ===========================================================
\section{Órdenes parciales y totales}
% ===========================================================

\subsection{Preórdenes y órdenes parciales}

\begin{definicion}{Preorden}{preorden}
Una relación $\leq$ sobre $A$ es un \textbf{preorden} (o \textbf{cuasiorden}) si es reflexiva y transitiva. El par $(A, \leq)$ se llama un \textbf{conjunto proordenado}.
\end{definicion}

\begin{definicion}{Orden parcial}{ordparcial}
Una relación $\leq$ sobre $A$ es un \textbf{orden parcial} (o \textbf{orden parcial no estricto}) si es:
\begin{enumerate}[leftmargin=2em, label=(\roman*)]
  \item \textbf{Reflexiva:} $\forall a,\; a \leq a$.
  \item \textbf{Antisimétrica:} $\forall a,b,\; (a \leq b \wedge b \leq a) \to a = b$.
  \item \textbf{Transitiva:} $\forall a,b,c,\; (a \leq b \wedge b \leq c) \to a \leq c$.
\end{enumerate}
El par $(A, \leq)$ se llama un \textbf{conjunto parcialmente ordenado} o \textbf{poset}.

La relación \textbf{estricta asociada} es $a < b \iff a \leq b \wedge a \neq b$.
\end{definicion}

\begin{definicion}{Orden parcial estricto}{ordparcest}
Una relación $<$ sobre $A$ es un \textbf{orden parcial estricto} si es:
\begin{enumerate}[leftmargin=2em, label=(\roman*)]
  \item \textbf{Irreflexiva:} $\forall a,\; \neg(a < a)$.
  \item \textbf{Transitiva:} $\forall a,b,c,\; (a < b \wedge b < c) \to a < c$.
\end{enumerate}
Nótese que la transitividad y la irreflexividad implican la asimetría.
\end{definicion}

\begin{proposicion}{Correspondencia entre órdenes estrictos y no estrictos}{estrictonoestricto}
Existe una correspondencia biunívoca entre órdenes parciales no estrictos y órdenes parciales estrictos sobre un mismo conjunto $A$:
\[
  a \leq b \;\iff\; a < b \vee a = b.
\]
Si $\leq$ es un orden parcial no estricto, la relación $<$ definida así es un orden parcial estricto; y recíprocamente.
\end{proposicion}

\begin{proof}
($\Rightarrow$) Sean $a < b$ y $b < c$. Entonces $a \leq b$ y $b \leq c$, luego $a \leq c$ por transitividad. Además, $a \neq b$ y $b \neq c$; si fuera $a = c$, entonces $a \leq b \leq a$ y por antisimetría $a = b$, contradicción. Luego $a < c$. La irreflexividad de $<$ es inmediata: $a \leq a$ junto con la antisimetría solo da $a = a$, no $a < a$.

($\Leftarrow$) Sea $\leq$ la relación definida por $a \leq b \iff a < b \vee a = b$. Reflexividad: $a = a$, luego $a \leq a$. Antisimetría: si $a \leq b$ y $b \leq a$, entonces o bien $a = b$ (y estamos listos) o bien $a < b$ y ($b < a$ o $b = a$). Si $a < b$ y $b < a$, por transitividad $a < a$, contradicción con la irreflexividad. Luego $a = b$. Transitividad: análoga.
\end{proof}

\begin{ejemplo}{Posets fundamentales}{}
\begin{enumerate}[leftmargin=2em, label=(\alph*)]
  \item $(\mathbb{N}, \leq)$ con el orden usual. Es un orden parcial (de hecho, total). Los elementos $2, 3$ son comparables: $2 < 3$.

  \item $(\mathcal{P}(\{1,2,3\}), \subseteq)$. Es un orden parcial pero no total: $\{1,2\}$ y $\{1,3\}$ son incomparables (ninguno contiene al otro). El diagrama de Hasse de este poset tiene la forma de un cubo.

  \item $(\mathbb{Z}^+, \mid)$ con la divisibilidad. Es un orden parcial. Los elementos $4$ y $6$ son incomparables ($4 \nmid 6$ y $6 \nmid 4$), pero $2 \mid 4$ y $2 \mid 6$, de modo que $2$ es una cota inferior común.

  \item $(\mathbb{R}, \leq)$ con el orden usual. Es un orden total pero no un bien-orden: el subconjunto $(0,1)$ no tiene elemento mínimo.
\end{enumerate}
\end{ejemplo}

\begin{definicion}{Elementos notables en un poset}{notables}
Sea $(A, \leq)$ un poset. Un elemento $a \in A$ es:
\begin{itemize}[leftmargin=2em]
  \item \textbf{Mínimo} (o \textbf{menor elemento}): $\forall b \in A,\; a \leq b$.
  \item \textbf{Máximo} (o \textbf{mayor elemento}): $\forall b \in A,\; b \leq a$.
  \item \textbf{Minimal}: $\nexists b \in A,\; b < a$.
  \item \textbf{Maximal}: $\nexists b \in A,\; a < b$.
  \item \textbf{Cota inferior} de $S \subseteq A$: $\forall b \in S,\; a \leq b$.
  \item \textbf{Ínfimo} (o \textbf{mayor cota inferior}) de $S$: cota inferior $a$ tal que $b \leq a$ para toda cota inferior $b$.
  \item \textbf{Supremo} (o \textbf{menor cota superior}) de $S$: simétricamente.
\end{itemize}
\end{definicion}

\begin{observacion}{Mínimo vs.\ minimal}{minvsmin}
Todo mínimo es minimal, pero no todo minimal es mínimo. En $(\mathcal{P}(\{1,2,3\}) \setminus \{\emptyset\}, \subseteq)$, los conjuntos $\{1\}$, $\{2\}$ y $\{3\}$ son todos minimales, pero ninguno es mínimo (porque son incomparables entre sí). En $(\mathcal{P}(\{1,2,3\}), \subseteq)$, el conjunto $\emptyset$ es el único mínimo.
\end{observacion}

\subsection{Órdenes totales (órdenes lineales)}

\begin{definicion}{Orden total (lineal)}{ordtotal}
Un orden parcial $\leq$ sobre $A$ es un \textbf{orden total} (o \textbf{lineal}) si satisface adicionalmente la propiedad de \textbf{totalidad}:
\[
  \forall a,b \in A,\quad a \leq b \;\vee\; b \leq a.
\]
El par $(A, \leq)$ se llama \textbf{cadena} o \textbf{conjunto totalmente ordenado}.

De manera equivalente, en la versión estricta, $<$ satisface la \textbf{tricotomía}:
\[
  \forall a,b \in A,\quad a < b \;\vee\; a = b \;\vee\; b < a.
\]
\end{definicion}

\begin{ejemplo}{Órdenes totales y órdenes no totales}{}
\begin{enumerate}[leftmargin=2em, label=(\alph*)]
  \item $(\mathbb{N}, \leq)$, $(\mathbb{Z}, \leq)$, $(\mathbb{Q}, \leq)$ y $(\mathbb{R}, \leq)$ son órdenes totales.

  \item $(\mathcal{P}(X), \subseteq)$ para $|X| \geq 2$ \emph{no} es un orden total. Por ejemplo, $\{a\}$ y $\{b\}$ ($a \neq b$) son incomparables.

  \item El \textbf{orden lexicográfico} en $\mathbb{Z}^2$: $(a_1, a_2) \leq_{\mathrm{lex}} (b_1, b_2)$ si $a_1 < b_1$, o bien $a_1 = b_1$ y $a_2 \leq b_2$. Es un orden total sobre $\mathbb{Z}^2$.

  \item El \textbf{orden lexicográfico} en las cadenas finitas de $\{0,1\}^*$ extiende la idea del diccionario y es un orden total (si las cadenas tienen longitud fija) o un preorden (si pueden tener distinta longitud).
\end{enumerate}
\end{ejemplo}

\begin{proposicion}{Caracterización de cadenas en posets}{cadenas}
Sea $(A, \leq)$ un poset. Un subconjunto $C \subseteq A$ es una \textbf{cadena en $A$} si la restricción de $\leq$ a $C$ es un orden total. Equivalentemente, cualesquiera dos elementos de $C$ son comparables.
\end{proposicion}

\begin{ejemplo}{Lema de Zorn como herramienta}{}
Recordando el Lema de Zorn de la Clase~2: si $(A, \leq)$ es un poset no vacío tal que toda cadena en $A$ tiene una cota superior en $A$, entonces $A$ tiene un elemento maximal. Este lema es equivalente al Axioma de Elección y se usará con frecuencia para demostrar la existencia de bien-órdenes.
\end{ejemplo}

% ===========================================================
\section{Conjuntos bien ordenados}
% ===========================================================

\subsection{Definición y primeras propiedades}

\begin{definicion}{Conjunto bien ordenado}{biorden}
Un orden total $(A, <)$ es un \textbf{bien-orden} (o \textbf{buena ordenación}) si satisface la propiedad de \textbf{buen orden}: todo subconjunto no vacío de $A$ tiene un \textbf{elemento mínimo}. Formalmente:
\[
  \forall S \subseteq A,\; \bigl(S \neq \emptyset \to \exists m \in S\,\forall s \in S\,(m \leq s)\bigr).
\]
El par $(A, <)$ se llama un \textbf{conjunto bien ordenado}.
\end{definicion}

\begin{observacion}{El elemento mínimo es único}{minunico}
En un orden total, el elemento mínimo de un subconjunto no vacío, cuando existe, es \emph{único}. En efecto, si $m$ y $m'$ son ambos mínimos de $S$, entonces $m \leq m'$ y $m' \leq m$, luego $m = m'$ por antisimetría.
\end{observacion}

\begin{ejemplo}{Conjuntos bien ordenados y no bien ordenados}{}
\begin{enumerate}[leftmargin=2em, label=(\alph*)]
  \item $(\mathbb{N}, \leq)$ es un conjunto bien ordenado: todo subconjunto no vacío de $\mathbb{N}$ tiene un mínimo. Este es el \emph{principio del buen orden de los naturales}, equivalente a la inducción matemática ordinaria.

  \item $(\mathbb{Z}, \leq)$ \emph{no} es un bien-orden: el subconjunto $\mathbb{Z}$ mismo no tiene mínimo (para cualquier $n \in \mathbb{Z}$, se tiene $n-1 < n$).

  \item $(\mathbb{Q} \cap [0,1], \leq)$ \emph{no} es un bien-orden: el subconjunto $(0,1) \cap \mathbb{Q}$ no tiene mínimo.

  \item $(\mathbb{R}_{\geq 0}, \leq)$ \emph{no} es un bien-orden: el subconjunto $(0, +\infty)$ no tiene mínimo.

  \item El conjunto $\{0, 1, 2, \ldots, \omega, \omega+1\}$ con el orden que extiende $\mathbb{N}$ poniendo todos los naturales antes de $\omega$ y $\omega < \omega+1$ es un bien-orden. (Veremos esto con precisión cuando definamos los ordinales.)

  \item \textbf{Orden antiléxico en $\mathbb{N}^2$:} $(m_1, n_1) \prec (m_2, n_2)$ si $m_1 + n_1 < m_2 + n_2$, o bien $m_1 + n_1 = m_2 + n_2$ y $m_1 < m_2$. Se puede verificar que $(\mathbb{N}^2, \prec)$ es un bien-orden. Los primeros elementos son:
  \[
    (0,0) \prec (1,0) \prec (0,1) \prec (2,0) \prec (1,1) \prec (0,2) \prec (3,0) \prec \cdots
  \]
\end{enumerate}
\end{ejemplo}

\begin{proposicion}{El bien orden implica buen fundamento}{bienfund}
Todo conjunto bien ordenado $(A, <)$ es \textbf{bien fundado}: no existen cadenas descendentes infinitas en $A$, es decir, no existe ninguna sucesión $(a_n)_{n \in \mathbb{N}}$ con $a_0 > a_1 > a_2 > \cdots$
\end{proposicion}

\begin{proof}
Supongamos que existe una cadena descendente $a_0 > a_1 > a_2 > \cdots$\ El conjunto $S = \{a_n \mid n \in \mathbb{N}\}$ es un subconjunto no vacío de $A$. Por la propiedad de buen orden, $S$ tiene un elemento mínimo $m = a_k$ para algún $k$. Pero entonces $a_{k+1} < a_k = m$ y $a_{k+1} \in S$, lo que contradice que $m$ sea el mínimo de $S$.
\end{proof}

\begin{proposicion}{El bien orden implica orden total}{biord-total}
Si $(A, <)$ es un bien-orden, entonces $(A, <)$ es un orden total.
\end{proposicion}

\begin{proof}
Sean $a, b \in A$ con $a \neq b$. El conjunto $\{a, b\}$ es no vacío, luego tiene un mínimo: o $a < b$ o $b < a$. En cualquier caso, $a$ y $b$ son comparables.
\end{proof}

\begin{teorema}{Principio del mínimo y buen orden}{minbien}
Sea $(A, <)$ un orden total. Las siguientes afirmaciones son equivalentes:
\begin{enumerate}[leftmargin=2em, label=(\roman*)]
  \item $(A, <)$ es un bien-orden.
  \item \textbf{Principio del mínimo:} toda propiedad $\varphi$ que se cumple para algún elemento de $A$ se cumple para algún elemento mínimo con respecto a esa propiedad, es decir:
  \[
    \exists a \in A\,\varphi(a) \;\implies\; \exists m \in A\,\bigl(\varphi(m) \wedge \forall b < m,\, \neg\varphi(b)\bigr).
  \]
  \item \textbf{Principio de inducción:} si $P \subseteq A$ satisface
  \[
    \forall a \in A,\; \bigl(\forall b < a,\; b \in P\bigr) \to a \in P,
  \]
  entonces $P = A$.
\end{enumerate}
\end{teorema}

\begin{proof}
\textbf{(i) $\Rightarrow$ (ii):} Dado $a \in A$ con $\varphi(a)$, el conjunto $S = \{x \in A \mid \varphi(x)\}$ es no vacío (contiene a $a$). Como $(A,<)$ es un bien-orden, $S$ tiene un mínimo $m$, que satisface $\varphi(m)$ y $\neg\varphi(b)$ para todo $b < m$.

\textbf{(ii) $\Rightarrow$ (iii):} Sea $P \subseteq A$ satisfaciendo la hipótesis de inducción. Supongamos que $P \neq A$, es decir, $A \setminus P \neq \emptyset$. Sea $\varphi(a) = (a \notin P)$. Por (ii), existe $m$ con $m \notin P$ y $b \in P$ para todo $b < m$. Pero entonces la hipótesis inductiva aplicada a $m$ daría $m \in P$, contradicción.

\textbf{(iii) $\Rightarrow$ (i):} Sea $S \subseteq A$ no vacío; queremos mostrar que $S$ tiene un mínimo. Definamos $P = \{a \in A \mid a \notin S\} \cup \{a \in A \mid a \text{ es el mínimo de } S\}$. Un argumento cuidadoso, usando la propiedad (iii) con $P = A \setminus S$ y el hecho de que $S \neq \emptyset$, muestra que el primer elemento de $S$ que «se alcanza» es el mínimo. Formalmente: si $P = \{a \in A \mid \forall b \leq a,\, b \notin S\}$, la hipótesis (iii) implica que o $P = A$ (lo que sería vacíamente posible solo si $S = \emptyset$) o existe un mínimo en $S$.
\end{proof}

\begin{observacion}{Equivalencia con inducción ordinaria}{equivind}
Para $A = \mathbb{N}$, la equivalencia del Teorema~\ref{teo:minbien} recupera la equivalencia clásica entre:
\begin{itemize}[leftmargin=2em]
  \item El principio del buen orden de $\mathbb{N}$ (todo subconjunto no vacío tiene mínimo).
  \item El principio de inducción matemática fuerte (si la verdad de $P(n)$ se sigue de $P(k)$ para todo $k < n$, entonces $P$ vale para todos los naturales).
\end{itemize}
La inducción transfinita generaliza exactamente este esquema a ordinales arbitrarios.
\end{observacion}

\subsection{Ejemplos desarrollados de bien-órdenes}

\begin{ejemplo}{El bien-orden de $\mathbb{N} \times \mathbb{N}$}{}
Definimos el \textbf{orden lexicográfico} en $\mathbb{N} \times \mathbb{N}$:
\[
  (m_1, n_1) <_{\mathrm{lex}} (m_2, n_2) \;\iff\; m_1 < m_2, \text{ o } (m_1 = m_2 \text{ y } n_1 < n_2).
\]
\textbf{Afirmación:} $(\mathbb{N} \times \mathbb{N}, <_{\mathrm{lex}})$ es un bien-orden.

\textbf{Demostración:} Sea $S \subseteq \mathbb{N} \times \mathbb{N}$ no vacío. Sea $S_1 = \{m \in \mathbb{N} \mid \exists n, (m,n) \in S\}$. Como $S_1 \subseteq \mathbb{N}$ es no vacío, tiene un mínimo $m_0$. Sea $S_2 = \{n \in \mathbb{N} \mid (m_0, n) \in S\}$. Como $S_2 \subseteq \mathbb{N}$ es no vacío, tiene un mínimo $n_0$. Entonces $(m_0, n_0)$ es el elemento mínimo de $S$ en el orden lexicográfico.
\end{ejemplo}

\begin{ejemplo}{El bien-orden en cadenas de longitud fija}{}
Sea $\Sigma = \{0, 1\}$ y sea $\Sigma^n = \{0,1\}^n$ el conjunto de cadenas binarias de longitud $n$. El orden lexicográfico en $\Sigma^n$ (comparando coordenada a coordenada) es un bien-orden. Además, el orden lexicográfico en $\Sigma^* = \bigcup_{n \in \mathbb{N}} \Sigma^n$ puede extenderse a un bien-orden de $\Sigma^*$ ordenando primero por longitud y luego lexicográficamente.
\end{ejemplo}

% ===========================================================
\section{Segmentos iniciales e isomorfismos de orden}
% ===========================================================

\subsection{Segmentos iniciales}

\begin{definicion}{Segmento inicial}{seginicial}
Sea $(A, <)$ un conjunto ordenado. Un subconjunto $I \subseteq A$ es un \textbf{segmento inicial} (o \textbf{sección inicial}) de $A$ si es \emph{cerrado hacia abajo}:
\[
  \forall a \in I,\; \forall b \in A,\; b < a \;\implies\; b \in I.
\]
Un segmento inicial $I$ es \textbf{propio} si $I \neq A$.
\end{definicion}

\begin{ejemplo}{Segmentos iniciales en distintos órdenes}{}
\begin{enumerate}[leftmargin=2em, label=(\alph*)]
  \item En $(\mathbb{N}, <)$, los segmentos iniciales son exactamente $\emptyset$ y los conjuntos $\{0, 1, \ldots, n-1\}$ para $n \in \mathbb{N}$, junto con $\mathbb{N}$ mismo.

  \item En $(\mathbb{R}, <)$, los segmentos iniciales son los intervalos $(-\infty, a)$ y $(-\infty, a]$ para $a \in \mathbb{R}$, junto con $\emptyset$ y $\mathbb{R}$.

  \item En $(\mathcal{P}(\{1,2,3\}), \subsetneq)$ (como orden parcial), el conjunto $\{\emptyset, \{1\}, \{2\}\}$ \emph{no} es un segmento inicial porque $\{1,2\} \supset \{1\}$ pero $\{1,2\} \notin \{\emptyset,\{1\},\{2\}\}$. En cambio, $\{\emptyset, \{1\}, \{2\}, \{1,2\}\}$ sí lo es.
\end{enumerate}
\end{ejemplo}

\begin{definicion}{Segmento inicial generado por un elemento}{seginelem}
Sea $(A, <)$ un conjunto ordenado y $a \in A$. El \textbf{segmento inicial generado por $a$} (o \textbf{sección de $a$}) es el conjunto
\[
  s(a) \;=\; \mathrm{pred}(a) \;=\; \{b \in A \mid b < a\}.
\]
\end{definicion}

\begin{observacion}{Los segmentos iniciales propios en bien-órdenes}{propbien}
En un bien-orden $(A, <)$, los segmentos iniciales propios son exactamente los conjuntos de la forma $s(a) = \{b \in A \mid b < a\}$ para algún $a \in A$. Esto se demuestra a continuación.
\end{observacion}

\begin{teorema}{Caracterización de segmentos iniciales en bien-órdenes}{caract-segin}
Sea $(A, <)$ un bien-orden. Un subconjunto propio $I \subsetneq A$ es un segmento inicial de $A$ si y solo si $I = s(a)$ para algún $a \in A$ (el único $a$ tal que $I = s(a)$ es $a = \min(A \setminus I)$).
\end{teorema}

\begin{proof}
\textbf{($\Leftarrow$)} $s(a)$ es un segmento inicial: si $b \in s(a)$ y $c < b$, entonces $c < b < a$ luego $c < a$ por transitividad, de modo que $c \in s(a)$. Y $s(a) \neq A$ porque $a \notin s(a)$.

\textbf{($\Rightarrow$)} Sea $I \subsetneq A$ un segmento inicial propio. El conjunto $A \setminus I$ es no vacío; como $(A,<)$ es un bien-orden, tiene un mínimo $a = \min(A \setminus I)$.

\emph{Afirmamos que $I = s(a)$.}
\begin{itemize}[leftmargin=2em]
  \item $I \subseteq s(a)$: Si $b \in I$, entonces $b \notin A \setminus I$, luego $b \neq a$ y $b \not\geq a$ (pues $a$ es el mínimo de $A \setminus I$). Por totalidad del orden, $b < a$, es decir, $b \in s(a)$.
  \item $s(a) \subseteq I$: Si $b < a$, entonces $b \notin A \setminus I$ (pues $a = \min(A \setminus I)$), luego $b \in I$.
\end{itemize}
Por tanto $I = s(a)$. La unicidad de $a$ es clara: si $s(a) = s(a')$ con $a \neq a'$, entonces (suponiendo $a < a'$) tendríamos $a \in s(a') = s(a)$, lo que daría $a < a$, contradicción.
\end{proof}

\begin{corolario}{Ningún bien-orden es isomorfo a un segmento inicial propio suyo}{nobien-iso}
Sea $(A, <)$ un bien-orden. Entonces $A$ no es isomorfo (como orden) a ningún segmento inicial propio $s(a)$ con $a \in A$.
\end{corolario}

\begin{proof}
Ver Teorema~\ref{teo:unicidadiso} más adelante.
\end{proof}

\subsection{Isomorfismos de orden}

\begin{definicion}{Homomorfismo e isomorfismo de orden}{isoord}
Sean $(A, <_A)$ y $(B, <_B)$ dos conjuntos ordenados. Una función $f: A \to B$ es:
\begin{itemize}[leftmargin=2em]
  \item Un \textbf{homomorfismo de orden} (o \textbf{función estrictamente creciente}) si preserva el orden: $a <_A a' \implies f(a) <_B f(a')$.
  \item Un \textbf{isomorfismo de orden} si es una biyección y preserva el orden en ambas direcciones:
  \[
    a <_A a' \;\iff\; f(a) <_B f(a').
  \]
\end{itemize}
Se escribe $(A, <_A) \cong (B, <_B)$ cuando existe un isomorfismo de orden entre $A$ y $B$.
\end{definicion}

\begin{observacion}{Todo homomorfismo estrictamente creciente es inyectivo}{homoinj}
Si $f: A \to B$ es estrictamente creciente y $a \neq a'$ en $A$, entonces o $a < a'$ (luego $f(a) < f(a')$ y $f(a) \neq f(a')$) o $a' < a$ (simétricamente). Luego $f$ es inyectiva.
\end{observacion}

\begin{proposicion}{Inversa de un isomorfismo de orden}{inviso}
Si $f: (A, <_A) \to (B, <_B)$ es un isomorfismo de orden, entonces su inversa $f^{-1}: B \to A$ también es un isomorfismo de orden.
\end{proposicion}

\begin{proof}
Como $f$ es biyección, $f^{-1}$ existe y es biyección. Sean $b, b' \in B$ con $b <_B b'$. Sean $a = f^{-1}(b)$ y $a' = f^{-1}(b')$. Por la propiedad de $f$: $a <_A a'$ (de lo contrario $f(a) \geq_B f(a')$, contradiciendo $b <_B b'$). Luego $f^{-1}(b) <_A f^{-1}(b')$.
\end{proof}

\begin{ejemplo}{Isomorfismos de orden entre órdenes conocidos}{}
\begin{enumerate}[leftmargin=2em, label=(\alph*)]
  \item $(\mathbb{N}, <) \cong (\mathbb{N}_{\text{par}}, <)$ donde $\mathbb{N}_{\text{par}} = \{0, 2, 4, \ldots\}$, mediante $f(n) = 2n$.

  \item $(\mathbb{Q} \cap (0,1), <) \cong (\mathbb{Q}, <)$: la función $f(q) = \tan(\pi q - \pi/2)$ es un isomorfismo de orden (de $\mathbb{Q} \cap (0,1)$ en su imagen, que es densa en $\mathbb{R}$; pero aquí se usa una biyección explícita que preserva el orden racional, cuya construcción requiere la teoría de DLO).

  \item $(\{0,1,2\}, <) \ncong (\{0,1,2,3\}, <)$: el primer conjunto tiene $3$ elementos y el segundo $4$; como un isomorfismo es una biyección, no pueden ser isomorfos.

  \item $(\omega, <)$ (naturales con orden usual) $\ncong$ $(\omega + 1, <)$ (que incluye un elemento $\omega$ mayor que todos los naturales): el segundo tiene un elemento máximo, el primero no; un isomorfismo preservaría la propiedad de tener máximo, luego no pueden ser isomorfos.
\end{enumerate}
\end{ejemplo}

\subsection{Teoremas fundamentales sobre isomorfismos de bien-órdenes}

Los siguientes teoremas son los resultados más importantes de esta clase y prepararan directamente la definición de número ordinal.

\begin{teorema}{Rigidez de los bien-órdenes}{rigidez}
Sea $(A, <)$ un bien-orden. El único isomorfismo de orden de $A$ en sí mismo (automorfismo) es la identidad.
\end{teorema}

\begin{proof}
Sea $f: A \to A$ un isomorfismo de orden. Supongamos que $f \neq \mathrm{id}_A$. Sea $S = \{a \in A \mid f(a) \neq a\}$, que por hipótesis es no vacío. Como $(A,<)$ es un bien-orden, $S$ tiene un mínimo $m$.

Como $f$ es isomorfismo, $m$ y $f(m)$ son comparables. Distinguimos casos:
\begin{itemize}[leftmargin=2em]
  \item \textbf{Caso $f(m) < m$:} Sea $b = f(m) < m$. Como $b < m$, por minimalidad de $m$, $b \notin S$, es decir, $f(b) = b$. Pero $f(f(m)) = f(b) = b = f(m)$, y como $f$ es inyectiva, $f(m) = m$, contradicción.
  \item \textbf{Caso $f(m) > m$:} Sea $b = f(m) > m$. Tenemos $f(m) = b$ y aplicando $f$: $f(b) = f(f(m))$. Como $f$ preserva el orden y $m < b$, $f(m) < f(b)$, es decir, $b < f(b)$, lo que significa $f(b) > b$, luego $f(b) \neq b$ y $b \in S$. Pero $b > m$, contradiciendo que $m$ es el mínimo de $S$.
\end{itemize}
En ambos casos llegamos a una contradicción. Luego $S = \emptyset$ y $f = \mathrm{id}_A$.
\end{proof}

\begin{teorema}{Unicidad del isomorfismo entre bien-órdenes}{unicidadiso}
Sean $(A, <_A)$ y $(B, <_B)$ bien-órdenes. Existe \textbf{a lo sumo un} isomorfismo de orden $f: A \to B$.
\end{teorema}

\begin{proof}
Si $f, g: A \to B$ son isomorfismos de orden, entonces $g^{-1} \circ f: A \to A$ es un automorfismo del bien-orden $A$. Por el Teorema~\ref{teo:rigidez}, $g^{-1} \circ f = \mathrm{id}_A$, luego $f = g$.
\end{proof}

\begin{teorema}{Tricotomía de bien-órdenes}{tricobien}
Sean $(A, <_A)$ y $(B, <_B)$ bien-órdenes. Se cumple exactamente una de las siguientes:
\begin{enumerate}[leftmargin=2em, label=(\roman*)]
  \item $(A, <_A) \cong (B, <_B)$ (son isomorfos).
  \item $(A, <_A) \cong (s_B(b), <_B)$ para algún $b \in B$ (es isomorfo a un segmento inicial propio de $B$).
  \item $(B, <_B) \cong (s_A(a), <_A)$ para algún $a \in A$ (es isomorfo a un segmento inicial propio de $A$).
\end{enumerate}
\end{teorema}

\begin{proof}
Consideremos el conjunto de todas las funciones de orden entre segmentos iniciales de $A$ y segmentos iniciales de $B$:
\[
  \mathcal{F} = \{f: s_A(a) \to s_B(b) \mid f \text{ es isomorfismo de orden, } a \in A \cup \{A\},\, b \in B \cup \{B\}\}.
\]
La unión $F = \bigcup_{f \in \mathcal{F}} f$ (como conjunto de pares) es una función (pues dos isomorfismos en $\mathcal{F}$ que coincidan en un punto coinciden en todo su dominio, por la unicidad). Sea $D = \mathrm{dom}(F)$ y $R = \mathrm{ran}(F)$.

\textbf{$D$ es un segmento inicial de $A$:} Si $a' \in D$ y $a <_A a'$, entonces $a \in \mathrm{dom}(f)$ para cualquier $f \in \mathcal{F}$ con $a' \in \mathrm{dom}(f)$, luego $a \in D$.

\textbf{$R$ es un segmento inicial de $B$:} Simétricamente.

\textbf{$F: D \to R$ es un isomorfismo de orden.}

Por el Teorema~\ref{teo:caract-segin}, $D = A$ o $D = s_A(a_0)$ para algún $a_0$; análogamente $R = B$ o $R = s_B(b_0)$. Si $D = s_A(a_0)$ y $R = s_B(b_0)$, entonces $F$ se extiende mapeando $a_0 \mapsto b_0$, contradiciendo la maximalidad. Luego se cumple exactamente una de las tres opciones del enunciado.
\end{proof}

\begin{corolario}{Ningún bien-orden es isomorfo a un segmento inicial propio suyo}{nobienisopropio}
Si $(A, <)$ es un bien-orden y $a \in A$, entonces $A \ncong s(a)$.
\end{corolario}

\begin{proof}
Si $f: A \to s(a)$ fuera un isomorfismo, como $s(a) \subsetneq A$, definiríamos $g: A \to A$ como $g = \iota \circ f$ donde $\iota: s(a) \hookrightarrow A$ es la inclusión. Entonces $g$ sería un homomorfismo estrictamente creciente de $A$ en $A$ con imagen contenida en $s(a)$. En particular $g(a) < a$. Pero por inducción sobre el bien-orden (aplicando el principio del mínimo al conjunto $\{x \in A \mid g(x) < x\}$ y notando que si $g(x) < x$ para el mínimo tal $x$, obtenemos $g(g(x)) < g(x) < x$, contradiciendo la minimalidad de $x$) se puede mostrar que cualquier tal función satisface $g(x) \geq x$ para todo $x$. Contradicción.
\end{proof}

% ===========================================================
\section{Inducción transfinita: introducción preliminar}
% ===========================================================

\subsection{Motivación: más allá de los naturales}

La inducción matemática ordinaria en $\mathbb{N}$ afirma que si una propiedad $P$ es verdadera para $0$ y se hereda al sucesor (de $P(n)$ se sigue $P(n+1)$), entonces $P$ es verdadera para todos los naturales. Esta es una consecuencia directa del buen orden de $\mathbb{N}$.

El \textbf{principio de inducción transfinita} generaliza esta idea a cualquier conjunto bien ordenado, y en particular a los números ordinales que definiremos en la Clase~4. Lo enunciamos aquí en su versión general.

\subsection{El principio de inducción transfinita}

\begin{teorema}{Principio de inducción transfinita}{induct-transfinita}
Sea $(A, <)$ un bien-orden y sea $P \subseteq A$ un subconjunto tal que
\[
  \forall a \in A,\quad \bigl(\forall b \in A,\; b < a \implies b \in P\bigr) \;\implies\; a \in P.
\]
Entonces $P = A$.
\end{teorema}

\begin{proof}
Este es precisamente el principio (iii) del Teorema~\ref{teo:minbien}, cuya equivalencia con el buen orden ya demostramos. La hipótesis dice que $P$ contiene a $a$ siempre que $P$ contiene a todos los predecesores de $a$. Si $P \neq A$, el conjunto $A \setminus P$ es no vacío y tiene un mínimo $m$. Por la hipótesis de inducción aplicada a $m$: todo $b < m$ satisface $b \in P$ (pues $m$ es el mínimo de $A \setminus P$). Pero entonces la condición del enunciado da $m \in P$, contradicción.
\end{proof}

\begin{observacion}{Tres formas de la inducción transfinita}{tresformas}
Al igual que en la inducción ordinaria, la inducción transfinita suele presentarse en tres formas, dependiendo del tipo de conjunto bien ordenado:
\begin{enumerate}[leftmargin=2em, label=(\arabic*)]
  \item \textbf{Forma básica (válida para cualquier bien-orden):} La enunciada en el Teorema~\ref{teo:induct-transfinita}.

  \item \textbf{Forma para ordinales con distinción de tipos:} Cuando el bien-orden es el de los ordinales, los elementos se clasifican en \emph{ordinales sucesor} ($\alpha = \beta + 1$) y \emph{ordinales límite} (no son el sucesor de ningún otro). En este caso, la inducción transfinita tiene tres pasos:
  \begin{itemize}[leftmargin=2em]
    \item \textbf{Caso base:} Demostrar $P(0)$.
    \item \textbf{Paso sucesor:} Demostrar que $P(\alpha) \implies P(\alpha+1)$.
    \item \textbf{Paso límite:} Demostrar que si $P(\beta)$ para todo $\beta < \lambda$ (donde $\lambda$ es un ordinal límite), entonces $P(\lambda)$.
  \end{itemize}

  \item \textbf{Recursión transfinita:} Permite \emph{definir} funciones $f: A \to C$ en un conjunto bien ordenado $A$ especificando el valor $f(a)$ en términos de los valores $\{f(b) \mid b < a\}$ ya calculados. Esto se tratará con detalle en la Clase~4.
\end{enumerate}
\end{observacion}

\begin{ejemplo}{Aplicación de la inducción transfinita: propiedades de $\mathbb{N}$}{}
Demostraremos que $\forall n \in \mathbb{N},\; n + 0 = n$ usando inducción transfinita (que para $\mathbb{N}$ coincide con la inducción fuerte).

Sea $P = \{n \in \mathbb{N} \mid n + 0 = n\}$. Supongamos que $P$ contiene a todos los predecesores de $n$ (es decir, a todos los $k < n$); debemos mostrar que $n \in P$.
\begin{itemize}[leftmargin=2em]
  \item Si $n = 0$: $0 + 0 = 0$ por definición de la suma (caso base).
  \item Si $n = k+1$ para algún $k$: tenemos $k \in P$ por hipótesis, es decir, $k + 0 = k$. Entonces $(k+1) + 0 = k + 0 + 1 = k + 1$ (usando la recursión de la suma).
\end{itemize}
Por inducción transfinita, $P = \mathbb{N}$.
\end{ejemplo}

\begin{ejemplo}{Primer ejemplo transfinito: ordenar $\mathbb{N} \cup \{\omega\}$}{}
Consideremos el conjunto $A = \mathbb{N} \cup \{\omega\}$ con el orden $n < \omega$ para todo $n \in \mathbb{N}$ y el orden usual en $\mathbb{N}$. Verifiquemos que $(A, <)$ es un bien-orden.

Sea $S \subseteq A$ no vacío. Si $S \cap \mathbb{N} \neq \emptyset$, entonces $S \cap \mathbb{N}$ tiene un mínimo $m$ en $\mathbb{N}$ (pues $(\mathbb{N}, <)$ es un bien-orden); como todo $n \in \mathbb{N}$ satisface $n < \omega$, el elemento $m$ es también el mínimo de $S$ en $A$. Si $S \cap \mathbb{N} = \emptyset$, entonces $S = \{\omega\}$, cuyo único elemento es $\omega$, que es trivialmente el mínimo.

Ahora realizamos inducción transfinita para mostrar que toda propiedad hereditaria se cumple en $A$. Sea $P \subseteq A$ con la propiedad: para todo $a \in A$, si $P$ contiene a todos los predecesores de $a$ entonces $a \in P$.
\begin{itemize}[leftmargin=2em]
  \item Para $0 \in \mathbb{N}$: $0$ no tiene predecesores, luego la hipótesis es vacuamente cierta, y $0 \in P$.
  \item Para $n+1 \in \mathbb{N}$: los predecesores de $n+1$ son $0,1,\ldots,n$. Si todos están en $P$, en particular $n \in P$; si la hipótesis dice que de $P(n)$ se sigue $P(n+1)$, concluimos.
  \item Para $\omega$: los predecesores de $\omega$ son todos los naturales. Si $\mathbb{N} \subseteq P$, la hipótesis da $\omega \in P$.
\end{itemize}
Luego $P = A$. Este es un ejemplo concreto de inducción transfinita que va \emph{más allá} de los naturales.
\end{ejemplo}

\subsection{Recursión transfinita (enunciado)}

\begin{teorema}{Principio de recursión transfinita}{recursion}
Sea $(A, <)$ un bien-orden y sea $G$ una función que a cada conjunto de la forma $\{(b, f(b)) \mid b < a\}$ (la «historia» de $f$ hasta $a$) le asigna un valor en un conjunto dado $C$. Entonces existe una \textbf{única} función $f: A \to C$ tal que
\[
  f(a) = G\bigl(\{(b, f(b)) \mid b < a\}\bigr) \quad \text{para todo } a \in A.
\]
\end{teorema}

\begin{proof}
La demostración completa usa inducción transfinita para construir la función $f$ por etapas, y se estudiará en la Clase~4 junto con su aplicación a la aritmética de ordinales.
\end{proof}

\begin{importante}
La inducción transfinita y la recursión transfinita son las herramientas fundamentales de la teoría de ordinales. Permiten definir la suma, el producto y la potencia de ordinales de manera recursiva, extender construcciones de los naturales a todos los ordinales, y demostrar resultados sobre la jerarquía del universo de conjuntos $V = \bigcup_{\alpha} V_\alpha$.
\end{importante}

% ===========================================================
\section{El Teorema del Buen Orden revisitado}
% ===========================================================

\begin{teorema}{Teorema del Buen Orden (Zermelo, 1904)}{teobienord}
Todo conjunto puede ser bien ordenado. Formalmente, para todo conjunto $A$ existe una relación $<$ sobre $A$ tal que $(A, <)$ es un bien-orden.
\end{teorema}

\begin{proof}
La demostración usa el Axioma de Elección en su forma equivalente (se vio en la Clase~2: AC $\iff$ TBO). Presentamos la idea principal.

Sea $A$ un conjunto. Por el Axioma de Elección, existe una función de elección $f: \mathcal{P}(A) \setminus \{\emptyset\} \to A$ con $f(S) \in S$ para todo $S \neq \emptyset$.

Se define inductivamente la enumeración: $a_0 = f(A)$, $a_1 = f(A \setminus \{a_0\})$, y en general $a_\alpha = f(A \setminus \{a_\beta \mid \beta < \alpha\})$ (para ordinales $\alpha$ tales que $A \setminus \{a_\beta \mid \beta < \alpha\} \neq \emptyset$). Este proceso termina en algún ordinal, dando una enumeración bien ordenada de $A$. El detalle técnico requiere la teoría de ordinales que desarrollaremos en la Clase~4.
\end{proof}

\begin{observacion}{Consecuencias del TBO para la teoría de órdenes}{consTBO}
El Teorema del Buen Orden implica que todo orden parcial está contenido en un orden total (se puede extender cualquier orden parcial a un orden total), y en particular, todo conjunto tiene al menos un bien-orden. Sin embargo, salvo para conjuntos muy específicos, no existe un procedimiento efectivo para construir dicho bien-orden.
\end{observacion}

% ===========================================================
\section{Resumen y conexiones hacia la siguiente clase}
% ===========================================================

En esta clase hemos construido el marco teórico para la definición de los números ordinales. Los resultados más importantes que debemos recordar son:

\begin{importante}
\begin{enumerate}[leftmargin=2em, label=\textbf{\arabic*.}]
  \item \textbf{Bien-orden $\Leftrightarrow$ Principio del mínimo $\Leftrightarrow$ Inducción transfinita} (Teorema~\ref{teo:minbien}).

  \item \textbf{Los segmentos iniciales propios de un bien-orden son exactamente los $s(a)$} (Teorema~\ref{teo:caract-segin}).

  \item \textbf{Ningún bien-orden es isomorfo a uno de sus segmentos iniciales propios} (Corolario~\ref{cor:nobienisopropio}).

  \item \textbf{Unicidad del isomorfismo:} entre dos bien-órdenes, si existe un isomorfismo de orden, es único (Teorema~\ref{teo:unicidadiso}).

  \item \textbf{Tricotomía de bien-órdenes:} dos bien-órdenes cualesquiera son comparables (uno es isomorfo a un segmento inicial propio del otro, o son isomorfos) (Teorema~\ref{teo:tricobien}).
\end{enumerate}
\end{importante}

En la \textbf{Clase~4} usaremos estas propiedades para definir los \emph{números ordinales de von Neumann}: cada ordinal es el conjunto de todos los ordinales menores, y los bien-órdenes dan a cada conjunto un «tipo ordinal» que lo caracteriza salvo isomorfismo.

% ===========================================================
\section{Problemas y tareas}
% ===========================================================

\begin{enumerate}[leftmargin=2em, label=\textbf{\arabic*.}]

  \item \textbf{(Propiedades de relaciones.)}
  Sea $A = \{1, 2, 3, 4\}$ y sea $R = \{(1,1),(1,2),(2,3),(1,3),(3,4),(2,4),(1,4),(2,2),(3,3),(4,4)\}$.
  \begin{enumerate}[label=(\alph*)]
    \item Verifica, directamente desde la definición, que $R$ es reflexiva, antisimétrica y transitiva. Concluye que $(A, R)$ es un poset.
    \item ¿Es $R$ un orden total? Justifica.
    \item Dibuja el diagrama de Hasse de $(A, R)$.
    \item Determina el mínimo, el máximo, los elementos minimales y los maximales (si existen) de $A$.
  \end{enumerate}

  \item \textbf{(Poset de divisibilidad.)}
  Considera el conjunto $D_{30} = \{1, 2, 3, 5, 6, 10, 15, 30\}$ (divisores de $30$) con la relación de divisibilidad $\mid$.
  \begin{enumerate}[label=(\alph*)]
    \item Verifica que $(D_{30}, \mid)$ es un poset.
    \item Dibuja el diagrama de Hasse de $(D_{30}, \mid)$.
    \item Determina, para cada par $\{a,b\} \subseteq D_{30}$, el ínfimo (máximo común divisor) y el supremo (mínimo común múltiplo) en el poset. ¿Es $(D_{30}, \mid)$ un reticulado? (\emph{Recuerda: un reticulado es un poset donde todo par de elementos tiene ínfimo y supremo.})
    \item ¿Es $(D_{30}, \mid)$ isomorfo (como orden) a $(\mathcal{P}(\{p_1, p_2, p_3\}), \subseteq)$ donde $p_1 = 2$, $p_2 = 3$, $p_3 = 5$ son los factores primos de $30$? Exhibe un isomorfismo o demuestra que no existe.
  \end{enumerate}

  \item \textbf{(Cadenas y anticiclos.)}
  \begin{enumerate}[label=(\alph*)]
    \item Sea $(A, \leq)$ un poset. Una \textbf{anticadena} en $A$ es un subconjunto $S \subseteq A$ tal que cualesquiera dos elementos de $S$ son incomparables. Encuentra una anticadena de tamaño máximo en $(\mathcal{P}(\{1,2,3,4\}), \subseteq)$.
    \item Enuncia y demuestra el \textbf{Teorema de Dilworth}: en cualquier poset finito, el tamaño máximo de una anticadena es igual al número mínimo de cadenas necesarias para cubrir el poset.
    \item Aplica el Teorema de Dilworth al poset $(D_{30}, \mid)$ del problema anterior.
  \end{enumerate}

  \item \textbf{(Bien-órdenes: verificación directa.)}
  \begin{enumerate}[label=(\alph*)]
    \item Demuestra que el orden lexicográfico en $\mathbb{N}^n$ (cadenas de $n$ naturales, con orden lexicográfico coordinada a coordinada) es un bien-orden, por inducción sobre $n$.
    \item Demuestra que la unión $A \cup B$ de dos bien-órdenes disjuntos (donde todo elemento de $A$ precede a todo elemento de $B$) es un bien-orden. ¿Es necesaria la condición de disjunción?
    \item Sea $f: (A, <_A) \to (B, <_B)$ un isomorfismo de orden. Demuestra que si $(B, <_B)$ es un bien-orden, entonces $(A, <_A)$ también lo es.
  \end{enumerate}

  \item \textbf{(Segmentos iniciales.)}
  \begin{enumerate}[label=(\alph*)]
    \item Sea $(A, <)$ un bien-orden y sea $I \subseteq A$ un segmento inicial. Demuestra que $(I, <\!\restriction_I)$ también es un bien-orden (donde $<\!\restriction_I$ es la restricción del orden a $I$).
    \item Sean $I$ y $J$ dos segmentos iniciales de un bien-orden $(A, <)$. Demuestra que $I \subseteq J$ o $J \subseteq I$ (es decir, los segmentos iniciales de un bien-orden forman una cadena bajo inclusión).
    \item Demuestra que la intersección de cualquier familia de segmentos iniciales es un segmento inicial, y que la unión también lo es.
    \item Usa el inciso anterior para demostrar que $I = \bigcap\{J \mid J \text{ es segmento inicial y } a \in J\} = s(a) \cup \{a\}$ para todo $a \in A$.
  \end{enumerate}

  \item \textbf{(Isomorfismos de orden.)}
  \begin{enumerate}[label=(\alph*)]
    \item Demuestra que la relación $\cong$ (ser isomorfos como órdenes) es una relación de equivalencia en la clase de conjuntos bien ordenados.
    \item Sea $f: (A, <_A) \to (B, <_B)$ un homomorfismo de orden inyectivo entre bien-órdenes. Demuestra que $f(A)$ es un segmento inicial de $B$ o $f(A) = B$.
    \item Usa el Teorema de Tricotomía para demostrar que si $(A, <_A)$ y $(B, <_B)$ son bien-órdenes finitos con $|A| = |B|$, entonces $A \cong B$.
  \end{enumerate}

  \item \textbf{(Inducción transfinita: aplicaciones.)}
  \begin{enumerate}[label=(\alph*)]
    \item Sea $(A, <)$ un bien-orden y sea $f: A \to A$ un homomorfismo de orden (no necesariamente biyectivo). Demuestra por inducción transfinita que $f(a) \geq a$ para todo $a \in A$.
    \item Usa el resultado anterior para probar el Corolario~\ref{cor:nobienisopropio}: ningún bien-orden es isomorfo a uno de sus segmentos iniciales propios.
    \item Sea $P \subseteq \mathbb{N}$ tal que: (i) $0 \in P$; (ii) si $n \in P$ entonces $n+1 \in P$; (iii) si $n > 0$ y $k \in P$ para todo $k < n$, entonces $n \in P$. Demuestra que $P = \mathbb{N}$ usando el Teorema~\ref{teo:induct-transfinita}.
  \end{enumerate}

  \item \textbf{(El orden de Cantor-Bendixson.)}
  Sea $(X, \tau)$ un espacio topológico. El \textbf{conjunto derivado} $X'$ de $X$ es el conjunto de todos los puntos de acumulación de $X$. Define inductivamente:
  \[
    X^{(0)} = X, \quad X^{(1)} = X', \quad X^{(\alpha+1)} = \bigl(X^{(\alpha)}\bigr)', \quad X^{(\lambda)} = \bigcap_{\alpha < \lambda} X^{(\alpha)} \text{ para } \lambda \text{ límite.}
  \]
  \begin{enumerate}[label=(\alph*)]
    \item Verifica que para $X = [0,1]$ con la topología usual, $X^{(\alpha)} = [0,1]$ para todo ordinal $\alpha$. ¿Por qué no se estabiliza?
    \item Sea $X = \{1/n \mid n \in \mathbb{N}^+\} \cup \{0\}$ con la topología del subespacio de $\mathbb{R}$. Calcula $X^{(0)}, X^{(1)}, X^{(2)}$ y $X^{(\alpha)}$ para $\alpha \geq 2$.
    \item Explica por qué este proceso de derivación es un ejemplo concreto de recursión transfinita.
  \end{enumerate}

  \item \textbf{(Construcción de un bien-orden explícito.)}
  Sea $A = \mathbb{Q} \cap [0,1]$. El Teorema del Buen Orden garantiza que $A$ puede ser bien ordenado, pero no nos da un bien-orden explícito. No obstante:
  \begin{enumerate}[label=(\alph*)]
    \item Demuestra que existe una biyección $g: \mathbb{N} \to \mathbb{Q} \cap [0,1]$ (es decir, $\mathbb{Q} \cap [0,1]$ es numerable). Puedes usar la diagonalización de Cantor.
    \item Usa $g$ para definir un bien-orden en $\mathbb{Q} \cap [0,1]$: $q \prec q' \iff g^{-1}(q) < g^{-1}(q')$ en $\mathbb{N}$. Demuestra que $(\mathbb{Q} \cap [0,1], \prec)$ es un bien-orden.
    \item ¿Es el bien-orden $\prec$ compatible con el orden usual $\leq$ de los racionales? Justifica.
  \end{enumerate}

  \item \textbf{(Proyecto: el tipo de orden y los ordinales.)}
  Investiga y redacta un resumen de 2--3 páginas sobre los siguientes puntos, anticipando la Clase~4:
  \begin{itemize}[leftmargin=2em]
    \item La definición de \textbf{número ordinal} de von Neumann: $\alpha$ es un ordinal si es un conjunto transitivo bien ordenado por $\in$.
    \item Cómo el Teorema de Tricotomía garantiza que cada clase de isomorfismo de bien-órdenes contiene exactamente un ordinal de von Neumann.
    \item Los primeros ordinales: $0 = \emptyset$, $1 = \{0\}$, $2 = \{0,1\}$, $\ldots$, $\omega = \mathbb{N}$, $\omega+1$, $\omega \cdot 2$, $\omega^2$, $\omega^\omega$, $\varepsilon_0$.
    \item Qué significa que $\omega$ sea el «primer ordinal infinito» y en qué sentido $(\mathbb{N}, <)$ y $(\omega, \in)$ son isomorfos.
  \end{itemize}
  \textbf{Fuentes sugeridas:} \cite{jech2003}, \cite{kunen2011}, \cite{enderton1977}.

\end{enumerate}

% ===========================================================
\section*{Referencias bibliográficas de esta clase}
% ===========================================================

Los temas tratados en esta clase se desarrollan con mayor profundidad y detalle en las siguientes fuentes principales:

\begin{itemize}[leftmargin=2em]
  \item \textbf{Teoría de conjuntos y órdenes (tratamiento avanzado):} \cite{jech2003} contiene la presentación canónica de bien-órdenes, ordinales y el principio de inducción transfinita; \cite{kunen2011} es preciso y exhaustivo en los aspectos formales de la axiomática.

  \item \textbf{Introducción a la teoría de conjuntos:} \cite{enderton1977} ofrece una presentación accesible y rigurosa de relaciones, órdenes y ordinales, ideal para el nivel de este curso; \cite{halmos1960} es clásico e intuitivo.

  \item \textbf{Órdenes lineales y bien-órdenes:} \cite{rosenstein1982} es una monografía dedicada íntegramente a los órdenes lineales, con una discusión profunda de los tipos de orden y los ordinales; \cite{sierpinski1958} cubre con detalle la teoría clásica de los números ordinales y cardinales.

  \item \textbf{Inducción transfinita y recursión:} \cite{devlin1993} presenta la recursión transfinita de forma clara y con aplicaciones; \cite{hrbacek1999} es un texto introductorio muy recomendable para el estudiante que se acerca por primera vez a la teoría transfinita.

  \item \textbf{Contexto histórico:} \cite{cantor1883} es el artículo seminal de Cantor donde se introducen los números ordinales transfinitos; \cite{kanamori2003} ofrece una perspectiva histórica y técnica del desarrollo de la jerarquía de ordinales y la inducción transfinita.
\end{itemize}

% ===========================================================
\chapter{Clase 4: Funciones, Equipotencia e Infinitud — Primeros Resultados de Cantor}
% ===========================================================

La noción de \emph{cardinalidad} —el «tamaño» de un conjunto— adquiere su forma definitiva cuando se formaliza mediante \textbf{biyecciones}. La idea central, debida a Georg Cantor en la década de 1870, es radical en su simplicidad: dos conjuntos tienen el mismo número de elementos si y solo si existe una correspondencia uno a uno entre ellos, independientemente de que los conjuntos sean finitos o infinitos. Esta sesión formaliza dicha idea, estudia sus consecuencias inmediatas —la equipotencia como relación de equivalencia y el Teorema de Cantor–Bernstein–Schröder— y culmina con el descubrimiento más sorprendente de Cantor: existen \emph{distintos tamaños de infinito}, en el sentido de que $\mathbb{R}$ no puede ponerse en correspondencia biyectiva con $\mathbb{N}$.

Todo el desarrollo sigue la tradición de la teoría de conjuntos axiomática ZFC tal como se estableció en la Clase~2. En particular, una función es, por definición, un conjunto de pares ordenados que satisface ciertas condiciones \cite{enderton1977, jech2003, kunen2011}.

% ===========================================================
\section{Funciones: definición formal y tipos}
% ===========================================================

\subsection{Funciones como relaciones}

\begin{definicion}{Función}{funcion}
Una \textbf{función} (o \textbf{aplicación}) de $A$ en $B$ es una relación $f \subseteq A \times B$ tal que:
\begin{enumerate}[leftmargin=2em, label=(\roman*)]
  \item \textbf{Totalidad:} $\forall a \in A,\; \exists b \in B,\; (a,b) \in f$.
  \item \textbf{Unicidad:} $\forall a \in A,\; \forall b, b' \in B,\; \bigl[(a,b) \in f \wedge (a,b') \in f \to b = b'\bigr]$.
\end{enumerate}
Se escribe $f: A \to B$ y, cuando $(a,b) \in f$, se escribe $f(a) = b$. El conjunto $A$ se llama el \textbf{dominio} de $f$ y $B$ el \textbf{codominio}. La \textbf{imagen} (o \textbf{rango}) de $f$ es
\[
  \mathrm{Im}(f) \;=\; f(A) \;=\; \{f(a) \mid a \in A\} \;=\; \{b \in B \mid \exists a \in A,\, f(a) = b\}.
\]
\end{definicion}

\begin{observacion}{Las funciones son conjuntos}{funcconjunto}
En ZFC, toda función $f: A \to B$ es literalmente un conjunto: el conjunto de pares de Kuratowski $\{(a, f(a)) \mid a \in A\}$. La totalidad y la unicidad garantizan que este conjunto «asigna» exactamente un elemento de $B$ a cada elemento de $A$. El conjunto de todas las funciones de $A$ en $B$ se denota $B^A$ o ${}^{A}B$, y su existencia se garantiza mediante el Axioma de Potencia y el Esquema de Separación (pues $B^A \subseteq \mathcal{P}(A \times B)$).
\end{observacion}

\begin{definicion}{Imagen directa e imagen inversa}{imagenes}
Sea $f: A \to B$. Para $S \subseteq A$ y $T \subseteq B$, se definen:
\begin{itemize}[leftmargin=2em]
  \item \textbf{Imagen directa:} $f(S) = \{f(a) \mid a \in S\} \subseteq B$.
  \item \textbf{Preimagen} (o \textbf{imagen inversa}): $f^{-1}(T) = \{a \in A \mid f(a) \in T\} \subseteq A$.
\end{itemize}
\end{definicion}

\begin{ejemplo}{Imágenes directas e inversas}{}
Sea $f: \mathbb{Z} \to \mathbb{Z}$ definida por $f(n) = n^2$.
\begin{enumerate}[leftmargin=2em, label=(\alph*)]
  \item $f(\{-2, -1, 0, 1, 2\}) = \{0, 1, 4\}$.
  \item $f(\mathbb{Z}) = \{0, 1, 4, 9, 16, \ldots\} = \{n^2 \mid n \in \mathbb{Z}\}$, el conjunto de cuadrados perfectos no negativos.
  \item $f^{-1}(\{4\}) = \{-2, 2\}$.
  \item $f^{-1}(\{3\}) = \emptyset$ (pues $3$ no es cuadrado perfecto).
  \item $f^{-1}(\{0, 1, 4\}) = \{-2, -1, 0, 1, 2\}$.
\end{enumerate}
Nótese que $f^{-1}(\{4\})$ tiene dos elementos: $f$ no es inyectiva.
\end{ejemplo}

\begin{proposicion}{Propiedades de las imágenes}{propimag}
Sea $f: A \to B$, $S, S' \subseteq A$ y $T, T' \subseteq B$. Entonces:
\begin{enumerate}[leftmargin=2em, label=(\roman*)]
  \item $f(S \cup S') = f(S) \cup f(S')$.
  \item $f(S \cap S') \subseteq f(S) \cap f(S')$ \quad (la igualdad no siempre se cumple).
  \item $f^{-1}(T \cup T') = f^{-1}(T) \cup f^{-1}(T')$.
  \item $f^{-1}(T \cap T') = f^{-1}(T) \cap f^{-1}(T')$.
  \item $f^{-1}(B \setminus T) = A \setminus f^{-1}(T)$.
  \item $S \subseteq f^{-1}(f(S))$; la igualdad se cumple si $f$ es inyectiva.
  \item $f(f^{-1}(T)) \subseteq T$; la igualdad se cumple si $f$ es sobreyectiva.
\end{enumerate}
\end{proposicion}

\begin{proof}
Demostramos (i) y (ii); las demás son análogas.

\textbf{(i)} Sea $b \in f(S \cup S')$. Entonces existe $a \in S \cup S'$ con $f(a) = b$. Si $a \in S$, entonces $b \in f(S)$; si $a \in S'$, entonces $b \in f(S')$. En cualquier caso, $b \in f(S) \cup f(S')$. La otra inclusión es evidente.

\textbf{(ii)} Si $b \in f(S \cap S')$, existe $a \in S \cap S'$ con $f(a) = b$; entonces $a \in S$ y $a \in S'$, luego $b \in f(S)$ y $b \in f(S')$. Para ver que la inclusión puede ser estricta, tomemos $A = \{1, 2\}$, $B = \{0\}$, $f(1) = f(2) = 0$, $S = \{1\}$, $S' = \{2\}$. Entonces $f(S \cap S') = f(\emptyset) = \emptyset$ pero $f(S) \cap f(S') = \{0\} \cap \{0\} = \{0\}$.
\end{proof}

\subsection{Inyectividad, sobreyectividad y biyectividad}

\begin{definicion}{Función inyectiva}{inyectiva}
Una función $f: A \to B$ es \textbf{inyectiva} (o \textbf{uno a uno}) si
\[
  \forall a, a' \in A,\quad f(a) = f(a') \;\implies\; a = a'.
\]
Equivalentemente, $a \neq a'$ implica $f(a) \neq f(a')$.
\end{definicion}

\begin{definicion}{Función sobreyectiva}{sobreyectiva}
Una función $f: A \to B$ es \textbf{sobreyectiva} (o \textbf{suprayectiva}, o \emph{sobre}) si
\[
  \forall b \in B,\quad \exists a \in A,\quad f(a) = b,
\]
es decir, $\mathrm{Im}(f) = B$.
\end{definicion}

\begin{definicion}{Función biyectiva}{biyectiva}
Una función $f: A \to B$ es \textbf{biyectiva} (o una \textbf{biyección}, o \textbf{correspondencia uno a uno}) si es simultáneamente inyectiva y sobreyectiva.
\end{definicion}

\begin{observacion}{Intuición geométrica}{intuicionbiy}
\begin{itemize}[leftmargin=2em]
  \item \emph{Inyectiva:} distintas entradas producen distintas salidas; «no colapsa elementos».
  \item \emph{Sobreyectiva:} todo elemento del codominio es «alcanzado»; ningún elemento de $B$ queda sin preimagen.
  \item \emph{Biyectiva:} empareja perfectamente cada elemento de $A$ con exactamente uno de $B$ y viceversa; es el concepto preciso de «misma cantidad de elementos».
\end{itemize}
\end{observacion}

\begin{ejemplo}{Clasificación de funciones}{}
\begin{enumerate}[leftmargin=2em, label=(\alph*)]
  \item $f: \mathbb{N} \to \mathbb{N}$, $f(n) = 2n$. Es \textbf{inyectiva} (si $2n = 2m$ entonces $n = m$) pero \textbf{no sobreyectiva} ($1 \notin \mathrm{Im}(f)$).

  \item $g: \mathbb{Z} \to \mathbb{Z}$, $g(n) = n^2$. Es \textbf{no inyectiva} ($g(-2) = g(2) = 4$) y \textbf{no sobreyectiva} ($-1 \notin \mathrm{Im}(g)$).

  \item $h: \mathbb{R} \to \mathbb{R}^+$, $h(x) = e^x$. Es \textbf{inyectiva} (la exponencial es estrictamente creciente) y \textbf{sobreyectiva} (todo $y > 0$ tiene un logaritmo), luego es \textbf{biyectiva}.

  \item La función \textbf{identidad} $\mathrm{id}_A: A \to A$, $\mathrm{id}_A(a) = a$ es siempre biyectiva.

  \item Sea $A = \{1, 2, 3\}$, $B = \{a, b, c\}$. La función $f(1) = a$, $f(2) = a$, $f(3) = b$ es \textbf{no inyectiva} y \textbf{no sobreyectiva} ($c \notin \mathrm{Im}(f)$).

  \item La función \textbf{de Cantor} $\langle \cdot, \cdot \rangle: \mathbb{N} \times \mathbb{N} \to \mathbb{N}$, dada por
  \[
    \langle m, n \rangle = \frac{(m+n)(m+n+1)}{2} + m,
  \]
  es biyectiva. Esto demuestra que $|\mathbb{N} \times \mathbb{N}| = |\mathbb{N}|$, un hecho aparentemente paradójico.
\end{enumerate}
\end{ejemplo}

\begin{ejemplo}{La función de paridad}{paridad}
Definamos $f: \mathbb{N} \to \{0,1\}$ por $f(n) = n \bmod 2$ (el resto al dividir por 2). Entonces:
\begin{itemize}[leftmargin=2em]
  \item $f$ es \textbf{sobreyectiva}: $f(0) = 0$ y $f(1) = 1$.
  \item $f$ \textbf{no} es inyectiva: $f(0) = f(2) = 0$.
\end{itemize}
Las preimágenes son los conjuntos de pares e impares: $f^{-1}(\{0\}) = \{0, 2, 4, \ldots\}$ y $f^{-1}(\{1\}) = \{1, 3, 5, \ldots\}$.
\end{ejemplo}

\subsection{Composición e inversas}

\begin{definicion}{Composición de funciones}{composicion}
Sean $f: A \to B$ y $g: B \to C$. La \textbf{composición} $g \circ f: A \to C$ se define por
\[
  (g \circ f)(a) = g(f(a)) \qquad \forall a \in A.
\]
Como conjunto: $g \circ f = \{(a, c) \in A \times C \mid \exists b \in B,\, (a,b) \in f \wedge (b,c) \in g\}$.
\end{definicion}

\begin{proposicion}{Composición y tipo de función}{compoinysupr}
Sea $f: A \to B$ y $g: B \to C$.
\begin{enumerate}[leftmargin=2em, label=(\roman*)]
  \item Si $f$ y $g$ son inyectivas, entonces $g \circ f$ es inyectiva.
  \item Si $f$ y $g$ son sobreyectivas, entonces $g \circ f$ es sobreyectiva.
  \item Si $f$ y $g$ son biyectivas, entonces $g \circ f$ es biyectiva.
  \item Si $g \circ f$ es inyectiva, entonces $f$ es inyectiva.
  \item Si $g \circ f$ es sobreyectiva, entonces $g$ es sobreyectiva.
\end{enumerate}
\end{proposicion}

\begin{proof}
\textbf{(i)} Si $(g \circ f)(a) = (g \circ f)(a')$, entonces $g(f(a)) = g(f(a'))$. Como $g$ es inyectiva, $f(a) = f(a')$. Como $f$ es inyectiva, $a = a'$.

\textbf{(ii)} Sea $c \in C$. Como $g$ es sobreyectiva, $\exists b \in B$ con $g(b) = c$. Como $f$ es sobreyectiva, $\exists a \in A$ con $f(a) = b$. Entonces $(g \circ f)(a) = g(b) = c$.

\textbf{(iii)} Sigue de (i) y (ii).

\textbf{(iv)} Si $f(a) = f(a')$, entonces $g(f(a)) = g(f(a'))$, es decir, $(g \circ f)(a) = (g \circ f)(a')$. Como $g \circ f$ es inyectiva, $a = a'$.

\textbf{(v)} Sea $c \in C$. Como $g \circ f$ es sobreyectiva, $\exists a \in A$ con $(g \circ f)(a) = c$, es decir, $g(f(a)) = c$. Tomando $b = f(a) \in B$, tenemos $g(b) = c$.
\end{proof}

\begin{definicion}{Función inversa}{inversa}
Sea $f: A \to B$ biyectiva. La \textbf{función inversa} $f^{-1}: B \to A$ se define por
\[
  f^{-1}(b) = a \quad \iff \quad f(a) = b.
\]
La unicidad de $a$ está garantizada por la inyectividad de $f$; la existencia, por la sobreyectividad.
\end{definicion}

\begin{proposicion}{Propiedades de la inversa}{propinversa}
Sea $f: A \to B$ biyectiva. Entonces:
\begin{enumerate}[leftmargin=2em, label=(\roman*)]
  \item $f^{-1}$ es biyectiva y $(f^{-1})^{-1} = f$.
  \item $f^{-1} \circ f = \mathrm{id}_A$ y $f \circ f^{-1} = \mathrm{id}_B$.
  \item Si $g: B \to C$ también es biyectiva, entonces $(g \circ f)^{-1} = f^{-1} \circ g^{-1}$.
\end{enumerate}
\end{proposicion}

\begin{proof}
\textbf{(i)} $f^{-1}$ es inyectiva: si $f^{-1}(b) = f^{-1}(b')$, sea $a$ ese valor común; entonces $f(a) = b$ y $f(a) = b'$, luego $b = b'$. $f^{-1}$ es sobreyectiva: para todo $a \in A$, $f(a) \in B$ y $f^{-1}(f(a)) = a$.

\textbf{(ii)} $(f^{-1} \circ f)(a) = f^{-1}(f(a)) = a$ por definición de $f^{-1}$. Análogamente $(f \circ f^{-1})(b) = f(f^{-1}(b)) = b$.

\textbf{(iii)} $(g \circ f)(a) = c \iff g(f(a)) = c \iff f(a) = g^{-1}(c) \iff a = f^{-1}(g^{-1}(c))$.
\end{proof}

% ===========================================================
\section{Conjuntos finitos e infinitos}
% ===========================================================

\subsection{Conjuntos finitos y la numeración por naturales}

\begin{definicion}{Segmentos iniciales de $\mathbb{N}$}{segnaturales}
Para cada $n \in \mathbb{N}$, definimos el \textbf{segmento inicial}
\[
  [n] \;=\; \{k \in \mathbb{N} \mid k < n\} \;=\; \{0, 1, 2, \ldots, n-1\}.
\]
En particular, $[0] = \emptyset$ y $[n]$ tiene exactamente $n$ elementos.
\end{definicion}

\begin{definicion}{Conjunto finito}{conjfinito}
Un conjunto $A$ es \textbf{finito} si existe $n \in \mathbb{N}$ y una biyección $f: [n] \to A$. En ese caso, el \textbf{cardinal} de $A$ es $|A| = n$. Un conjunto que no es finito se dice \textbf{infinito}.
\end{definicion}

\begin{observacion}{El cardinal está bien definido}{cardinalbien}
Para que la Definición~\ref{def:conjfinito} sea coherente, hay que verificar que si existen biyecciones $f: [m] \to A$ y $g: [n] \to A$, entonces $m = n$. Esto es el llamado \textbf{Principio de las Casillas} (o lema de Dedekind–Cantor para conjuntos finitos) y se prueba por inducción sobre $\min(m,n)$. No lo demostraremos en detalle aquí, pero el resultado es esencial para la consistencia de la noción de cardinalidad finita.
\end{observacion}

\begin{proposicion}{Propiedades de los conjuntos finitos}{propfinitos}
\begin{enumerate}[leftmargin=2em, label=(\roman*)]
  \item Todo subconjunto de un conjunto finito es finito.
  \item La unión finita de conjuntos finitos es finita. En particular, si $|A| = m$ y $|B| = n$ y $A \cap B = \emptyset$, entonces $|A \cup B| = m + n$.
  \item Si $|A| = m$ y $|B| = n$, entonces $|A \times B| = m \cdot n$.
  \item Si $|A| = n$, entonces $|\mathcal{P}(A)| = 2^n$.
  \item No existe biyección entre $[m]$ y $[n]$ para $m \neq n$ (principio de las casillas).
\end{enumerate}
\end{proposicion}

\subsection{La definición de Dedekind de conjunto infinito}

La noción de «infinito» puede también caracterizarse desde adentro, sin referencia explícita a $\mathbb{N}$:

\begin{definicion}{Conjunto Dedekind-infinito}{dedinfinito}
Un conjunto $A$ es \textbf{Dedekind-infinito} si existe una biyección $f: A \to B$ donde $B \subsetneq A$ (un subconjunto \textbf{propio} de $A$). Equivalentemente, $A$ es Dedekind-infinito si existe una función inyectiva $f: A \to A$ que \textbf{no} es sobreyectiva.

Un conjunto que no es Dedekind-infinito se llama \textbf{Dedekind-finito}.
\end{definicion}

\begin{observacion}{Origen histórico}{dedekindhistoria}
Esta definición aparece en el trabajo de Richard Dedekind «Was sind und was sollen die Zahlen?» (1888), donde la usa precisamente para \emph{definir} los conjuntos infinitos de manera intrínseca, sin depender de los naturales \cite{dedekind1888}. La idea es que un conjunto es infinito si y solo si tiene «la misma cantidad» de elementos que alguna parte propia suya; algo imposible en los conjuntos finitos.
\end{observacion}

\begin{ejemplo}{$\mathbb{N}$ es Dedekind-infinito}{ndedinfinito}
La función $f: \mathbb{N} \to \mathbb{N}$ dada por $f(n) = n + 1$ es inyectiva (si $n+1 = m+1$ entonces $n = m$) pero no sobreyectiva ($0 \notin \mathrm{Im}(f)$). Su imagen es $\mathbb{N} \setminus \{0\} = \{1, 2, 3, \ldots\} \subsetneq \mathbb{N}$.

Más aún, $f$ establece una biyección entre $\mathbb{N}$ y el subconjunto propio $\{1, 2, 3, \ldots\}$, lo que confirma que $\mathbb{N}$ es Dedekind-infinito. Este fenómeno fue la fuente de la famosa paradoja informal del «Hotel de Hilbert».
\end{ejemplo}

\begin{ejemplo}{El Hotel de Hilbert (ilustración intuitiva)}{hilbert}
Imaginemos un hotel con habitaciones numeradas $0, 1, 2, 3, \ldots$ (una por cada natural), todas ocupadas. Si llega un nuevo huésped:
\begin{itemize}[leftmargin=2em]
  \item Movemos al huésped de la habitación $n$ a la habitación $n+1$ (para todo $n \in \mathbb{N}$).
  \item La habitación $0$ queda libre para el nuevo huésped.
\end{itemize}
Este «movimiento» corresponde exactamente a la función $f(n) = n+1$: es una biyección de $\mathbb{N}$ en $\mathbb{N} \setminus \{0\}$. Formalmente, reacomodamos el hotel estableciendo la correspondencia
\[
  \{0,1,2,3,\ldots\} \;\longleftrightarrow\; \{1,2,3,4,\ldots\},
\]
dejando libre la habitación $0$ para el nuevo cliente. Si llegan infinitos nuevos huéspedes (numerados $0, 1, 2, \ldots$), movemos al huésped de la habitación $n$ a la $2n+1$ (impar), dejando todas las pares libres. Esta es la idea detrás de muchas pruebas de numerabilidad.
\end{ejemplo}

\begin{proposicion}{Conjuntos finitos son Dedekind-finitos}{finDedfinito}
Todo conjunto finito es Dedekind-finito; es decir, ningún conjunto finito es biyectable con un subconjunto propio suyo.
\end{proposicion}

\begin{proof}
Por inducción en $n = |A|$. Si $|A| = 0$, entonces $A = \emptyset$ y no tiene subconjuntos propios. Supongamos el resultado cierto para todos los conjuntos de cardinal menor que $n$, y sea $|A| = n \geq 1$. Supongamos que $f: A \to B$ es una biyección con $B \subsetneq A$. Sea $a_0 \in A \setminus B$. Definimos $A' = A \setminus \{a_0\}$, que tiene $n-1$ elementos. La restricción de $f$ da una inyección de $A$ en $A'$ (pues $B \subseteq A' \cup B \subseteq A$). En particular, $|f(A)| = |A| = n$ pero $f(A) = B \subseteq A' \setminus \{a_0\}$, que tiene a lo sumo $n-1$ elementos. Esto contradice el principio de las casillas. (Una prueba detallada requiere el principio de las casillas, que a su vez se demuestra por inducción.)
\end{proof}

\begin{proposicion}{Relación entre infinitud e infinitud de Dedekind}{infDed}
Bajo el Axioma de Elección (AC), las siguientes condiciones sobre un conjunto $A$ son equivalentes:
\begin{enumerate}[leftmargin=2em, label=(\roman*)]
  \item $A$ es infinito (no finito).
  \item $A$ es Dedekind-infinito.
  \item Existe una función inyectiva $\mathbb{N} \to A$.
\end{enumerate}
Sin el Axioma de Elección, puede existir un conjunto infinito que sea Dedekind-finito. \cite{herrlich2006}
\end{proposicion}

\begin{proof}[Esquema de la prueba]
\textbf{(iii) $\Rightarrow$ (ii):} Si $\iota: \mathbb{N} \to A$ es inyectiva, definimos $f: A \to A$ por
\[
  f(a) = \begin{cases} \iota(\iota^{-1}(a) + 1) & \text{si } a \in \mathrm{Im}(\iota), \\ a & \text{si } a \notin \mathrm{Im}(\iota). \end{cases}
\]
Entonces $f$ es inyectiva y $\iota(0) \notin \mathrm{Im}(f)$, luego $A$ es Dedekind-infinito.

\textbf{(ii) $\Rightarrow$ (iii):} Si $f: A \to A$ es inyectiva y $a_0 \in A \setminus \mathrm{Im}(f)$, definimos $\iota: \mathbb{N} \to A$ por $\iota(0) = a_0$ e $\iota(n+1) = f(\iota(n))$. Por inducción, $\iota$ es inyectiva.

\textbf{(i) $\Rightarrow$ (iii):} Usando AC (de forma controlada), un conjunto infinito $A$ contiene subconjuntos de todo tamaño finito, y se puede construir por inducción una función inyectiva de $\mathbb{N}$ en $A$.

\textbf{(iii) $\Rightarrow$ (i):} Si existiera una biyección $h: [n] \to A$ y una inyección $\iota: \mathbb{N} \to A$, la composición $h^{-1} \circ \iota: \mathbb{N} \to [n]$ sería inyectiva, lo que contradice el principio de las casillas.
\end{proof}

% ===========================================================
\section{Equipotencia como relación de equivalencia}
% ===========================================================

\subsection{Definición y primeras propiedades}

\begin{definicion}{Equipotencia}{equipotencia}
Dos conjuntos $A$ y $B$ son \textbf{equipotentes} (o tienen el mismo \textbf{cardinal}) si existe una biyección $f: A \to B$. Se escribe $A \sim B$ o $|A| = |B|$.
\end{definicion}

\begin{teorema}{La equipotencia es una relación de equivalencia}{equivequipot}
La relación $\sim$ de equipotencia satisface:
\begin{enumerate}[leftmargin=2em, label=(\roman*)]
  \item \textbf{Reflexividad:} $A \sim A$ para todo conjunto $A$.
  \item \textbf{Simetría:} Si $A \sim B$, entonces $B \sim A$.
  \item \textbf{Transitividad:} Si $A \sim B$ y $B \sim C$, entonces $A \sim C$.
\end{enumerate}
\end{teorema}

\begin{proof}
\textbf{(i)} La identidad $\mathrm{id}_A: A \to A$ es una biyección.

\textbf{(ii)} Si $f: A \to B$ es biyectiva, entonces $f^{-1}: B \to A$ también lo es (Proposición~\ref{prop:propinversa}).

\textbf{(iii)} Si $f: A \to B$ y $g: B \to C$ son biyectivas, entonces $g \circ f: A \to C$ es biyectiva (Proposición~\ref{prop:compoinysupr}).
\end{proof}

\begin{observacion}{Clases de equivalencia y cardinalidad}{clasecardinal}
Formalmente, la «clase de equivalencia» de $A$ bajo $\sim$ sería el conjunto $\{B \mid B \sim A\}$. Esto no es un conjunto en ZFC (sería demasiado grande, una clase propia). Por ello, la teoría de cardinales introduce los \textbf{números cardinales} como representantes canónicos: en ZFC, el cardinal de $A$ se define como el ordinal mínimo equipotente a $A$ (ordinal de Hartogs). Para los propósitos de esta clase, trabajaremos con la equipotencia directamente sin necesidad de este refinamiento.
\end{observacion}

\begin{ejemplo}{Biyecciones entre conjuntos infinitos}{}
\begin{enumerate}[leftmargin=2em, label=(\alph*)]
  \item $\mathbb{N} \sim \mathbb{Z}$: la función $f: \mathbb{N} \to \mathbb{Z}$ definida por
  \[
    f(n) = \begin{cases} n/2 & \text{si } n \text{ es par}, \\ -(n+1)/2 & \text{si } n \text{ es impar}, \end{cases}
  \]
  es una biyección. Explícitamente: $f(0) = 0$, $f(1) = -1$, $f(2) = 1$, $f(3) = -2$, $f(4) = 2$, $\ldots$

  \item $\mathbb{N} \sim \mathbb{N} \times \mathbb{N}$: la función de emparejamiento de Cantor $\langle m, n \rangle = \frac{(m+n)(m+n+1)}{2} + m$ es biyectiva. Los primeros valores son:
  \[
    \langle 0,0\rangle = 0,\quad \langle 1,0\rangle = 1,\quad \langle 0,1\rangle = 2,\quad \langle 2,0\rangle = 3,\quad \langle 1,1\rangle = 4,\quad \langle 0,2\rangle = 5,\ldots
  \]

  \item $(0,1) \sim \mathbb{R}$: la función $f: (0,1) \to \mathbb{R}$, $f(x) = \tan\!\bigl(\pi(x - \tfrac{1}{2})\bigr)$, es una biyección.

  \item $[0,1] \sim (0,1)$: puede construirse una biyección aunque los dos intervalos no sean idénticos. Una construcción explícita es:
  \[
    f(x) = \begin{cases}
      1/(n+2) & \text{si } x = 1/n \text{ para algún } n \in \mathbb{N}^+, \\
      0 & \text{si } x = 0, \\
      x & \text{en otro caso.}
    \end{cases}
  \]
  Verificar que $f: [0,1] \to (0,1)$ es biyectiva es un buen ejercicio.
\end{enumerate}
\end{ejemplo}

\subsection{Comparación de cardinales: la relación $\leq$}

\begin{definicion}{Comparación de cardinales}{cardleq}
Se dice que $|A| \leq |B|$ (el cardinal de $A$ es \textbf{menor o igual} que el de $B$) si existe una función inyectiva $f: A \to B$. Se escribe $A \preceq B$.

Se dice que $|A| < |B|$ (el cardinal de $A$ es \textbf{estrictamente menor}) si $A \preceq B$ pero $A \not\sim B$.
\end{definicion}

\begin{observacion}{Reflexividad y transitividad de $\preceq$}{reflextransprec}
La relación $\preceq$ es reflexiva (la identidad es inyectiva) y transitiva (la composición de inyecciones es inyectiva). La antisimetría —que $A \preceq B$ y $B \preceq A$ impliquen $A \sim B$— es precisamente el contenido del Teorema de Cantor–Bernstein–Schröder que demostraremos en la siguiente sección.
\end{observacion}

% ===========================================================
\section{Teorema de Cantor–Bernstein–Schröder}
% ===========================================================

\subsection{Enunciado}

\begin{teorema}{Cantor–Bernstein–Schröder}{CBS}
Sean $A$ y $B$ conjuntos. Si existe una inyección $f: A \to B$ y una inyección $g: B \to A$, entonces existe una biyección $h: A \to B$. En notación de cardinales:
\[
  |A| \leq |B| \wedge |B| \leq |A| \;\implies\; |A| = |B|.
\]
\end{teorema}

\begin{observacion}{Historia del teorema}{historiaCBS}
Este resultado fue enunciado por Cantor pero no demostrado satisfactoriamente por él. La primera demostración correcta es debida a Felix Bernstein (1898, en su tesis doctoral) y a Ernst Schröder (1896, aunque con un error subsanado). Existen múltiples demostraciones; presentamos la clásica debida a la idea de los «cadenas de Bernstein» \cite{jech2003, halmos1960, hrbacek1999}.
\end{observacion}

\subsection{Demostración}

\begin{proof}[Demostración del Teorema~\ref{teo:CBS}]
Sean $f: A \to B$ y $g: B \to A$ inyecciones. Construiremos una biyección $h: A \to B$.

\medskip
\textbf{Paso 1: Definición de las cadenas.}

Dado $x \in A \cup B$, definimos la \textbf{cadena} de $x$ como la sucesión (finita o infinita hacia atrás) obtenida aplicando $f^{-1}$ y $g^{-1}$ alternativamente, en tanto sea posible:
\[
  \ldots \to g^{-1}(x) \to f^{-1}(g^{-1}(x)) \to \cdots \to x.
\]
Más precisamente: si $x \in A$, entonces (si existe) $f(x) \in B$, y (si existe) $g^{-1}(x) \in B$, luego (si existe) $f^{-1}(g^{-1}(x)) \in A$, y así sucesivamente. Cada elemento de $A \cup B$ pertenece a exactamente una cadena.

\medskip
\textbf{Paso 2: Tipos de cadenas.}

Cada cadena es de uno de tres tipos:
\begin{itemize}[leftmargin=2em]
  \item \textbf{Tipo A:} la cadena es infinita y el «primer» elemento (el de más a la izquierda) pertenece a $A$.
  \item \textbf{Tipo B:} la cadena es infinita y el primer elemento pertenece a $B$.
  \item \textbf{Tipo $\infty$:} la cadena es doblemente infinita (sin extremo izquierdo).
\end{itemize}

\medskip
\textbf{Paso 3: Definición de $h$.}

Definimos $h: A \to B$ por:
\[
  h(x) = \begin{cases}
    f(x) & \text{si } x \text{ está en una cadena de tipo A o de tipo } \infty, \\
    g^{-1}(x) & \text{si } x \text{ está en una cadena de tipo B.}
  \end{cases}
\]

\medskip
\textbf{Paso 4: $h$ está bien definida.}

Si $x$ está en una cadena de tipo B, entonces $x \in g(\mathrm{Im}(f)) \subseteq \mathrm{Im}(g)$, y por ello $g^{-1}(x)$ existe. En los otros casos, usamos $f(x)$, que siempre está definido.

\medskip
\textbf{Paso 5: $h$ es inyectiva.}

Elementos en cadenas distintas tienen imágenes distintas porque $f$ y $g^{-1}$ son inyecciones. Dentro de la misma cadena, distintos elementos de $A$ producen distintas imágenes por la misma razón.

\medskip
\textbf{Paso 6: $h$ es sobreyectiva.}

Sea $y \in B$. Analizamos el tipo de la cadena de $y$:
\begin{itemize}[leftmargin=2em]
  \item Si $y$ está en una cadena de tipo A o $\infty$: en esa cadena, el predecesor inmediato de $y$ (por $f$) pertenece a $A$; llámalo $x = f^{-1}(y)$. Entonces $h(x) = f(x) = y$.
  \item Si $y$ está en una cadena de tipo B: $g(y) \in A$ está en la misma cadena, y $h(g(y)) = g^{-1}(g(y)) = y$.
\end{itemize}
En todo caso, $y$ tiene preimagen bajo $h$.
\end{proof}

\begin{corolario}{$\preceq$ es antisimétrica}{prec-antisim}
La relación $\preceq$ (de comparación de cardinales por inyecciones) satisface la antisimetría: $A \preceq B$ y $B \preceq A$ implican $A \sim B$.
\end{corolario}

\begin{ejemplo}{Aplicación del Teorema CBS}{aplicCBS}
\begin{enumerate}[leftmargin=2em, label=(\alph*)]
  \item \textbf{$[0,1] \sim (0,1)$:} Tenemos la inyección obvia $\iota: (0,1) \hookrightarrow [0,1]$ (inclusión). Para la otra dirección, la función $f: [0,1] \to (0,1)$ definida por $f(x) = \frac{x+1}{3}$ es inyectiva (pues $\frac{x+1}{3} \in (\frac{1}{3}, \frac{2}{3}) \subset (0,1)$). Por CBS, $[0,1] \sim (0,1)$.

  \item \textbf{$\mathbb{N} \sim 2\mathbb{N}$:} Tenemos $\iota: 2\mathbb{N} \hookrightarrow \mathbb{N}$ (inclusión) y $f: \mathbb{N} \to 2\mathbb{N}$, $f(n) = 2n$ (inyectiva). Por CBS, $\mathbb{N} \sim 2\mathbb{N}$ (también demostrable directamente).

  \item \textbf{$\mathcal{P}(\mathbb{N}) \sim \mathbb{R}$:} Se construyen inyecciones en ambas direcciones. Una dirección: la representación binaria de reales en $(0,1)$ da una inyección de $(0,1)$ (salvo un conjunto numerable) en $\{0,1\}^{\mathbb{N}} \sim \mathcal{P}(\mathbb{N})$; la otra: la representación binaria da una inyección de $\mathcal{P}(\mathbb{N})$ en $[0,1]$. Los detalles se completan con CBS.
\end{enumerate}
\end{ejemplo}

\subsection{Variante constructiva: demostración de Knaster-Tarski}

Existe otra elegante demostración del Teorema CBS que utiliza el Teorema del Punto Fijo de Knaster–Tarski:

\begin{proof}[Demostración alternativa (Knaster--Tarski)]
Sean $f: A \to B$ y $g: B \to A$ inyecciones. Consideremos la función $\Phi: \mathcal{P}(A) \to \mathcal{P}(A)$ definida por
\[
  \Phi(X) \;=\; A \setminus g\!\bigl(B \setminus f(X)\bigr).
\]
Se verifica que $\Phi$ es \textbf{monótona}: si $X \subseteq X'$, entonces $f(X) \subseteq f(X')$, luego $B \setminus f(X) \supseteq B \setminus f(X')$, luego $g(B \setminus f(X)) \supseteq g(B \setminus f(X'))$, luego $\Phi(X) \subseteq \Phi(X')$.

Por el Teorema del Punto Fijo de Knaster–Tarski (el retículo $\mathcal{P}(A)$ es completo), $\Phi$ tiene un \textbf{punto fijo}: existe $C \subseteq A$ con $\Phi(C) = C$, es decir,
\[
  C = A \setminus g\!\bigl(B \setminus f(C)\bigr).
\]
Definimos $h: A \to B$ por
\[
  h(a) = \begin{cases} f(a) & \text{si } a \in C, \\ g^{-1}(a) & \text{si } a \in A \setminus C. \end{cases}
\]
De la ecuación del punto fijo: $A \setminus C = g(B \setminus f(C))$, de modo que $g^{-1}$ está bien definida en $A \setminus C$ con imagen $B \setminus f(C)$. Se comprueba que $h$ es biyectiva: su imagen es $f(C) \cup (B \setminus f(C)) = B$.
\end{proof}

% ===========================================================
\section{Conjuntos numerables y no numerables}
% ===========================================================

\subsection{Conjuntos numerables}

\begin{definicion}{Conjunto numerable}{numerable}
Un conjunto $A$ es:
\begin{itemize}[leftmargin=2em]
  \item \textbf{Finito} si $A \sim [n]$ para algún $n \in \mathbb{N}$.
  \item \textbf{Infinito numerable} (o \textbf{contablemente infinito}) si $A \sim \mathbb{N}$.
  \item \textbf{Numerable} (o \textbf{a lo sumo numerable}) si es finito o infinito numerable, es decir, si $|A| \leq |\mathbb{N}|$.
  \item \textbf{No numerable} (o \textbf{incontable}) si $|\mathbb{N}| < |A|$, es decir, si $A$ no es numerable.
\end{itemize}
\end{definicion}

\begin{observacion}{Terminología alternativa}{terminonumerable}
En la literatura en inglés: \emph{countably infinite} para $A \sim \mathbb{N}$, \emph{countable} para a lo sumo numerable, \emph{uncountable} para no numerable. En español se usa también «contable» como sinónimo de «numerable». El cardinal de $\mathbb{N}$ se denota $\aleph_0$ (aleph-cero), el primero de los \textbf{números cardinales transfinitos} introducidos por Cantor.
\end{observacion}

\begin{teorema}{Caracterizaciones de la numerabilidad}{carnum}
Para un conjunto infinito $A$, las siguientes condiciones son equivalentes:
\begin{enumerate}[leftmargin=2em, label=(\roman*)]
  \item $A \sim \mathbb{N}$.
  \item Existe una enumeración de $A$: una sucesión $a_0, a_1, a_2, \ldots$ tal que $A = \{a_0, a_1, a_2, \ldots\}$ y todos los $a_i$ son distintos.
  \item Existe una función sobreyectiva $\mathbb{N} \to A$.
  \item Existe una función inyectiva $A \to \mathbb{N}$.
\end{enumerate}
\end{teorema}

\begin{proof}
\textbf{(i) $\Leftrightarrow$ (ii):} Una biyección $f: \mathbb{N} \to A$ es exactamente una enumeración $a_n = f(n)$.

\textbf{(i) $\Rightarrow$ (iii):} Toda biyección es sobreyectiva.

\textbf{(iii) $\Rightarrow$ (i):} Sea $s: \mathbb{N} \to A$ sobreyectiva. Definamos $f: A \to \mathbb{N}$ por $f(a) = \min\{n \in \mathbb{N} \mid s(n) = a\}$. Esta $f$ es inyectiva. Como $A$ es infinito, $f(A)$ es un subconjunto infinito de $\mathbb{N}$, y todo subconjunto infinito de $\mathbb{N}$ es equipotente a $\mathbb{N}$ (por la función que enumera sus elementos en orden creciente). Componiendo, $A \sim \mathbb{N}$.

\textbf{(i) $\Rightarrow$ (iv):} Toda biyección $A \to \mathbb{N}$ es inyectiva.

\textbf{(iv) $\Rightarrow$ (i):} Si $\iota: A \hookrightarrow \mathbb{N}$ es inyectiva, entonces $A \sim \iota(A) \subseteq \mathbb{N}$. Como $A$ es infinito y $\iota(A)$ también lo es, $\iota(A)$ es un subconjunto infinito de $\mathbb{N}$, equipotente a $\mathbb{N}$.
\end{proof}

\begin{proposicion}{Todo subconjunto infinito de $\mathbb{N}$ es numerable}{subinfN}
Si $S \subseteq \mathbb{N}$ es infinito, entonces $S \sim \mathbb{N}$.
\end{proposicion}

\begin{proof}
Definimos $f: \mathbb{N} \to S$ por recursión: $f(0) = \min S$ y $f(n+1) = \min(S \setminus \{f(0), \ldots, f(n)\})$. Como $S$ es infinito, este mínimo siempre existe. La función $f$ es estrictamente creciente (y en particular inyectiva) y su imagen es $S$ (si algún $s \in S$ no estuviera en $\mathrm{Im}(f)$, habría una cantidad infinita de elementos de $S$ menores que $s$, lo que es imposible). Luego $f$ es una biyección de $\mathbb{N}$ en $S$.
\end{proof}

\begin{teorema}{Numerabilidad de $\mathbb{Z}$, $\mathbb{Q}$ y uniones numerables}{numZQ}
\begin{enumerate}[leftmargin=2em, label=(\roman*)]
  \item $\mathbb{Z}$ es numerable.
  \item $\mathbb{N} \times \mathbb{N}$ es numerable.
  \item Si $A$ y $B$ son numerables, $A \cup B$ es numerable.
  \item Una unión numerable de conjuntos numerables es numerable: si $(A_n)_{n \in \mathbb{N}}$ es una familia de conjuntos numerables, $\bigcup_{n \in \mathbb{N}} A_n$ es numerable.
  \item $\mathbb{Q}$ es numerable.
\end{enumerate}
\end{teorema}

\begin{proof}
\textbf{(i)} Ya construimos explícitamente la biyección $f: \mathbb{N} \to \mathbb{Z}$.

\textbf{(ii)} La función de emparejamiento de Cantor $\langle m, n \rangle: \mathbb{N} \times \mathbb{N} \to \mathbb{N}$ es biyectiva. Una segunda prueba es por el argumento de la diagonal: enumeramos $\mathbb{N} \times \mathbb{N}$ recorriendo las diagonales $\{(m,n) \mid m + n = k\}$ para $k = 0, 1, 2, \ldots$:
\[
  (0,0),\ (1,0),\ (0,1),\ (2,0),\ (1,1),\ (0,2),\ (3,0),\ \ldots
\]

\textbf{(iii)} Sean $a_0, a_1, a_2, \ldots$ y $b_0, b_1, b_2, \ldots$ enumeraciones de $A$ y $B$ (sin perder generalidad, suponemos $A \cap B = \emptyset$). Entonces $a_0, b_0, a_1, b_1, a_2, b_2, \ldots$ es una enumeración de $A \cup B$.

\textbf{(iv)} Para cada $n$, sea $a_{n,0}, a_{n,1}, a_{n,2}, \ldots$ una enumeración de $A_n$. Definimos una enumeración de $\bigcup A_n$ recorriendo la tabla $(a_{n,k})_{n,k \in \mathbb{N}}$ en diagonal (método de Cantor): ordenamos los pares $(n,k)$ según $n+k$ y luego por $n$, y enumeramos $a_{n,k}$ en ese orden, omitiendo repeticiones. Por (ii), esto da una enumeración.

\textbf{(v)} Tenemos $\mathbb{Q} = \bigcup_{n \in \mathbb{Z}} \{p/n \mid p \in \mathbb{Z}, \gcd(p,n) = 1\}$ (unión numerable de conjuntos numerables). Por (iv), $\mathbb{Q}$ es numerable.

Una demostración más directa: la función $f: \mathbb{Z} \times \mathbb{N}^+ \to \mathbb{Q}$, $f(p,q) = p/q$, es sobreyectiva. Como $\mathbb{Z} \times \mathbb{N}^+$ es numerable (por (ii) y (i)), $\mathbb{Q}$ lo es.
\end{proof}

\begin{ejemplo}{El argumento diagonal para $\mathbb{N} \times \mathbb{N}$}{diagonalNN}
El argumento diagonal de Cantor para $\mathbb{N} \times \mathbb{N}$ puede visualizarse como una tabla infinita:
\[
\begin{array}{cccccc}
(0,0) & (0,1) & (0,2) & (0,3) & (0,4) & \cdots \\
(1,0) & (1,1) & (1,2) & (1,3) & \cdots  & \\
(2,0) & (2,1) & (2,2) & \cdots  &  & \\
(3,0) & (3,1) & \cdots  &  &  & \\
\vdots &  &  &  &  &
\end{array}
\]
Recorremos las antidiagonales $\{(i,j) \mid i+j = k\}$ para $k = 0, 1, 2, \ldots$, obteniendo la enumeración:
\[
(0,0),\ (1,0),\ (0,1),\ (2,0),\ (1,1),\ (0,2),\ (3,0),\ (2,1),\ (1,2),\ (0,3),\ \ldots
\]
\end{ejemplo}

\begin{corolario}{El conjunto de polinomios con coeficientes enteros es numerable}{polnum}
El conjunto $\mathbb{Z}[x]$ de todos los polinomios con coeficientes enteros es numerable.
\end{corolario}

\begin{proof}
Para cada $n \in \mathbb{N}$, el conjunto $\mathbb{Z}_n[x]$ de polinomios de grado $\leq n$ está en biyección con $\mathbb{Z}^{n+1}$ (los coeficientes), que es numerable. Entonces $\mathbb{Z}[x] = \bigcup_{n \geq 0} \mathbb{Z}_n[x]$ es unión numerable de conjuntos numerables, luego numerable.
\end{proof}

\begin{corolario}{Los números algebraicos son numerables}{algnum}
El conjunto $\mathbb{A}$ de los \textbf{números algebraicos} (raíces de polinomios con coeficientes enteros) es numerable.
\end{corolario}

\begin{proof}
Cada polinomio $p \in \mathbb{Z}[x]$ de grado $n$ tiene a lo sumo $n$ raíces complejas. El conjunto total de raíces es $\mathbb{A} = \bigcup_{p \in \mathbb{Z}[x]} \{\text{raíces de } p\}$, unión numerable de conjuntos finitos, luego numerable.
\end{proof}

\begin{observacion}{Números trascendentes}{trascendentes}
Como veremos en la siguiente sección, $\mathbb{R}$ es no numerable. Dado que $\mathbb{A}$ es numerable y $\mathbb{A} \subset \mathbb{R}$, el conjunto $\mathbb{R} \setminus \mathbb{A}$ de los \textbf{números trascendentes} (como $\pi$ y $e$) es también no numerable; en particular, no puede ser vacío. Esto es una demostración de existencia puramente conjuntista: los trascendentes existen (y son «la mayoría» de los reales) sin necesidad de exhibir ningún ejemplo concreto.
\end{observacion}

% ===========================================================
\section{Conjuntos no numerables: el argumento diagonal de Cantor}
% ===========================================================

\subsection{La no numerabilidad de $\mathbb{R}$}

El resultado siguiente es uno de los teoremas más famosos de las matemáticas:

\begin{teorema}{No numerabilidad de $\mathbb{R}$}{nonum}
El conjunto $\mathbb{R}$ de los números reales no es numerable: $|\mathbb{R}| > |\mathbb{N}|$.
\end{teorema}

\begin{proof}[Demostración (argumento diagonal de Cantor, 1891)]
Basta demostrar que el intervalo $(0,1)$ no es numerable (pues $(0,1) \sim \mathbb{R}$ por la biyección $x \mapsto \tan(\pi(x - \frac{1}{2}))$, y si $(0,1)$ fuera numerable también lo sería $\mathbb{R}$).

Supongamos, para obtener una contradicción, que $(0,1)$ \textbf{es} numerable. Entonces existe una enumeración $x_0, x_1, x_2, \ldots$ de todos los elementos de $(0,1)$. Escribimos cada $x_i$ en su expansión decimal:
\[
  x_i = 0.d_{i,0}\, d_{i,1}\, d_{i,2}\, d_{i,3} \cdots
\]
donde $d_{i,k} \in \{0, 1, \ldots, 9\}$ son los dígitos decimales de $x_i$.

Construimos ahora un número $y = 0.e_0\, e_1\, e_2\, e_3\,\cdots \in (0,1)$ definiendo:
\[
  e_k = \begin{cases} 5 & \text{si } d_{k,k} \neq 5, \\ 6 & \text{si } d_{k,k} = 5. \end{cases}
\]
(O cualquier elección que garantice $e_k \neq d_{k,k}$ y evite los dígitos $0$ y $9$ para esquivar ambigüedades en la representación decimal.)

Afirmamos que $y \in (0,1)$ y que $y \neq x_k$ para todo $k \in \mathbb{N}$. En efecto, $y$ y $x_k$ difieren en el $k$-ésimo dígito decimal (en la posición $k$): $e_k \neq d_{k,k}$. Como los dígitos elegidos están en $\{5, 6\}$, no hay ambigüedad de representación (no se presentan los casos $0.999\ldots = 1.000\ldots$), así que esta diferencia de dígito garantiza $y \neq x_k$.

Pero entonces $y \in (0,1)$ y $y \notin \{x_0, x_1, x_2, \ldots\}$, lo cual contradice que la enumeración sea de todos los elementos de $(0,1)$. Por tanto, $(0,1)$ no es numerable.
\end{proof}

\begin{ejemplo}{Construcción explícita de la diagonal}{diagexplicita}
Ilustramos el argumento con una lista hipotética. Supongamos que la enumeración comienza:
\begin{align*}
  x_0 &= 0.\mathbf{3}14159265\ldots \\
  x_1 &= 0.7\mathbf{1}8281828\ldots \\
  x_2 &= 0.57\mathbf{7}21566\ldots \\
  x_3 &= 0.141\mathbf{5}9265\ldots \\
  x_4 &= 0.2718\mathbf{2}8182\ldots \\
  x_5 &= 0.61803\mathbf{3}9887\ldots \\
  \vdots
\end{align*}
Los dígitos de la diagonal son $d_{0,0} = 3$, $d_{1,1} = 1$, $d_{2,2} = 7$, $d_{3,3} = 5$, $d_{4,4} = 2$, $d_{5,5} = 3$, $\ldots$

Aplicando la regla $e_k = 5$ si $d_{k,k} \neq 5$, $e_k = 6$ si $d_{k,k} = 5$:
\[
  y = 0.\underbrace{5}_{e_0}\underbrace{5}_{e_1}\underbrace{5}_{e_2}\underbrace{6}_{e_3}\underbrace{5}_{e_4}\underbrace{5}_{e_5}\cdots = 0.555655\ldots
\]
Este número $y$ difiere de $x_k$ en al menos el $k$-ésimo dígito para cada $k$, luego $y \neq x_k$ para todo $k$.
\end{ejemplo}

\begin{observacion}{Esencia del argumento diagonal}{esenciadiag}
El argumento diagonal es mucho más general que su aplicación a $\mathbb{R}$: es un método para construir, dada cualquier lista de objetos de un cierto tipo, un objeto del mismo tipo que no aparece en la lista. Sus variantes aparecen en:
\begin{itemize}[leftmargin=2em]
  \item El Teorema de Cantor: $|A| < |\mathcal{P}(A)|$ para todo conjunto $A$ (Sección~\ref{sec:cantorpot}).
  \item La prueba de incompletitud de Gödel.
  \item La indecidibilidad del problema de la parada (halting problem) de Turing.
  \item La construcción de números de Liouville (números trascendentes).
\end{itemize}
\end{observacion}

\subsection{Distintos infinitos: $\aleph_0$ y $\mathfrak{c}$}

\begin{definicion}{Los cardinales $\aleph_0$ y $\mathfrak{c}$}{alephnota}
El \textbf{cardinal} $\aleph_0 = |\mathbb{N}|$ es el primer cardinal infinito. El cardinal $\mathfrak{c} = |\mathbb{R}|$ se llama la \textbf{potencia del continuo}. El Teorema~\ref{teo:nonum} establece que $\aleph_0 < \mathfrak{c}$.
\end{definicion}

\begin{proposicion}{Valor de $\mathfrak{c}$}{valorc}
$\mathfrak{c} = |\mathbb{R}| = |\mathcal{P}(\mathbb{N})| = |2^{\mathbb{N}}| = 2^{\aleph_0}$, donde $2^{\mathbb{N}}$ denota el conjunto de todas las funciones $\mathbb{N} \to \{0,1\}$.
\end{proposicion}

\begin{proof}
Construimos biyecciones $\mathbb{R} \sim \mathcal{P}(\mathbb{N}) \sim \{0,1\}^{\mathbb{N}}$.

\textbf{$\{0,1\}^{\mathbb{N}} \sim \mathcal{P}(\mathbb{N})$:} La biyección es $f \mapsto f^{-1}(\{1\}) = \{n \in \mathbb{N} \mid f(n) = 1\}$, que identifica cada función $\{0,1\}^{\mathbb{N}}$ con su conjunto «soporte».

\textbf{$\mathcal{P}(\mathbb{N}) \preceq [0,1]$:} A cada $S \subseteq \mathbb{N}$ le asociamos el número $\sum_{n \in S} 3^{-(n+1)}$ (expansión en base 3 con dígitos 0 y 1). Esta función es inyectiva (de hecho, su imagen es el conjunto de Cantor ternario).

\textbf{$[0,1] \preceq \mathcal{P}(\mathbb{N})$:} Cada $x \in [0,1]$ tiene una expansión binaria $x = \sum_{n=0}^\infty b_n 2^{-(n+1)}$, $b_n \in \{0,1\}$. La función $x \mapsto \{n \mid b_n = 1\}$ es inyectiva (salvo los números diádicos con dos expansiones, pero se puede elegir siempre la que termina en 0).

Por CBS, $\mathcal{P}(\mathbb{N}) \sim [0,1] \sim \mathbb{R}$.
\end{proof}

% ===========================================================
\section{El Teorema de Cantor: estrictamente más grande}\label{sec:cantorpot}
% ===========================================================

\begin{teorema}{Teorema de Cantor}{cantorpot}
Para todo conjunto $A$, $|A| < |\mathcal{P}(A)|$. Es decir, el conjunto potencia siempre tiene cardinal estrictamente mayor que el conjunto original.
\end{teorema}

\begin{proof}
\textbf{Parte 1: $|A| \leq |\mathcal{P}(A)|$.} La función $a \mapsto \{a\}$ es una inyección de $A$ en $\mathcal{P}(A)$.

\textbf{Parte 2: $|A| \neq |\mathcal{P}(A)|$.} Supongamos que existe una biyección $f: A \to \mathcal{P}(A)$. Definimos el subconjunto \textbf{diagonal}:
\[
  D = \{a \in A \mid a \notin f(a)\}.
\]
Como $D \subseteq A$, tenemos $D \in \mathcal{P}(A)$. Como $f$ es sobreyectiva, existe $a_0 \in A$ con $f(a_0) = D$. Entonces:
\begin{itemize}[leftmargin=2em]
  \item Si $a_0 \in D$, entonces por definición de $D$: $a_0 \notin f(a_0) = D$. Contradicción.
  \item Si $a_0 \notin D$, entonces $a_0 \notin f(a_0) = D$, luego por la definición de $D$: $a_0 \in D$. Contradicción.
\end{itemize}
En cualquier caso se obtiene una contradicción. Por tanto, no puede existir tal biyección.
\end{proof}

\begin{corolario}{Jerarquía infinita de cardinales}{jerarquiacard}
Existe una sucesión estrictamente creciente de cardinales infinitos:
\[
  \aleph_0 \;=\; |\mathbb{N}| \;<\; |\mathcal{P}(\mathbb{N})| \;<\; |\mathcal{P}(\mathcal{P}(\mathbb{N}))| \;<\; \cdots
\]
En particular, no existe un «conjunto de todos los cardinales» ni un «infinito máximo».
\end{corolario}

\begin{ejemplo}{El teorema de Cantor para conjuntos finitos}{cantorfinito}
Para $A = \{1, 2, 3\}$, $|A| = 3$ y $|\mathcal{P}(A)| = 2^3 = 8 > 3$. El subconjunto diagonal en la prueba del Teorema~\ref{teo:cantorpot} depende de la función $f$ elegida. Por ejemplo, si $f(1) = \{2, 3\}$, $f(2) = \{1\}$, $f(3) = \{1, 2, 3\}$, entonces:
\begin{align*}
  D &= \{a \in A \mid a \notin f(a)\}.
\end{align*}
Verificamos cada elemento:
\begin{itemize}[leftmargin=2em]
  \item $1 \notin f(1) = \{2,3\}$: \textbf{verdadero}, luego $1 \in D$.
  \item $2 \notin f(2) = \{1\}$: \textbf{verdadero}, luego $2 \in D$.
  \item $3 \notin f(3) = \{1,2,3\}$: \textbf{falso} ($3 \in \{1,2,3\}$), luego $3 \notin D$.
\end{itemize}
Por tanto $D = \{1, 2\}$. El conjunto $D = \{1, 2\}$ no es ninguno de los valores $f(1), f(2), f(3)$, lo que confirma que $f$ no puede ser sobreyectiva.
\end{ejemplo}

\begin{observacion}{La Hipótesis del Continuo}{hipotesisCH}
El Teorema de Cantor establece $\aleph_0 < \mathfrak{c} = 2^{\aleph_0}$. La pregunta natural es: ¿existe algún cardinal $\kappa$ con $\aleph_0 < \kappa < \mathfrak{c}$? La respuesta es que esta pregunta —conocida como la \textbf{Hipótesis del Continuo} (HC): «no existe tal $\kappa$», i.e., $\mathfrak{c} = \aleph_1$— es \textbf{independiente de ZFC}: Gödel demostró en 1940 que HC es consistente con ZFC, y Cohen demostró en 1963 que su negación también lo es \cite{godel1940, cohen1966}. Este es uno de los resultados más profundos de la lógica matemática del siglo XX.
\end{observacion}

% ===========================================================
\section{Resumen y panorama}
% ===========================================================

\begin{importante}
\textbf{Resumen de los resultados principales de esta clase:}
\begin{enumerate}[leftmargin=2em, label=\textbf{\arabic*.}]
  \item Una función $f: A \to B$ es inyectiva si no colapsa elementos distintos; sobreyectiva si alcanza todo el codominio; biyectiva si es ambas cosas a la vez.
  \item La equipotencia ($A \sim B$: existe una biyección) es una relación de equivalencia.
  \item El Teorema de Cantor–Bernstein–Schröder: $|A| \leq |B|$ y $|B| \leq |A|$ implican $|A| = |B|$.
  \item $\mathbb{Z}$, $\mathbb{Q}$, $\mathbb{Z}[x]$ y $\mathbb{A}$ (los algebraicos) son todos numerables: $|\cdot| = \aleph_0$.
  \item $\mathbb{R}$ no es numerable (argumento diagonal de Cantor): $|\mathbb{R}| = \mathfrak{c} > \aleph_0$.
  \item Para todo $A$: $|A| < |\mathcal{P}(A)|$ (Teorema de Cantor), lo que genera una jerarquía infinita de cardinalidades.
  \item En particular: $\aleph_0 < \mathfrak{c} = 2^{\aleph_0}$.
\end{enumerate}
\end{importante}

La teoría de cardinalidades continuará en la Clase~5, donde introduciremos los números cardinales de von Neumann, la aritmética cardinal (suma, producto y potencia de cardinales) y exploraremos las consecuencias del Axioma de Elección sobre la comparabilidad de cardinales.

% ===========================================================
\section{Problemas y tareas}
% ===========================================================

\begin{enumerate}[leftmargin=2.5em, label=\textbf{\arabic*.}]

  \item \textbf{(Tipos de funciones.)}
  \begin{enumerate}[label=(\alph*)]
    \item Para cada una de las siguientes funciones $f: \mathbb{R} \to \mathbb{R}$, determine si es inyectiva, sobreyectiva, biyectiva, o ninguna de las anteriores. Justifique formalmente.
    \begin{enumerate}[label=(\roman*)]
      \item $f(x) = x^3$.
      \item $f(x) = x^2$.
      \item $f(x) = e^x$.
      \item $f(x) = \sin(x)$.
      \item $f(x) = x^3 - x$.
    \end{enumerate}
    \item Dé un ejemplo de una función $f: \mathbb{N} \to \mathbb{N}$ que sea sobreyectiva pero no inyectiva. ¿Puede ocurrir lo contrario (inyectiva pero no sobreyectiva)?
    \item Si $f: A \to B$ es biyectiva, demuestre directamente desde la definición (sin usar la Proposición~\ref{prop:compoinysupr}) que $f^{-1}: B \to A$ es biyectiva.
  \end{enumerate}

  \item \textbf{(Imágenes directas e inversas.)}
  \begin{enumerate}[label=(\alph*)]
    \item Sea $f: \mathbb{R} \to \mathbb{R}$, $f(x) = x^2$. Calcule:
    \begin{itemize}
      \item $f([-3, 2])$, $f^{-1}([1, 4])$, $f^{-1}(\{-1\})$.
    \end{itemize}
    \item Demuestre que para $f: A \to B$ y $S, S' \subseteq A$:
    \[
      f(S \cap S') = f(S) \cap f(S') \iff f \text{ es inyectiva}.
    \]
    \item Demuestre que $f(f^{-1}(T)) = T$ para todo $T \subseteq B$ si y solo si $f$ es sobreyectiva.
    \item (Desafío) Encuentre conjuntos $A, B$ y una función $f: A \to B$ tal que $f(A \setminus S) \neq f(A) \setminus f(S)$ para algún $S \subseteq A$.
  \end{enumerate}

  \item \textbf{(Dedekind-infinitud.)}
  \begin{enumerate}[label=(\alph*)]
    \item Demuestre que $\mathbb{Z}$ es Dedekind-infinito exhibiendo explícitamente una biyección entre $\mathbb{Z}$ y un subconjunto propio suyo.
    \item Demuestre que si $A$ es Dedekind-infinito y $A \subseteq B$, entonces $B$ es Dedekind-infinito.
    \item Sea $A$ un conjunto y $\iota: \mathbb{N} \to A$ una función inyectiva. Construya explícitamente una función inyectiva $f: A \to A$ que no sea sobreyectiva, demostrando que $A$ es Dedekind-infinito.
    \item (Investigación) Enuncie y explique por qué la equivalencia «infinito $\Leftrightarrow$ Dedekind-infinito» requiere el Axioma de Elección Numerable (una forma debilitada del Axioma de Elección). Consulte \cite{herrlich2006}.
  \end{enumerate}

  \item \textbf{(Equipotencia.)}
  \begin{enumerate}[label=(\alph*)]
    \item Demuestre directamente (sin CBS) que $[0,1] \sim [0,2]$ y que $[0,1] \sim [a,b]$ para cualquier $a < b$.
    \item Demuestre que $(0, +\infty) \sim \mathbb{R}$ construyendo una biyección explícita.
    \item Demuestre que $\mathbb{N} \sim \mathbb{N} \setminus \{0, 1, 2, \ldots, k\}$ para todo $k \in \mathbb{N}$.
    \item Sea $A = \{f: \mathbb{N} \to \{0,1,2\} \mid f \text{ es eventualmente } 0\}$ (funciones de $\mathbb{N}$ en $\{0,1,2\}$ que son $0$ salvo finitos valores). Determine si $A \sim \mathbb{N}$, $A \sim \mathbb{R}$, o ninguna de las dos. Justifique.
    \item Demuestre que $\mathbb{N}^{\mathbb{N}}$ (el conjunto de todas las sucesiones de naturales) cumple $|\mathbb{N}^{\mathbb{N}}| = |\mathbb{R}|$.
  \end{enumerate}

  \item \textbf{(Teorema de Cantor–Bernstein–Schröder.)}
  \begin{enumerate}[label=(\alph*)]
    \item Use CBS para demostrar que $[0,1] \sim [0,1)$ sin construir una biyección explícita. (Sugerencia: construya inyecciones en ambas direcciones.)
    \item Demuestre que el conjunto de los irracionales en $(0,1)$ es equipotente a $\mathbb{R}$.
    \item Sea $A \subseteq B \subseteq C$ con $A \sim C$. Demuestre que $A \sim B \sim C$. (Este es el «Teorema de Dedekind–Bernstein».)
    \item (Desafío) Demuestre que si $f: A \to B$ es sobreyectiva, entonces $|B| \leq |A|$. ¿Requiere AC este resultado?
  \end{enumerate}

  \item \textbf{(Numerabilidad.)}
  \begin{enumerate}[label=(\alph*)]
    \item Demuestre que el conjunto de los números racionales en $(0,1)$ es numerable, construyendo explícitamente una enumeración.
    \item Demuestre que el conjunto de todos los pares $(p, q)$ con $p, q \in \mathbb{Z}$ y $\gcd(p,q) = 1$ es numerable.
    \item Demuestre que el conjunto de todas las sucesiones finitas de $\{0,1\}$ (cadenas binarias finitas, incluyendo la cadena vacía) es numerable.
    \item Demuestre que el conjunto de todos los subconjuntos \emph{finitos} de $\mathbb{N}$ es numerable, pero el de todos los subconjuntos de $\mathbb{N}$ no lo es.
    \item Sea $A = \{r \in \mathbb{R} \mid r \text{ es solución de un polinomio de grado } \leq 5 \text{ con coeficientes en } \mathbb{Z}\}$. Determine si $A$ es numerable.
  \end{enumerate}

  \item \textbf{(Argumento diagonal.)}
  \begin{enumerate}[label=(\alph*)]
    \item Demuestre que el conjunto $\{0,1\}^{\mathbb{N}}$ de todas las sucesiones binarias infinitas no es numerable, usando el argumento diagonal directamente sobre sucesiones (sin hacer referencia a $\mathbb{R}$).
    \item Deduzca a partir del punto (a) que $(0,1)$ no es numerable, usando la representación binaria de los reales. ¿Qué hay que hacer con los números diádicos?
    \item (Paradoja de Berry) Considere la siguiente «definición»: «el menor número natural que no puede describirse en menos de veinte palabras en español». ¿Por qué esta definición es paradójica? Relaciones esto con el argumento diagonal.
    \item (Generalización) Demuestre el siguiente resultado: para cualquier conjunto $A$ y cualquier enumeración $f_0, f_1, f_2, \ldots$ de funciones $f_i: A \to \{0,1\}$, existe una función $g: A \to \{0,1\}$ que no aparece en la lista. ¿Qué condición se necesita sobre $A$?
  \end{enumerate}

  \item \textbf{(Teorema de Cantor: $|A| < |\mathcal{P}(A)|$.)}
  \begin{enumerate}[label=(\alph*)]
    \item Demuestre el Teorema de Cantor para $A = \mathbb{N}$: construya explícitamente un subconjunto $D \subseteq \mathbb{N}$ que no sea la imagen de ningún elemento bajo una función hipotética $f: \mathbb{N} \to \mathcal{P}(\mathbb{N})$ sobreyectiva.
    \item Demuestre que no existe ninguna función sobreyectiva de $A$ en $\mathcal{P}(A)$, sin suponer que $f$ sea inyectiva. (El teorema de Cantor solo requiere la ausencia de sobreyección, no la ausencia de inyección.)
    \item Use el Teorema de Cantor para probar que no existe un «conjunto universal» en ZFC, es decir, no puede existir un conjunto $V$ tal que $A \in V$ para todo conjunto $A$.
    \item Deduzca que la clase propia $\mathbf{V}$ de todos los conjuntos no puede ser un conjunto.
  \end{enumerate}

  \item \textbf{(Proyecto: la Hipótesis del Continuo.)}
  Investigue y redacte un ensayo breve (2--3 páginas) sobre los siguientes puntos:
  \begin{itemize}[leftmargin=2em]
    \item Enuncie la Hipótesis del Continuo: $\nexists \kappa$ con $\aleph_0 < \kappa < \mathfrak{c}$, equivalentemente $\mathfrak{c} = \aleph_1$.
    \item Explique en qué sentido Gödel (1940) demostró que HC es \emph{consistente} con ZFC, usando el constructo de los «conjuntos construibles» $\mathbf{L}$.
    \item Explique en qué sentido Cohen (1963) demostró que la negación de HC también es consistente con ZFC, usando la técnica del \emph{forcing}.
    \item ¿Qué significa que HC sea \emph{independiente} de ZFC? ¿Es esto un defecto de ZFC?
    \item ¿Cuál es el estado actual de la investigación sobre la verdad de HC?
  \end{itemize}
  \textbf{Fuentes sugeridas:} \cite{cohen1966, jech2003, kunen2011, kanamori2003}.

\end{enumerate}

% ===========================================================
\section*{Referencias bibliográficas de esta clase}
% ===========================================================

Los resultados desarrollados en esta clase se encuentran tratados con rigor y detalle en las siguientes fuentes:

\begin{itemize}[leftmargin=2em]
  \item \textbf{Tratamientos estándar de funciones, equipotencia y cardinales:} \cite{enderton1977} ofrece una presentación muy clara y rigurosa de funciones, equipotencia y los primeros cardinales transfinitos, accesible al nivel de este curso; \cite{halmos1960} es un clásico indispensable; \cite{hrbacek1999} es un texto introductorio completo y recomendado.

  \item \textbf{Teoría de conjuntos avanzada:} \cite{jech2003} es la referencia canónica para todos los aspectos de la teoría de cardinales y ordinales; \cite{kunen2011} es preciso y exhaustivo en los aspectos formales.

  \item \textbf{Teorema de Cantor–Bernstein–Schröder:} El resultado aparece en su forma moderna en \cite{enderton1977} y \cite{jech2003}; la demostración vía el Teorema del Punto Fijo de Knaster–Tarski se puede encontrar en \cite{devlin1993} y en muchos textos de álgebra de retículos.

  \item \textbf{Argumento diagonal y no numerabilidad de $\mathbb{R}$:} El artículo original de Cantor es \cite{cantor1891}; una presentación moderna y accesible se encuentra en \cite{enderton1977} y \cite{sierpinski1958}.

  \item \textbf{Conjuntos Dedekind-infinitos y Axioma de Elección:} \cite{herrlich2006} es una monografía dedicada íntegramente a las consecuencias y patologías del Axioma de Elección; \cite{moore1982} ofrece el contexto histórico.

  \item \textbf{Hipótesis del Continuo:} \cite{cohen1966} es el libro clásico de Cohen donde se presenta el forcing; \cite{godel1940} es el trabajo original de Gödel; \cite{kanamori2003} sitúa estos resultados en el panorama más amplio de los cardinales grandes.

  \item \textbf{Definición de Dedekind de conjunto infinito:} \cite{dedekind1888} es la fuente histórica original donde Dedekind define los conjuntos infinitos mediante la reflexividad.
\end{itemize}

% ===========================================================
\chapter{Clase 5: Números Ordinales — Definición, Jerarquía y Aritmética Elemental}
% ===========================================================

Los \textbf{números ordinales} constituyen uno de los logros más profundos de la teoría de conjuntos moderna. Si los números cardinales miden \emph{cuántos} elementos tiene un conjunto, los números ordinales codifican algo más sutil: el \emph{tipo de orden} de un conjunto bien ordenado. Así, mientras que el cardinal $\aleph_0$ mide el «tamaño» de $\mathbb{N}$, el ordinal $\omega$ captura la estructura específica del orden natural de $\mathbb{N}$: un elemento mínimo, cada elemento con un sucesor inmediato, y ningún elemento máximo.

La idea de representar un número ordinal como el conjunto de todos los ordinales que lo preceden fue formulada por \textbf{John von Neumann} en 1923 \cite{vonneumann1923}. Esta representación canónica tiene la virtud de hacer que el orden entre ordinales coincida exactamente con la pertenencia: $\alpha < \beta$ si y solo si $\alpha \in \beta$. El resultado es una clase absolutamente estándar, libremente comparable y axiomáticamente transparente.

Esta clase construye la teoría de los números ordinales desde sus cimientos. Comenzamos con la definición formal de ordinal en ZFC, clasificamos los ordinales en sucesor y límite, introducimos $\omega$ como el primer ordinal transfinito, desarrollamos la inducción y la recursión transfinita en toda su generalidad, y concluimos con la aritmética ordinal, destacando en qué medida difiere de la aritmética de los naturales \cite{jech2003, kunen2011, enderton1977, hrbacek1999}.

% ===========================================================
\section{Ordinales de von Neumann}
% ===========================================================

\subsection{Motivación: tipos de orden bien fundado}

En la Clase~3 se demostró que todo conjunto bien ordenado $(W, <)$ admite un único isomorfismo con su «tipo ordinal»; la tarea ahora es elegir un representante canónico para cada tipo de orden. Informalmente:

\begin{itemize}[leftmargin=2em]
  \item El tipo de orden del conjunto vacío es $0$.
  \item El tipo de orden de un conjunto de un solo elemento es $1$.
  \item El tipo de orden de $\{0, 1, 2, \ldots, n-1\}$ (con el orden usual) es $n$.
  \item El tipo de orden de $\mathbb{N}$ (con el orden usual) es $\omega$.
\end{itemize}

La idea de von Neumann es hacer que cada tipo de orden \emph{sea} el conjunto de todos los tipos de orden que lo preceden. De este modo, el ordinal $3$ es el conjunto $\{0, 1, 2\}$, es decir, el conjunto formado por los tres ordinales anteriores, que son precisamente $0 = \emptyset$, $1 = \{\emptyset\}$ y $2 = \{\emptyset, \{\emptyset\}\}$.

\subsection{Conjuntos transitivos}

\begin{definicion}{Conjunto transitivo}{conjtrans}
Un conjunto $T$ se dice \textbf{transitivo} si todo elemento de un elemento de $T$ es también un elemento de $T$:
\[
  \forall x\,\forall y\;\bigl(x \in y \wedge y \in T \;\to\; x \in T\bigr),
\]
o equivalentemente, $\bigcup T \subseteq T$, es decir, todo elemento de $T$ es a su vez un subconjunto de $T$.
\end{definicion}

\begin{ejemplo}{Conjuntos transitivos}{ejemplotrans}
\begin{enumerate}[leftmargin=2em, label=(\alph*)]
  \item $\emptyset$ es transitivo vacuamente.
  \item $\{\emptyset\}$ es transitivo: su único elemento es $\emptyset$, y $\emptyset \subseteq \{\emptyset\}$.
  \item $\{\emptyset, \{\emptyset\}\}$ es transitivo: $\bigcup\{\emptyset,\{\emptyset\}\} = \emptyset \cup \{\emptyset\} = \{\emptyset\} \subseteq \{\emptyset, \{\emptyset\}\}$.
  \item $\{\{\emptyset\}\}$ \textbf{no} es transitivo: $\emptyset \in \{\emptyset\} \in \{\{\emptyset\}\}$ pero $\emptyset \notin \{\{\emptyset\}\}$.
  \item Todo número natural de von Neumann es transitivo (se demuestra por inducción).
\end{enumerate}
\end{ejemplo}

\subsection{Definición formal de número ordinal}

\begin{definicion}{Número ordinal (von Neumann)}{ordinal}
Un conjunto $\alpha$ es un \textbf{número ordinal} (o simplemente un \textbf{ordinal}) si satisface simultáneamente:
\begin{enumerate}[leftmargin=2em, label=(\roman*)]
  \item $\alpha$ es \textbf{transitivo}: $\forall\beta\,\forall\gamma\;(\gamma \in \beta \in \alpha \to \gamma \in \alpha)$.
  \item $\alpha$ está \textbf{bien ordenado por} $\in$: la relación $\in \upharpoonright \alpha$ es un bien-orden sobre $\alpha$.
\end{enumerate}
Se escribe $\mathbf{ON}$ (o $\mathrm{Ord}$) para denotar la clase de todos los ordinales.
\end{definicion}

\begin{observacion}{Condición equivalente}{ordequiv}
Como consecuencia de la Regularidad y la transitividad, la condición (ii) puede reformularse: $\in$ es una relación de \emph{orden total estricto} sobre $\alpha$ (tricotomía + transitividad + irreflexividad), y todo subconjunto no vacío de $\alpha$ tiene un $\in$-mínimo. En presencia de transitividad, la irreflexividad de $\in$ (es decir, $\alpha \notin \alpha$) se sigue del Axioma de Regularidad, y la relación de orden total se reduce a la tricotomía.
\end{observacion}

\subsection{Los primeros ordinales}

\begin{ejemplo}{Los primeros ordinales de von Neumann}{primerosordes}
Construimos los ordinales finitos de modo explícito:
\begin{align*}
  0 &= \emptyset,\\
  1 &= \{0\} = \{\emptyset\},\\
  2 &= \{0,1\} = \{\emptyset,\{\emptyset\}\},\\
  3 &= \{0,1,2\} = \{\emptyset,\{\emptyset\},\{\emptyset,\{\emptyset\}\}\},\\
  n+1 &= \{0,1,\ldots,n\} = n \cup \{n\}.
\end{align*}
En cada caso, el ordinal $n$ es transitivo (pues sus elementos son los ordinales $0, 1, \ldots, n-1$, cada uno subconjunto de $n$) y $\in$ es un bien-orden sobre $n$ (el orden es $0 \in 1 \in 2 \in \cdots \in n$, y cualquier subconjunto no vacío de $n$ tiene un mínimo).

Nótese que $n = |\{0, 1, \ldots, n-1\}|$, o sea, el cardinal del ordinal $n$ como conjunto es $n$. Esto ilustra la conveniente dualidad ordinal-cardinal para los naturales.
\end{ejemplo}

\subsection{La clase $\mathbf{ON}$ y sus propiedades fundamentales}

\begin{teorema}{Propiedades básicas de los ordinales}{propord}
Sean $\alpha, \beta, \gamma$ ordinales. Entonces:
\begin{enumerate}[leftmargin=2em, label=(\roman*)]
  \item $\alpha \notin \alpha$ \quad (irreflexividad).
  \item Si $\alpha \in \beta$ y $\beta \in \gamma$, entonces $\alpha \in \gamma$ \quad (transitividad de $\in$).
  \item $\alpha \in \beta$, $\alpha = \beta$ o $\beta \in \alpha$ \quad (tricotomía: comparabilidad total).
  \item Si $\alpha \in \beta$, entonces $\alpha \subsetneq \beta$ \quad (todo elemento es subconjunto propio).
  \item Si $A \subseteq \mathbf{ON}$ es un conjunto de ordinales, entonces $\bigcup A$ es un ordinal y $\bigcup A = \sup A$.
  \item Todo subconjunto no vacío de $\mathbf{ON}$ tiene un elemento $\in$-mínimo.
  \item $\mathbf{ON}$ es una clase propia (no es un conjunto).
\end{enumerate}
\end{teorema}

\begin{proof}
Probamos las partes más relevantes.

\textbf{(i)} Por el Axioma de Regularidad, ningún conjunto puede ser elemento de sí mismo: $\alpha \notin \alpha$.

\textbf{(ii)} Si $\alpha \in \beta$ y $\beta \in \gamma$, como $\gamma$ es transitivo, $\alpha \in \gamma$.

\textbf{(iii)} La tricotomía requiere un argumento más elaborado. Se define $\gamma = \alpha \cap \beta$ y se muestra que $\gamma$ es un ordinal, que $\gamma \in \alpha$ o $\gamma = \alpha$, y que $\gamma \in \beta$ o $\gamma = \beta$. Del análisis de los cuatro casos posibles, solo pueden ocurrir tres: $\alpha \in \beta$, $\alpha = \beta$, o $\beta \in \alpha$ (los casos $\alpha \in \beta$ y $\beta \in \alpha$ simultáneos conducirían a $\alpha \in \alpha$ por (ii), violando (i)).

\textbf{(iv)} Si $\alpha \in \beta$, entonces $\alpha \subseteq \beta$ por transitividad de $\beta$ (ya que $\alpha$ es elemento de $\beta$, todo elemento de $\alpha$ está en $\beta$). La inclusión es propia pues $\alpha \in \beta$ pero $\alpha \notin \alpha$.

\textbf{(vii)} Si $\mathbf{ON}$ fuera un conjunto, como es transitivo y bien ordenado por $\in$, sería él mismo un ordinal, luego $\mathbf{ON} \in \mathbf{ON}$, contradiciendo (i). Esta es la \textbf{paradoja de Burali-Forti}.
\end{proof}

\begin{importante}
La identificación $\alpha < \beta \iff \alpha \in \beta$ es la clave de la representación de von Neumann. Gracias a ella, el orden entre ordinales es simplemente la relación de pertenencia, y el ordinal $\alpha$ es \emph{literalmente} el conjunto de todos los ordinales menores que $\alpha$:
\[
  \alpha = \{\beta \in \mathbf{ON} : \beta < \alpha\}.
\]
\end{importante}

% ===========================================================
\section{Ordinales sucesor y ordinales límite}
% ===========================================================

\subsection{El ordinal sucesor}

\begin{definicion}{Ordinal sucesor}{ordsucesor}
Para cada ordinal $\alpha$, su \textbf{sucesor} es el ordinal
\[
  S(\alpha) = \alpha^{+} = \alpha \cup \{\alpha\}.
\]
Un ordinal $\beta$ se llama \textbf{ordinal sucesor} si existe un ordinal $\alpha$ tal que $\beta = S(\alpha)$; en ese caso $\alpha$ se llama el \textbf{predecesor} de $\beta$.
\end{definicion}

\begin{proposicion}{El sucesor es un ordinal}{sucesordinal}
Para todo ordinal $\alpha$, el conjunto $\alpha \cup \{\alpha\}$ es un ordinal y es el menor ordinal estrictamente mayor que $\alpha$.
\end{proposicion}

\begin{proof}
\textbf{Transitividad:} Sea $x \in y \in \alpha \cup \{\alpha\}$. Si $y \in \alpha$, entonces $x \in y \in \alpha$ y por transitividad de $\alpha$, $x \in \alpha \subseteq \alpha \cup \{\alpha\}$. Si $y = \alpha$, entonces $x \in \alpha \subseteq \alpha \cup \{\alpha\}$.

\textbf{Bien-orden por $\in$:} Cualquier subconjunto no vacío $T \subseteq \alpha \cup \{\alpha\}$ es de la forma $T = T' \cup \{\alpha\}$ con $T' \subseteq \alpha$, o bien $T \subseteq \alpha$. Si $T' \neq \emptyset$, el mínimo de $T'$ (que existe porque $\alpha$ es ordinal) es también el mínimo de $T$. Si $T' = \emptyset$, entonces $T = \{\alpha\}$ y su mínimo es $\alpha$.

\textbf{Es el menor sucesor:} Claramente $\alpha \in \alpha \cup \{\alpha\}$, así que $\alpha < S(\alpha)$. Si $\alpha < \gamma < S(\alpha) = \alpha \cup \{\alpha\}$, entonces $\gamma \in \alpha \cup \{\alpha\}$, es decir $\gamma \in \alpha$ o $\gamma = \alpha$. En el primer caso $\gamma < \alpha$, contradicción. En el segundo $\gamma = \alpha$, contradicción con $\alpha < \gamma$.
\end{proof}

\begin{ejemplo}{Los primeros ordinales sucesor}{ej-sucesor}
\begin{align*}
  S(0) &= 0 \cup \{0\} = \emptyset \cup \{\emptyset\} = \{\emptyset\} = 1,\\
  S(1) &= 1 \cup \{1\} = \{\emptyset\} \cup \{\{\emptyset\}\} = \{\emptyset, \{\emptyset\}\} = 2,\\
  S(2) &= 2 \cup \{2\} = \{\emptyset, \{\emptyset\}\} \cup \{\{\emptyset, \{\emptyset\}\}\} = \{\emptyset, \{\emptyset\}, \{\emptyset, \{\emptyset\}\}\} = 3,\\
  S(n) &= n \cup \{n\} = n+1 \quad\text{(para todo natural } n\text{)}.
\end{align*}
\end{ejemplo}

\subsection{Ordinales límite}

\begin{definicion}{Ordinal límite}{ordlimite}
Un ordinal $\lambda$ es un \textbf{ordinal límite} si $\lambda \neq 0$ y $\lambda$ no es un ordinal sucesor, es decir, no existe ningún ordinal $\alpha$ con $\lambda = S(\alpha)$.

Equivalentemente, $\lambda$ es límite si y solo si $\lambda = \bigcup \lambda = \sup\{\alpha : \alpha < \lambda\}$, es decir, $\lambda$ es igual a la unión de todos sus elementos.
\end{definicion}

\begin{observacion}{Tricotomía de ordinales}{tricordinal}
Cada ordinal $\alpha$ pertenece exactamente a una de tres categorías:
\begin{enumerate}[leftmargin=2em, label=(\roman*)]
  \item $\alpha = 0 = \emptyset$: el ordinal cero.
  \item $\alpha = S(\beta)$ para algún ordinal $\beta$: ordinal sucesor.
  \item $\alpha$ es un ordinal límite: $\alpha \neq 0$ y $\alpha = \bigcup \alpha$.
\end{enumerate}
Esta tricotomía es la base de todos los argumentos por inducción y recursión transfinita.
\end{observacion}

\begin{ejemplo}{Identificación de tipos de ordinales}{ej-tipos}
\begin{itemize}[leftmargin=2em]
  \item $0 = \emptyset$: es el ordinal cero.
  \item $1, 2, 3, 4, \ldots, n, \ldots$: todos son ordinales sucesor ($n = S(n-1)$).
  \item $\omega = \{0, 1, 2, 3, \ldots\}$: es el primer ordinal límite (véase la sección siguiente).
  \item $\omega + 1 = S(\omega)$: ordinal sucesor.
  \item $\omega \cdot 2 = \omega + \omega$: ordinal límite.
  \item $\omega^2, \omega^3, \ldots$: ordinales límite.
  \item $\omega^\omega$: ordinal límite.
  \item $\varepsilon_0 = \sup\{\omega, \omega^\omega, \omega^{\omega^\omega}, \ldots\}$: ordinal límite.
\end{itemize}
\end{ejemplo}

% ===========================================================
\section{\texorpdfstring{$\omega$}{omega} como primer ordinal infinito}
% ===========================================================

\subsection{El Axioma del Infinito y la existencia de $\omega$}

Los axiomas ZFC presentados en la Clase~2 garantizan la existencia de conjuntos infinitos mediante el \textbf{Axioma del Infinito}:
\[
  \exists I\;\Bigl(\emptyset \in I \;\wedge\; \forall x \in I\;\bigl(x \cup \{x\} \in I\bigr)\Bigr).
\]
Este axioma asegura que existe un conjunto que contiene $0, 1, 2, 3, \ldots$ (todos los naturales de von Neumann). A partir de él, se define:

\begin{definicion}{El ordinal $\omega$}{defomega}
El ordinal $\omega$ es el menor conjunto inductivo, es decir,
\[
  \omega = \bigcap\{I : I \text{ es inductivo}\} = \bigcap\bigl\{I : \emptyset \in I \;\wedge\; \forall x \in I\,(x \cup \{x\} \in I)\bigr\}.
\]
Equivalentemente, $\omega$ es el conjunto de todos los ordinales finitos (ordinales que se obtienen mediante sucesor a partir de $0$ en un número finito de pasos).
\end{definicion}

\begin{teorema}{$\omega$ es un ordinal límite}{omegaordlimite}
El conjunto $\omega$ es un ordinal límite. Es decir:
\begin{enumerate}[leftmargin=2em, label=(\roman*)]
  \item $\omega$ es un ordinal.
  \item $\omega \neq 0$.
  \item $\omega$ no es sucesor de ningún ordinal.
  \item $\omega = \bigcup \omega$.
\end{enumerate}
\end{teorema}

\begin{proof}
\textbf{(i) Transitividad:} Si $n \in \omega$ (es decir, $n$ es un natural de von Neumann), entonces $n$ es un ordinal finito, luego $n \subseteq \omega$ (todos los elementos de $n$ son naturales, pues $n = \{0,1,\ldots,n-1\}$).

\textbf{(i) Bien-orden:} Todo subconjunto no vacío $S \subseteq \omega$ tiene un elemento con el menor número de predecesores; ese es el mínimo con respecto a $\in$.

\textbf{(ii)} $0 = \emptyset \in \omega$, luego $\omega \neq \emptyset = 0$.

\textbf{(iii)} Si $\omega = S(\alpha) = \alpha \cup \{\alpha\}$ para algún $\alpha \in \omega$, entonces $\alpha \in \omega$ y $S(\alpha) \in \omega$ (por ser $\omega$ inductivo). Pero $\omega = S(\alpha)$ implicaría $\omega \in \omega$, lo que viola la Regularidad.

\textbf{(iv)} $\bigcup \omega = \bigcup_{n \in \omega} n = \omega$ (pues $\bigcup \omega \subseteq \omega$ por transitividad, y para cada $n \in \omega$, $n \in S(n) \subseteq \omega$, luego $n \in \bigcup \omega$).
\end{proof}

\begin{proposicion}{$\omega$ es el menor ordinal infinito}{omegaminimo}
Si $\lambda$ es un ordinal límite, entonces $\omega \leq \lambda$. En particular, si $\alpha < \omega$, entonces $\alpha$ es un ordinal finito.
\end{proposicion}

\begin{proof}
Sea $\lambda$ cualquier ordinal límite. Si $\omega \not\leq \lambda$, entonces $\lambda < \omega$, luego $\lambda \in \omega$, lo que significa que $\lambda$ es un natural de von Neumann y por tanto tiene un predecesor, contradiciendo que $\lambda$ es límite. Luego $\omega \leq \lambda$.
\end{proof}

\subsection{Los naturales como ordinales finitos}

\begin{proposicion}{Identificación con los naturales}{natordfinitos}
Los ordinales finitos son exactamente los elementos de $\omega$: $\{0, 1, 2, 3, \ldots\} = \omega$. En particular, los axiomas de Peano son válidos para $\omega$ con la función sucesor $S(\alpha) = \alpha \cup \{\alpha\}$:
\begin{enumerate}[leftmargin=2em, label=(\roman*)]
  \item $0 \in \omega$.
  \item $\forall n \in \omega,\; S(n) \in \omega$.
  \item $\forall n \in \omega,\; S(n) \neq 0$.
  \item $\forall m, n \in \omega,\; S(m) = S(n) \to m = n$.
  \item \textbf{Inducción:} $\forall A \subseteq \omega\;\bigl[0 \in A \;\wedge\; \forall n \in \omega\,(n \in A \to S(n) \in A)\bigr] \to A = \omega$.
\end{enumerate}
\end{proposicion}

\begin{ejemplo}{La estructura de $\omega$}{estructuraomega}
El ordinal $\omega$ como conjunto es:
\[
  \omega = \{0, 1, 2, 3, 4, \ldots\} = \bigl\{\emptyset,\, \{\emptyset\},\, \{\emptyset,\{\emptyset\}\},\, \{\emptyset,\{\emptyset\},\{\emptyset,\{\emptyset\}\}\},\, \ldots\bigr\}.
\]
El orden dado por $\in$ coincide con el orden usual de los naturales: $0 < 1 < 2 < 3 < \cdots$. El tipo de orden de este conjunto bien ordenado es precisamente $\omega$.

Comparemos $\omega$ con los ordinales que lo suceden:
\[
  \omega < \omega+1 < \omega+2 < \cdots < \omega+\omega < \cdots
\]
donde $\omega + 1 = S(\omega) = \omega \cup \{\omega\}$, etc.
\end{ejemplo}

% ===========================================================
\section{Inducción y Recursión Transfinita}
% ===========================================================

\subsection{Principio de inducción transfinita}

La inducción ordinaria (sobre los naturales) es el caso $\alpha \in \omega$ de la inducción transfinita. La versión transfinita extiende este principio a toda la clase $\mathbf{ON}$.

\begin{teorema}{Inducción transfinita}{inducciontransfinita}
Sea $P(\alpha)$ una propiedad sobre ordinales. Supongamos que para todo ordinal $\alpha$ se cumple:
\[
  \bigl[\,\forall \beta < \alpha,\; P(\beta)\,\bigr] \;\Longrightarrow\; P(\alpha).
\]
(Es decir, si $P$ vale para todos los ordinales menores que $\alpha$, entonces vale para $\alpha$.) Entonces $P(\alpha)$ vale para todo ordinal $\alpha$.
\end{teorema}

\begin{proof}
Supongamos, por reducción al absurdo, que existe algún ordinal para el que $P$ falla. Sea $A = \{\alpha \in \mathbf{ON} : \neg P(\alpha)\} \neq \emptyset$. Como $\mathbf{ON}$ es bien ordenado por $\in$, el conjunto $A$ tiene un elemento $\in$-mínimo $\alpha_0$. Por mínimalidad, $P(\beta)$ vale para todo $\beta < \alpha_0$. La hipótesis del teorema implica entonces $P(\alpha_0)$, contradicción.
\end{proof}

\begin{observacion}{La hipótesis de inducción es fuerte}{hipfuerte}
A diferencia de la inducción ordinaria, donde se asume solo $P(n)$ para demostrar $P(n+1)$, la inducción transfinita asume $P(\beta)$ para \emph{todos} los $\beta < \alpha$. Esta versión se conoce como \textbf{inducción fuerte} (o \textbf{inducción completa}). En el contexto transfinito, la distinción entre inducción débil y fuerte desaparece: la condición «$P$ vale para todos los predecesores» es la única natural, porque los ordinales límite no tienen predecesor inmediato.
\end{observacion}

\subsection{Formulación en casos: sucesor y límite}

En la práctica, la inducción transfinita se aplica verificando tres casos:

\begin{importante}
\begin{teorema}{Inducción transfinita por casos}{inducciontresCasos}
Sea $P(\alpha)$ una propiedad sobre ordinales. $P(\alpha)$ vale para todo ordinal $\alpha$ si y solo si:
\begin{enumerate}[leftmargin=2em, label=(\Roman*)]
  \item \textbf{Caso base:} $P(0)$.
  \item \textbf{Caso sucesor:} Para todo ordinal $\alpha$, si $P(\alpha)$ entonces $P(\alpha+1)$.
  \item \textbf{Caso límite:} Para todo ordinal límite $\lambda$, si $P(\beta)$ para todo $\beta < \lambda$, entonces $P(\lambda)$.
\end{enumerate}
\end{teorema}
\end{importante}

\begin{proof}
La suficiencia se deduce del Teorema~\ref{teo:inducciontransfinita}: si $\alpha$ es un ordinal y $P$ vale para todos los $\beta < \alpha$, distinguimos tres subcasos según la tricotomía. Si $\alpha = 0$, aplicamos (I). Si $\alpha = \beta + 1$ para algún $\beta$, entonces $\beta < \alpha$, luego $P(\beta)$ y por (II), $P(\alpha)$. Si $\alpha$ es límite, (III) da $P(\alpha)$.

La necesidad es trivial: (I) es el caso $\alpha = 0$; (II) es el caso $\alpha$ sucesor con un solo predecesor inmediato; (III) es el caso $\alpha$ límite.
\end{proof}

\begin{ejemplo}{Inducción transfinita: todo ordinal es transitivo}{ejinductrans}
Demostraremos por inducción transfinita que todo ordinal es un conjunto transitivo.

Sea $P(\alpha)$: «$\alpha$ es transitivo».

\textbf{Caso base:} $P(0)$: $\emptyset$ es transitivo vacuamente. \checkmark

\textbf{Caso sucesor:} Asumamos $P(\alpha)$, es decir, $\alpha$ es transitivo. Queremos ver que $\alpha + 1 = \alpha \cup \{\alpha\}$ es transitivo. Sea $x \in y \in \alpha \cup \{\alpha\}$.
\begin{itemize}
  \item Si $y \in \alpha$: como $\alpha$ es transitivo y $x \in y \in \alpha$, se sigue $x \in \alpha \subseteq \alpha \cup \{\alpha\}$.
  \item Si $y = \alpha$: entonces $x \in \alpha \subseteq \alpha \cup \{\alpha\}$.
\end{itemize}
Luego $\alpha \cup \{\alpha\}$ es transitivo. \checkmark

\textbf{Caso límite:} Sea $\lambda$ un ordinal límite y supongamos que $P(\beta)$ vale para todo $\beta < \lambda$. Sea $x \in y \in \lambda = \bigcup_{\beta < \lambda} \beta$. Entonces existe $\beta_0 < \lambda$ con $y \in \beta_0$. Por transitividad de $\beta_0$, $x \in \beta_0 \subseteq \lambda$. Luego $\lambda$ es transitivo. \checkmark

Por el Teorema~\ref{teo:inducciontresCasos}, todo ordinal es transitivo.
\end{ejemplo}

\subsection{Recursión transfinita}

La recursión transfinita es el análogo del principio de definición por recursión sobre los naturales, extendido a toda la clase de ordinales. Permite definir funciones cuyo dominio es $\mathbf{ON}$ (o un segmento inicial de $\mathbf{ON}$) especificando:
\begin{itemize}
  \item el valor en $0$,
  \item el valor en $\alpha + 1$ en función del valor en $\alpha$,
  \item el valor en un límite $\lambda$ en función de los valores en todos los $\beta < \lambda$.

\end{itemize}

\begin{teorema}{Recursión transfinita (forma general)}{recursiongeneral}
Sea $G: V \to V$ un operador de clase (dado por una fórmula de ZFC). Entonces existe una única función de clase $F: \mathbf{ON} \to V$ tal que, para todo ordinal $\alpha$:
\[
  F(\alpha) = G\bigl(F \upharpoonright \alpha\bigr),
\]
donde $F \upharpoonright \alpha$ denota la restricción de $F$ a $\{0, 1, \ldots, \beta, \ldots : \beta < \alpha\}$.
\end{teorema}

\begin{proof}[Esquema de demostración]
La unicidad se establece por inducción transfinita: si $F$ y $F'$ son dos funciones que satisfacen la ecuación de recursión, entonces el conjunto $\{\alpha : F(\alpha) \neq F'(\alpha)\}$ tiene un mínimo $\alpha_0$. Pero $F \upharpoonright \alpha_0 = F' \upharpoonright \alpha_0$ por mínimalidad, luego $F(\alpha_0) = G(F \upharpoonright \alpha_0) = G(F' \upharpoonright \alpha_0) = F'(\alpha_0)$, contradicción.

La existencia se prueba mostrando que para cada ordinal $\alpha$ existe una única función $f_\alpha: \alpha \to V$ con $f_\alpha(\beta) = G(f_\alpha \upharpoonright \beta)$ para todo $\beta < \alpha$. Estas funciones locales son compatibles y su unión define $F$.
\end{proof}

\begin{observacion}{Formulación por casos de la recursión}{recporc}
En la práctica, la recursión transfinita se especifica en tres casos:
\begin{align*}
  F(0) &= a_0 & &\text{(valor inicial)},\\
  F(\alpha + 1) &= h\bigl(F(\alpha)\bigr) & &\text{(paso sucesor, para alguna función } h\text{)},\\
  F(\lambda) &= \sup_{\beta < \lambda} F(\beta) & &\text{(paso límite, para } \lambda \text{ límite)}.
\end{align*}
El paso límite puede tomar otras formas (unión, ínfimo, etc.), según el contexto.
\end{observacion}

\begin{ejemplo}{Recursión transfinita: la jerarquía acumulativa $V_\alpha$}{ejrecursion}
La \textbf{jerarquía acumulativa} de von Neumann se define por recursión transfinita:
\begin{align*}
  V_0 &= \emptyset,\\
  V_{\alpha+1} &= \mathcal{P}(V_\alpha),\\
  V_\lambda &= \bigcup_{\alpha < \lambda} V_\alpha \quad\text{(para } \lambda \text{ límite)}.
\end{align*}
Los primeros niveles son:
\begin{align*}
  V_0 &= \emptyset,\\
  V_1 &= \{\emptyset\},\\
  V_2 &= \{\emptyset, \{\emptyset\}\},\\
  V_3 &= \{\emptyset, \{\emptyset\}, \{\{\emptyset\}\}, \{\emptyset, \{\emptyset\}\}\},\\
  V_\omega &= \bigcup_{n \in \omega} V_n = \text{todos los conjuntos hereditariamente finitos}.
\end{align*}
El \textbf{Teorema de Acumulación}: todo conjunto $x$ pertenece a algún $V_\alpha$ (equivalentemente, $V = \bigcup_{\alpha \in \mathbf{ON}} V_\alpha$). El ordinal mínimo $\alpha$ con $x \in V_\alpha$ es el \textbf{rango} de $x$.
\end{ejemplo}

% ===========================================================
\section{Aritmética Ordinal}
% ===========================================================

La aritmética ordinal define suma, producto y potencia para ordinales mediante recursión transfinita. A diferencia de la aritmética natural, estas operaciones \textbf{no son conmutativas}, como veremos mediante ejemplos concretos.

\subsection{Suma ordinal}

\begin{definicion}{Suma ordinal}{sumaordinal}
Para cada ordinal $\alpha$, la \textbf{suma} $\alpha + \beta$ se define por recursión transfinita en $\beta$:
\begin{align*}
  \alpha + 0 &= \alpha,\\
  \alpha + (\beta + 1) &= (\alpha + \beta) + 1,\\
  \alpha + \lambda &= \sup_{\beta < \lambda}(\alpha + \beta) \quad\text{(para } \lambda \text{ límite)}.
\end{align*}
\end{definicion}

\begin{ejemplo}{Cálculo de sumas ordinales}{ejsumas}
Calculamos algunas sumas con $\omega$.

\textbf{$1 + \omega$:} Por definición, $1 + \omega = \sup_{n < \omega}(1 + n)$.
\begin{align*}
  1 + 0 &= 1,\\
  1 + 1 &= 2,\\
  1 + 2 &= 3,\\
  &\vdots\\
  1 + n &= n+1 \;\text{ para cada } n \in \omega.
\end{align*}
Luego $1 + \omega = \sup\{1, 2, 3, \ldots\} = \omega$.

\textbf{$\omega + 1$:} Por definición directa, $\omega + 1 = (\omega + 0) + 1 = \omega + 1$. Como conjunto, $\omega + 1 = S(\omega) = \omega \cup \{\omega\}$, que contiene a $\omega$ como elemento.

\textbf{Comparación:}
\[
  1 + \omega = \omega \neq \omega + 1.
\]
Por tanto, \textbf{la suma ordinal no es conmutativa}.

\textbf{$\omega + \omega$:} $\omega + \omega = \sup_{n < \omega}(\omega + n) = \sup\{\omega, \omega+1, \omega+2, \ldots\}$. Este ordinal se denota también $\omega \cdot 2$.
\end{ejemplo}

\begin{proposicion}{Propiedades de la suma ordinal}{propsuma}
Para todos los ordinales $\alpha, \beta, \gamma$:
\begin{enumerate}[leftmargin=2em, label=(\roman*)]
  \item \textbf{Asociatividad:} $(\alpha + \beta) + \gamma = \alpha + (\beta + \gamma)$.
  \item \textbf{Elemento neutro:} $\alpha + 0 = 0 + \alpha = \alpha$.
  \item \textbf{Monotonía (derecha):} $\beta < \gamma \implies \alpha + \beta < \alpha + \gamma$.
  \item \textbf{Monotonía (izquierda):} $\beta < \gamma \implies \beta + \alpha \leq \gamma + \alpha$ \quad (puede ser igualdad).
  \item \textbf{No conmutatividad:} En general, $\alpha + \beta \neq \beta + \alpha$.
\end{enumerate}
\end{proposicion}

\begin{proof}
\textbf{(i)} Por inducción transfinita en $\gamma$.

\textit{Base:} $(\alpha+\beta)+0 = \alpha+\beta = \alpha+(\beta+0)$.

\textit{Sucesor:} $(\alpha+\beta)+(\gamma+1) = ((\alpha+\beta)+\gamma)+1 = (\alpha+(\beta+\gamma))+1 = \alpha+((\beta+\gamma)+1) = \alpha+(\beta+(\gamma+1))$.

\textit{Límite:} $(\alpha+\beta)+\lambda = \sup_{\gamma < \lambda}((\alpha+\beta)+\gamma) = \sup_{\gamma < \lambda}(\alpha+(\beta+\gamma)) = \alpha + \sup_{\gamma < \lambda}(\beta+\gamma) = \alpha+(\beta+\lambda)$.

\textbf{(ii)} $\alpha + 0 = \alpha$ por definición. Para $0 + \alpha$, se prueba por inducción transfinita en $\alpha$: $0+0 = 0$; $0+(\alpha+1) = (0+\alpha)+1 = \alpha+1$; $0+\lambda = \sup_{\beta < \lambda}(0+\beta) = \sup_{\beta < \lambda} \beta = \lambda$.

\textbf{(iii) y (v)} son consecuencia directa de las definiciones y se dejan como ejercicio (ver Problemas).
\end{proof}

\begin{ejemplo}{No conmutatividad de la suma: interpretación geométrica}{noconmsuma}
La suma ordinal $\alpha + \beta$ puede visualizarse como yuxtaposición de tipos de orden: primero se ponen los elementos con tipo de orden $\alpha$ y luego los de tipo $\beta$.

\begin{itemize}[leftmargin=2em]
  \item $1 + \omega$: un punto, seguido de una cadena de tipo $\omega$. El punto queda absorbido «al principio»; el tipo resultante es indistinguible de $\omega$. (No hay elemento máximo, y el conjunto es isomorfo a $\mathbb{N}$.)
  \item $\omega + 1$: una cadena de tipo $\omega$ seguida de un punto extra. Ahora el tipo resultante tiene un elemento máximo, lo que lo hace distinto de $\omega$.
\end{itemize}

Así, $1 + \omega = \omega < \omega + 1$: añadir un punto al principio de una cadena infinita no cambia el tipo de orden, pero añadirlo al final sí lo cambia.
\end{ejemplo}

\subsection{Producto ordinal}

\begin{definicion}{Producto ordinal}{productordinal}
Para cada ordinal $\alpha$, el \textbf{producto} $\alpha \cdot \beta$ se define por recursión transfinita en $\beta$:
\begin{align*}
  \alpha \cdot 0 &= 0,\\
  \alpha \cdot (\beta + 1) &= (\alpha \cdot \beta) + \alpha,\\
  \alpha \cdot \lambda &= \sup_{\beta < \lambda}(\alpha \cdot \beta) \quad\text{(para } \lambda \text{ límite)}.
\end{align*}
\end{definicion}

\begin{ejemplo}{Cálculo de productos ordinales}{ejproductos}
\textbf{$2 \cdot \omega$:} $2 \cdot \omega = \sup_{n < \omega}(2 \cdot n) = \sup\{0, 2, 4, 6, \ldots\} = \omega$.

\textbf{$\omega \cdot 2$:} $\omega \cdot 2 = (\omega \cdot 1) + \omega = \omega + \omega$.

\textbf{Comparación:}
\[
  2 \cdot \omega = \omega \neq \omega + \omega = \omega \cdot 2.
\]
Por tanto, \textbf{el producto ordinal no es conmutativo}.

\textbf{$\omega \cdot \omega = \omega^2$:} $\omega^2 = \sup_{n < \omega}(\omega \cdot n)$. Este ordinal tiene el tipo de orden del conjunto $\omega \times \omega$ ordenado lexicográficamente.
\end{ejemplo}

\begin{proposicion}{Propiedades del producto ordinal}{propprod}
Para todos los ordinales $\alpha, \beta, \gamma$:
\begin{enumerate}[leftmargin=2em, label=(\roman*)]
  \item \textbf{Asociatividad:} $(\alpha \cdot \beta) \cdot \gamma = \alpha \cdot (\beta \cdot \gamma)$.
  \item $\alpha \cdot 0 = 0$ y $\alpha \cdot 1 = \alpha$.
  \item \textbf{Distributividad izquierda:} $\alpha \cdot (\beta + \gamma) = (\alpha \cdot \beta) + (\alpha \cdot \gamma)$.
  \item \textbf{Distributividad derecha falla:} En general, $(\alpha + \beta) \cdot \gamma \neq (\alpha \cdot \gamma) + (\beta \cdot \gamma)$.
  \item \textbf{No conmutatividad:} En general, $\alpha \cdot \beta \neq \beta \cdot \alpha$.
  \item \textbf{Anulación:} $0 \cdot \alpha = 0$ para todo $\alpha$.
\end{enumerate}
\end{proposicion}

\begin{proof}
\textbf{(iii)} Por inducción en $\gamma$.

\textit{Base:} $\alpha \cdot (\beta + 0) = \alpha \cdot \beta = (\alpha \cdot \beta) + 0 = (\alpha \cdot \beta) + (\alpha \cdot 0)$.

\textit{Sucesor:} $\alpha \cdot (\beta + (\gamma+1)) = \alpha \cdot ((\beta+\gamma)+1) = (\alpha \cdot (\beta+\gamma)) + \alpha$. Por hipótesis de inducción: $= ((\alpha\cdot\beta)+(\alpha\cdot\gamma))+\alpha = (\alpha\cdot\beta)+((\alpha\cdot\gamma)+\alpha) = (\alpha\cdot\beta)+(\alpha\cdot(\gamma+1))$.

\textit{Límite:} $\alpha \cdot (\beta+\lambda) = \alpha \cdot \sup_{\gamma<\lambda}(\beta+\gamma) = \sup_{\gamma<\lambda}(\alpha \cdot (\beta+\gamma)) = \sup_{\gamma<\lambda}((\alpha\cdot\beta)+(\alpha\cdot\gamma)) = (\alpha\cdot\beta) + \sup_{\gamma<\lambda}(\alpha\cdot\gamma) = (\alpha\cdot\beta)+(\alpha\cdot\lambda)$.

\textbf{(iv)} Contraejemplo: $(1+1)\cdot\omega = 2\cdot\omega = \omega$, pero $(1\cdot\omega)+(1\cdot\omega) = \omega+\omega \neq \omega$.
\end{proof}

\begin{ejemplo}{No conmutatividad del producto: interpretación}{noconmprod}
El producto ordinal $\alpha \cdot \beta$ puede interpretarse como el tipo de orden del conjunto $\beta \times \alpha$ ordenado lexicográficamente (primero por la coordenada izquierda $\beta$, luego por $\alpha$):
\[
  \alpha \cdot \beta \;\cong\; \beta \times \alpha \;\text{ con orden lexicográfico.}
\]
\begin{itemize}[leftmargin=2em]
  \item $2 \cdot \omega$: tipo de orden de $\omega \times 2 = \{(0,0),(0,1),(1,0),(1,1),(2,0),(2,1),\ldots\}$ con orden lexicográfico. El orden es $(0,0),(0,1),(1,0),(1,1),(2,0),(2,1),\ldots$, que es isomorfo a $\omega$.
  \item $\omega \cdot 2$: tipo de orden de $2 \times \omega$ con orden lexicográfico: $(0,0),(0,1),(0,2),\ldots,(1,0),(1,1),(1,2),\ldots$, que es isomorfo a $\omega + \omega$.
\end{itemize}
La diferencia ilustra que el orden en que se «copian» las cadenas importa.
\end{ejemplo}

\subsection{Potencia ordinal}

\begin{definicion}{Potencia ordinal}{potenciaordinal}
Para cada ordinal $\alpha$, la \textbf{potencia} $\alpha^\beta$ se define por recursión transfinita en $\beta$:
\begin{align*}
  \alpha^0 &= 1,\\
  \alpha^{\beta+1} &= \alpha^\beta \cdot \alpha,\\
  \alpha^\lambda &= \sup_{\beta < \lambda} \alpha^\beta \quad\text{(para } \lambda \text{ límite)}.
\end{align*}
\end{definicion}

\begin{ejemplo}{Potencias de $\omega$}{ejpotenciasomega}
\begin{align*}
  \omega^0 &= 1,\\
  \omega^1 &= \omega,\\
  \omega^2 &= \omega \cdot \omega = \sup_{n < \omega} \omega \cdot n = \omega^2\;\text{ (tipo de orden de } \omega \times \omega\text{)},\\
  \omega^3 &= \omega^2 \cdot \omega,\\
  \omega^\omega &= \sup_{n < \omega} \omega^n.
\end{align*}
El ordinal $\omega^\omega$ es el tipo de orden de los polinomios con coeficientes naturales ordenados lexicográficamente. Es un ordinal límite mayor que todos los $\omega^n$ ($n \in \omega$).

La \textbf{primera épsilon-ordinal} $\varepsilon_0$ se define como el menor ordinal $\varepsilon$ tal que $\omega^\varepsilon = \varepsilon$:
\[
  \varepsilon_0 = \sup\bigl\{\omega, \omega^\omega, \omega^{\omega^\omega}, \ldots\bigr\}.
\]
\end{ejemplo}

\subsection{La forma normal de Cantor}

\begin{teorema}{Forma normal de Cantor}{formanormal}
Todo ordinal $\alpha > 0$ se escribe de manera única en la \textbf{forma normal de Cantor}:
\[
  \alpha = \omega^{\beta_1} \cdot c_1 + \omega^{\beta_2} \cdot c_2 + \cdots + \omega^{\beta_k} \cdot c_k,
\]
donde $k \geq 1$, $\alpha \geq \beta_1 > \beta_2 > \cdots > \beta_k$, y $c_1, c_2, \ldots, c_k$ son enteros positivos.
\end{teorema}

\begin{ejemplo}{Forma normal de Cantor}{ejcantor}
\begin{align*}
  5 &= \omega^0 \cdot 5,\\
  \omega &= \omega^1,\\
  \omega + 3 &= \omega^1 \cdot 1 + \omega^0 \cdot 3,\\
  \omega^2 + \omega \cdot 4 + 2 &= \omega^2 \cdot 1 + \omega^1 \cdot 4 + \omega^0 \cdot 2,\\
  \omega^\omega &= \omega^\omega \cdot 1.
\end{align*}
La forma normal de Cantor generaliza la representación decimal de los enteros al reino transfinito. Sin embargo, el exponente $\beta_1$ puede ser en sí mismo un ordinal transfinito; en ese caso la representación no es «finita» en el mismo sentido.
\end{ejemplo}

\subsection{Tabla comparativa: aritmética ordinal vs.\ aritmética natural}

\begin{center}
\renewcommand{\arraystretch}{1.4}
\begin{tabular}{lcc}
\hline
\textbf{Propiedad} & \textbf{$\mathbb{N}$} & \textbf{Ordinales} \\
\hline
Conmutatividad de $+$ & \checkmark & $\times$ \\
Conmutatividad de $\cdot$ & \checkmark & $\times$ \\
Asociatividad de $+$ & \checkmark & \checkmark \\
Asociatividad de $\cdot$ & \checkmark & \checkmark \\
Distributividad izquierda: $a(b+c) = ab+ac$ & \checkmark & \checkmark \\
Distributividad derecha: $(a+b)c = ac+bc$ & \checkmark & $\times$ \\
Cancelación izquierda de $+$: $a+b = a+c \Rightarrow b=c$ & \checkmark & \checkmark \\
Cancelación derecha de $+$: $b+a = c+a \Rightarrow b=c$ & \checkmark & $\times$ \\
Ley de absorción: $n + \omega = \omega$ para $n$ finito & $\times$ & \checkmark \\
\hline
\end{tabular}
\end{center}

\begin{observacion}{Absorción a la izquierda}{absorcion}
La ley $n + \omega = \omega$ para $n \in \omega$ es un fenómeno exclusivamente transfinito: un número finito de pasos queda «absorbido» por la cadena infinita que viene a continuación. En cambio, $\omega + n > \omega$ para $n > 0$, ilustrando la asimetría fundamental de la suma ordinal.
\end{observacion}

% ===========================================================
\section{Segmentos iniciales y tipo de orden}
% ===========================================================

\begin{definicion}{Segmento inicial propio}{segmentoinicial}
Sea $(A, <)$ un conjunto bien ordenado y $a \in A$. El \textbf{segmento inicial propio} determinado por $a$ es
\[
  \text{seg}(A, a) = \{x \in A : x < a\}.
\]
\end{definicion}

\begin{teorema}{Representación de bien-órdenes por ordinales}{represbienord}
Para todo conjunto bien ordenado $(W, <)$ existe un único ordinal $\alpha$ (el \textbf{tipo de orden} de $W$) tal que $(W, <) \cong (\alpha, \in)$. El isomorfismo es también único.
\end{teorema}

\begin{proof}[Esquema]
Se define por recursión transfinita: para cada $w \in W$, se asigna $f(w) = \{f(x) : x < w\}$ (el conjunto de los tipos de orden de los segmentos iniciales propios). Se verifica que $f: W \to \mathbf{ON}$ es un isomorfismo con $\text{Im}(f)$, que resulta ser un ordinal. La unicidad se sigue de que dos isomorfismos entre bien-órdenes deben coincidir.
\end{proof}

\begin{corolario}{Comparabilidad de bien-órdenes}{compbienord}
Dados dos conjuntos bien ordenados $(A, <_A)$ y $(B, <_B)$, exactamente una de las siguientes alternativas se cumple:
\begin{enumerate}[leftmargin=2em, label=(\roman*)]
  \item $(A, <_A) \cong (B, <_B)$,
  \item $(A, <_A) \cong (B', <_B)$ para algún segmento inicial propio $B'$ de $B$,
  \item $(A', <_A) \cong (B, <_B)$ para algún segmento inicial propio $A'$ de $A$.
\end{enumerate}
\end{corolario}

% ===========================================================
\section{Resumen y panorama}
% ===========================================================

\begin{importante}
\textbf{Resumen de los resultados principales de esta clase:}
\begin{enumerate}[leftmargin=2em, label=\textbf{\arabic*.}]
  \item Un \textbf{ordinal} (en el sentido de von Neumann) es un conjunto transitivo bien ordenado por $\in$. El ordinal $\alpha$ es el conjunto de todos los ordinales menores que él: $\alpha = \{\beta : \beta < \alpha\}$.
  \item Los ordinales se clasifican en tres tipos: el ordinal $0 = \emptyset$, los \textbf{ordinales sucesor} $\alpha + 1 = \alpha \cup \{\alpha\}$, y los \textbf{ordinales límite} $\lambda = \bigcup \lambda$ (distintos de 0).
  \item El ordinal $\omega$ es el primer ordinal infinito y el primer ordinal límite. Sus elementos son exactamente los números naturales de von Neumann $0, 1, 2, 3, \ldots$
  \item La \textbf{inducción transfinita} extiende la inducción ordinaria: una propiedad que se preserva bajo sucesor y bajo límites vale para todos los ordinales.
  \item La \textbf{recursión transfinita} permite definir funciones en $\mathbf{ON}$ especificando el valor en $0$, el paso sucesor y el paso límite.
  \item La \textbf{suma ordinal} ($\alpha + \beta$), el \textbf{producto ordinal} ($\alpha \cdot \beta$) y la \textbf{potencia ordinal} ($\alpha^\beta$) se definen por recursión transfinita y generalizan la aritmética natural.
  \item La aritmética ordinal es \textbf{no conmutativa}: $1 + \omega = \omega \neq \omega + 1$; $2 \cdot \omega = \omega \neq \omega + \omega = \omega \cdot 2$.
  \item La aritmética ordinal es \textbf{asociativa}, la suma y el producto son \textbf{distributivos por la izquierda}, pero \textbf{no por la derecha}.
  \item Todo ordinal admite una única \textbf{forma normal de Cantor} $\omega^{\beta_1}c_1 + \cdots + \omega^{\beta_k}c_k$.
\end{enumerate}
\end{importante}

En la Clase~6 abordaremos los \textbf{números cardinales} desde el punto de vista ordinal: definiremos los cardinales como ordinales iniciales (los ordinales que no son equipotentes a ningún ordinal menor), estudiaremos la sucesión de alephs $\aleph_0, \aleph_1, \aleph_2, \ldots$, y desarrollaremos la aritmética cardinal, que tiene su propio conjunto de propiedades, distintas de la aritmética ordinal.

% ===========================================================
\section{Problemas y tareas}
% ===========================================================

\begin{enumerate}[leftmargin=2.5em, label=\textbf{\arabic*.}]

  \item \textbf{(Conjuntos transitivos.)}
  \begin{enumerate}[label=(\alph*)]
    \item Determine cuáles de los siguientes conjuntos son transitivos. Justifique.
    \begin{enumerate}[label=(\roman*)]
      \item $\emptyset$.
      \item $\{\emptyset, \{\emptyset\}\}$.
      \item $\{\{\emptyset\}, \{\{\emptyset\}\}\}$.
      \item $\{\emptyset, \{\emptyset\}, \{\emptyset, \{\emptyset\}\}\}$.
      \item $\mathcal{P}(\mathcal{P}(\emptyset))$.
    \end{enumerate}
    \item Demuestre que la unión de dos conjuntos transitivos es transitiva.
    \item Demuestre que la intersección de dos conjuntos transitivos es transitiva.
    \item Demuestre que $\mathcal{P}(T)$ es transitivo siempre que $T$ lo sea.
    \item (Desafío) Demuestre que para todo conjunto $X$ existe el menor conjunto transitivo que contiene a $X$ (llamado la \emph{clausura transitiva} de $X$, denotada $\mathrm{TC}(X)$). Demuestre que $\mathrm{TC}(X) = X \cup \bigcup X \cup \bigcup\bigcup X \cup \cdots$
  \end{enumerate}

  \item \textbf{(Los primeros ordinales.)}
  \begin{enumerate}[label=(\alph*)]
    \item Escriba explícitamente como conjuntos los ordinales $0, 1, 2, 3, 4$.
    \item Verifique que $4 = \{0, 1, 2, 3\}$ es un conjunto transitivo bien ordenado por $\in$.
    \item Demuestre por inducción (ordinaria, sobre $n \in \mathbb{N}$) que el ordinal $n$ (de von Neumann) tiene exactamente $n$ elementos.
    \item Demuestre que si $\alpha$ es un ordinal, entonces $\alpha \notin \alpha$ (use el Axioma de Regularidad).
    \item Demuestre que si $\alpha$ y $\beta$ son ordinales con $\alpha \subsetneq \beta$, entonces $\alpha \in \beta$.
  \end{enumerate}

  \item \textbf{(Ordinales sucesor y límite.)}
  \begin{enumerate}[label=(\alph*)]
    \item Demuestre que $\alpha \cup \{\alpha\}$ es un ordinal para todo ordinal $\alpha$.
    \item Demuestre que $\omega$ no es un ordinal sucesor. (Sugerencia: si $\omega = \beta \cup \{\beta\}$, entonces $\beta \in \omega$, luego $\beta$ es un natural, lo que llevaría a $\omega = \beta+1$ siendo un natural, absurdo.)
    \item Demuestre directamente desde la definición que $\omega = \bigcup \omega$.
    \item Clasifique los siguientes ordinales como $0$, sucesor o límite, justificando: $\omega, \omega+1, \omega+\omega, \omega\cdot\omega, \omega^{\omega}, \omega^{\omega}+1$.
    \item Demuestre que si $\lambda$ es un ordinal límite, entonces para todo $\alpha < \lambda$ se tiene $\alpha + 1 < \lambda$.
  \end{enumerate}

  \item \textbf{(Inducción transfinita.)}
  \begin{enumerate}[label=(\alph*)]
    \item Demuestre por inducción transfinita que todo ordinal finito es igual a $0$ o es sucesor de otro ordinal finito.
    \item Demuestre por inducción transfinita que para todo ordinal $\alpha$ se tiene $\alpha \subseteq \beta$ o $\beta \subseteq \alpha$ (comparabilidad de subconjuntos de ordinales).
    \item Demuestre por inducción transfinita que $0 + \alpha = \alpha$ para todo ordinal $\alpha$.
    \item Demuestre por inducción transfinita que la función $f: \mathbf{ON} \to \mathbf{ON}$ definida por $f(\alpha) = \omega + \alpha$ es estrictamente creciente.
    \item (Desafío) Demuestre por inducción transfinita que el conjunto de ordinales límite $\leq \alpha$ tiene cardinalidad $|\alpha|$ para todo ordinal infinito $\alpha$.
  \end{enumerate}

  \item \textbf{(Aritmética ordinal — Suma.)}
  \begin{enumerate}[label=(\alph*)]
    \item Calcule, justificando paso a paso:
    \begin{enumerate}[label=(\roman*)]
      \item $3 + \omega$.
      \item $\omega + 3$.
      \item $\omega + \omega$.
      \item $\omega \cdot 2 + \omega$.
      \item $(\omega + 1) + (\omega + 1)$.
    \end{enumerate}
    \item Demuestre que $n + \omega = \omega$ para todo $n \in \omega$ y que $\omega + n > \omega$ para todo $n > 0$.
    \item Demuestre que la suma ordinal es estrictamente monótona por la derecha: $\beta < \gamma \Rightarrow \alpha + \beta < \alpha + \gamma$.
    \item Muestre con un contraejemplo que la suma ordinal \emph{no} es estrictamente monótona por la izquierda, es decir, que $\beta < \gamma$ no implica $\beta + \alpha < \gamma + \alpha$ en general.
    \item Demuestre la ley de cancelación izquierda: $\alpha + \beta = \alpha + \gamma \Rightarrow \beta = \gamma$.
    \item Muestre con un contraejemplo que la cancelación derecha falla: exhiba $\alpha, \beta, \gamma$ con $\beta + \alpha = \gamma + \alpha$ pero $\beta \neq \gamma$.
  \end{enumerate}

  \item \textbf{(Aritmética ordinal — Producto.)}
  \begin{enumerate}[label=(\alph*)]
    \item Calcule, justificando:
    \begin{enumerate}[label=(\roman*)]
      \item $3 \cdot \omega$.
      \item $\omega \cdot 3$.
      \item $\omega \cdot \omega$.
      \item $(\omega + 1) \cdot 2$.
      \item $2 \cdot (\omega + 1)$.
    \end{enumerate}
    \item Demuestre que $n \cdot \omega = \omega$ para todo $n \geq 1$ finito.
    \item Demuestre la distributividad izquierda: $\alpha(\beta + \gamma) = \alpha\beta + \alpha\gamma$.
    \item Muestre con un contraejemplo que la distributividad derecha falla: $(1+1)\omega \neq 1\cdot\omega + 1\cdot\omega$.
    \item Demuestre que $\omega \cdot n = \omega \cdot m \Rightarrow n = m$ para $n, m \in \omega \setminus \{0\}$.
  \end{enumerate}

  \item \textbf{(Forma normal de Cantor.)}
  \begin{enumerate}[label=(\alph*)]
    \item Escriba en forma normal de Cantor los siguientes ordinales:
    \begin{enumerate}[label=(\roman*)]
      \item $\omega^3 + \omega^2 \cdot 5 + \omega \cdot 3 + 7$.
      \item $\omega^\omega + \omega^3 + 1$.
      \item $\omega^{\omega^2}$.
    \end{enumerate}
    \item Dados $\alpha = \omega^2 \cdot 3 + \omega \cdot 2 + 1$ y $\beta = \omega^2 + \omega \cdot 4 + 5$, calcule $\alpha + \beta$ y $\beta + \alpha$ usando la forma normal.
    \item Calcule $\alpha \cdot \omega$ para $\alpha = \omega^2 \cdot 3 + \omega \cdot 2 + 1$.
    \item (Desafío) Demuestre el Teorema de Cantor de la Forma Normal por inducción transfinita sobre $\alpha$.
  \end{enumerate}

  \item \textbf{(Recursión transfinita.)}
  \begin{enumerate}[label=(\alph*)]
    \item La \textbf{sucesión de Veblen} $(\varphi_0, \varphi_1, \varphi_2, \ldots)$ comienza con $\varphi_0(\alpha) = \omega^\alpha$. Calcule $\varphi_0(0), \varphi_0(1), \varphi_0(\omega), \varphi_0(\omega+1)$.
    \item Defina por recursión transfinita la función $f: \omega \to \omega$ dada por $f(0) = 1$, $f(n+1) = f(n) + f(n)$. Calcule $f(0), f(1), \ldots, f(4)$ y determine una fórmula cerrada.
    \item Defina los \textbf{rangos de von Neumann} $V_\alpha$ por recursión transfinita y calcule $V_0, V_1, V_2, V_3, V_\omega$.
    \item Demuestre por inducción transfinita que $V_\alpha \subseteq V_\beta$ siempre que $\alpha \leq \beta$.
    \item (Proyecto) Estudie la definición de la función $\alpha \mapsto \omega_\alpha$ (los cardinales alephs) por recursión transfinita en los ordinales. Enuncie la propiedad característica de $\omega_{\alpha+1}$ (ser el menor cardinal mayor que $\omega_\alpha$) y de $\omega_\lambda$ para $\lambda$ límite.
  \end{enumerate}

  \item \textbf{(Paradoja de Burali-Forti.)}
  \begin{enumerate}[label=(\alph*)]
    \item Enuncie la paradoja de Burali-Forti: si $\mathbf{ON}$ fuera un conjunto, sería un ordinal, luego $\mathbf{ON} \in \mathbf{ON}$, contradiciendo el Axioma de Regularidad.
    \item Explique cómo ZFC resuelve esta paradoja: $\mathbf{ON}$ es una \emph{clase propia}, no un conjunto.
    \item Relacione este fenómeno con la paradoja de Russell: en ambos casos, la «colección demasiado grande» no puede ser un conjunto en ZFC.
  \end{enumerate}

  \item \textbf{(Investigación: ordinales y tipos de orden.)}
  Consulte \cite{sierpinski1958} o \cite{rosenstein1982} y redacte una exposición de 1--2 páginas sobre los siguientes puntos:
  \begin{itemize}[leftmargin=2em]
    \item ¿Cuántos tipos de orden distintos (no isomorfos como conjuntos bien ordenados) existen sobre un conjunto de $n$ elementos?
    \item Describa el tipo de orden $\eta$ del conjunto de los racionales con el orden usual. ¿Es $\eta$ un ordinal? ¿Por qué?
    \item Enuncie el Teorema de Cantor sobre la unicidad del tipo de orden del conjunto denso, sin extremos y numerable.
  \end{itemize}

\end{enumerate}

% ===========================================================
\section*{Referencias bibliográficas de esta clase}
% ===========================================================

Los resultados desarrollados en esta clase se encuentran tratados con rigor y detalle en las siguientes fuentes:

\begin{itemize}[leftmargin=2em]
  \item \textbf{Definición de ordinal y representación de von Neumann:} La fuente histórica original es \cite{vonneumann1923}, en la que von Neumann introduce los ordinales como conjuntos bien ordenados por la pertenencia. Una presentación moderna accesible se encuentra en \cite{enderton1977}; los desarrollos exhaustivos están en \cite{jech2003} y \cite{kunen2011}.

  \item \textbf{Inducción y recursión transfinita:} \cite{enderton1977} contiene una exposición cuidadosa de la recursión transfinita y sus aplicaciones; \cite{hrbacek1999} es un texto introductorio completo; \cite{suppes1972} presenta una formulación axiomática detallada.

  \item \textbf{Aritmética ordinal:} La presentación clásica de la aritmética ordinal y la forma normal de Cantor se encuentra en \cite{sierpinski1958}, que sigue siendo una referencia fundamental; \cite{levy1979} y \cite{jech2003} contienen tratamientos modernos completos.

  \item \textbf{Tipos de orden y bien-órdenes:} \cite{rosenstein1982} es una monografía dedicada a los órdenes lineales, incluyendo los bien-órdenes y sus tipos ordinales; \cite{cantor1883} es el artículo histórico en el que Cantor introduce los ordinales transfinitos.

  \item \textbf{La jerarquía acumulativa $V_\alpha$:} \cite{kunen2011} y \cite{devlin1993} contienen exposiciones detalladas; \cite{jech2003} la desarrolla en el contexto más amplio del universo de von Neumann.

  \item \textbf{Paradoja de Burali-Forti y clases propias:} La paradoja original es de Burali-Forti (1897); una discusión moderna de las clases propias en ZFC se encuentra en \cite{jech2003} y \cite{mendelson2015}.

  \item \textbf{Texto introductorio recomendado:} \cite{halmos1960} proporciona una primera lectura informal y motivadora; \cite{hrbacek1999} y \cite{enderton1977} son los textos de nivel intermedio más recomendados para este curso.
\end{itemize}

% ===========================================================
\chapter{Clase 6: Aritmética Ordinal Avanzada y la Jerarquía de Ordinales Infinitos}
% ===========================================================

La aritmética ordinal elemental —suma, producto y potencia— fue presentada en la Clase~5 mediante definiciones recursivas. En esta clase completamos ese programa con dos grandes logros: la \textbf{demostración completa del Teorema de la Forma Normal de Cantor} y el estudio del universo de ordinales que la potenciación $\alpha \mapsto \omega^\alpha$ genera más allá de $\omega^\omega$. El personaje central es el ordinal $\varepsilon_0$: el primer punto fijo de la función $\alpha \mapsto \omega^\alpha$, un ordinal que ni siquiera la forma normal de Cantor puede representar de manera finita sin aludirse a sí mismo.

Junto con la aritmética avanzada, introducimos la noción de \textbf{cofinalidad}, que permite clasificar los ordinales límite según el «ritmo» con que se aproximan a su propio valor, y que resulta esencial en la teoría cardinal y en la teoría de grandes cardinales \cite{jech2003, kanamori2003}.

La exposición sigue los lineamientos clásicos de \cite{sierpinski1958, enderton1977, hrbacek1999, levy1979}, con una atención especial a los detalles de los argumentos por inducción y recursión transfinita.

% ===========================================================
\section{Potenciación ordinal: propiedades y cálculos avanzados}
% ===========================================================

\subsection{Recordatorio de la definición}

Recordemos que la \textbf{potencia ordinal} $\alpha^\beta$ se define por recursión transfinita en $\beta$:

\begin{definicion}{Potencia ordinal}{potordrecord}
Para ordinales $\alpha$ y $\beta$:
\begin{align*}
  \alpha^0 &= 1, \\
  \alpha^{\beta+1} &= \alpha^\beta \cdot \alpha, \\
  \alpha^\lambda &= \sup_{\beta < \lambda} \alpha^\beta \quad \text{(para } \lambda \text{ límite, } \alpha > 0\text{)}.
\end{align*}
Se establece además $0^\lambda = 0$ para todo $\lambda > 0$.
\end{definicion}

\begin{observacion}{El caso $\alpha = 0$}{cerobase}
La convención $0^\lambda = 0$ para $\lambda > 0$ es consistente con el caso límite: $\sup_{\beta < \lambda} 0^\beta = \sup\{0\} = 0$. El caso $0^0 = 1$ sigue la convención universal que hace la recursión coherente.
\end{observacion}

\subsection{Propiedades fundamentales de la potencia ordinal}

\begin{proposicion}{Leyes de los exponentes ordinales}{leyesexp}
Para todos los ordinales $\alpha, \beta, \gamma$:
\begin{enumerate}[leftmargin=2em, label=(\roman*)]
  \item $\alpha^{\beta+\gamma} = \alpha^\beta \cdot \alpha^\gamma$.
  \item $(\alpha^\beta)^\gamma = \alpha^{\beta \cdot \gamma}$.
  \item Si $1 < \alpha$ y $\beta < \gamma$, entonces $\alpha^\beta < \alpha^\gamma$.
  \item Si $\alpha \leq \beta$ y $\gamma > 0$, entonces $\alpha^\gamma \leq \beta^\gamma$.
  \item $\alpha^1 = \alpha$ y $1^\alpha = 1$.
\end{enumerate}
\end{proposicion}

\begin{proof}
Probamos (i) y (ii) por inducción transfinita; (iii)--(v) siguen de las mismas técnicas.

\textbf{(i)} Inducción transfinita en $\gamma$.

\textit{Base:} $\alpha^{\beta+0} = \alpha^\beta = \alpha^\beta \cdot 1 = \alpha^\beta \cdot \alpha^0$. \checkmark

\textit{Paso sucesor:} Asuma $\alpha^{\beta+\gamma} = \alpha^\beta \cdot \alpha^\gamma$. Entonces
\[
  \alpha^{\beta+(\gamma+1)} = \alpha^{(\beta+\gamma)+1} = \alpha^{\beta+\gamma} \cdot \alpha
  = (\alpha^\beta \cdot \alpha^\gamma) \cdot \alpha = \alpha^\beta \cdot (\alpha^\gamma \cdot \alpha)
  = \alpha^\beta \cdot \alpha^{\gamma+1}. \checkmark
\]

\textit{Paso límite:} Sea $\lambda$ un ordinal límite. Entonces $\beta + \lambda = \sup_{\gamma<\lambda}(\beta+\gamma)$ es también un límite, y
\[
  \alpha^{\beta+\lambda}
  = \sup_{\gamma < \lambda} \alpha^{\beta+\gamma}
  \stackrel{\text{HI}}{=} \sup_{\gamma < \lambda} (\alpha^\beta \cdot \alpha^\gamma)
  = \alpha^\beta \cdot \sup_{\gamma < \lambda} \alpha^\gamma
  = \alpha^\beta \cdot \alpha^\lambda. \checkmark
\]
(El penúltimo paso usa que la multiplicación por $\alpha^\beta$ conmuta con el supremo cuando $\alpha^\beta > 0$, hecho que se verifica por inducción.)

\textbf{(ii)} Inducción transfinita en $\gamma$.

\textit{Base:} $(\alpha^\beta)^0 = 1 = \alpha^0 = \alpha^{\beta \cdot 0}$. \checkmark

\textit{Paso sucesor:} $(\alpha^\beta)^{\gamma+1} = (\alpha^\beta)^\gamma \cdot \alpha^\beta
\stackrel{\text{HI}}{=} \alpha^{\beta\gamma} \cdot \alpha^\beta = \alpha^{\beta\gamma + \beta}
= \alpha^{\beta(\gamma+1)}$. \checkmark

\textit{Paso límite:} $(\alpha^\beta)^\lambda = \sup_{\gamma<\lambda}(\alpha^\beta)^\gamma
\stackrel{\text{HI}}{=} \sup_{\gamma<\lambda} \alpha^{\beta\gamma}
= \alpha^{\sup_{\gamma<\lambda}\beta\gamma} = \alpha^{\beta\lambda}$. \checkmark
\end{proof}

\begin{observacion}{Las leyes de exponentes fallan sin orden}{noconmpot}
La ley $(\alpha \cdot \beta)^\gamma = \alpha^\gamma \cdot \beta^\gamma$ \emph{no} vale en general. Por ejemplo, con $\alpha = \omega$, $\beta = \omega$ y $\gamma = 2$:
\[
  (\omega \cdot \omega)^2 = \omega^2 \cdot \omega^2 = \omega^4,
\]
pero aplicando la fórmula candidata se obtendría $\omega^2 \cdot \omega^2 = \omega^4$, que en este caso coincide. Sin embargo, si tomamos $\alpha = 2$, $\beta = \omega$ y $\gamma = \omega$:
\[
  (2 \cdot \omega)^\omega = \omega^\omega,
  \quad\text{mientras que}\quad
  2^\omega \cdot \omega^\omega = \omega \cdot \omega^\omega = \omega^{1+\omega} = \omega^\omega.
\]
Aquí coinciden por absorción, pero la identidad en cuestión falla de manera no trivial para otros valores \cite{sierpinski1958}.
\end{observacion}

\subsection{Las potencias de $\omega$: cálculos detallados}

\begin{ejemplo}{Las potencias de $\omega$ paso a paso}{potomegadet}
\begin{align*}
  \omega^0 &= 1, \\
  \omega^1 &= \omega, \\
  \omega^2 &= \omega^1 \cdot \omega = \omega \cdot \omega = \sup_{n<\omega}(\omega \cdot n)
    = \sup\{0,\omega,\omega+\omega,\omega+\omega+\omega,\ldots\}, \\
  \omega^3 &= \omega^2 \cdot \omega = \sup_{n<\omega}(\omega^2 \cdot n), \\
  \omega^n &= \underbrace{\omega \cdot \omega \cdots \omega}_{n\text{ veces}}
    \quad(n \in \omega), \\
  \omega^\omega &= \sup_{n<\omega} \omega^n.
\end{align*}

El ordinal $\omega^2$ tiene el tipo de orden del conjunto
$\omega \times \omega = \{(i,j) : i,j \in \omega\}$
con el orden lexicográfico (primero por $i$, luego por $j$):
\[
  (0,0) < (0,1) < (0,2) < \cdots < (1,0) < (1,1) < \cdots < (2,0) < \cdots
\]

El ordinal $\omega^\omega$ es el supremo de todos $\omega^n$, un ordinal límite que no es alcanzado por ninguna potencia finita de $\omega$. Representa el tipo de orden de las sucesiones $(a_0, a_1, a_2, \ldots)$ con valores en $\omega$, casi todas nulas, bajo el orden lexicográfico inverso.
\end{ejemplo}

\begin{ejemplo}{Potencias mixtas}{potmixtas}
\begin{align*}
  2^\omega &= \sup_{n<\omega} 2^n = \sup\{1, 2, 4, 8, \ldots\} = \omega, \\
  3^\omega &= \sup_{n<\omega} 3^n = \omega, \\
  n^\omega &= \omega \quad \text{para todo } 1 < n < \omega, \\
  \omega^{\omega+1} &= \omega^\omega \cdot \omega = \sup_{n<\omega}(\omega^\omega \cdot n), \\
  \omega^{\omega \cdot 2} &= (\omega^\omega)^2 = \omega^\omega \cdot \omega^\omega, \\
  \omega^{\omega^2} &= \sup_{n<\omega} \omega^{\omega \cdot n}.
\end{align*}

En particular, $n^\omega = \omega$ para todo entero $n \geq 2$: la potencia $n^\omega$ es absorbida por $\omega$ porque se trata del supremo de la sucesión de enteros $n^0 < n^1 < n^2 < \cdots$ que sigue siendo contable y coinicial con $\omega$.
\end{ejemplo}

% ===========================================================
\section{El Teorema de la Forma Normal de Cantor}
% ===========================================================

\subsection{Enunciado preciso}

El Teorema de la Forma Normal de Cantor es el análogo ordinal de la representación en una base numérica. Todo número natural tiene una única representación decimal; todo ordinal positivo tiene una única representación como «polinomio en $\omega$».

\begin{teorema}{Forma Normal de Cantor}{fncantor}
Para todo ordinal $\alpha > 0$ existe una única representación
\[
  \alpha = \omega^{\beta_1} \cdot c_1 + \omega^{\beta_2} \cdot c_2 + \cdots + \omega^{\beta_k} \cdot c_k,
\]
donde:
\begin{enumerate}[leftmargin=2em, label=(\roman*)]
  \item $k \geq 1$.
  \item $\beta_1 > \beta_2 > \cdots > \beta_k$ son ordinales (llamados los \textbf{exponentes}).
  \item $c_1, c_2, \ldots, c_k$ son enteros positivos (llamados los \textbf{coeficientes}).
  \item $\beta_1 \leq \alpha$ (el exponente más alto no supera al propio ordinal).
\end{enumerate}
\end{teorema}

\begin{proof}
La demostración consta de dos partes: \emph{existencia} y \emph{unicidad}.

\medskip
\textbf{Existencia.} Procederemos por inducción transfinita en $\alpha$.

Sea $\alpha > 0$. Consideremos el conjunto
\[
  E = \{\beta \in \mathbf{ON} : \omega^\beta \leq \alpha\}.
\]
El conjunto $E$ es no vacío (pues $0 \in E$: $\omega^0 = 1 \leq \alpha$) y está acotado superiormente por $\alpha$ (si $\omega^\beta > \alpha$ para todo $\beta > \alpha$, basta verificar que $\omega^{\alpha+1} > \alpha$, lo cual se prueba por inducción sobre $\alpha$). Por tanto $E$ tiene un supremo $\beta_1 = \sup E$, y como $E$ es un conjunto de ordinales con cota, $\beta_1 \in \mathbf{ON}$.

\textit{Afirmamos que $\beta_1 \in E$, es decir, $\omega^{\beta_1} \leq \alpha$.}
Si $\beta_1 = 0$, entonces $\omega^0 = 1 \leq \alpha$. Si $\beta_1 = \gamma + 1$ para algún $\gamma$, entonces $\gamma \in E$, es decir $\omega^\gamma \leq \alpha < \omega^{\gamma+1} = \omega^{\beta_1}$, contradiciendo $\beta_1 = \sup E$ (ya que $\beta_1 \notin E$ si $\omega^{\beta_1} > \alpha$). Luego $\beta_1$ es límite o $\beta_1 = 0$, y en ambos casos $\beta_1 \in E$.

Así $\omega^{\beta_1} \leq \alpha$. Como $\beta_1+1 \notin E$, tenemos $\omega^{\beta_1+1} > \alpha$, es decir $\omega^{\beta_1} \cdot \omega > \alpha$. Sea
\[
  c_1 = \max\{c \in \omega : \omega^{\beta_1} \cdot c \leq \alpha\}.
\]
Este máximo existe porque $\omega^{\beta_1} \cdot \omega > \alpha$, así que el conjunto es finito y no vacío (contiene al menos a $1$). Define $\alpha_1 = \alpha - \omega^{\beta_1} \cdot c_1$ (diferencia ordinal por la derecha), es decir, el único ordinal tal que $\omega^{\beta_1} \cdot c_1 + \alpha_1 = \alpha$.

Si $\alpha_1 = 0$, hemos terminado: $\alpha = \omega^{\beta_1} \cdot c_1$.

Si $\alpha_1 > 0$, nótese que $\alpha_1 < \omega^{\beta_1}$ (pues si $\alpha_1 \geq \omega^{\beta_1}$, entonces $c_1$ no sería el máximo). Por hipótesis de inducción, $\alpha_1$ tiene una forma normal
\[
  \alpha_1 = \omega^{\beta_2} \cdot c_2 + \cdots + \omega^{\beta_k} \cdot c_k,
\]
con $\beta_2 < \beta_1$ (ya que $\alpha_1 < \omega^{\beta_1}$ implica que todos los exponentes de $\alpha_1$ son estrictamente menores que $\beta_1$). Concatenando, obtenemos la forma normal de $\alpha$.

\medskip
\textbf{Unicidad.} Supongamos que
\[
  \alpha = \omega^{\beta_1} c_1 + \cdots + \omega^{\beta_k} c_k
  = \omega^{\gamma_1} d_1 + \cdots + \omega^{\gamma_m} d_m,
\]
con $\beta_1 > \cdots > \beta_k$ y $\gamma_1 > \cdots > \gamma_m$. Procedemos por inducción transfinita en $\alpha$.

Primero observamos que el mayor exponente en la forma normal de $\alpha$ es el mayor ordinal $\delta$ con $\omega^\delta \leq \alpha$, es decir $\beta_1 = \gamma_1$ (ambos son el supremo del conjunto $E$ anterior). Así $\beta_1 = \gamma_1$.

A continuación, el coeficiente $c_1$ es el mayor entero positivo con $\omega^{\beta_1} \cdot c_1 \leq \alpha$, así que $c_1 = d_1$.

Luego $\omega^{\beta_2} c_2 + \cdots = \omega^{\gamma_2} d_2 + \cdots = \alpha - \omega^{\beta_1}c_1$, y la unicidad de este ordinal (que es estrictamente menor que $\alpha$) sigue por hipótesis de inducción.
\end{proof}

\subsection{Ejemplos de cálculo de la forma normal}

\begin{ejemplo}{Forma normal de ordinales concretos}{ejfncantor}
Calculamos la forma normal de Cantor paso a paso.

\textbf{(a) $\alpha = 1427$.} El mayor $\beta$ con $\omega^\beta \leq 1427$ es $\beta_1 = 0$ (ya que $\omega^1 = \omega > 1427$).
\[
  1427 = \omega^0 \cdot 1427.
\]
En general, todo entero positivo $n$ tiene forma normal $\omega^0 \cdot n$.

\textbf{(b) $\alpha = \omega^3 + \omega^2 \cdot 5 + \omega \cdot 3 + 7$.} Esta ya es la forma normal con exponentes $3 > 2 > 1 > 0$ y coeficientes $1, 5, 3, 7$.

\textbf{(c) $\alpha = \omega^4 + \omega \cdot 2 + 1$.} Los exponentes son $4 > 1 > 0$ y los coeficientes $1, 2, 1$.

\textbf{(d) $\alpha = \omega^\omega$.} La forma normal es $\omega^\omega \cdot 1$: un solo término con exponente $\omega$ y coeficiente $1$. Nótese que aquí el exponente $\beta_1 = \omega$ es un ordinal infinito.

\textbf{(e) $\alpha = \omega^{\omega^2} + \omega^3 \cdot 4 + 2$.} Exponentes: $\omega^2 > 3 > 0$, coeficientes: $1, 4, 2$.

\textbf{(f) $\alpha = \omega^{\omega^\omega}$.} La forma normal es $\omega^{\omega^\omega} \cdot 1$. El exponente $\omega^\omega$ es en sí un ordinal transfinito con su propia forma normal $\omega^\omega \cdot 1$.
\end{ejemplo}

\begin{ejemplo}{Aritmética con la forma normal}{aritfncantor}
Sean $\alpha = \omega^3 \cdot 2 + \omega \cdot 5 + 3$ y $\beta = \omega^2 \cdot 4 + \omega \cdot 7 + 1$.

\textbf{Suma $\alpha + \beta$:} En la suma ordinal, los términos de $\alpha$ con exponente $\leq \beta_1(\beta) = 2$ quedan absorbidos por $\beta$:
\begin{align*}
  \alpha + \beta
  &= (\omega^3 \cdot 2 + \omega \cdot 5 + 3) + (\omega^2 \cdot 4 + \omega \cdot 7 + 1) \\
  &= \omega^3 \cdot 2 + (\omega \cdot 5 + 3 + \omega^2 \cdot 4) + \omega \cdot 7 + 1 \\
  &= \omega^3 \cdot 2 + \omega^2 \cdot 4 + \omega \cdot 7 + 1.
\end{align*}
(El término $\omega \cdot 5 + 3$ es absorbido por $\omega^2 \cdot 4$, pues $\omega \cdot 5 + 3 < \omega^2 \leq \omega^2 \cdot 4$.)

\textbf{Suma $\beta + \alpha$:} Aquí el término más alto de $\alpha$ es $\omega^3 > \omega^2$:
\[
  \beta + \alpha = \omega^3 \cdot 2 + \omega \cdot 5 + 3 = \alpha.
\]
Esto refleja que $\beta < \omega^3$ y $\omega^3$ absorbe a $\beta$: $\beta + \omega^3 = \omega^3$, y por tanto $\beta + (\omega^3 \cdot 2 + \cdots) = \omega^3 \cdot 2 + \cdots = \alpha$.

\textbf{Producto $\alpha \cdot \omega$:} Para cualquier $\alpha = \omega^{\beta_1} c_1 + (\text{términos menores})$, se tiene $\alpha \cdot \omega = \omega^{\beta_1+1}$, pues
\[
  \alpha \cdot \omega = \sup_{n<\omega}(\alpha \cdot n)
  = \sup_{n<\omega}(\omega^{\beta_1} c_1 n + \cdots)
  = \omega^{\beta_1+1}.
\]
En nuestro caso, $\alpha \cdot \omega = \omega^4$.
\end{ejemplo}

\begin{corolario}{Algoritmo de suma con la forma normal}{algsuma}
Si $\alpha = \omega^{\beta_1} c_1 + \cdots + \omega^{\beta_k} c_k$ y $\gamma = \omega^{\delta_1} d_1 + \cdots + \omega^{\delta_m} d_m$ con $\beta_1 \geq \cdots \geq \beta_k$ y $\delta_1 \geq \cdots \geq \delta_m$, entonces para calcular $\alpha + \gamma$:
\begin{enumerate}[leftmargin=2em, label=(\arabic*)]
  \item Se eliminan de $\alpha$ todos los términos $\omega^{\beta_i} c_i$ con $\beta_i < \delta_1$ (quedan absorbidos).
  \item Si el mayor exponente restante de $\alpha$ coincide con $\delta_1$, se suman los coeficientes correspondientes.
  \item Se concatenan los términos restantes de $\alpha$ con todos los de $\gamma$.
\end{enumerate}
\end{corolario}

\begin{proposicion}{El exponente máximo y la forma normal}{expmaxfn}
Sea $\alpha > 0$ con forma normal $\alpha = \omega^{\beta_1} c_1 + \cdots$. Entonces:
\begin{enumerate}[leftmargin=2em, label=(\roman*)]
  \item $\alpha \geq \omega^{\beta_1}$.
  \item $\alpha < \omega^{\beta_1+1}$.
  \item $\alpha$ es un ordinal límite si y solo si $\beta_k > 0$ (el último exponente es positivo) o $k \geq 2$ con $\beta_k = 0$... más precisamente, $\alpha$ es límite si y solo si el término de menor exponente tiene exponente $> 0$.
  \item $\alpha = \omega^{\beta_1} \cdot c_1$ si y solo si $k = 1$.
\end{enumerate}
\end{proposicion}

\begin{proof}
\textbf{(i)} y \textbf{(ii)} son consecuencia directa de la construcción de $\beta_1$ en la prueba del teorema.

\textbf{(iii)} $\alpha$ es sucesor $\iff$ $\alpha = \alpha' + 1$ para algún $\alpha'$, lo que en la forma normal ocurre $\iff$ $\beta_k = 0$ (el último término es $\omega^0 \cdot c_k = c_k$), en cuyo caso $\alpha = (\omega^{\beta_1} c_1 + \cdots + \omega^{\beta_{k-1}} c_{k-1} + (c_k - 1)) + 1$.
\end{proof}

% ===========================================================
\section{Puntos fijos de \texorpdfstring{$\alpha \mapsto \omega^\alpha$}{alpha -> omega^alpha} y el ordinal \texorpdfstring{$\varepsilon_0$}{ε₀}}
% ===========================================================

\subsection{La función \texorpdfstring{$\alpha \mapsto \omega^\alpha$}{α ↦ ωᵅ} y sus propiedades}

La función $\Phi: \mathbf{ON} \to \mathbf{ON}$ definida por $\Phi(\alpha) = \omega^\alpha$ es un objeto fundamental en la teoría de ordinales. Antes de estudiar sus puntos fijos, sistematizamos sus propiedades.

\begin{proposicion}{Propiedades de $\alpha \mapsto \omega^\alpha$}{propsomegaalpha}
La función $\Phi(\alpha) = \omega^\alpha$ satisface:
\begin{enumerate}[leftmargin=2em, label=(\roman*)]
  \item Es \textbf{estrictamente creciente}: $\alpha < \beta \Rightarrow \omega^\alpha < \omega^\beta$.
  \item Es \textbf{continua por la derecha}: para todo ordinal límite $\lambda$, $\omega^\lambda = \sup_{\alpha < \lambda} \omega^\alpha$.
  \item Para todo $\alpha$ finito, $\omega^\alpha < \omega$.
  \item Para todo $\alpha \geq 1$, $\alpha < \omega^\alpha$.
  \item La imagen de $\Phi$ es una subclase propia de $\mathbf{ON}$.
\end{enumerate}
\end{proposicion}

\begin{proof}
\textbf{(i)} Por inducción en $\beta$ se demuestra que $\omega^\alpha < \omega^{\alpha+1}$ y que la función es creciente; luego por transitividad es estrictamente creciente.

\textbf{(ii)} Es la definición misma de $\omega^\lambda$ para $\lambda$ límite.

\textbf{(iv)} Por inducción transfinita en $\alpha$.
\textit{Base $\alpha=1$}: $1 < \omega = \omega^1$. \checkmark
\textit{Sucesor}: Si $\alpha < \omega^\alpha$, entonces $\alpha + 1 \leq \omega^\alpha < \omega^\alpha \cdot \omega = \omega^{\alpha+1}$. \checkmark
\textit{Límite}: Si $\alpha$ es límite y $\beta < \alpha$ para todo $\beta$ implica $\beta < \omega^\beta \leq \omega^\alpha$, luego $\alpha = \sup_{\beta < \alpha} \beta \leq \omega^\alpha$. Y la igualdad no ocurre si $\alpha$ no es punto fijo (lo cual demostraremos abajo). \checkmark
\end{proof}

\subsection{La existencia de puntos fijos: definición de $\varepsilon_0$}

\begin{definicion}{Punto fijo de $\alpha \mapsto \omega^\alpha$}{puntofijo}
Un ordinal $\varepsilon$ es un \textbf{punto fijo} de la función $\alpha \mapsto \omega^\alpha$ si
\[
  \omega^\varepsilon = \varepsilon.
\]
Un ordinal así recibe el nombre de \textbf{ordinal épsilon}.
\end{definicion}

\begin{importante}
¿Por qué $\omega$, $\omega^2$, $\omega^\omega$, etc., no son puntos fijos? Verifiquemos:
\begin{align*}
  \omega^1 &= \omega \quad\Rightarrow\quad \omega^{\omega^1} = \omega^\omega \neq \omega.\\
  \omega^{\omega} &\quad\Rightarrow\quad \omega^{\omega^\omega} = \omega^{\omega^\omega} \neq \omega^\omega \quad\text{(pues $\omega^\omega < \omega^{\omega^\omega}$)}.
\end{align*}
Ningún ordinal de la forma $\omega^n$, $\omega^\omega$, $\omega^{\omega^\omega}$, etc., satisface la ecuación $\omega^\varepsilon = \varepsilon$. El primer punto fijo requiere un proceso iterado «hasta el infinito».
\end{importante}

\begin{teorema}{Existencia y definición de $\varepsilon_0$}{defeps0}
Existe un ordinal $\varepsilon_0$ tal que $\omega^{\varepsilon_0} = \varepsilon_0$, y es el \textbf{mínimo} ordinal con esta propiedad. Está dado explícitamente por
\[
  \varepsilon_0 = \sup\bigl\{\omega,\; \omega^\omega,\; \omega^{\omega^\omega},\; \omega^{\omega^{\omega^\omega}},\; \ldots\bigr\}
  = \sup_{n < \omega} \varphi_n,
\]
donde $\varphi_0 = \omega$ y $\varphi_{n+1} = \omega^{\varphi_n}$ para todo $n \in \omega$.
\end{teorema}

\begin{proof}
Definimos por recursión ordinaria:
\[
  \varphi_0 = \omega, \qquad \varphi_{n+1} = \omega^{\varphi_n}.
\]
La sucesión $\varphi_0 < \varphi_1 < \varphi_2 < \cdots$ es estrictamente creciente (por la propiedad (i) de la proposición anterior, ya que $\varphi_n < \omega^{\varphi_n} = \varphi_{n+1}$). Sea $\varepsilon_0 = \sup_{n<\omega} \varphi_n$.

\textbf{$\varepsilon_0$ es un punto fijo:} Como $\varepsilon_0$ es el supremo de una sucesión de tipo $\omega$, es un ordinal límite. Aplicando la continuidad de $\alpha \mapsto \omega^\alpha$:
\[
  \omega^{\varepsilon_0}
  = \omega^{\sup_n \varphi_n}
  = \sup_n \omega^{\varphi_n}
  = \sup_n \varphi_{n+1}
  = \sup_{n \geq 1} \varphi_n
  = \varepsilon_0. \checkmark
\]
(El último paso es porque $\sup_{n \geq 1} \varphi_n = \sup_{n \geq 0} \varphi_n = \varepsilon_0$, pues $\varphi_0 < \varphi_1$.)

\textbf{$\varepsilon_0$ es el mínimo:} Sea $\varepsilon$ cualquier punto fijo de $\alpha \mapsto \omega^\alpha$. Como $\varepsilon \geq \omega = \omega^0$ (pues $\omega^1 = \omega \leq \omega^\varepsilon = \varepsilon$, y $\varepsilon \geq 1$), tenemos $\varphi_0 = \omega \leq \varepsilon$. Por inducción: si $\varphi_n \leq \varepsilon$, entonces $\varphi_{n+1} = \omega^{\varphi_n} \leq \omega^\varepsilon = \varepsilon$. Luego $\varphi_n \leq \varepsilon$ para todo $n$, y por tanto $\varepsilon_0 = \sup_n \varphi_n \leq \varepsilon$.
\end{proof}

\subsection{Representación de $\varepsilon_0$ como torre de exponentes}

\begin{ejemplo}{La torre exponencial de $\varepsilon_0$}{torreexp}
La sucesión que converge a $\varepsilon_0$ puede escribirse explícitamente como torres de exponentes:
\begin{align*}
  \varphi_0 &= \omega, \\
  \varphi_1 &= \omega^\omega, \\
  \varphi_2 &= \omega^{\omega^\omega}, \\
  \varphi_3 &= \omega^{\omega^{\omega^\omega}}, \\
  \varphi_n &= \underbrace{\omega^{\omega^{\cdot^{\cdot^{\cdot^\omega}}}}}_{n+1\;\omega\text{'s}}.
\end{align*}
El límite $\varepsilon_0 = \sup_n \varphi_n$ puede imaginarse como la «torre infinita»:
\[
  \varepsilon_0 = \omega^{\omega^{\omega^{\omega^{\cdots}}}}.
\]
Por supuesto, esta «expresión» no es formalmente una notación finita; refleja el hecho de que $\varepsilon_0$ no puede ser escrito como una expresión finita con $\omega$ y las operaciones $+$, $\cdot$, y la potenciación, sin referirse a sí mismo. Esto contrasta con todos los ordinales que tienen forma normal de Cantor con exponentes finitos.
\end{ejemplo}

\subsection{La forma normal de Cantor y $\varepsilon_0$}

\begin{observacion}{El límite de la forma normal de Cantor}{fnylimite}
Si $\alpha < \varepsilon_0$, entonces la forma normal de Cantor de $\alpha$ involucra exponentes $\beta_1 < \varepsilon_0$, cuyos exponentes son a su vez $< \varepsilon_0$, y así sucesivamente; pero la cadena de exponentes es \emph{finita}. Esto es: para $\alpha < \varepsilon_0$, la forma normal de Cantor puede expandirse de manera \emph{completamente finita}, en el sentido de que todos los exponentes que aparecen son también $<\varepsilon_0$ y el proceso termina en un número finito de pasos.

Por el contrario, $\varepsilon_0$ mismo no puede representarse por la forma normal de Cantor sin usar un exponente igual a $\varepsilon_0$:
\[
  \varepsilon_0 = \omega^{\varepsilon_0} \cdot 1.
\]
Esto no es circular (la forma normal existe) pero sí implica que la cadena de expansión de exponentes no termina; el proceso de expandir el exponente siempre produce de nuevo $\varepsilon_0$.
\end{observacion}

\begin{importante}
$\varepsilon_0$ es el primer ordinal para el cual la forma normal de Cantor es «autorreferente»: $\varepsilon_0 = \omega^{\varepsilon_0}$. Todos los ordinales $\alpha < \varepsilon_0$ tienen representaciones en forma normal de Cantor que se resuelven en términos de ordinales estrictamente menores a través de una cadena finita de expansiones.
\end{importante}

% ===========================================================
\section{La jerarquía de ordinales épsilon}
% ===========================================================

\subsection{Ordinales épsilon más allá de $\varepsilon_0$}

La existencia de $\varepsilon_0$ es solo el comienzo. Una vez construido $\varepsilon_0$, podemos preguntar: ¿existen más puntos fijos de $\alpha \mapsto \omega^\alpha$? La respuesta es afirmativa, y de hecho forman una clase propia.

\begin{teorema}{Existencia de $\varepsilon_{\alpha}$ para todo ordinal $\alpha$}{existepsalpha}
Para cada ordinal $\alpha$ existe un único ordinal $\varepsilon_\alpha$ tal que:
\begin{enumerate}[leftmargin=2em, label=(\roman*)]
  \item $\omega^{\varepsilon_\alpha} = \varepsilon_\alpha$ \quad ($\varepsilon_\alpha$ es punto fijo de $\alpha \mapsto \omega^\alpha$).
  \item $\varepsilon_\alpha$ es el $(\alpha+1)$-ésimo punto fijo en orden creciente \quad (es decir, es el mínimo punto fijo mayor que todos $\varepsilon_\beta$ con $\beta < \alpha$).
\end{enumerate}
La colección $\{\varepsilon_\alpha : \alpha \in \mathbf{ON}\}$ es la \textbf{jerarquía épsilon} o \textbf{clase de ordinales épsilon}.
\end{teorema}

\begin{proof}[Construcción]
Definimos la jerarquía por recursión transfinita en $\alpha$.

\textbf{Caso base:} $\varepsilon_0$ se define como en el Teorema~\ref{teo:defeps0}.

\textbf{Caso sucesor:} Dado $\varepsilon_\alpha$, definimos $\varepsilon_{\alpha+1}$ como el mínimo punto fijo de $\beta \mapsto \omega^\beta$ mayor que $\varepsilon_\alpha$. Explícitamente:
\[
  \varepsilon_{\alpha+1} = \sup\{\varepsilon_\alpha + 1,\; \omega^{\varepsilon_\alpha+1},\; \omega^{\omega^{\varepsilon_\alpha+1}},\; \ldots\}.
\]
El mismo argumento de la prueba del Teorema~\ref{teo:defeps0} muestra que este supremo es un punto fijo.

\textbf{Caso límite:} Para $\lambda$ un ordinal límite,
\[
  \varepsilon_\lambda = \sup_{\alpha < \lambda} \varepsilon_\alpha.
\]
Verificamos que $\varepsilon_\lambda$ es un punto fijo:
\[
  \omega^{\varepsilon_\lambda}
  = \omega^{\sup_{\alpha < \lambda}\varepsilon_\alpha}
  = \sup_{\alpha < \lambda} \omega^{\varepsilon_\alpha}
  = \sup_{\alpha < \lambda} \varepsilon_\alpha
  = \varepsilon_\lambda. \checkmark
\]
\end{proof}

\subsection{Los primeros ordinales épsilon}

\begin{ejemplo}{Los primeros $\varepsilon_\alpha$}{primeroseps}
\begin{align*}
  \varepsilon_0 &= \sup\{\omega, \omega^\omega, \omega^{\omega^\omega}, \ldots\}
    &&\text{(primer punto fijo de $\alpha \mapsto \omega^\alpha$),} \\
  \varepsilon_1 &= \sup\{\varepsilon_0+1,\; \omega^{\varepsilon_0+1},\; \omega^{\omega^{\varepsilon_0+1}},\; \ldots\}
    &&\text{(segundo punto fijo),} \\
  \varepsilon_2 &= \sup\{\varepsilon_1+1,\; \omega^{\varepsilon_1+1},\; \ldots\}
    &&\text{(tercer punto fijo),} \\
  \varepsilon_\omega &= \sup_{n < \omega} \varepsilon_n
    &&\text{(primer punto fijo límite de la jerarquía $\varepsilon$),} \\
  \varepsilon_{\varepsilon_0} &
    &&\text{(un ordinal épsilon cuyo índice es él mismo un ordinal épsilon).}
\end{align*}
La jerarquía épsilon crece sin cota: para cada $\alpha$, $\varepsilon_\alpha < \varepsilon_{\alpha+1}$ estrictamente.
\end{ejemplo}

\begin{proposicion}{Propiedades de la jerarquía épsilon}{propseps}
\begin{enumerate}[leftmargin=2em, label=(\roman*)]
  \item La función $\alpha \mapsto \varepsilon_\alpha$ es estrictamente creciente y continua.
  \item Para todo $\alpha$, $\varepsilon_\alpha$ es un ordinal límite.
  \item Para todo $\alpha$, $\alpha < \varepsilon_\alpha$.
  \item Para todo $\alpha$ y todo $\beta < \varepsilon_\alpha$, se tiene $\omega^\beta < \varepsilon_\alpha$.
  \item La clase de todos los ordinales épsilon es una clase propia.
\end{enumerate}
\end{proposicion}

\begin{proof}
\textbf{(ii)} Si $\varepsilon_\alpha$ fuera sucesor, $\varepsilon_\alpha = \gamma + 1$, entonces $\omega^{\varepsilon_\alpha} = \omega^{\gamma+1} = \omega^\gamma \cdot \omega > \gamma + 1 = \varepsilon_\alpha$ para $\gamma$ infinito, contradicción.

\textbf{(iii)} Por inducción: $\varepsilon_0 \geq \omega > 0$, y si $\alpha < \varepsilon_\alpha$ entonces $\alpha + 1 \leq \varepsilon_\alpha < \varepsilon_{\alpha+1}$.

\textbf{(v)} Si la clase fuera un conjunto $E$, sea $\varepsilon = \sup E$. Entonces $\omega^\varepsilon = \sup_{\alpha \in E} \omega^{\varepsilon_\alpha} = \sup_{\alpha \in E} \varepsilon_\alpha = \varepsilon$, luego $\varepsilon$ sería un ordinal épsilon no en $E$, contradicción.
\end{proof}

\subsection{La función de Veblen y la jerarquía más allá de $\varepsilon$}

\begin{observacion}{La función $\varphi$ de Veblen}{veblen}
La jerarquía épsilon puede verse como los puntos fijos de $\varphi_0(\alpha) = \omega^\alpha$. Si definimos $\varphi_1(\alpha) = \varepsilon_\alpha$ (el $\alpha$-ésimo punto fijo de $\varphi_0$), podemos luego considerar $\varphi_2(\alpha) = $ el $\alpha$-ésimo punto fijo de $\varphi_1$, etc. Esta es la \textbf{función de Veblen}, que produce ordinales cada vez mayores:
\[
  \varphi_0(\alpha) = \omega^\alpha, \quad
  \varphi_1(\alpha) = \varepsilon_\alpha, \quad
  \varphi_2(\alpha) = \text{$\alpha$-ésimo punto fijo de $\varphi_1$}, \quad \ldots
\]
El ordinal $\Gamma_0 = \varphi_{\Gamma_0}(0)$ (primer punto fijo del proceso de Veblen a través de todos sus niveles) es la llamada \textbf{constante de Feferman-Schütte}, que aparece en la teoría de la prueba como cota de la prueba de sistemas como la aritmética predicativa. Esto ilustra cómo la jerarquía de ordinales transfinitos tiene profundas conexiones con la teoría de la demostración \cite{pohlers1989, smorynski1991}.
\end{observacion}

% ===========================================================
\section{Cofinalidad de un ordinal}
% ===========================================================

\subsection{Motivación: ordinales límite de «distintos ritmos»}

Los ordinales límite no son todos iguales desde el punto de vista de cómo son «alcanzados» por sucesiones de ordinales menores. Compárense:

\begin{itemize}[leftmargin=2em]
  \item $\omega = \sup\{0, 1, 2, 3, \ldots\}$: alcanzado por una sucesión de tipo $\omega$ (enumerable, bien ordenada por su tipo).
  \item $\omega_1$ (el primer ordinal no numerable): no puede alcanzarse por ninguna sucesión de tipo $\omega$; se necesita una sucesión de longitud $\omega_1$ para «agotar» $\omega_1$.
  \item $\omega_\omega = \sup\{\omega_0, \omega_1, \omega_2, \ldots\}$: aunque es un ordinal límite extremadamente grande, es alcanzado por una sucesión de tipo $\omega$ (la sucesión $\omega_0, \omega_1, \omega_2, \ldots$).
\end{itemize}

Esta observación motiva la siguiente definición.

\subsection{Definición formal de cofinalidad}

\begin{definicion}{Función cofinala y cofinalidad}{cofinalidad}
Sea $\lambda$ un ordinal límite.
\begin{enumerate}[leftmargin=2em, label=(\roman*)]
  \item Una función $f: \alpha \to \lambda$ es \textbf{cofinal en $\lambda$} si para todo $\xi < \lambda$ existe $i < \alpha$ con $f(i) \geq \xi$ (equivalentemente, $\sup_{\iota < \alpha} f(\iota) = \lambda$, aunque $\lambda$ no se alcanza como valor).
  \item La \textbf{cofinalidad} de $\lambda$, denotada $\mathrm{cf}(\lambda)$, es el menor ordinal $\alpha$ tal que existe una función cofinal $f: \alpha \to \lambda$.
\end{enumerate}
Para ordinales sucesor y para $0$ se define $\mathrm{cf}(\lambda) = 1$ si $\lambda > 0$ (son alcanzados en un solo paso) y $\mathrm{cf}(0) = 0$.
\end{definicion}

\begin{observacion}{Reformulación equivalente}{cofequiv}
$\mathrm{cf}(\lambda)$ es el menor cardinal $\kappa$ tal que $\lambda$ es unión de $\kappa$ muchos conjuntos de ordinales, cada uno acotado en $\lambda$. También puede definirse como el mínimo longitud de una sucesión estrictamente creciente $\langle \alpha_i \rangle_{i < \kappa}$ de ordinales $< \lambda$ cuyo supremo es $\lambda$.
\end{observacion}

\subsection{Ejemplos de cofinalidad}

\begin{ejemplo}{Cofinalidad de los primeros ordinales}{cofprimerosord}
\begin{enumerate}[leftmargin=2em, label=(\alph*)]
  \item $\mathrm{cf}(\omega) = \omega$: la sucesión $0, 1, 2, 3, \ldots$ es cofinal en $\omega$ y tiene longitud $\omega$. No hay sucesión finita cofinal en $\omega$ (toda sucesión finita tiene un máximo, que no puede ser cota superior de todos los naturales).

  \item $\mathrm{cf}(\omega + \omega) = \omega$: la sucesión $\omega, \omega+1, \omega+2, \ldots$ es cofinal en $\omega+\omega$ con longitud $\omega$.

  \item $\mathrm{cf}(\omega^\omega) = \omega$: la sucesión $\omega^0, \omega^1, \omega^2, \ldots$ es cofinal en $\omega^\omega$ con longitud $\omega$.

  \item $\mathrm{cf}(\varepsilon_0) = \omega$: la sucesión $\omega, \omega^\omega, \omega^{\omega^\omega}, \ldots$ es cofinal en $\varepsilon_0$ con longitud $\omega$.

  \item $\mathrm{cf}(\omega_1) = \omega_1$: no existe ninguna sucesión de longitud $< \omega_1$ cofinal en $\omega_1$ (esto es equivalente a afirmar que $\omega_1$ no es numerable, lo cual es su definición).
\end{enumerate}
\end{ejemplo}

\begin{definicion}{Ordinal regular y ordinal singular}{regular}
Un ordinal límite $\lambda$ es:
\begin{itemize}[leftmargin=2em]
  \item \textbf{Regular} si $\mathrm{cf}(\lambda) = \lambda$.
  \item \textbf{Singular} si $\mathrm{cf}(\lambda) < \lambda$.
\end{itemize}
\end{definicion}

\begin{ejemplo}{Ordinales regulares y singulares}{regsingjemplo}
\begin{itemize}[leftmargin=2em]
  \item $\omega$ es regular: $\mathrm{cf}(\omega) = \omega$.
  \item $\omega + \omega$ es singular: $\mathrm{cf}(\omega + \omega) = \omega < \omega + \omega$.
  \item $\omega^\omega$ es singular: $\mathrm{cf}(\omega^\omega) = \omega < \omega^\omega$.
  \item $\varepsilon_0$ es singular: $\mathrm{cf}(\varepsilon_0) = \omega < \varepsilon_0$.
  \item $\omega_1$ (el primer ordinal no numerable) es regular: $\mathrm{cf}(\omega_1) = \omega_1$. (Esto se probará formalmente al estudiar cardinales.)
  \item $\omega_\omega = \sup_{n < \omega} \omega_n$ es singular: $\mathrm{cf}(\omega_\omega) = \omega < \omega_\omega$.
\end{itemize}
\end{ejemplo}

\subsection{Propiedades básicas de la cofinalidad}

\begin{proposicion}{La cofinalidad es un cardinal regular}{cofregular}
Para todo ordinal límite $\lambda$, $\mathrm{cf}(\lambda)$ es un cardinal regular, es decir, $\mathrm{cf}(\mathrm{cf}(\lambda)) = \mathrm{cf}(\lambda)$.
\end{proposicion}

\begin{proof}
Sea $\kappa = \mathrm{cf}(\lambda)$ y sea $f: \kappa \to \lambda$ cofinal. Si $\mathrm{cf}(\kappa) = \mu < \kappa$, sea $g: \mu \to \kappa$ cofinal. Entonces $f \circ g: \mu \to \lambda$ sería cofinal (pues para todo $\xi < \lambda$ existe $i < \kappa$ con $f(i) \geq \xi$, y existe $j < \mu$ con $g(j) \geq i$, luego $f(g(j)) \geq f(i) \geq \xi$). Esto contradice la minimalidad de $\kappa$. Luego $\mathrm{cf}(\kappa) = \kappa$.
\end{proof}

\begin{proposicion}{Cofinalidad y potencias de $\omega$}{cofpotomega}
\begin{enumerate}[leftmargin=2em, label=(\roman*)]
  \item $\mathrm{cf}(\omega^\alpha) = \omega$ para todo $\alpha \geq 1$ finito.
  \item $\mathrm{cf}(\omega^\omega) = \omega$.
  \item $\mathrm{cf}(\varepsilon_\alpha) = \omega$ para todo $\alpha$ numerable.
\end{enumerate}
\end{proposicion}

\begin{proof}
\textbf{(i)} Para $\alpha = n \geq 1$ finito, la sucesión $(\omega^{n-1} \cdot k)_{k < \omega}$ (o más directamente $(\omega^{n-1} \cdot k)_{k \in \omega}$) es cofinal en $\omega^n$ y tiene longitud $\omega$.

\textbf{(ii)} La sucesión $\omega^n$ para $n \in \omega$ es cofinal en $\omega^\omega$.

\textbf{(iii)} Para $\alpha$ numerable, la construcción de $\varepsilon_\alpha$ involucra un supremo de una sucesión de tipo $\omega$ o $\omega \cdot \alpha$ (enumerable), por lo que $\mathrm{cf}(\varepsilon_\alpha) = \omega$.
\end{proof}

\begin{observacion}{Cofinalidad e inaccessibilidad}{cofinacc}
Un cardinal $\kappa > \omega$ que es regular y es límite fuerte (es decir, $2^\lambda < \kappa$ para todo $\lambda < \kappa$) recibe el nombre de \textbf{cardinal fuertemente inaccesible}. Estos cardinales no pueden demostrarse que existan en ZFC; su existencia es una hipótesis adicional. La cofinalidad es el primer paso para entender este tipo de fenómenos en la jerarquía cardinal \cite{kanamori2003, jech2003}.
\end{observacion}

\subsection{La cofinalidad y la forma normal de Cantor}

\begin{proposicion}{Cofinalidad desde la forma normal}{cofinfn}
Sea $\alpha > 0$ con forma normal $\alpha = \omega^{\beta_1} c_1 + \cdots + \omega^{\beta_k} c_k$.
\begin{enumerate}[leftmargin=2em, label=(\roman*)]
  \item Si $\beta_k = 0$ (el último término es finito), entonces $\alpha$ es sucesor y $\mathrm{cf}(\alpha) = 1$.
  \item Si $\beta_k > 0$ y $\beta_k$ es sucesor, entonces $\mathrm{cf}(\alpha) = \omega$.
  \item Si $\beta_k > 0$ y $\beta_k$ es un ordinal límite, entonces $\mathrm{cf}(\alpha) = \mathrm{cf}(\beta_k)$.
  \item En particular, para $\alpha = \omega^\beta$ con $\beta$ límite, $\mathrm{cf}(\omega^\beta) = \mathrm{cf}(\beta)$.
\end{enumerate}
\end{proposicion}

\begin{proof}
\textbf{(i)} Ya probado: $\alpha$ es sucesor si y solo si termina en un coeficiente finito con exponente $0$.

\textbf{(ii)} Si $\beta_k = \gamma + 1$, entonces $\omega^{\beta_k} = \omega^{\gamma+1} = \omega^\gamma \cdot \omega$ es un ordinal de cofinalidad $\omega$ (la sucesión $\omega^\gamma \cdot n$ para $n \in \omega$ es cofinal), y el término $\omega^{\beta_k} \cdot c_k$ tiene también cofinalidad $\omega$. Como este es el término dominante (el «final»), $\mathrm{cf}(\alpha) = \omega$.

\textbf{(iii) y (iv)} Si $\beta_k$ es límite, $\omega^{\beta_k} = \sup_{\gamma < \beta_k} \omega^\gamma$, y la cofinalidad de $\omega^{\beta_k}$ coincide con la de $\beta_k$.
\end{proof}

% ===========================================================
\section{Visión panorámica de la jerarquía ordinal}
% ===========================================================

\subsection{Estratos de la jerarquía}

La siguiente descripción sintetiza los distintos «estratos» de la jerarquía ordinal que hemos estudiado:

\begin{center}
\renewcommand{\arraystretch}{1.6}
\begin{tabular}{lll}
\hline
\textbf{Ordinal} & \textbf{Descripción} & \textbf{Cofinalidad} \\
\hline
$0, 1, 2, \ldots, n, \ldots$ & Ordinales finitos & $1$ (sucesores), $0$ ($=0$) \\
$\omega$ & Primer ordinal infinito & $\omega$ (regular) \\
$\omega+1, \omega+2, \ldots$ & Sucesores de $\omega$ & $1$ \\
$\omega \cdot 2, \omega \cdot 3, \ldots$ & Múltiplos de $\omega$ & $\omega$ \\
$\omega^2, \omega^3, \ldots, \omega^n$ & Potencias finitas de $\omega$ & $\omega$ \\
$\omega^\omega$ & Primera potencia transfinita & $\omega$ \\
$\omega^{\omega^\omega}$ & Torre de tres $\omega$'s & $\omega$ \\
$\varepsilon_0$ & Primer punto fijo de $\alpha \mapsto \omega^\alpha$ & $\omega$ \\
$\varepsilon_1, \varepsilon_2, \ldots$ & Jerarquía épsilon & $\omega$ \\
$\varepsilon_\omega$ & Primer épsilon-límite & $\omega$ \\
$\omega_1$ & Primer ordinal no numerable & $\omega_1$ (regular) \\
$\omega_\omega$ & El $\omega$-ésimo cardinal infinito & $\omega$ (singular) \\
\hline
\end{tabular}
\end{center}

\subsection{Comparación: forma normal y sus límites}

\begin{importante}
La forma normal de Cantor proporciona una representación única para cada ordinal, pero tiene un límite preciso: es completamente efectiva (sin «autorreferencia») solo para los ordinales estrictamente menores que $\varepsilon_0$. Los ordinales $\geq \varepsilon_0$ tienen formas normales que involucran a sí mismos como exponentes.

Esto tiene consecuencias en la teoría de la prueba: el sistema de aritmética de Peano ($\mathrm{PA}$) puede demostrar la consistencia de sistemas más débiles usando ordinales $< \varepsilon_0$ para medir «rangos de prueba». Sin embargo, la consistencia de $\mathrm{PA}$ mismo requiere ordinales $\geq \varepsilon_0$ (Teorema de Gentzen, 1936 \cite{gentzen1936}).
\end{importante}

\begin{ejemplo}{Notación de Knuth y $\varepsilon_0$}{knuth}
La «notación de flecha de Knuth» $\uparrow\uparrow$ para torres exponenciales en los naturales es un reflejo informal de la construcción que lleva a $\varepsilon_0$:
\[
  2 \uparrow\uparrow 3 = 2^{2^2} = 16,
  \quad
  2 \uparrow\uparrow 4 = 2^{2^{2^2}} = 2^{65536}.
\]
El límite ordinal $\varepsilon_0$ corresponde a la «$\omega$-torre» $\omega \uparrow\uparrow \omega$ en el universo ordinal, un límite que los números naturales solo pueden aproximar finítamente.
\end{ejemplo}

% ===========================================================
\section{Resumen de la clase}
% ===========================================================

\begin{importante}
\textbf{Resultados principales de la Clase~6:}
\begin{enumerate}[leftmargin=2em, label=\textbf{\arabic*.}]
  \item La \textbf{potenciación ordinal} $\alpha^\beta$ satisface leyes de exponentes análogas a las naturales: $\alpha^{\beta+\gamma} = \alpha^\beta \cdot \alpha^\gamma$ y $(\alpha^\beta)^\gamma = \alpha^{\beta\gamma}$. Sin embargo, $(\alpha\beta)^\gamma \neq \alpha^\gamma \beta^\gamma$ en general.

  \item El \textbf{Teorema de la Forma Normal de Cantor}: todo ordinal $\alpha > 0$ se escribe de manera única como $\alpha = \omega^{\beta_1}c_1 + \cdots + \omega^{\beta_k}c_k$ con $\beta_1 > \cdots > \beta_k$ y $c_i \in \mathbb{N}^+$.

  \item La función $\alpha \mapsto \omega^\alpha$ es estrictamente creciente, continua (para límites), y satisface $\alpha < \omega^\alpha$ para todo $\alpha \geq 1$.

  \item El ordinal $\varepsilon_0 = \sup\{\omega, \omega^\omega, \omega^{\omega^\omega}, \ldots\}$ es el \textbf{mínimo punto fijo} de $\alpha \mapsto \omega^\alpha$, es decir, $\omega^{\varepsilon_0} = \varepsilon_0$.

  \item La \textbf{jerarquía épsilon} $\{\varepsilon_\alpha : \alpha \in \mathbf{ON}\}$ forma una clase propia de puntos fijos de $\alpha \mapsto \omega^\alpha$, definida por recursión transfinita.

  \item La \textbf{cofinalidad} $\mathrm{cf}(\lambda)$ de un ordinal límite $\lambda$ es el mínimo tipo de orden de una sucesión cofinal en $\lambda$. Un ordinal límite es \textbf{regular} si $\mathrm{cf}(\lambda) = \lambda$, y \textbf{singular} en caso contrario.

  \item $\mathrm{cf}(\omega) = \omega$ (regular); $\mathrm{cf}(\omega^\omega) = \omega$ (singular); $\mathrm{cf}(\varepsilon_0) = \omega$ (singular). El primer cardinal regular no numerable es $\omega_1$.

  \item La forma normal de Cantor es «completa» para ordinales $< \varepsilon_0$; para ordinales $\geq \varepsilon_0$ involucra autoexponentes, lo que conecta con la teoría de la prueba (Teorema de Gentzen).
\end{enumerate}
\end{importante}

En la Clase~7 estudiaremos los \textbf{números cardinales} desde el punto de vista ordinal: definiremos los cardinales como ordinales iniciales, los \textbf{alephs} $\aleph_0, \aleph_1, \aleph_2, \ldots$, y desarrollaremos la aritmética cardinal, que presenta sus propias reglas y sorpresas.

% ===========================================================
\section{Problemas y tareas}
% ===========================================================

\begin{enumerate}[leftmargin=2.5em, label=\textbf{\arabic*.}]

  \item \textbf{(Potenciación ordinal: cálculos básicos.)}
  \begin{enumerate}[label=(\alph*)]
    \item Calcule los siguientes ordinales, justificando cada paso mediante la definición recursiva:
    \begin{enumerate}[label=(\roman*)]
      \item $\omega^2 \cdot \omega$.
      \item $(\omega+1)^2$.
      \item $(\omega+1)^\omega$.
      \item $2^{\omega+1}$.
      \item $\omega^{\omega \cdot 2}$.
    \end{enumerate}
    \item Demuestre que $\omega^{\omega+1} = \omega^\omega \cdot \omega$.
    \item Demuestre que $(\omega+1)^\omega = \omega^\omega$ sin usar la ley $(\alpha\beta)^\gamma = \alpha^\gamma \beta^\gamma$.
    \item Calcule $(\omega \cdot 2)^\omega$ y $(2 \cdot \omega)^\omega$. ¿Son iguales? Justifique.
  \end{enumerate}

  \item \textbf{(Leyes de exponentes ordinales.)}
  \begin{enumerate}[label=(\alph*)]
    \item Demuestre por inducción transfinita en $\gamma$ que $\alpha^{\beta+\gamma} = \alpha^\beta \cdot \alpha^\gamma$.
    \item Demuestre por inducción transfinita en $\gamma$ que $(\alpha^\beta)^\gamma = \alpha^{\beta\gamma}$.
    \item Encuentre un contraejemplo a la ley $(\alpha \cdot \beta)^\gamma = \alpha^\gamma \cdot \beta^\gamma$ con $\alpha, \beta, \gamma$ ordinales infinitos.
    \item ¿Es cierto que $\alpha^\beta \cdot \alpha^\gamma = \alpha^{\beta+\gamma}$ cuando $\alpha = 0$? Considere todos los casos ($\beta, \gamma$ cero, sucesor, límite).
  \end{enumerate}

  \item \textbf{(Forma Normal de Cantor.)}
  \begin{enumerate}[label=(\alph*)]
    \item Escriba en forma normal de Cantor los siguientes ordinales:
    \begin{enumerate}[label=(\roman*)]
      \item $\omega^5 + \omega^3 \cdot 7 + \omega^2 + \omega \cdot 4 + 11$.
      \item $\omega^{\omega+3} + \omega^\omega \cdot 6 + \omega^4$.
      \item $\omega^{\omega^2 + \omega} + \omega^{\omega^2} \cdot 3$.
    \end{enumerate}
    \item Dados $\alpha = \omega^3 \cdot 4 + \omega \cdot 2 + 5$ y $\beta = \omega^3 \cdot 2 + \omega^2 + 3$, calcule:
    \begin{enumerate}[label=(\roman*)]
      \item $\alpha + \beta$.
      \item $\beta + \alpha$.
      \item $\alpha \cdot \omega$.
      \item $\alpha \cdot 3$.
    \end{enumerate}
    \item Demuestre que si $\alpha = \omega^{\beta_1}c_1 + \cdots + \omega^{\beta_k}c_k$ es la forma normal y $\beta_k \geq 1$, entonces $\alpha$ es un ordinal límite.
    \item (Desafío) Complete la prueba de la unicidad en el Teorema de la Forma Normal de Cantor para el caso en que $\alpha$ tiene dos o más términos.
    \item (Desafío) Demuestre que la suma ordinal de dos ordinales en forma normal puede calcularse mediante el algoritmo del Corolario~\ref{cor:algsuma}.
  \end{enumerate}

  \item \textbf{(El ordinal $\varepsilon_0$.)}
  \begin{enumerate}[label=(\alph*)]
    \item Calcule los primeros cinco términos de la sucesión $\varphi_0 = \omega$, $\varphi_{n+1} = \omega^{\varphi_n}$. Escriba $\varphi_2, \varphi_3$ en términos de torres exponenciales.
    \item Demuestre que $\varphi_n < \varphi_{n+1}$ para todo $n \in \omega$.
    \item Demuestre directamente desde la definición que $\varepsilon_0 = \sup_{n \in \omega} \varphi_n$ satisface $\omega^{\varepsilon_0} = \varepsilon_0$.
    \item Demuestre que $\varepsilon_0$ es el \emph{mínimo} ordinal $\varepsilon$ con $\omega^\varepsilon = \varepsilon$, sin suponer que ya conocemos ningún punto fijo.
    \item (Desafío) Demuestre que para todo $n \in \omega$ se tiene $\varphi_n < \varepsilon_0$, y que todo ordinal $\alpha < \varepsilon_0$ satisface $\alpha < \varphi_n$ para algún $n$.
  \end{enumerate}

  \item \textbf{(Jerarquía épsilon.)}
  \begin{enumerate}[label=(\alph*)]
    \item Describa (sin demostración formal) cómo se construye $\varepsilon_1$ a partir de $\varepsilon_0$ por un proceso análogo al de $\varepsilon_0$ a partir de $\omega$.
    \item Demuestre que $\varepsilon_\omega = \sup_{n < \omega} \varepsilon_n$ satisface $\omega^{\varepsilon_\omega} = \varepsilon_\omega$.
    \item ¿Es $\varepsilon_0 + 1$ un punto fijo de $\alpha \mapsto \omega^\alpha$? Justifique.
    \item Demuestre que entre cualesquiera dos ordinales épsilon consecutivos $\varepsilon_\alpha < \varepsilon_{\alpha+1}$ existen ordinales que no son puntos fijos de $\alpha \mapsto \omega^\alpha$. Exhiba uno explícitamente.
    \item (Desafío) Demuestre que la clase de todos los ordinales épsilon es una clase propia.
  \end{enumerate}

  \item \textbf{(Cofinalidad.)}
  \begin{enumerate}[label=(\alph*)]
    \item Determine $\mathrm{cf}(\alpha)$ para cada uno de los siguientes ordinales, justificando su respuesta exhibiendo una sucesión cofinal:
    \begin{enumerate}[label=(\roman*)]
      \item $\omega \cdot \omega$.
      \item $\omega^3$.
      \item $\omega^\omega + \omega^3$.
      \item $\varepsilon_0$.
      \item $\varepsilon_\omega$.
    \end{enumerate}
    \item Demuestre que $\mathrm{cf}(\omega) = \omega$ mostrando que no existe una función cofinal $f: n \to \omega$ con $n$ finito.
    \item Demuestre que $\mathrm{cf}(\lambda) \leq \lambda$ para todo ordinal límite $\lambda$.
    \item Demuestre que $\mathrm{cf}(\mathrm{cf}(\lambda)) = \mathrm{cf}(\lambda)$ (la cofinalidad de la cofinalidad es idempotente).
    \item Sea $f: \alpha \to \beta$ estrictamente creciente y cofinal. Demuestre que $\mathrm{cf}(\alpha) = \mathrm{cf}(\beta)$.
  \end{enumerate}

  \item \textbf{(Ordinales regulares y singulares.)}
  \begin{enumerate}[label=(\alph*)]
    \item ¿Cuáles de los siguientes ordinales son regulares? $\omega$, $\omega+\omega$, $\omega^\omega$, $\varepsilon_0$, $\omega_1$. Justifique.
    \item Demuestre que $\omega$ es el único ordinal límite regular que es numerable.
    \item Enuncie (sin demostrar) el resultado que afirma que todo cardinal acessible $\omega_\alpha$ con $\alpha$ sucesor es regular.
    \item (Investigación) Consulte \cite{jech2003} y describa brevemente qué es un cardinal débilmente inaccesible y por qué su existencia no puede probarse en ZFC.
  \end{enumerate}

  \item \textbf{(Conexión con la teoría de la prueba.)}
  \begin{enumerate}[label=(\alph*)]
    \item Enuncie el Teorema de Gentzen (1936) que afirma que la consistencia de la aritmética de Peano es equivalente a la buena fundación de los ordinales $< \varepsilon_0$ bajo el sistema de notación de Cantor.
    \item Explique informalmente por qué $\varepsilon_0$ actúa como «medida de la fuerza» de la aritmética de Peano en la escala transfinita.
    \item (Proyecto) Consulte \cite{pohlers1989} o \cite{smorynski1991} y redacte un resumen de 1--2 páginas sobre el \emph{análisis de la prueba} (\textit{proof-theoretic ordinals}) para los sistemas formales más comunes.
  \end{enumerate}

  \item \textbf{(Síntesis general.)}
  \begin{enumerate}[label=(\alph*)]
    \item Ubique cada uno de los siguientes ordinales en la jerarquía $\omega < \omega^\omega < \varepsilon_0 < \varepsilon_1 < \varepsilon_\omega < \cdots$, indicando con precisión entre cuáles de estos hitos se encuentra:
    \begin{enumerate}[label=(\roman*)]
      \item $\omega^{\omega^3} + \omega^5$.
      \item $\omega^{\varepsilon_0}$.
      \item $\varepsilon_0 \cdot 2$.
      \item $\varepsilon_0 + \varepsilon_0$.
      \item $\varepsilon_0^{\varepsilon_0}$.
    \end{enumerate}
    \item ¿Cuál es la forma normal de Cantor de $\varepsilon_0 + \omega^{\varepsilon_0+1}$?
    \item Determine $\mathrm{cf}(\varepsilon_0^{\varepsilon_0})$.
    \item (Desafío) Determine el menor ordinal $\alpha$ tal que $\varepsilon_\alpha = \alpha$ (un «punto fijo de la jerarquía épsilon»). ¿Es este ordinal igual a alguno de los $\varepsilon_\beta$?
  \end{enumerate}

  \item \textbf{(Investigación y exploración.)}
  Consulte al menos dos de las referencias indicadas al final de esta clase y redacte una exposición de 1--2 páginas sobre \textbf{uno} de los siguientes temas:
  \begin{enumerate}[label=(\alph*)]
    \item La \textbf{función de Veblen} $\varphi_\alpha(\beta)$: defínala, calcule sus primeros valores y describa los ordinales $\Gamma_0$ (constante de Feferman-Schütte).
    \item Los \textbf{ordinales de Church-Kleene}: los ordinales recursivos y la notación $\omega_1^{CK}$ (el primer ordinal no recursivo).
    \item La \textbf{hipótesis del continuo} y su relación con los cardinales regulares y singulares.
  \end{enumerate}

\end{enumerate}

% ===========================================================
\section*{Referencias bibliográficas de esta clase}
% ===========================================================

Los resultados desarrollados en esta clase se encuentran tratados con rigor y detalle en las siguientes fuentes:

\begin{itemize}[leftmargin=2em]
  \item \textbf{Forma Normal de Cantor:} El resultado original se encuentra en \cite{cantor1897}, artículo en el que Cantor establece la representación en «base $\omega$» de los ordinales transfinitos. Presentaciones modernas accesibles se hallan en \cite{enderton1977}, \cite{hrbacek1999} y \cite{devlin1993}. La prueba completa con unicidad aparece en \cite{sierpinski1958} y en \cite{levy1979}.

  \item \textbf{Potenciación ordinal y aritmética avanzada:} La fuente más sistemática sigue siendo \cite{sierpinski1958}; los textos modernos de referencia son \cite{jech2003} (capítulo 2) y \cite{kunen2011} (capítulo I). Para una visión más introductoria, véase \cite{suppes1972}.

  \item \textbf{El ordinal $\varepsilon_0$ y los ordinales épsilon:} La jerarquía épsilon se estudia en detalle en \cite{jech2003} y \cite{levy1979}. La conexión con la función de Veblen se describe en \cite{pohlers1989}; la introducción más asequible es \cite{hrbacek1999}.

  \item \textbf{Función de Veblen y jerarquía más allá de $\varepsilon_0$:} La función de Veblen fue introducida por Oswald Veblen en 1908; su tratamiento moderno aparece en \cite{pohlers1989} y \cite{smorynski1991}.

  \item \textbf{Cofinalidad:} La noción de cofinalidad se trata en todos los textos avanzados de teoría de conjuntos; los mejores desarrollos para este nivel están en \cite{jech2003} (sección 3) y \cite{kunen2011}. Una presentación más introductoria se encuentra en \cite{hrbacek1999} y \cite{devlin1993}.

  \item \textbf{Conexión con teoría de la prueba y el Teorema de Gentzen:} El resultado de Gentzen (1936) sobre la consistencia de la aritmética de Peano mediante $\varepsilon_0$ se expone en \cite{gentzen1936}; presentaciones modernas accesibles aparecen en \cite{smorynski1991} y en el primer capítulo de \cite{pohlers1989}.

  \item \textbf{Grandes cardinales y cofinalidad:} La conexión entre cofinalidad, regularidad e inaccesibilidad se desarrolla en \cite{kanamori2003}, cuyo capítulo inicial sirve como referencia estándar.

  \item \textbf{Texto introductorio recomendado:} Para una primera lectura, \cite{halmos1960} y \cite{enderton1977} son las referencias de entrada más amigables; \cite{hrbacek1999} es el texto de nivel intermedio más completo para cubrir todos los temas de esta clase.
\end{itemize}

% ===========================================================
\chapter{Clase 7: Números Cardinales Transfinitos — Alephs y Teorema de Cantor}
% ===========================================================

Los \textbf{números cardinales} dan respuesta formal a la pregunta más antigua de la aritmética infinita: ¿cuándo dos conjuntos tienen el mismo «tamaño»? Georg Cantor respondió en los años 1870–1890 con la noción de equipotencia y con la radical observación de que no todos los conjuntos infinitos son equipotentes. Las clases 4 y 5 sentaron las bases: equipotencia, la desigualdad de Cantor-Schröder-Bernstein y los ordinales de von Neumann. En esta clase se corona ese programa con tres logros fundamentales:

\begin{enumerate}[leftmargin=2em, label=(\roman*)]
  \item \textbf{Cardinales como ordinales iniciales}: se define formalmente el cardinal de un conjunto —dentro de ZFC con el Axioma de Elección— como el menor ordinal equipotente a él, y se estudia la estructura de la clase de cardinales.
  \item \textbf{La jerarquía aleph}: se construye la sucesión transfinita $\aleph_0, \aleph_1, \aleph_2, \ldots, \aleph_\omega, \ldots$ de cardinales infinitos y se prueba que cada cardinal infinito es un $\aleph$.
  \item \textbf{El Teorema de Cantor y los beth numbers}: se demuestra que para todo conjunto $A$ se cumple $|A| < |\mathcal{P}(A)|$, se introduce la jerarquía $\beth_0, \beth_1, \beth_2, \ldots$ y se contrasta con la jerarquía aleph, anticipando la Hipótesis del Continuo.
\end{enumerate}

Al final se hace una primera presentación de los \textbf{cardinales inaccesibles}, que escapan a la jerarquía construida por las operaciones cardinales ordinarias y constituyen el punto de entrada a la teoría de grandes cardinales \cite{kanamori2003, jech2003}.

La exposición sigue los textos clásicos \cite{jech2003, kunen2011, enderton1977, hrbacek1999, levy1979, sierpinski1958}, con énfasis en la motivación intuitiva y en los detalles formales de cada demostración.

% ===========================================================
\section{Cardinales como ordinales iniciales}
% ===========================================================

\subsection{Motivación: el problema del representante canónico}

La Clase~4 introdujo la \emph{equipotencia}: dos conjuntos $A$ y $B$ son equipotentes, escrito $A \sim B$, si existe una biyección $f \colon A \to B$. La equipotencia es una relación de equivalencia sobre la clase universal; las clases de equivalencia son las «cardinalidades». Sin embargo, las clases de equivalencia no son conjuntos en ZFC (la clase de todos los conjuntos con dos elementos, por ejemplo, no es un conjunto). Para trabajar con cardinales como objetos matemáticos ordinarios, necesitamos elegir un representante canónico en cada clase.

\begin{itemize}[leftmargin=2em]
  \item Para los conjuntos \emph{finitos} la solución es natural: el representante de la clase de un conjunto con $n$ elementos es el número natural $n = \{0,1,\ldots,n-1\}$.
  \item Para los conjuntos \emph{infinitos} la solución requiere el Axioma de Elección: el representante de la clase de $A$ es el \emph{menor ordinal} equipotente a $A$.
\end{itemize}

\subsection{Ordinales iniciales}

\begin{definicion}{Ordinal inicial}{ordinalinicial}
Un ordinal $\kappa$ se denomina \textbf{ordinal inicial} (o \textbf{cardinal de von Neumann}) si no es equipotente a ningún ordinal $\alpha < \kappa$:
\[
  \forall \alpha < \kappa \;:\; \alpha \not\sim \kappa.
\]
Equivalentemente, $\kappa$ es un ordinal inicial si $|\kappa| \neq |\alpha|$ para todo $\alpha < \kappa$.
\end{definicion}

\begin{observacion}{Todo natural es un ordinal inicial}{natinicial}
Los ordinales finitos $0, 1, 2, 3, \ldots$ son todos ordinales iniciales: el ordinal $n$ tiene exactamente $n$ elementos, y ningún ordinal menor que $n$ tiene $n$ elementos.
\end{observacion}

\begin{definicion}{Número cardinal (en ZFC)}{numcardinalzfc}
Dado un conjunto $A$, su \textbf{número cardinal} $|A|$ (también denotado $\mathrm{card}(A)$) es el único ordinal inicial equipotente a $A$. Formalmente, asumiendo el Axioma de Elección:
\[
  |A| := \min\bigl\{\alpha \in \mathbf{ON} \mid \alpha \sim A\bigr\}.
\]
Este mínimo existe porque, por el Teorema del Bien-Orden (equivalente al AC), todo conjunto puede ser bien ordenado y, en consecuencia, es equipotente a algún ordinal.
\end{definicion}

\begin{observacion}{Sin el Axioma de Elección}{sinchoice}
Sin el Axioma de Elección no todo conjunto puede ser bien ordenado, y la definición anterior deja de funcionar. Existen formulaciones alternativas de cardinal (por ejemplo, usando rangos de Scott), pero pierden muchas propiedades agradables. En estas notas siempre trabajamos en ZFC.
\end{observacion}

\begin{ejemplo}{Los primeros cardinales}{primcardinals}
\begin{enumerate}[leftmargin=2em, label=(\alph*)]
  \item $|\emptyset| = 0$, $|\{a\}| = 1$, $|\{a,b\}| = 2$ para $a \neq b$.
  \item $|\mathbb{N}| = \omega$. En efecto, $\mathbb{N} \sim \omega$ (la función identidad es una biyección), y ningún ordinal finito $n$ es equipotente a $\omega$ (pues $\omega$ es infinito). Por tanto $\omega$ es el cardinal de $\mathbb{N}$.
  \item $|\mathbb{Z}| = \omega$. La función $n \mapsto \begin{cases} 2n & n \geq 0 \\ -2n-1 & n < 0 \end{cases}$ es una biyección $\mathbb{Z} \to \mathbb{N}$.
  \item $|\mathbb{Q}| = \omega$. El argumento diagonal de Cantor muestra que $\mathbb{Q}$ es numerable.
\end{enumerate}
\end{ejemplo}

\begin{proposicion}{Propiedades básicas de cardinales}{propcardbasicas}
Sean $A$, $B$ conjuntos. Entonces:
\begin{enumerate}[leftmargin=2em, label=(\roman*)]
  \item $|A| = |B|$ si y solo si $A \sim B$.
  \item $|A| \leq |B|$ (definido como: existe una inyección $A \hookrightarrow B$) si y solo si $|A| \leq_{\mathrm{ON}} |B|$ (orden usual de ordinales).
  \item $|A| < |B|$ si y solo si $|A| \leq |B|$ y $|A| \neq |B|$.
  \item Si $A \subseteq B$, entonces $|A| \leq |B|$.
\end{enumerate}
\end{proposicion}

\begin{proof}
(i) Si $|A| = |B|$, ambos son el mismo ordinal inicial $\kappa$; las biyecciones $A \to \kappa$ y $B \to \kappa$ componen para dar $A \sim B$. Recíprocamente, si $A \sim B$, el menor ordinal equipotente a $A$ también lo es a $B$, de modo que $|A| = |B|$. (ii)--(iv) son consecuencias directas de la definición y del teorema de Cantor-Schröder-Bernstein \cite{enderton1977}.
\end{proof}

\subsection{La clase de los cardinales infinitos}

\begin{definicion}{Cardinal infinito}{cardinfinito}
Un cardinal $\kappa$ es \textbf{infinito} si $\omega \leq \kappa$, es decir, si contiene a $\omega$ como subconjunto (o equivalentemente, si $\kappa$ no es un natural).
\end{definicion}

\begin{teorema}{Caracterización de los cardinales infinitos}{carinfinito}
Un ordinal $\kappa$ es un cardinal infinito si y solo si $\kappa$ es un ordinal inicial y $\kappa \geq \omega$.
\end{teorema}

\begin{proof}
Por definición, $\kappa$ es cardinal infinito si y solo si es un ordinal inicial mayor o igual que $\omega$. Los ordinales iniciales finitos son exactamente $0, 1, 2, \ldots$; el primer ordinal inicial infinito es $\omega$.
\end{proof}

% ===========================================================
\section{La sucesión de Alephs}
% ===========================================================

\subsection{Construcción de la jerarquía aleph}

La teoría ordinal permite construir, por recursión transfinita, la sucesión de todos los cardinales infinitos. Esta sucesión se indexa por los ordinales y produce los llamados \textbf{números aleph} (en referencia a la primera letra del alfabeto hebreo, $\aleph$).

\begin{definicion}{Números aleph}{numsaleph}
La función $\alpha \mapsto \aleph_\alpha$ se define por recursión transfinita en $\alpha$:
\begin{align*}
  \aleph_0 &:= \omega, \\
  \aleph_{\alpha+1} &:= \omega_{\alpha+1} := \text{el menor cardinal mayor que } \aleph_\alpha, \\
  \aleph_\lambda &:= \sup_{\alpha < \lambda} \aleph_\alpha \quad \text{para } \lambda \text{ un ordinal límite.}
\end{align*}
Cada $\aleph_\alpha$ es también denotado $\omega_\alpha$ cuando se enfatiza su naturaleza ordinal.
\end{definicion}

\begin{observacion}{Existencia del sucesor cardinal}{existsucc}
El «menor cardinal mayor que $\kappa$» existe siempre. En efecto, por el Teorema del Bien-Orden, $\mathcal{P}(\kappa)$ puede ser bien ordenado, luego es equipotente a algún ordinal $\lambda > \kappa$. El conjunto de ordinales iniciales estrictamente mayores que $\kappa$ es no vacío, y su mínimo es el \emph{sucesor cardinal} de $\kappa$, denotado $\kappa^+$.
\end{observacion}

\begin{definicion}{Cardinal sucesor y cardinal límite}{cardsuclim}
Sea $\kappa$ un cardinal infinito.
\begin{itemize}[leftmargin=2em]
  \item El \textbf{sucesor cardinal} de $\kappa$, denotado $\kappa^+$, es el menor cardinal estrictamente mayor que $\kappa$.
  \item Un cardinal infinito $\lambda$ es un \textbf{cardinal límite} si no es el sucesor cardinal de ningún cardinal, es decir, si para todo cardinal $\kappa < \lambda$ se tiene $\kappa^+ < \lambda$.
\end{itemize}
\end{definicion}

\begin{ejemplo}{Los primeros alephs}{primerosaleph}
Calculamos explícitamente los primeros valores:
\begin{enumerate}[leftmargin=2em, label=(\alph*)]
  \item $\aleph_0 = \omega$. Es el cardinal de $\mathbb{N}$, $\mathbb{Z}$, $\mathbb{Q}$, y de cualquier conjunto numerable infinito.
  \item $\aleph_1 = \omega_1$. Es el primer cardinal no numerable. Por definición, $\aleph_1 = (\aleph_0)^+$ es el menor ordinal inicial no numerable. Contiene a todos los ordinales numerables.
  \item $\aleph_2 = \omega_2 = (\aleph_1)^+$. El segundo cardinal no numerable.
  \item $\aleph_n$ para cada $n \in \mathbb{N}$: el $n$-ésimo cardinal no numerable (el $(n-1)$-ésimo sucesor de $\aleph_0$).
  \item $\aleph_\omega = \sup_{n < \omega} \aleph_n$. Es el primer cardinal límite no numerable, igual al supremo de $\aleph_0, \aleph_1, \aleph_2, \ldots$
  \item $\aleph_{\omega+1} = (\aleph_\omega)^+$.
  \item $\aleph_{\omega \cdot 2}, \aleph_{\omega^2}, \aleph_{\omega^\omega}, \aleph_{\varepsilon_0}, \ldots$ continúan la jerarquía.
\end{enumerate}
\end{ejemplo}

\subsection{Todo cardinal infinito es un aleph}

El siguiente teorema, clave para la teoría cardinal, afirma que los $\aleph_\alpha$ agotan \emph{todos} los cardinales infinitos.

\begin{teorema}{Los alephs son todos los cardinales infinitos}{alephtodos}
Para todo cardinal infinito $\kappa$ existe un único ordinal $\alpha$ tal que $\kappa = \aleph_\alpha$.
\end{teorema}

\begin{proof}
Sea $\kappa$ un cardinal infinito. Definimos $\alpha$ como el tipo de orden del conjunto bien ordenado
\[
  S := \{\lambda \in \mathbf{Card} \mid \lambda \text{ es un cardinal infinito y } \lambda < \kappa\},
\]
donde el orden es el usual de los ordinales. $S$ es un conjunto (no una clase propia) porque es un subconjunto de $\kappa$.

Afirmamos que $\alpha$ es el ordinal inicial del tipo de orden de $S$, y que la función $\beta \mapsto \aleph_\beta$ es un isomorfismo de orden de $\alpha$ en $S \cup \{\kappa\}$. Esto se verifica por inducción transfinita en $\alpha$: el primer cardinal infinito mayor que todos los elementos de $S$ de índice $< \beta$ debe ser $\aleph_\beta$ por la definición recursiva. Por tanto $\kappa = \aleph_\alpha$ \cite{jech2003, kunen2011}.
\end{proof}

\begin{corolario}{Unicidad del índice aleph}{unicidadaleph}
La función $\alpha \mapsto \aleph_\alpha$ es una biyección estrictamente creciente entre la clase de todos los ordinales $\mathbf{ON}$ y la clase de todos los cardinales infinitos.
\end{corolario}

\subsection{Aritmética cardinal: sumas y productos}

\begin{proposicion}{Suma y producto de cardinales infinitos}{aritcardinfmain}
Sea $\kappa$ un cardinal infinito y $\lambda$ un cardinal con $1 \leq \lambda \leq \kappa$. Entonces:
\[
  \kappa + \lambda = \kappa \cdot \lambda = \kappa.
\]
En particular, $\kappa + \kappa = \kappa$ y $\kappa \cdot \kappa = \kappa$.
\end{proposicion}

\begin{proof}
Probamos $\kappa \cdot \kappa = \kappa$ por inducción transfinita en $\kappa$. Esto es el Teorema de Hessenberg; su prueba utiliza que $\kappa \times \kappa$ puede ser bien ordenado con el orden lexicográfico por pares, y que este bien-orden tiene tipo ordinal $\kappa$ \cite{jech2003, enderton1977}. La igualdad $\kappa + \kappa = \kappa$ se sigue de $\kappa \cdot \kappa = \kappa$.
\end{proof}

\begin{ejemplo}{Cardinalidades conocidas}{cardejemplos}
\begin{enumerate}[leftmargin=2em, label=(\alph*)]
  \item $|\mathbb{N} \times \mathbb{N}| = \aleph_0 \cdot \aleph_0 = \aleph_0$.
  \item $\bigcup_{n \in \mathbb{N}} \mathbb{N}^n$ (unión de todos los $n$-uplos de naturales) tiene cardinal $\aleph_0$.
  \item $|\mathbb{R}|$ no es $\aleph_0$ (por el argumento diagonal de Cantor) pero ¿cuál $\aleph$ es? La Hipótesis del Continuo (HC) afirma que $|\mathbb{R}| = \aleph_1$; la negación de HC es consistente con ZFC \cite{cohen1966, godel1940}.
\end{enumerate}
\end{ejemplo}

% ===========================================================
\section{El Teorema de Cantor}
% ===========================================================

\subsection{Enunciado y primer comentario}

El \textbf{Teorema de Cantor} —publicado en 1891 \cite{cantor1891}— es uno de los resultados más elegantes y fértiles de la matemática. Afirma que el conjunto de partes de cualquier conjunto es estrictamente más grande que el conjunto original; en particular, no existe un conjunto «más grande que todos» ni hay un cardinal máximo.

\begin{teorema}{Teorema de Cantor}{teocantor}
Para todo conjunto $A$:
\[
  |A| < |\mathcal{P}(A)|.
\]
Más precisamente: existe una inyección $A \hookrightarrow \mathcal{P}(A)$ pero no existe ninguna sobreyección $A \twoheadrightarrow \mathcal{P}(A)$.
\end{teorema}

\begin{proof}
\textbf{Paso 1: $|A| \leq |\mathcal{P}(A)|$.} La función $a \mapsto \{a\}$ es una inyección $A \hookrightarrow \mathcal{P}(A)$.

\textbf{Paso 2: $|\mathcal{P}(A)| \not\leq |A|$}, es decir, no existe ninguna sobreyección $f \colon A \twoheadrightarrow \mathcal{P}(A)$.

Sea $f \colon A \to \mathcal{P}(A)$ una función cualquiera. Definimos:
\[
  D := \bigl\{a \in A \mid a \notin f(a)\bigr\}.
\]
Claramente $D \subseteq A$, luego $D \in \mathcal{P}(A)$. Afirmamos que $D \notin \mathrm{Im}(f)$.

Supongamos por reducción al absurdo que $D = f(d)$ para algún $d \in A$. Entonces:
\[
  d \in D \iff d \notin f(d) \iff d \notin D.
\]
Contradicción. Por tanto $D \notin \mathrm{Im}(f)$, y $f$ no es sobreyectiva. Como $f$ fue arbitraria, ninguna función $A \to \mathcal{P}(A)$ es sobreyectiva.

Combinando los dos pasos: $|A| \leq |\mathcal{P}(A)|$ y $|A| \neq |\mathcal{P}(A)|$, luego $|A| < |\mathcal{P}(A)|$.
\end{proof}

\begin{observacion}{La diagonal de Cantor}{diagcantor}
El conjunto $D = \{a \in A \mid a \notin f(a)\}$ es la famosa \emph{diagonal de Cantor}. Esta construcción —una función que deliberadamente «evade» a $f$— es el arquetipo del argumento diagonal, que reaparece en el Teorema de la Incompletitud de Gödel, en el Problema de la Parada de Turing, en la insolubilidad del problema de la correspondencia de Post, y en muchos resultados de límite en lógica y computación \cite{enderton2001, mendelson2015}.
\end{observacion}

\begin{corolario}{No hay conjunto de todos los cardinales}{noconjcard}
No existe ningún conjunto que contenga a todos los cardinales. En particular, la clase $\mathbf{Card}$ de todos los cardinales es una clase propia.
\end{corolario}

\begin{proof}
Supongamos que existiera un conjunto $X$ que contiene a todos los cardinales. Sea $\kappa = |\bigcup X|$. Entonces $\kappa$ es un cardinal, luego $\kappa \in X \subseteq \bigcup X$. Pero entonces $\kappa \leq |\bigcup X| = \kappa$, y en particular $\kappa \in \bigcup X$ implica que hay un $A \in X$ con $\kappa \in A$. Por el Teorema de Cantor $|\mathcal{P}(\bigcup X)| > \kappa$, y $\mathcal{P}(\bigcup X)$ es equipotente a algún cardinal $\mu > \kappa$ que debería estar en $X$, pero $\mu > |\bigcup X|$, contradicción.

Una forma más directa: si $C$ fuera el conjunto de todos los cardinales, $\bigcup C$ sería un conjunto, y $|\mathcal{P}(\bigcup C)|$ sería un cardinal mayor que todo cardinal en $C$, contradicción \cite{jech2003}.
\end{proof}

\begin{corolario}{Sucesor de todo cardinal}{succexiste}
Para todo cardinal $\kappa$ existe un cardinal estrictamente mayor: su sucesor cardinal $\kappa^+$. En particular, la jerarquía de cardinales no tiene cima.
\end{corolario}

\begin{proof}
Por el Teorema de Cantor, $|\mathcal{P}(\kappa)| > \kappa$. El menor ordinal inicial mayor que $\kappa$ es $\kappa^+$, y satisface $\kappa < \kappa^+ \leq |\mathcal{P}(\kappa)|$.
\end{proof}

\subsection{El Teorema de Cantor para conjuntos concretos}

\begin{ejemplo}{Cantor aplicado a $\mathbb{N}$}{cantorN}
El conjunto $\mathcal{P}(\mathbb{N})$ no es numerable. Concretamente, toda función $f \colon \mathbb{N} \to \mathcal{P}(\mathbb{N})$ omite el conjunto
\[
  D = \{n \in \mathbb{N} \mid n \notin f(n)\}.
\]

Visualicemos $f$ como una tabla de 0s y 1s, donde la fila $n$ indica qué naturales pertenecen a $f(n)$:

\medskip
\begin{center}
\begin{tabular}{c|ccccccc}
 & 0 & 1 & 2 & 3 & 4 & 5 & $\cdots$ \\
\hline
$f(0)$ & \textbf{1} & 0 & 1 & 0 & 1 & 0 & $\cdots$ \\
$f(1)$ & 0 & \textbf{0} & 0 & 1 & 1 & 0 & $\cdots$ \\
$f(2)$ & 1 & 1 & \textbf{1} & 0 & 0 & 1 & $\cdots$ \\
$f(3)$ & 0 & 0 & 0 & \textbf{1} & 0 & 0 & $\cdots$ \\
$f(4)$ & 1 & 1 & 1 & 1 & \textbf{0} & 1 & $\cdots$ \\
\vdots & & & & & & & $\ddots$
\end{tabular}
\end{center}
\medskip

La diagonal marcada en negrita es $1, 0, 1, 1, 0, \ldots$ El conjunto $D$ invierte cada bit de la diagonal: $D$ contiene $1$ (porque $f(1)$ no lo contiene), $D$ no contiene $0$ (porque $f(0)$ sí lo contiene), etc. Así $D \neq f(n)$ para todo $n$.
\end{ejemplo}

\begin{ejemplo}{Cantor aplicado a $\mathbb{R}$}{cantorR}
De manera similar, $|\mathbb{R}| = |\mathcal{P}(\mathbb{N})|$ (se prueba construyendo biyecciones explícitas vía expansiones binarias). Por tanto $|\mathbb{R}| > \aleph_0$. Los reales no son numerables, como demostró Cantor en 1874 \cite{cantor1883} mediante el argumento de los intervalos encajados, y reiteró en 1891 \cite{cantor1891} con el argumento diagonal.
\end{ejemplo}

\subsection{Aritmética de la exponenciación cardinal}

\begin{definicion}{Exponenciación cardinal}{expcard}
Para cardinales $\kappa$ y $\lambda$, la \textbf{potencia cardinal} $\kappa^\lambda$ se define como el cardinal del conjunto de todas las funciones de $\lambda$ en $\kappa$:
\[
  \kappa^\lambda := |\kappa^\lambda| = \bigl|\{f \mid f \colon \lambda \to \kappa\}\bigr|.
\]
\end{definicion}

\begin{observacion}{Relación con el conjunto de partes}{expcardpow}
En particular, $2^\kappa = |\mathcal{P}(\kappa)|$, pues cada subconjunto de $\kappa$ se identifica con su función característica $\kappa \to 2 = \{0,1\}$. Por tanto el Teorema de Cantor afirma:
\[
  \kappa < 2^\kappa \quad \text{para todo cardinal } \kappa.
\]
\end{observacion}

\begin{proposicion}{Propiedades de la exponenciación cardinal}{propexpcard}
Para cardinales $\kappa, \lambda, \mu$:
\begin{enumerate}[leftmargin=2em, label=(\roman*)]
  \item $\kappa^{\lambda + \mu} = \kappa^\lambda \cdot \kappa^\mu$.
  \item $(\kappa^\lambda)^\mu = \kappa^{\lambda \cdot \mu}$.
  \item $(\kappa \cdot \lambda)^\mu = \kappa^\mu \cdot \lambda^\mu$.
  \item Si $\kappa \leq \lambda$, entonces $\kappa^\mu \leq \lambda^\mu$.
  \item $2^{\aleph_0} \geq \aleph_1$ (por el Teorema de Cantor).
\end{enumerate}
\end{proposicion}

% ===========================================================
\section{Beth Numbers}
% ===========================================================

\subsection{Definición y primeros valores}

Los \textbf{beth numbers} (del hebreo $\beth$, la segunda letra del alfabeto) forman una jerarquía de cardinales que crece «maximalmente» mediante potencias de 2, partiendo de $\aleph_0$.

\begin{definicion}{Beth numbers}{bethnumbers}
La función $\alpha \mapsto \beth_\alpha$ se define por recursión transfinita:
\begin{align*}
  \beth_0 &:= \aleph_0 = \omega, \\
  \beth_{\alpha+1} &:= 2^{\beth_\alpha}, \\
  \beth_\lambda &:= \sup_{\alpha < \lambda} \beth_\alpha \quad \text{para } \lambda \text{ límite.}
\end{align*}
\end{definicion}

\begin{ejemplo}{Los primeros beth numbers}{primerosbeths}
\begin{enumerate}[leftmargin=2em, label=(\alph*)]
  \item $\beth_0 = \aleph_0$: el cardinal de los naturales.
  \item $\beth_1 = 2^{\aleph_0}$: el cardinal del continuo, $|\mathbb{R}| = |\mathcal{P}(\mathbb{N})|$.
  \item $\beth_2 = 2^{\beth_1} = 2^{2^{\aleph_0}}$: el cardinal de $\mathcal{P}(\mathbb{R})$, también igual al cardinal del conjunto de todas las funciones $\mathbb{R} \to \mathbb{R}$.
  \item $\beth_3 = 2^{\beth_2}$: el cardinal de $\mathcal{P}(\mathcal{P}(\mathbb{R}))$.
  \item $\beth_\omega = \sup_{n < \omega} \beth_n$: el primer beth límite.
\end{enumerate}
\end{ejemplo}

\begin{observacion}{La jerarquía beth es estrictamente creciente}{bethcrece}
Por el Teorema de Cantor, $\beth_\alpha < 2^{\beth_\alpha} = \beth_{\alpha+1}$ para todo $\alpha$, de modo que la sucesión $(\beth_\alpha)_{\alpha \in \mathbf{ON}}$ es estrictamente creciente.
\end{observacion}

\subsection{La jerarquía de von Neumann y los beth numbers}

La jerarquía de conjuntos $V_\alpha$ (jerarquía von Neumann) está directamente relacionada con los beth numbers:

\begin{proposicion}{Cardinal de $V_\alpha$}{cardvonneumann}
Para ordinales $\alpha$ finitos:
\begin{align*}
  |V_0| &= 0, \\
  |V_{n+1}| &= 2^{|V_n|}, \\
  |V_\omega| &= \aleph_0.
\end{align*}
Para ordinales transfinitos: $|V_{\omega + \alpha}| = \beth_\alpha$. En particular, $|V_{\omega+1}| = \beth_1$, $|V_{\omega+2}| = \beth_2$, etc.
\end{proposicion}

\subsection{Comparación entre alephs y beths}

La relación entre las dos jerarquías es el corazón de la teoría cardinal moderna:

\begin{importante}
Por el Teorema de Cantor, $\beth_\alpha \geq \aleph_\alpha$ para todo $\alpha$. La igualdad $\beth_\alpha = \aleph_\alpha$ para todos los $\alpha$ es la \textbf{Hipótesis Generalizada del Continuo} (GCH). La igualdad $\beth_1 = \aleph_1$ es la \textbf{Hipótesis del Continuo} (CH) de Cantor.
\end{importante}

\begin{ejemplo}{Desigualdades entre alephs y beths}{alephbeth}
\begin{enumerate}[leftmargin=2em, label=(\alph*)]
  \item Siempre: $\aleph_0 = \beth_0$.
  \item Siempre: $\aleph_1 \leq \beth_1 = 2^{\aleph_0}$.
  \item \textit{Si CH}: $\aleph_1 = \beth_1 = 2^{\aleph_0}$.
  \item \textit{Si $\neg$CH}: $\aleph_1 < \beth_1 = 2^{\aleph_0}$; de hecho, $2^{\aleph_0}$ puede ser cualquier $\aleph_\alpha$ con $\alpha$ cofinal regular con $\omega$ (Teorema de Easton \cite{jech2003}).
  \item Siempre: $\aleph_\omega \leq \beth_\omega$ y $\beth_\omega$ es un cardinal límite.
\end{enumerate}
\end{ejemplo}

\begin{teorema}{König}{teokoning}
Para todo cardinal infinito $\kappa$: $\mathrm{cf}(2^\kappa) > \kappa$, donde $\mathrm{cf}$ denota la cofinalidad. En particular, $\mathrm{cf}(2^{\aleph_0}) > \aleph_0$, es decir, $2^{\aleph_0} \neq \aleph_\omega$.
\end{teorema}

\begin{observacion}{Consecuencia inmediata}{koningcons}
El Teorema de König impone restricciones a los posibles valores de $2^{\aleph_0}$. Por ejemplo, $2^{\aleph_0} \neq \aleph_\omega$ (ya que $\mathrm{cf}(\aleph_\omega) = \omega$), ni puede ser $\aleph_{\omega_1}$ (ya que $\mathrm{cf}(\aleph_{\omega_1}) = \omega_1 \leq \aleph_0^+$); pero podría ser $\aleph_1, \aleph_2, \aleph_{\omega+1}$, etc.
\end{observacion}

% ===========================================================
\section{Cardinales Inaccesibles: una primera noción}
% ===========================================================

\subsection{Motivación}

La jerarquía aleph nos da todos los cardinales infinitos; la aritmética cardinal (suma, producto, potencia) genera nuevos cardinales a partir de los dados. Sin embargo, hay una pregunta natural: ¿existe algún cardinal que sea «tan grande» que no pueda alcanzarse desde abajo mediante ninguna de estas operaciones?

\begin{definicion}{Cardinal débilmente inaccesible}{carddebilinaccesible}
Un cardinal $\kappa > \aleph_0$ es \textbf{débilmente inaccesible} si:
\begin{enumerate}[leftmargin=2em, label=(\roman*)]
  \item $\kappa$ es un cardinal límite: no es sucesor de ningún cardinal.
  \item $\kappa$ tiene cofinalidad regular: $\mathrm{cf}(\kappa) = \kappa$ (es decir, ninguna sucesión de longitud $< \kappa$ cofinal en $\kappa$ existe).
\end{enumerate}
\end{definicion}

\begin{definicion}{Cardinal fuertemente inaccesible}{cardfuerteinaccesible}
Un cardinal $\kappa > \aleph_0$ es \textbf{fuertemente inaccesible} si:
\begin{enumerate}[leftmargin=2em, label=(\roman*)]
  \item $\kappa$ es débilmente inaccesible.
  \item Para todo cardinal $\lambda < \kappa$: $2^\lambda < \kappa$.
\end{enumerate}
Bajo GCH, las nociones débil y fuerte coinciden.
\end{definicion}

\begin{observacion}{¿Por qué «inaccesible»?}{inaccmotiv}
$\kappa$ es fuertemente inaccesible porque es imposible «alcanzarlo» desde abajo:
\begin{itemize}[leftmargin=2em]
  \item No puede alcanzarse por suma/producto de $< \kappa$ cardinales menores (por regularidad).
  \item No puede alcanzarse por potenciación $2^\lambda$ con $\lambda < \kappa$.
  \item No es el sucesor de ningún cardinal menor.
\end{itemize}
Los cardinales inaccesibles son los «universos» de la teoría de conjuntos: si $\kappa$ es fuertemente inaccesible, entonces $V_\kappa$ es un modelo de ZFC \cite{kunen2011, kanamori2003}.
\end{observacion}

\begin{observacion}{Independencia de ZFC}{indepzfc}
La existencia de cardinales inaccesibles no puede probarse desde ZFC (si ZFC es consistente): tal existencia es independiente de los axiomas de Zermelo-Fraenkel. Es una hipótesis de \emph{gran cardinal} que, junto con muchas otras hipótesis progresivamente más fuertes (cardinales de Mahlo, cardinales medibles, cardinales supercompactos, etc.), forma la escala jerárquica de la \emph{teoría de grandes cardinales} \cite{kanamori2003}.
\end{observacion}

\begin{ejemplo}{Los alephs no son inaccesibles (en general)}{alephnosinac}
\begin{enumerate}[leftmargin=2em, label=(\alph*)]
  \item $\aleph_1 = \aleph_0^+$ es un cardinal sucesor, luego no es límite: no es débilmente inaccesible.
  \item $\aleph_\omega$ es un cardinal límite, pero $\mathrm{cf}(\aleph_\omega) = \omega < \aleph_\omega$: no es regular, luego no es débilmente inaccesible.
  \item $\aleph_{\omega_1}$ es un cardinal límite con $\mathrm{cf}(\aleph_{\omega_1}) = \omega_1$: es débilmente inaccesible (bajo ciertos supuestos). Si además $2^\lambda < \aleph_{\omega_1}$ para todo $\lambda < \aleph_{\omega_1}$, es fuertemente inaccesible.
  \item Si $\kappa$ es fuertemente inaccesible, entonces $\kappa = \aleph_\kappa = \beth_\kappa$; es un punto fijo de ambas jerarquías.
\end{enumerate}
\end{ejemplo}

% ===========================================================
\section{La Hipótesis del Continuo: anticipación}
% ===========================================================

\subsection{Formulación de la HC y de la HGC}

\begin{importante}
\textbf{Hipótesis del Continuo (HC):} $2^{\aleph_0} = \aleph_1$, es decir, $\beth_1 = \aleph_1$.\\[0.3em]
\textbf{Hipótesis Generalizada del Continuo (HGC):} Para todo ordinal $\alpha$, $2^{\aleph_\alpha} = \aleph_{\alpha+1}$, es decir, $\beth_\alpha = \aleph_\alpha$ para todo $\alpha$.
\end{importante}

\begin{observacion}{Equivalencias de la HC}{hcequiv}
Las siguientes afirmaciones son equivalentes bajo ZFC:
\begin{enumerate}[leftmargin=2em, label=(\roman*)]
  \item $2^{\aleph_0} = \aleph_1$ (formulación de Cantor).
  \item No existe ningún cardinal $\kappa$ con $\aleph_0 < \kappa < 2^{\aleph_0}$ (no hay cardinalidades «intermedias» entre $\mathbb{N}$ y $\mathbb{R}$).
  \item $|\mathbb{R}| = \aleph_1$.
\end{enumerate}
\end{observacion}

\begin{observacion}{Estado metamatemático de la HC}{hcmeta}
La HC fue formulada por Cantor en 1878 y colocada por Hilbert como primer problema de su famosa lista de 1900.
\begin{itemize}[leftmargin=2em]
  \item En 1938, Gödel demostró que HC es \emph{consistente} con ZFC: no puede refutarse desde los axiomas de ZFC \cite{godel1940}.
  \item En 1963, Cohen demostró que $\neg$HC también es consistente con ZFC: no puede probarse desde los axiomas de ZFC \cite{cohen1966}.
\end{itemize}
Por tanto, la HC es \textbf{independiente} de ZFC: ninguna de las dos respuestas (HC o $\neg$HC) puede deducirse de ZFC. Este resultado fue calificado por muchos matemáticos como uno de los más profundos del siglo XX.
\end{observacion}

% ===========================================================
\section{Síntesis: diagrama de la jerarquía cardinal}
% ===========================================================

La siguiente tabla resume las relaciones entre los distintos niveles de la jerarquía cardinal, comparando alephs, beths y sus significados conjuntistas:

\begin{center}
\renewcommand{\arraystretch}{1.5}
\begin{tabular}{c|c|l}
$\aleph_\alpha$ & $\beth_\alpha$ & Significado \\ \hline
$\aleph_0$ & $\beth_0 = \aleph_0$ & $|\mathbb{N}|$, cardinal de los conjuntos numerables \\
$\aleph_1$ & $\beth_1 = 2^{\aleph_0}$ & Primer cardinal no numerable; $|\mathbb{R}|$ bajo CH \\
$\aleph_2$ & $\beth_2 = 2^{2^{\aleph_0}}$ & $|\mathcal{P}(\mathbb{R})|$ bajo CH \\
$\vdots$ & $\vdots$ & $\vdots$ \\
$\aleph_\omega$ & $\beth_\omega = \sup_n \beth_n$ & Primer cardinal límite no numerable \\
$\vdots$ & $\vdots$ & $\vdots$ \\
$\kappa$ inaccesible & $\kappa = \beth_\kappa = \aleph_\kappa$ & Punto fijo de ambas jerarquías; $V_\kappa \models \mathrm{ZFC}$
\end{tabular}
\end{center}

\begin{importante}
La relación $\aleph_\alpha \leq \beth_\alpha$ vale siempre. La igualdad para todo $\alpha$ es la HGC. La Teoría de Forcing (Cohen, 1963) muestra que los valores de $2^{\aleph_\alpha}$ pueden ser casi cualquier función no decreciente con cofinalidad regular, sujeta al Teorema de König.
\end{importante}

% ===========================================================
\section{Problemas y tareas}
% ===========================================================

Los siguientes ejercicios están diseñados para afianzar los conceptos de la clase. Se presentan en orden creciente de dificultad.

\begin{enumerate}[leftmargin=2em, label=\textbf{\arabic*.}]

\item \textbf{(Cardinales como ordinales iniciales.)}
Demuestre directamente, sin usar el Teorema del Bien-Orden, que $\omega$ es un ordinal inicial. \emph{Sugerencia}: muestre que ningún ordinal $n < \omega$ es equipotente a $\omega$.

\item \textbf{(Equipotencias básicas.)}
Pruebe que los siguientes conjuntos son todos de cardinal $\aleph_0$:
\begin{enumerate}[label=(\alph*)]
  \item El conjunto de los enteros $\mathbb{Z}$.
  \item El conjunto de los racionales $\mathbb{Q}$.
  \item El conjunto de polinomios con coeficientes enteros $\mathbb{Z}[x]$.
  \item La unión $\bigcup_{n \in \mathbb{N}} \mathbb{N}^n$ de todos los $n$-uplos de naturales.
\end{enumerate}

\item \textbf{(Conjuntos de cardinal $\beth_1$.)}
Construya biyecciones explícitas (o demuestre la equipotencia usando el Teorema de Cantor-Schröder-Bernstein) entre los siguientes conjuntos:
\begin{enumerate}[label=(\alph*)]
  \item $\mathcal{P}(\mathbb{N})$ y $[0,1]$.
  \item $\mathbb{R}$ y $\mathbb{R}^2$.
  \item El conjunto de sucesiones de naturales $\mathbb{N}^{\mathbb{N}}$ y $\mathbb{R}$.
  \item El conjunto de funciones continuas $f \colon \mathbb{R} \to \mathbb{R}$ y $\mathbb{R}$.
\end{enumerate}

\item \textbf{(Teorema de Cantor: variaciones.)}
\begin{enumerate}[label=(\alph*)]
  \item Pruebe que para cualquier función $f \colon \mathbb{N} \to \{0,1\}^{\mathbb{N}}$ (funciones de $\mathbb{N}$ a las sucesiones binarias), existe una sucesión binaria no en la imagen de $f$. Exhiba explícitamente esa sucesión en función de $f$.
  \item Demuestre que no existe ninguna función sobreyectiva $f \colon A \to 2^A$ para ningún conjunto $A$, donde $2^A$ denota el conjunto de todas las funciones $A \to \{0,1\}$.
  \item Pruebe que $|\mathbb{R}^{\mathbb{R}}| = 2^{\beth_1}$.
\end{enumerate}

\item \textbf{(Aritmética cardinal.)}
Calcule (justificando) los siguientes cardinales:
\begin{enumerate}[label=(\alph*)]
  \item $\aleph_0 + \aleph_1$.
  \item $\aleph_1 \cdot \aleph_2$.
  \item $\aleph_0^{\aleph_0}$.
  \item $\aleph_1^{\aleph_0}$.
  \item $2^{\aleph_1}$ (¿puede determinarse sin hipótesis adicionales?).
\end{enumerate}

\item \textbf{(Jerarquía aleph.)}
\begin{enumerate}[label=(\alph*)]
  \item Demuestre que $\aleph_\omega$ es un cardinal límite.
  \item Demuestre que $\mathrm{cf}(\aleph_\omega) = \aleph_0$ (es decir, $\aleph_\omega$ es un cardinal singular).
  \item Concluya que $\aleph_\omega$ no es débilmente inaccesible.
\end{enumerate}

\item \textbf{(Beth numbers.)}
\begin{enumerate}[label=(\alph*)]
  \item Demuestre que $\beth_\omega$ es un cardinal límite.
  \item Pruebe que $\aleph_\omega \leq \beth_\omega$. ¿Es siempre $\aleph_\omega = \beth_\omega$?
  \item Demuestre que si existe un cardinal fuertemente inaccesible $\kappa$, entonces $\kappa = \beth_\kappa$.
\end{enumerate}

\item \textbf{(Independencia de la HC: ideas intuitivas.)}
\begin{enumerate}[label=(\alph*)]
  \item Enuncia la Hipótesis del Continuo y explica en tus propias palabras qué significaría que fuera verdadera o falsa.
  \item Describe brevemente (sin entrar en detalles técnicos) la idea de Gödel para probar que CH no puede refutarse en ZFC \cite{godel1940}.
  \item Describe brevemente la idea de Cohen (forcing) para probar que CH no puede probarse en ZFC \cite{cohen1966}.
\end{enumerate}

\item \textbf{(Cardinales inaccesibles y ZFC.)}
\begin{enumerate}[label=(\alph*)]
  \item Sea $\kappa$ un cardinal fuertemente inaccesible. Demuestre que $V_\kappa$ satisface todos los axiomas de ZFC. (\emph{Sugerencia}: verifique axioma por axioma; use que $\kappa$ es regular y que $2^\lambda < \kappa$ para $\lambda < \kappa$.)
  \item Explique por qué la existencia de $\kappa$ no puede demostrarse en ZFC (asumiendo que ZFC es consistente).
\end{enumerate}

\item \textbf{(Problema de investigación.)}
La \emph{Hipótesis del Continuo en singular} es la afirmación $2^{\aleph_\omega} = \aleph_{\omega+1}$. ¿Es esta afirmación demostrable, refutable, o independiente de ZFC? Investiga el \emph{Teorema de Shelah} sobre la aritmética cardinal singular y explica qué restricciones impone al valor de $2^{\aleph_\omega}$.

\end{enumerate}

% ===========================================================
\section{Referencias de la Clase}
% ===========================================================

Los resultados de esta clase se encuentran desarrollados con todo detalle en las siguientes obras. Para una presentación axiomática completa de la teoría cardinal y la jerarquía aleph, los textos fundamentales son \cite{jech2003} y \cite{kunen2011}. Una introducción accesible y rigurosa está en \cite{enderton1977} y \cite{hrbacek1999}. Para la teoría de grandes cardinales y los cardinales inaccesibles, la referencia estándar es \cite{kanamori2003}. El artículo original de Cantor sobre el argumento diagonal es \cite{cantor1891}. Para la independencia de la HC, véanse el trabajo fundacional de Gödel \cite{godel1940} y el de Cohen \cite{cohen1966}. Para una exposición histórica de la teoría cardinal, \cite{sierpinski1958} y \cite{halmos1960} son lecturas muy recomendadas.

\nocite{cantor1897, vonneumann1923, levy1979, devlin1993, suppes1972}

% ===========================================================
\chapter{Clase 8: Aritmética Cardinal Transfinita}
% ===========================================================

La \textbf{aritmética de cardinales infinitos} es, en muchos sentidos, más simple que la aritmética de ordinales: la suma y el producto de dos cardinales infinitos cualesquiera colapsa al mayor de ellos. Sin embargo, la \textbf{potenciación cardinal} $\kappa^\lambda$ exhibe una riqueza genuina que escapa a este principio de absorción, y cuyo comportamiento general aún no está completamente determinado en ZFC. El teorema de König —una desigualdad estricta para sumas y productos de familias de cardinales— es la herramienta central que permite extraer consecuencias no triviales de la potenciación. Su corolario más sorprendente afirma que $2^{\aleph_0}$ no puede tener cofinalidad numerable: el continuo no puede ser $\aleph_\omega$, ni $\aleph_{\omega \cdot 2}$, ni ningún cardinal de cofinalidad $\omega$.

La clasificación de los cardinales infinitos en \textbf{regulares} y \textbf{singulares} según su cofinalidad es la noción transversal de esta clase. Los cardinales sucesores $\aleph_{\alpha+1}$ son siempre regulares; los cardinales límite como $\aleph_\omega$, $\aleph_{\omega^2}$, $\aleph_{\varepsilon_0}$ son singulares. Esta distinción tiene consecuencias profundas tanto en la aritmética cardinal elemental como en la teoría de grandes cardinales.

La Clase~6 introdujo la cofinalidad para ordinales y la Clase~7 presentó la jerarquía aleph y las primeras propiedades de la aritmética cardinal. Esta clase profundiza todos esos temas y los organiza en un marco coherente.

La exposición sigue los tratamientos clásicos de \cite{jech2003, kunen2011, enderton1977, hrbacek1999}, con especial atención a los detalles técnicos de las demostraciones y a la intuición geométrica y combinatoria detrás de cada resultado.

% ===========================================================
\section{Suma y producto de cardinales infinitos}
% ===========================================================

\subsection{Definiciones formales}

Recordamos que los números cardinales son ordinales iniciales, y que la aritmética cardinal difiere conceptualmente de la ordinal: la suma y el producto cardinales se definen mediante biyecciones entre conjuntos, no mediante concatenación ni iteración de orden.

\begin{definicion}{Suma cardinal}{sumacardinal}
Sean $\kappa$ y $\lambda$ cardinales. La \textbf{suma cardinal} $\kappa + \lambda$ es el cardinal del conjunto
\[
  (\kappa \times \{0\}) \cup (\lambda \times \{1\}),
\]
es decir, la unión disjunta de dos conjuntos de cardinalidades $\kappa$ y $\lambda$. Formalmente:
\[
  \kappa + \lambda := \bigl|(\kappa \times \{0\}) \cup (\lambda \times \{1\})\bigr|.
\]
\end{definicion}

\begin{definicion}{Producto cardinal}{productocardinal}
Sean $\kappa$ y $\lambda$ cardinales. El \textbf{producto cardinal} $\kappa \cdot \lambda$ es el cardinal del producto cartesiano:
\[
  \kappa \cdot \lambda := |\kappa \times \lambda|.
\]
\end{definicion}

\begin{observacion}{Diferencia con la aritmética ordinal}{diffcardord}
La suma ordinal $\omega + \omega$ y la suma cardinal $\aleph_0 + \aleph_0$ son objetos distintos. Como ordinal, $\omega + \omega \neq \omega$; como cardinal, $\aleph_0 + \aleph_0 = \aleph_0$. La razón es que la suma ordinal recuerda el \emph{orden} de la concatenación, mientras que la suma cardinal solo mide la \emph{cardinalidad} de la unión disjunta: ambas copias de $\mathbb{N}$ pueden entrelazarse en una sola copia de $\mathbb{N}$.
\end{observacion}

\subsection{La ley de absorción: todo colapsа al máximo}

El resultado central de la aritmética cardinal para cardinales infinitos es que ni la suma ni el producto producen cardinales nuevos: siempre se obtiene el máximo de los dos operandos.

\begin{teorema}{Ley de absorción cardinal}{leyabsorcion}
Sea $\kappa$ un cardinal infinito y $\lambda$ un cardinal con $1 \leq \lambda \leq \kappa$. Entonces:
\[
  \kappa + \lambda = \kappa \cdot \lambda = \kappa.
\]
En particular:
\[
  \kappa + \kappa = \kappa, \qquad \kappa \cdot \kappa = \kappa.
\]
\end{teorema}

La demostración del caso $\kappa \cdot \kappa = \kappa$ es el corazón del teorema. Presentamos la demostración clásica basada en el \textbf{Teorema de Hessenberg}, que utiliza un buen orden especial del producto $\kappa \times \kappa$.

\begin{proof}
\textbf{Paso 1: $\kappa \cdot \kappa = \kappa$ para todo cardinal infinito $\kappa$.}

Procedemos por inducción transfinita en $\kappa$. Tomamos como hipótesis de inducción que $\mu \cdot \mu = \mu$ para todo cardinal infinito $\mu < \kappa$.

\textbf{Base de la inducción.} Para $\kappa = \aleph_0$: el conjunto $\mathbb{N} \times \mathbb{N}$ puede numerarse mediante la diagonal de Cantor. Definimos el orden en $\omega \times \omega$ dado por
\[
  (\alpha_1, \beta_1) \prec (\alpha_2, \beta_2)
  \;\iff\;
  \max(\alpha_1,\beta_1) < \max(\alpha_2,\beta_2),
\]
y para igual máximo, por el orden lexicográfico. Bajo este orden, $\omega \times \omega$ tiene tipo ordinal $\omega$, lo que da una biyección $\omega \times \omega \to \omega$, es decir, $\aleph_0 \cdot \aleph_0 = \aleph_0$.

\textbf{Paso de inducción para $\kappa$ sucesor.} Si $\kappa = \mu^+$ para algún cardinal $\mu$, consideramos el orden de Hessenberg en $\kappa \times \kappa$:
\[
  (\alpha_1, \beta_1) \prec (\alpha_2, \beta_2)
  \;\iff\;
  \max(\alpha_1, \beta_1) < \max(\alpha_2, \beta_2),
\]
y para igual máximo, orden lexicográfico. Para cada $\gamma < \kappa$, el segmento inicial
\[
  \{(\alpha, \beta) \in \kappa \times \kappa : \max(\alpha,\beta) \leq \gamma\}
  = (\gamma+1) \times (\gamma+1)
\]
tiene cardinalidad $|\gamma+1|^2$. Si $\gamma < \kappa$, entonces $|\gamma| < \kappa$ y por hipótesis de inducción (si $|\gamma|$ es infinito) $|\gamma|^2 = |\gamma| < \kappa$; si $|\gamma|$ es finito, $|\gamma|^2$ es finito, luego $< \kappa$. En todo caso, cada segmento inicial propio tiene cardinalidad $< \kappa$. El tipo ordinal de $(\kappa \times \kappa, \prec)$ es entonces un ordinal $\leq \kappa$ cuyos segmentos iniciales propios tienen cardinalidad $< \kappa$, y como el tipo ordinal de un buen orden de cardinalidad $\kappa \cdot \kappa$ es al menos $\kappa$, se obtiene que el tipo ordinal es exactamente $\kappa$, lo que da $\kappa \cdot \kappa = \kappa$.

\textbf{Paso de inducción para $\kappa$ límite.} Sea $\kappa$ un cardinal límite infinito. Para cualquier $\mu < \kappa$ tenemos $\mu \cdot \mu = \mu < \kappa$ (por hipótesis de inducción o por ser finito). Entonces
\[
  \kappa \cdot \kappa = \left|\bigcup_{\mu < \kappa} \mu\right|^2 \leq \kappa \cdot \kappa,
\]
y de manera más precisa: toda función $\kappa \times \kappa \to \kappa$ puede construirse a partir de las biyecciones $\mu \times \mu \to \mu$ para $\mu < \kappa$ encadenadas coherentemente, dando $\kappa \cdot \kappa = \kappa$.

\textbf{Paso 2: $\kappa + \lambda = \kappa$ para $1 \leq \lambda \leq \kappa$.}

Se tiene
\[
  \kappa \leq \kappa + \lambda \leq \kappa + \kappa = 2 \cdot \kappa \leq \kappa \cdot \kappa = \kappa.
\]
Por tanto $\kappa + \lambda = \kappa$.

\textbf{Paso 3: $\kappa \cdot \lambda = \kappa$ para $1 \leq \lambda \leq \kappa$.}

Se tiene $\kappa \leq \kappa \cdot \lambda \leq \kappa \cdot \kappa = \kappa$, luego $\kappa \cdot \lambda = \kappa$ \cite{jech2003, enderton1977}.
\end{proof}

\begin{corolario}{Suma y producto de dos alephs}{sumaprodaleph}
Para cualquier par de cardinales infinitos $\kappa$ y $\lambda$:
\[
  \kappa + \lambda = \kappa \cdot \lambda = \max(\kappa, \lambda).
\]
\end{corolario}

\begin{proof}
Sin pérdida de generalidad supongamos $\kappa \leq \lambda$. Entonces $1 \leq \kappa \leq \lambda$, y por el Teorema~\ref{teo:leyabsorcion} aplicado con el papel de $\kappa$ jugado por $\lambda$ y el papel de $\lambda$ jugado por $\kappa$: $\lambda + \kappa = \lambda \cdot \kappa = \lambda = \max(\kappa, \lambda)$.
\end{proof}

\begin{ejemplo}{Cálculos de suma y producto}{ejsumaprod}
Calculamos algunos cardinales usando la ley de absorción.

\textbf{(a)} $\aleph_0 + \aleph_1 = \max(\aleph_0, \aleph_1) = \aleph_1$.

\textbf{(b)} $\aleph_3 + \aleph_3 + \aleph_3 = \aleph_3$ (tres sumandos; por inducción, cualquier suma finita de copias de $\aleph_3$ sigue siendo $\aleph_3$).

\textbf{(c)} $\aleph_0 \cdot \aleph_0 = \aleph_0$. Consecuencia: $|\mathbb{N} \times \mathbb{N}| = \aleph_0$, que es la base del hecho de que $\mathbb{Q}$ es numerable.

\textbf{(d)} $\aleph_2 \cdot \aleph_5 = \aleph_5$.

\textbf{(e)} $\aleph_\omega \cdot \aleph_\omega = \aleph_\omega$ (el cardinal $\aleph_\omega$ es infinito, luego $\aleph_\omega \cdot \aleph_\omega = \aleph_\omega$).

\textbf{(f)} El conjunto de todos los polinomios con coeficientes enteros tiene cardinal $\aleph_0$: en efecto, cada polinomio de grado $n$ está determinado por $n+1$ coeficientes enteros, luego el conjunto de polinomios de grado $n$ tiene cardinal $|\mathbb{Z}^{n+1}| = \aleph_0^{n+1} = \aleph_0$. La unión numerable de conjuntos de cardinal $\aleph_0$ tiene cardinal $\aleph_0 \cdot \aleph_0 = \aleph_0$.
\end{ejemplo}

\begin{ejemplo}{Unión numerable de conjuntos numerables}{unionnum}
Sea $\{A_n\}_{n \in \mathbb{N}}$ una familia numerable de conjuntos, cada uno de cardinalidad $\leq \aleph_0$. Entonces $\left|\bigcup_{n \in \mathbb{N}} A_n\right| \leq \aleph_0$.

\textit{Demostración.} Para cada $n$, fijamos una inyección $f_n \colon A_n \hookrightarrow \mathbb{N}$. Definimos $g \colon \bigcup_n A_n \to \mathbb{N} \times \mathbb{N}$ poniendo $g(a) = (n_0, f_{n_0}(a))$ donde $n_0 = \min\{n : a \in A_n\}$. Esta $g$ es inyectiva, luego $\left|\bigcup_n A_n\right| \leq |\mathbb{N} \times \mathbb{N}| = \aleph_0$.
\end{ejemplo}

\begin{ejemplo}{Suma infinita de cardinales iguales}{sumainfiguales}
Sea $\kappa$ un cardinal infinito y supongamos que sumamos $\kappa$ copias de $\kappa$:
\[
  \sum_{\alpha < \kappa} \kappa = \kappa \cdot \kappa = \kappa.
\]
Más generalmente, la suma cardinal de $\lambda$ cardinales iguales a $\kappa$ con $1 \leq \lambda \leq \kappa$ es:
\[
  \sum_{\alpha < \lambda} \kappa = \kappa \cdot \lambda = \max(\kappa, \lambda) \cdot 1 = \max(\kappa,\lambda) = \kappa.
\]
\end{ejemplo}

\subsection{Suma infinita de cardinales distintos}

La ley de absorción se extiende naturalmente a sumas infinitas indexadas por ordinales.

\begin{definicion}{Suma cardinal indexada}{sumacardindexada}
Sea $\{\kappa_i\}_{i \in I}$ una familia de cardinales indexada por un conjunto $I$. La \textbf{suma cardinal} $\sum_{i \in I} \kappa_i$ es el cardinal de la unión disjunta:
\[
  \sum_{i \in I} \kappa_i := \left|\bigsqcup_{i \in I} \kappa_i\right| = \left|\bigcup_{i \in I} (\kappa_i \times \{i\})\right|.
\]
\end{definicion}

\begin{proposicion}{Suma cardinal de una familia de cardinales}{sumafamcard}
Sea $\{\kappa_i\}_{i \in I}$ una familia de cardinales y sea $\kappa = \sup_{i \in I} \kappa_i$. Entonces:
\begin{enumerate}[leftmargin=2em, label=(\roman*)]
  \item $\sum_{i \in I} \kappa_i \leq |I| \cdot \kappa$.
  \item Si $|I| \leq \kappa$ y $\kappa_i \leq \kappa$ para todo $i \in I$, entonces $\sum_{i \in I} \kappa_i \leq \kappa$.
  \item Si existe $i_0 \in I$ con $\kappa_{i_0} = \kappa$ y $|I| \leq \kappa$, entonces $\sum_{i \in I} \kappa_i = \kappa$.
\end{enumerate}
\end{proposicion}

\begin{proof}
\textbf{(i)} Se tiene $\bigsqcup_{i \in I} \kappa_i \hookrightarrow I \times \kappa$ (por la inyección $(x, i) \mapsto (i, f_i(x))$ donde $f_i \colon \kappa_i \hookrightarrow \kappa$), luego $\sum_{i \in I} \kappa_i \leq |I \times \kappa| = |I| \cdot \kappa$.

\textbf{(ii)} Si $|I| \leq \kappa$, entonces $|I| \cdot \kappa = \kappa$, y por (i), $\sum_{i \in I} \kappa_i \leq \kappa$.

\textbf{(iii)} La cota inferior $\kappa \leq \sum_{i \in I} \kappa_i$ se sigue de $\kappa = \kappa_{i_0} \leq \sum_{i \in I}\kappa_i$; la cota superior es (ii).
\end{proof}

% ===========================================================
\section{Potenciación cardinal: \texorpdfstring{$\kappa^\lambda$}{κ\^{}λ}}
% ===========================================================

\subsection{Definición y ejemplos iniciales}

La potenciación cardinal es la operación donde la aritmética transfinita comienza a exhibir complejidad genuina.

\begin{definicion}{Potencia cardinal}{potenciacardinal}
Sean $\kappa$ y $\lambda$ cardinales. La \textbf{potencia cardinal} $\kappa^\lambda$ es el cardinal del conjunto de todas las funciones de $\lambda$ en $\kappa$:
\[
  \kappa^\lambda := |{}^\lambda\kappa| = |\{\,f \mid f \colon \lambda \to \kappa\,\}|.
\]
\end{definicion}

\begin{observacion}{Consistencia con la aritmética usual}{consisarith}
Para cardinales finitos $m$ y $n$, el conjunto ${}^n m$ de funciones de $n$ a $m$ tiene exactamente $m^n$ elementos, consistente con la definición usual. Para $\kappa^\lambda$ con $\kappa = 2$, obtenemos $2^\lambda = |\mathcal{P}(\lambda)|$ (identificando subconjuntos con funciones características), lo que conecta con el Teorema de Cantor: $2^\lambda > \lambda$ para todo $\lambda$.
\end{observacion}

\begin{proposicion}{Propiedades básicas de la potenciación cardinal}{propbasicpot}
Sean $\kappa$, $\lambda$, $\mu$ cardinales. Entonces:
\begin{enumerate}[leftmargin=2em, label=(\roman*)]
  \item $\kappa^{\lambda + \mu} = \kappa^\lambda \cdot \kappa^\mu$.
  \item $(\kappa \cdot \lambda)^\mu = \kappa^\mu \cdot \lambda^\mu$.
  \item $(\kappa^\lambda)^\mu = \kappa^{\lambda \cdot \mu}$.
  \item Si $\lambda \leq \mu$, entonces $\kappa^\lambda \leq \kappa^\mu$.
  \item Si $\kappa \leq \lambda$, entonces $\kappa^\mu \leq \lambda^\mu$.
  \item $\kappa^1 = \kappa$ y $\kappa^0 = 1$ y $1^\lambda = 1$.
\end{enumerate}
\end{proposicion}

\begin{proof}
\textbf{(i)} ${}^{\lambda+\mu}\kappa \cong {}^\lambda\kappa \times {}^\mu\kappa$ mediante la biyección $f \mapsto (f\!\restriction_\lambda, \, f\!\restriction_{[\lambda,\lambda+\mu)})$.

\textbf{(ii)} ${}^\mu(\kappa \times \lambda) \cong {}^\mu\kappa \times {}^\mu\lambda$ mediante $f \mapsto (\pi_1 \circ f, \pi_2 \circ f)$.

\textbf{(iii)} ${}^\mu({}^\lambda\kappa) \cong {}^{\lambda \cdot \mu}\kappa$ mediante la currificación: $(f \colon \mu \to {}^\lambda\kappa) \mapsto \bigl((\alpha, \beta) \mapsto f(\beta)(\alpha)\bigr)$.

\textbf{(iv)-(v)} Se obtienen exhibiendo inyecciones entre los conjuntos de funciones correspondientes.

\textbf{(vi)} Inmediato por las definiciones \cite{enderton1977, hrbacek1999}.
\end{proof}

\subsection{Valores concretos de potencias cardinales}

\begin{ejemplo}{Potencias de $\aleph_0$}{potenaleph0}
Calculamos $\aleph_0^\lambda$ para varios valores de $\lambda$.

\textbf{(a) $\aleph_0^n = \aleph_0$ para todo $n$ finito.} En efecto, $\aleph_0^n = \aleph_0 \cdot \aleph_0 \cdots \aleph_0 = \aleph_0$ por la ley de absorción.

\textbf{(b) $\aleph_0^{\aleph_0} = 2^{\aleph_0}$.} Probamos las dos desigualdades.
\begin{itemize}[leftmargin=2em]
  \item $2^{\aleph_0} \leq \aleph_0^{\aleph_0}$: la inyección ${}^\omega 2 \hookrightarrow {}^\omega \omega$ es la inclusión (toda función $\omega \to \{0,1\}$ es también función $\omega \to \omega$).
  \item $\aleph_0^{\aleph_0} \leq 2^{\aleph_0}$: cada función $f \colon \omega \to \omega$ puede codificarse mediante una función $\omega \to \{0,1\}$ (por ejemplo, codificando $f(n) \in \omega$ en binario, o usando la biyección $\omega \times \omega \leftrightarrow \omega$ para construir $F \colon \omega \to 2$). Formalmente, $|\omega^\omega| = |({}^\omega\omega)| \leq |{}^\omega 2^\omega| = |({}^\omega\omega)| \leq 2^{\aleph_0 \cdot \aleph_0} = 2^{\aleph_0}$.
\end{itemize}
Por el Teorema de Cantor-Schröder-Bernstein, $\aleph_0^{\aleph_0} = 2^{\aleph_0}$.

\textbf{(c)} En particular, $|\mathbb{R}| = 2^{\aleph_0} = \aleph_0^{\aleph_0}$: la cardinalidad del continuo coincide con la del espacio de sucesiones de naturales.
\end{ejemplo}

\begin{ejemplo}{Potencias de $\aleph_1$}{potenaleph1}
\textbf{(a) $\aleph_1^{\aleph_0}$:} Sabemos que $2^{\aleph_0} \leq \aleph_1^{\aleph_0}$. También, $\aleph_1^{\aleph_0} \leq (2^{\aleph_0})^{\aleph_0} = 2^{\aleph_0 \cdot \aleph_0} = 2^{\aleph_0}$. En consecuencia, $\aleph_1^{\aleph_0} = 2^{\aleph_0}$: el valor depende de si $\aleph_1 \leq 2^{\aleph_0}$ o $\aleph_1 > 2^{\aleph_0}$, pero la expresión $\aleph_1^{\aleph_0}$ siempre coincide con $2^{\aleph_0}$.

\textbf{(b)} Si se asume la Hipótesis del Continuo ($2^{\aleph_0} = \aleph_1$), entonces $\aleph_1^{\aleph_0} = 2^{\aleph_0} = \aleph_1$: el exponente $\aleph_0$ no cambia $\aleph_1$.

\textbf{(c) $\aleph_1^{\aleph_1}$:} Se tiene $2^{\aleph_1} \leq \aleph_1^{\aleph_1} \leq (2^{\aleph_1})^{\aleph_1} = 2^{\aleph_1 \cdot \aleph_1} = 2^{\aleph_1}$, luego $\aleph_1^{\aleph_1} = 2^{\aleph_1}$.
\end{ejemplo}

\begin{proposicion}{Fórmula general para $\kappa^\lambda$}{formulakaplam}
Sean $\kappa \geq 2$ y $\lambda \geq 1$ cardinales infinitos. Entonces:
\begin{enumerate}[leftmargin=2em, label=(\roman*)]
  \item Si $2^\lambda \geq \kappa$, entonces $\kappa^\lambda = 2^\lambda$.
  \item Si $2^\lambda < \kappa$, entonces $\kappa^\lambda = \kappa^{\mathrm{cf}(\kappa)}$ (solo depende de $\kappa$ y de $\mathrm{cf}(\kappa)$, no de $\lambda$).
\end{enumerate}
\end{proposicion}

\begin{proof}
\textbf{(i)} $2^\lambda \leq \kappa^\lambda \leq (2^\lambda)^\lambda = 2^{\lambda \cdot \lambda} = 2^\lambda$, luego $\kappa^\lambda = 2^\lambda$.

\textbf{(ii)} La demostración completa usa el \emph{Teorema de la Aritmética Cardinal de Easton} y resultados más avanzados sobre la aritmética de cardinales regulares y singulares, y queda fuera del alcance de esta clase \cite{jech2003, shelah1994}.
\end{proof}

\begin{ejemplo}{La relación $\kappa^\lambda$ y los beth numbers}{bethkaplam}
Recordemos que los \textbf{beth numbers} se definen por:
\[
  \beth_0 = \aleph_0, \qquad \beth_{\alpha+1} = 2^{\beth_\alpha}, \qquad \beth_\lambda = \sup_{\alpha < \lambda} \beth_\alpha.
\]
La Hipótesis del Continuo Generalizada (GCH) afirma $\aleph_\alpha = \beth_\alpha$ para todo $\alpha$. Bajo GCH:
\begin{enumerate}[leftmargin=2em, label=(\alph*)]
  \item $\aleph_\alpha^{\aleph_\beta} = \aleph_\alpha$ si $\beta < \alpha$ y $\mathrm{cf}(\aleph_\alpha) > \aleph_\beta$ (la potencia no supera $\kappa$).
  \item $\aleph_\alpha^{\aleph_\beta} = \aleph_{\alpha+1}$ en los demás casos en que $\aleph_\beta \leq \aleph_\alpha$.
\end{enumerate}
Sin GCH, el valor exacto de $2^{\aleph_0}$ —y en general de $\kappa^\lambda$— es independiente de ZFC.
\end{ejemplo}

% ===========================================================
\section{El Teorema de König}
% ===========================================================

\subsection{Formulación e intuición}

El \textbf{Teorema de König} (1905, \cite{konig1905}) es la desigualdad cardinal fundamental. Afirma que una suma de cardinales es estrictamente menor que el producto correspondiente de cardinales mayores. Su enunciado general es la versión transfinita de la afirmación intuitiva: «el producto de cosas más grandes supera la suma de cosas más pequeñas».

\begin{teorema}{Teorema de König}{teokoenig}
Sea $I$ un conjunto y sean $\{\kappa_i\}_{i \in I}$ y $\{\lambda_i\}_{i \in I}$ dos familias de cardinales tales que $\kappa_i < \lambda_i$ para todo $i \in I$. Entonces:
\[
  \sum_{i \in I} \kappa_i < \prod_{i \in I} \lambda_i.
\]
\end{teorema}

\begin{proof}
Necesitamos demostrar dos cosas:
\begin{enumerate}[leftmargin=2em, label=(\arabic*)]
  \item $\sum_{i \in I} \kappa_i \leq \prod_{i \in I} \lambda_i$, es decir, existe una inyección $\bigsqcup_{i \in I} \kappa_i \hookrightarrow \prod_{i \in I} \lambda_i$.
  \item Ninguna función $f \colon \bigsqcup_{i \in I} \kappa_i \to \prod_{i \in I} \lambda_i$ puede ser sobreyectiva.
\end{enumerate}

\textbf{Parte 1 (inyección).} Para cada $i \in I$ fijamos una inyección $\phi_i \colon \kappa_i \hookrightarrow \lambda_i$ (que existe porque $\kappa_i < \lambda_i$ implica $\kappa_i \leq \lambda_i$, es decir, hay una inyección de $\kappa_i$ en $\lambda_i$). Fijamos además un elemento distinguido $d_i \in \lambda_i$ para cada $i$. Definimos $\Phi \colon \bigsqcup_{i \in I} \kappa_i \to \prod_{i \in I} \lambda_i$ por:
\[
  \Phi(\alpha, j)(i) = \begin{cases} \phi_j(\alpha) & \text{si } i = j, \\ d_i & \text{si } i \neq j. \end{cases}
\]
Si $\Phi(\alpha, j) = \Phi(\beta, k)$, evaluando en $i = j$ obtenemos $\phi_j(\alpha) = \Phi(\alpha, j)(j) = \Phi(\beta, k)(j)$. Si $j \neq k$, entonces $\Phi(\beta, k)(j) = d_j$, por lo que $\phi_j(\alpha) = d_j$. Y evaluando en $i = k$, si $j \neq k$ se tendría $\phi_k(\beta) = \Phi(\beta,k)(k) = \Phi(\alpha,j)(k) = d_k$. Esto puede ajustarse para que ambas condiciones sean incompatibles si $j \neq k$ (eligiendo $d_j \notin \mathrm{Im}(\phi_j)$ cuando $\kappa_j < \lambda_j$ lo permite). Una construcción ligeramente diferente que garantiza la inyectividad usa una codificación del índice $j$ en la función, como por ejemplo:
\[
  \Phi'(\alpha, j)(i) = \begin{cases} \phi_j(\alpha) & \text{si } i = j, \\ 0_i & \text{si } i \neq j, \end{cases}
\]
donde $0_i$ es el elemento mínimo de $\lambda_i$ (como ordinal). Entonces $\Phi'(\alpha, j) = \Phi'(\beta, k)$ implica: evaluando en $j$ si $j \neq k$, $\phi_j(\alpha) = 0_j$ y evaluando en $k$, $\phi_k(\beta) = 0_k$; pero evaluando en $j$ del lado de $(\beta, k)$ obtenemos $0_j$ si $j \neq k$. La inyectividad se verifica con más detalle en \cite{jech2003}; la idea clave es que el «índice» $j$ queda registrado en el soporte de $\Phi'$.

Para un argumento limpio: fijamos para cada $j$ una función $e_j \colon \kappa_j \to \prod_{i} \lambda_i$ definida por $e_j(\alpha)(i) = \phi_j(\alpha)$ si $i = j$ y $e_j(\alpha)(i) = 0_i$ si $i \neq j$. Las imágenes de funciones distintas $e_j$ tienen «soporte» diferente (coordenada $j$ no nula), y dentro de la imagen de $e_j$ la inyectividad se hereda de $\phi_j$. Por tanto la función que manda $(\alpha, j) \mapsto e_j(\alpha)$ es inyectiva.

\textbf{Parte 2 (no sobreyección).} Sea $f \colon \bigsqcup_{i \in I} \kappa_i \to \prod_{i \in I} \lambda_i$ una función cualquiera. Debemos exhibir un elemento de $\prod_{i \in I} \lambda_i$ no en la imagen de $f$.

Para cada $i \in I$, consideramos la «sección $i$» de $f$ restringida a $\kappa_i$:
\[
  A_i := \{f(\alpha, i)(i) : \alpha \in \kappa_i\} \subseteq \lambda_i.
\]
El conjunto $A_i$ tiene cardinal $|A_i| \leq |\kappa_i| = \kappa_i < \lambda_i$. Por tanto $A_i \subsetneq \lambda_i$, y existe
\[
  \beta_i \in \lambda_i \setminus A_i.
\]
Definimos el elemento diagonal $g \in \prod_{i \in I} \lambda_i$ por $g(i) = \beta_i$ para cada $i \in I$.

Afirmamos que $g \notin \mathrm{Im}(f)$. En efecto, sea $(\alpha, j) \in \bigsqcup_{i \in I} \kappa_i$ cualquiera. Por definición, $f(\alpha, j)(j) \in A_j$, mientras que $g(j) = \beta_j \notin A_j$. Luego $f(\alpha, j) \neq g$ (difieren en la coordenada $j$). Como $(\alpha, j)$ fue arbitrario, $g \notin \mathrm{Im}(f)$.

Combinando las dos partes: $\bigsqcup_{i \in I} \kappa_i$ se inyecta en $\prod_{i \in I} \lambda_i$ pero no hay sobreyección, luego $\sum_i \kappa_i < \prod_i \lambda_i$ \cite{konig1905, jech2003, kunen2011}.
\end{proof}

\begin{observacion}{El argumento diagonal en König}{diagkoenig}
La segunda parte de la demostración es una versión del \emph{argumento diagonal de Cantor}: para cada $i$, elegimos $\beta_i$ fuera de la imagen de $f$ en la coordenada $i$. La función $g = (\beta_i)_{i \in I}$ «escapa» a $f$ porque difiere de cada $f(\alpha, i)$ en la coordenada $i$.
\end{observacion}

\subsection{Consecuencias fundamentales del Teorema de König}

El Teorema de König tiene tres corolarios inmediatos de gran importancia.

\begin{corolario}{König para dos cardinales}{koenigdos}
Sea $\kappa < \lambda$ y $\mu < \nu$ cardinales. Entonces $\kappa + \mu < \lambda \cdot \nu$.
\end{corolario}

\begin{proof}
Aplicar el Teorema de König con $I = \{0, 1\}$, $\kappa_0 = \kappa$, $\kappa_1 = \mu$, $\lambda_0 = \lambda$, $\lambda_1 = \nu$.
\end{proof}

\begin{corolario}{$\kappa < \kappa^{\mathrm{cf}(\kappa)}$ para cardinales infinitos}{kappacofkappa}
Para todo cardinal infinito $\kappa$ se tiene:
\[
  \kappa < \kappa^{\mathrm{cf}(\kappa)}.
\]
\end{corolario}

\begin{proof}
Sea $\delta = \mathrm{cf}(\kappa)$ y sea $f \colon \delta \to \kappa$ una función cofinal, es decir, $\sup_{\alpha < \delta} f(\alpha) = \kappa$. Para cada $\alpha < \delta$, la imagen $f(\alpha) < \kappa$. Considere la familia:
\[
  \kappa_\alpha = f(\alpha) \quad \text{y} \quad \lambda_\alpha = \kappa, \quad \text{para } \alpha < \delta.
\]
Entonces $\kappa_\alpha < \lambda_\alpha = \kappa$ y $\sum_{\alpha < \delta} \kappa_\alpha = \sup_{\alpha < \delta} f(\alpha) \cdot \delta = \kappa$ (pues $f$ es cofinal y $\delta \leq \kappa$). Por el Teorema de König:
\[
  \kappa = \sum_{\alpha < \delta} \kappa_\alpha < \prod_{\alpha < \delta} \lambda_\alpha = \kappa^\delta = \kappa^{\mathrm{cf}(\kappa)}.
\]
\end{proof}

\begin{corolario}{La cofinalidad de $2^\kappa$}{cofinalidadpot}
Para todo cardinal infinito $\kappa$, la cofinalidad de $2^\kappa$ satisface:
\[
  \mathrm{cf}(2^\kappa) > \kappa.
\]
En particular, $2^{\aleph_0}$ tiene cofinalidad $> \aleph_0$, es decir, $2^{\aleph_0} \neq \aleph_\omega$.
\end{corolario}

\begin{proof}
Supongamos por reducción al absurdo que $\mathrm{cf}(2^\kappa) \leq \kappa$. Entonces existe $\delta = \mathrm{cf}(2^\kappa) \leq \kappa$ y una sucesión cofinal $\{\mu_\alpha\}_{\alpha < \delta}$ de cardinales $\mu_\alpha < 2^\kappa$ con $\sup_{\alpha < \delta} \mu_\alpha = 2^\kappa$. Por el Corolario anterior aplicado a $\lambda = 2^\kappa$:
\[
  2^\kappa < (2^\kappa)^{\mathrm{cf}(2^\kappa)} \leq (2^\kappa)^\kappa = 2^{\kappa \cdot \kappa} = 2^\kappa.
\]
Contradicción. Luego $\mathrm{cf}(2^\kappa) > \kappa$.
\end{proof}

\begin{importante}
Este corolario es notable porque impone una restricción \emph{demostrable en ZFC} sobre el valor del continuo $2^{\aleph_0}$. Por ejemplo:
\begin{itemize}[leftmargin=2em]
  \item $2^{\aleph_0} \neq \aleph_\omega$ (pues $\mathrm{cf}(\aleph_\omega) = \aleph_0 \not> \aleph_0$).
  \item $2^{\aleph_0} \neq \aleph_{\omega^2}$ (pues $\mathrm{cf}(\aleph_{\omega^2}) = \aleph_0$).
  \item Sí es posible (consistente con ZFC) que $2^{\aleph_0} = \aleph_1$, $\aleph_2$, $\aleph_{\omega_1}$, etc., pues estos tienen cofinalidad $> \aleph_0$.
\end{itemize}
\end{importante}

\begin{ejemplo}{Aplicación del Teorema de König a $\aleph_\omega$}{koenigaomega}
Demostramos que $\aleph_\omega < \aleph_\omega^{\aleph_0}$ usando el Corolario~\ref{cor:kappacofkappa}.

Como $\mathrm{cf}(\aleph_\omega) = \aleph_0$ (lo demostraremos en la Sección~\ref{sec:cardinales-singulares}), el corolario afirma:
\[
  \aleph_\omega < \aleph_\omega^{\aleph_0}.
\]
Más cuantitativamente, el Teorema de König con $I = \omega$, $\kappa_n = \aleph_n$ y $\lambda_n = \aleph_\omega$ da:
\[
  \aleph_\omega = \sum_{n < \omega} \aleph_n < \prod_{n < \omega} \aleph_\omega = \aleph_\omega^{\aleph_0}.
\]
No puede determinarse en ZFC el valor exacto de $\aleph_\omega^{\aleph_0}$ más allá de estas cotas. El resultado de Shelah sobre la aritmética cardinal singular \cite{shelah1994} muestra que $\aleph_\omega^{\aleph_0} < \aleph_{\omega_4}$ bajo la hipótesis de que $\aleph_n < \aleph_\omega$ para todo $n < \omega$ (lo cual es evidente), pero este resultado es profundo e involucra la teoría de \emph{pcf} (possible cofinalities).
\end{ejemplo}

\begin{ejemplo}{König y la hipótesis del continuo generalizada}{koeniggch}
Bajo la Hipótesis del Continuo Generalizada (GCH), el Teorema de König toma la siguiente forma concreta. Si $\beta < \alpha$:
\[
  \aleph_\alpha^{\aleph_\beta} = \aleph_{\alpha+1}
  \quad \text{si } \mathrm{cf}(\aleph_\alpha) \leq \aleph_\beta,
  \qquad
  \aleph_\alpha^{\aleph_\beta} = \aleph_\alpha
  \quad \text{si } \mathrm{cf}(\aleph_\alpha) > \aleph_\beta.
\]
Esto sigue de la fórmula de Easton/König: si $\mathrm{cf}(\aleph_\alpha) \leq \aleph_\beta$, el Teorema de König garantiza $\aleph_\alpha < \aleph_\alpha^{\aleph_\beta}$, y bajo GCH este salto es exactamente $\aleph_{\alpha+1}$.
\end{ejemplo}

% ===========================================================
\section{Cofinalidad y cardinales regulares}
\label{sec:cofinalidad-regular}
% ===========================================================

\subsection{Cofinalidad cardinal: repaso y profundización}

La cofinalidad fue introducida para ordinales en la Clase~6. Aquí la estudiamos con énfasis en los cardinales infinitos, donde adquiere un significado combinatorio y aritmético preciso.

\begin{definicion}{Cofinalidad de un cardinal límite}{cofinalidadcardlim}
Sea $\kappa$ un cardinal infinito límite. La \textbf{cofinalidad} de $\kappa$, denotada $\mathrm{cf}(\kappa)$, es el mínimo ordinal $\delta$ tal que existe una función $f \colon \delta \to \kappa$ con $\sup_{\alpha < \delta} f(\alpha) = \kappa$:
\[
  \mathrm{cf}(\kappa) := \min\bigl\{\delta \in \mathbf{ON} : \exists\, f \colon \delta \to \kappa,\, \sup f = \kappa\bigr\}.
\]
Equivalentemente, $\mathrm{cf}(\kappa)$ es el menor cardinal $\delta$ tal que $\kappa$ es la unión de $\delta$ conjuntos de cardinal $< \kappa$.
\end{definicion}

\begin{proposicion}{Propiedades básicas de la cofinalidad cardinal}{propcofcard}
Sea $\kappa$ un cardinal infinito. Entonces:
\begin{enumerate}[leftmargin=2em, label=(\roman*)]
  \item $\mathrm{cf}(\kappa)$ es un cardinal infinito.
  \item $\mathrm{cf}(\kappa) \leq \kappa$.
  \item $\mathrm{cf}(\mathrm{cf}(\kappa)) = \mathrm{cf}(\kappa)$: la cofinalidad es idempotente.
  \item $\mathrm{cf}(\kappa^+) = \kappa^+$ para todo cardinal $\kappa$ (los cardinales sucesores son siempre regulares).
\end{enumerate}
\end{proposicion}

\begin{proof}
\textbf{(i)} $\mathrm{cf}(\kappa)$ es el tipo ordinal de una sucesión cofinal, que es un ordinal. Que es infinito se sigue de que ninguna sucesión finita puede ser cofinal en $\kappa$ (una sucesión finita tiene un máximo, y $\kappa$ no tiene máximo porque es un cardinal límite).

\textbf{(ii)} La función identidad $\mathrm{id} \colon \kappa \to \kappa$ es cofinal, luego $\mathrm{cf}(\kappa) \leq \kappa$.

\textbf{(iii)} Sea $\delta = \mathrm{cf}(\kappa)$ y $\mu = \mathrm{cf}(\delta)$. Si $g \colon \mu \to \delta$ es cofinal en $\delta$ y $f \colon \delta \to \kappa$ es cofinal en $\kappa$, entonces $f \circ g \colon \mu \to \kappa$ es cofinal en $\kappa$ (pues para cada $\xi < \kappa$ existe $\alpha < \delta$ con $f(\alpha) \geq \xi$, y existe $\beta < \mu$ con $g(\beta) \geq \alpha$, luego $f(g(\beta)) \geq f(\alpha) \geq \xi$). Así $\mathrm{cf}(\kappa) \leq \mu = \mathrm{cf}(\delta)$. Pero $\delta = \mathrm{cf}(\kappa) \leq \mathrm{cf}(\kappa) = \delta$, y la minimalidad de $\mathrm{cf}(\kappa)$ implica que $\mu = \delta$, es decir, $\mathrm{cf}(\delta) = \delta$.

\textbf{(iv)} Sea $\kappa^+$ el sucesor cardinal de $\kappa$. Supongamos que $\mathrm{cf}(\kappa^+) = \delta < \kappa^+$, con $f \colon \delta \to \kappa^+$ cofinal. Para cada $\alpha < \delta$, $f(\alpha) < \kappa^+$, luego $|f(\alpha)| \leq \kappa$. Por tanto:
\[
  \kappa^+ = \bigcup_{\alpha < \delta} f(\alpha)
\]
y cada $f(\alpha)$ tiene cardinal $\leq \kappa$. Entonces $|\kappa^+| \leq |\delta| \cdot \kappa = \delta \cdot \kappa$. Como $\delta \leq \kappa$ (pues $\delta < \kappa^+$ implica $\delta \leq \kappa$), tenemos $\delta \cdot \kappa = \kappa$, lo que daría $\kappa^+ \leq \kappa$, contradicción. Luego $\mathrm{cf}(\kappa^+) = \kappa^+$.
\end{proof}

\subsection{Cardinales regulares: definición y caracterización}

\begin{definicion}{Cardinal regular y singular}{cardreg}
Sea $\kappa$ un cardinal infinito.
\begin{itemize}[leftmargin=2em]
  \item $\kappa$ es \textbf{regular} si $\mathrm{cf}(\kappa) = \kappa$.
  \item $\kappa$ es \textbf{singular} si $\mathrm{cf}(\kappa) < \kappa$.
\end{itemize}
\end{definicion}

\begin{teorema}{Cardinales regulares: caracterización combinatoria}{caracregular}
Sea $\kappa$ un cardinal infinito. Son equivalentes:
\begin{enumerate}[leftmargin=2em, label=(\roman*)]
  \item $\kappa$ es regular, es decir, $\mathrm{cf}(\kappa) = \kappa$.
  \item $\kappa$ no es unión de menos de $\kappa$ conjuntos, cada uno de cardinalidad $< \kappa$. Formalmente: si $\kappa = \bigcup_{i \in I} A_i$ con $|A_i| < \kappa$ para todo $i$, entonces $|I| \geq \kappa$.
  \item No existe ninguna función $f \colon \lambda \to \kappa$ cofinal con $\lambda < \kappa$.
\end{enumerate}
\end{teorema}

\begin{proof}
Las tres condiciones son equivalentes por definición de cofinalidad: $\mathrm{cf}(\kappa) = \kappa$ significa exactamente que la menor longitud de una sucesión cofinal en $\kappa$ es $\kappa$, lo que equivale a (iii). Y (iii) equivale a (ii) porque si $\kappa = \bigcup_{i \in I} A_i$ con $|A_i| < \kappa$ y $|I| < \kappa$, entonces la función $f(i) = \sup A_i$ sería cofinal en $\kappa$ (o se puede ajustar) con dominio de cardinal $< \kappa$.
\end{proof}

\begin{corolario}{Todo cardinal sucesor es regular}{succregular}
Para todo cardinal $\kappa$, el sucesor $\kappa^+$ es regular.
\end{corolario}

\begin{proof}
Es la proposición~\ref{prop:propcofcard}(iv).
\end{proof}

\begin{ejemplo}{Los primeros cardinales regulares}{primreg}
\begin{enumerate}[leftmargin=2em, label=(\alph*)]
  \item $\aleph_0 = \omega$ es regular: ninguna sucesión finita es cofinal en $\omega$ (toda sucesión finita tiene un máximo). Así, $\mathrm{cf}(\aleph_0) = \aleph_0$.
  \item $\aleph_1 = \omega_1$ es regular: es el sucesor cardinal de $\aleph_0$, y los sucesores son regulares.
  \item $\aleph_2$ es regular: es el sucesor cardinal de $\aleph_1$.
  \item $\aleph_n$ es regular para todo $n < \omega$: todos son cardinales sucesores.
  \item $\aleph_{\omega+1}$ es regular: es el sucesor cardinal de $\aleph_\omega$.
  \item $\aleph_{\alpha+1}$ es regular para todo ordinal $\alpha$: todos los cardinales sucesores son regulares.
\end{enumerate}
\end{ejemplo}

\begin{ejemplo}{Verificación directa de la regularidad de $\aleph_1$}{verifregaleph1}
Comprobamos directamente que $\mathrm{cf}(\aleph_1) = \aleph_1$.

Supongamos que existe $f \colon \omega \to \omega_1$ cofinal (sucesión de tipo $\omega$ cofinal en $\omega_1$). Entonces $\omega_1 = \sup_{n<\omega} f(n)$, es decir, todo ordinal numerable está acotado por algún $f(n)$. Pero la unión $\bigcup_{n<\omega} f(n)$ (como conjunto de ordinales) tiene cardinal
\[
  \left|\bigcup_{n < \omega} f(n)\right| \leq \sum_{n<\omega} |f(n)| \leq \aleph_0 \cdot \aleph_0 = \aleph_0,
\]
ya que cada $f(n) < \omega_1$ es un ordinal numerable. Luego $\omega_1 \leq \aleph_0$, contradicción. Por tanto no existe sucesión numerable cofinal en $\omega_1$, y $\mathrm{cf}(\aleph_1) = \aleph_1$.
\end{ejemplo}

\begin{teorema}{El cardinal $\aleph_0$ y los cardinales sucesores agotan los cardinales regulares asequibles}{regzfc}
En ZFC, los únicos cardinales infinitos regulares cuya existencia puede probarse son:
\begin{itemize}[leftmargin=2em]
  \item $\aleph_0$, y
  \item los cardinales sucesores $\kappa^+$ para cada cardinal $\kappa$.
\end{itemize}
La existencia de un cardinal límite regular mayor que $\aleph_0$ (llamado \textbf{cardinal débilmente inaccesible}) no puede demostrarse en ZFC.
\end{teorema}

\begin{proof}
La afirmación sobre los sucesores es el Corolario~\ref{cor:succregular}. La afirmación sobre los inaccesibles es un resultado de independencia: si ZFC es consistente, también lo es ZFC + ``no existe ningún cardinal débilmente inaccesible'' \cite{godel1940, jech2003, kanamori2003}.
\end{proof}

% ===========================================================
\section{Cardinales singulares: definición, ejemplos y propiedades}
\label{sec:cardinales-singulares}
% ===========================================================

\subsection{Definición y primeros ejemplos}

\begin{definicion}{Cardinal singular}{cardsing}
Un cardinal infinito $\kappa$ es \textbf{singular} si existe una función cofinal $f \colon \delta \to \kappa$ con $\delta < \kappa$. Equivalentemente, $\kappa$ es singular si $\mathrm{cf}(\kappa) < \kappa$.

El cardinal $\delta = \mathrm{cf}(\kappa)$ recibe el nombre de \textbf{cofinalidad} de $\kappa$, y necesariamente es un cardinal regular (por la propiedad de idempotencia de la cofinalidad).
\end{definicion}

\begin{importante}
Para que un cardinal $\kappa$ sea singular, necesita:
\begin{enumerate}[leftmargin=2em, label=(\arabic*)]
  \item Ser un cardinal \emph{límite}: los cardinales sucesores son siempre regulares.
  \item Poder alcanzarse «desde abajo» mediante una sucesión de longitud $< \kappa$: existe $\{\kappa_\alpha\}_{\alpha < \delta}$ con $\delta < \kappa$ y $\kappa_\alpha < \kappa$ y $\sup_{\alpha<\delta} \kappa_\alpha = \kappa$.
\end{enumerate}
\end{importante}

\begin{ejemplo}{$\aleph_\omega$ es singular con $\mathrm{cf}(\aleph_\omega) = \aleph_0$}{alephomsing}
El cardinal $\aleph_\omega$ es el límite de la sucesión $\aleph_0 < \aleph_1 < \aleph_2 < \cdots$

Definimos $f \colon \omega \to \aleph_\omega$ por $f(n) = \aleph_n$. Entonces:
\[
  \sup_{n < \omega} f(n) = \sup_{n < \omega} \aleph_n = \aleph_\omega.
\]
Esta función $f$ es cofinal en $\aleph_\omega$ y tiene dominio $\omega = \aleph_0 < \aleph_\omega$. Luego $\mathrm{cf}(\aleph_\omega) \leq \aleph_0$.

Comprobamos que $\mathrm{cf}(\aleph_\omega) = \aleph_0$ (no puede ser finita): ninguna sucesión finita $\aleph_{n_1}, \ldots, \aleph_{n_k}$ es cofinal en $\aleph_\omega$, pues su máximo $\aleph_{\max(n_i)} < \aleph_\omega$.

Por tanto $\mathrm{cf}(\aleph_\omega) = \aleph_0 < \aleph_\omega$, y $\aleph_\omega$ es singular.
\end{ejemplo}

\begin{ejemplo}{Cardinales singulares de cofinalidad $\aleph_0$}{cardsingness0}
Todos los siguientes son cardinales singulares de cofinalidad $\aleph_0$:
\begin{enumerate}[leftmargin=2em, label=(\alph*)]
  \item $\aleph_\omega$: $\sup_{n<\omega} \aleph_n = \aleph_\omega$.
  \item $\aleph_{\omega \cdot 2} = \aleph_{\omega + \omega}$: $\sup_{n<\omega} \aleph_{\omega + n} = \aleph_{\omega \cdot 2}$.
  \item $\aleph_{\omega^2}$: $\sup_{n<\omega} \aleph_{\omega \cdot n} = \aleph_{\omega^2}$.
  \item $\aleph_{\omega^\omega}$: $\sup_{n<\omega} \aleph_{\omega^n} = \aleph_{\omega^\omega}$.
  \item $\aleph_{\varepsilon_0}$: $\sup_{n<\omega} \aleph_{\varphi_n} = \aleph_{\varepsilon_0}$ donde $\varphi_n$ es la sucesión que converge a $\varepsilon_0$.
\end{enumerate}
En general, $\aleph_\lambda$ es singular de cofinalidad $\aleph_0$ siempre que $\lambda$ sea un ordinal límite de cofinalidad $\omega$ (por ejemplo, $\lambda = \omega, \omega^2, \omega^\omega, \varepsilon_0$, etc.).
\end{ejemplo}

\begin{ejemplo}{Cardinal singular de cofinalidad $\aleph_1$}{cardsingness1}
Sea $\lambda = \omega_1$ (el primer ordinal no numerable). Entonces $\aleph_{\omega_1}$ es un cardinal singular con $\mathrm{cf}(\aleph_{\omega_1}) = \aleph_1$.

La función cofinal es $f \colon \omega_1 \to \aleph_{\omega_1}$ dada por $f(\alpha) = \aleph_\alpha$:
\[
  \sup_{\alpha < \omega_1} \aleph_\alpha = \aleph_{\omega_1}.
\]
Ninguna sucesión de longitud $< \omega_1$ (es decir, de longitud $\aleph_0$) puede ser cofinal en $\aleph_{\omega_1}$, pues la unión de $\aleph_0$ conjuntos de cardinalidad $\leq \aleph_\alpha$ para algún $\alpha < \omega_1$ tiene cardinalidad $\leq \aleph_0 \cdot \aleph_\alpha = \aleph_\alpha < \aleph_{\omega_1}$.

Por tanto $\mathrm{cf}(\aleph_{\omega_1}) = \aleph_1 < \aleph_{\omega_1}$, y $\aleph_{\omega_1}$ es singular.
\end{ejemplo}

\begin{ejemplo}{Caracterización de la cofinalidad de $\aleph_\alpha$ en función de $\alpha$}{cofinalidadindice}
La relación entre el índice $\alpha$ y la cofinalidad de $\aleph_\alpha$ es:
\begin{center}
\begin{tabular}{lll}
\textbf{$\alpha$} & \textbf{Tipo} & \textbf{$\mathrm{cf}(\aleph_\alpha)$} \\
\hline
$0$ & límite ($\omega = \aleph_0$) & $\aleph_0$ (regular) \\
$1, 2, 3, \ldots$ (sucesor finito) & sucesor & $\aleph_n$ (regular) \\
$\omega$ & límite, $\mathrm{cf}(\omega) = \omega$ & $\aleph_0$ (singular) \\
$\omega + 1$ & sucesor & $\aleph_{\omega+1}$ (regular) \\
$\omega \cdot 2$ & límite, $\mathrm{cf}(\omega \cdot 2) = \omega$ & $\aleph_0$ (singular) \\
$\omega_1$ & límite, $\mathrm{cf}(\omega_1) = \omega_1$ & $\aleph_1$ (singular) \\
$\omega_1 + 1$ & sucesor & $\aleph_{\omega_1+1}$ (regular) \\
\end{tabular}
\end{center}
\medskip
En general: $\mathrm{cf}(\aleph_\alpha) = \mathrm{cf}(\alpha)$ cuando $\alpha$ es un ordinal límite (la cofinalidad del cardinal es la cofinalidad del índice). Para $\alpha$ sucesor, $\aleph_\alpha$ es el sucesor cardinal de $\aleph_{\alpha-1}$ y es regular.
\end{ejemplo}

\subsection{Consecuencias de la singularidad en la aritmética cardinal}

\begin{teorema}{König aplicado a cardinales singulares}{koenig-singular}
Sea $\kappa$ un cardinal singular con $\mathrm{cf}(\kappa) = \delta$. Sea $f \colon \delta \to \kappa$ cofinal. Entonces:
\[
  \sum_{\alpha < \delta} f(\alpha) = \kappa \qquad \text{y} \qquad
  \kappa < \prod_{\alpha < \delta} \kappa = \kappa^\delta.
\]
En particular, $\kappa < \kappa^\delta = \kappa^{\mathrm{cf}(\kappa)}$.
\end{teorema}

\begin{proof}
La primera igualdad: como $f(\alpha) < \kappa$ y $\sup f(\alpha) = \kappa$, se tiene $\sum_{\alpha<\delta} f(\alpha) \leq \delta \cdot \kappa = \kappa$ (pues $\delta < \kappa$), y la desigualdad opuesta es obvia.

La desigualdad estricta es el Corolario~\ref{cor:kappacofkappa}, que es consecuencia del Teorema de König.
\end{proof}

\begin{corolario}{Restricción sobre $\aleph_\omega^{\aleph_0}$}{aomegacotas}
Se tiene la siguiente acotación para $\aleph_\omega^{\aleph_0}$:
\[
  \aleph_\omega < \aleph_\omega^{\aleph_0}.
\]
Además, por el Teorema de König (con $I = \omega$, $\kappa_n = \aleph_n < \lambda_n = \aleph_\omega^{\aleph_0}$):
\[
  \aleph_\omega = \sum_{n < \omega} \aleph_n < \prod_{n<\omega} \aleph_\omega^{\aleph_0} = \bigl(\aleph_\omega^{\aleph_0}\bigr)^{\aleph_0} = \aleph_\omega^{\aleph_0}.
\]
El valor exacto de $\aleph_\omega^{\aleph_0}$ depende de hipótesis adicionales sobre la aritmética cardinal singular. El \textbf{Teorema de Shelah} afirma que si $\aleph_\omega$ es fuerte-límite (es decir, $2^{\aleph_n} < \aleph_\omega$ para todo $n < \omega$), entonces $\aleph_\omega^{\aleph_0} < \aleph_{\omega_4}$ \cite{shelah1994}.
\end{corolario}

\begin{observacion}{Hipótesis de Cardinal Singular (SCH)}{sch}
La \textbf{Hipótesis de Cardinal Singular} (SCH) afirma: si $\kappa$ es un cardinal singular fuerte-límite, entonces $2^\kappa = \kappa^+$. Bajo GCH, SCH se cumple. Sin GCH, SCH no es necesariamente verdadera, pero es considerablemente más difícil de refutar que GCH; de hecho, refutar SCH requiere grandes cardinales. El resultado profundo de Shelah es que SCH es «casi siempre» verdadera para cardinales singulares de cofinalidad $\aleph_0$ \cite{shelah1994, jech2003}.
\end{observacion}

\subsection{Cardinales singulares y el Teorema de Silver}

\begin{teorema}{Teorema de Silver (enunciado)}{teosilver}
Sea $\kappa$ un cardinal singular de cofinalidad no numerable, es decir, $\mathrm{cf}(\kappa) > \aleph_0$. Si la Hipótesis del Continuo Generalizada vale para todos los cardinales $< \kappa$ (es decir, $2^\lambda = \lambda^+$ para todo cardinal infinito $\lambda < \kappa$), entonces $2^\kappa = \kappa^+$.
\end{teorema}

\begin{proof}[Referencia]
La demostración utiliza la teoría de \emph{ultrapotencias} y \emph{funciones regresivas} (funciones de Fodor). Se omite en este curso; véase \cite{jech2003} (Capítulo 8) o \cite{kunen2011} \cite{silver1975}.
\end{proof}

\begin{observacion}{Importancia del Teorema de Silver}{importsilver}
El Teorema de Silver fue sorprendente cuando fue demostrado en 1975 \cite{silver1975}: muestra que la aritmética de cardinales singulares no es completamente libre, sino que está parcialmente determinada por la aritmética de los cardinales más pequeños. Contrasta con el resultado de Easton (1970) \cite{easton1970}, que muestra que la función $\kappa \mapsto 2^\kappa$ para cardinales regulares puede tomar prácticamente cualquier valor consistente con el Teorema de König.
\end{observacion}

% ===========================================================
\section{Resumen: aritmética cardinal infinita}
% ===========================================================

La siguiente tabla recoge las identidades y desigualdades fundamentales establecidas en esta clase.

\begin{center}
\begin{tabular}{lll}
\textbf{Operación} & \textbf{Resultado} & \textbf{Condición} \\
\hline
$\kappa + \lambda$ & $\max(\kappa, \lambda)$ & $\kappa, \lambda$ infinitos \\
$\kappa \cdot \lambda$ & $\max(\kappa, \lambda)$ & $\kappa, \lambda$ infinitos \\
$\kappa^\lambda$ & $2^\lambda$ & $2^\lambda \geq \kappa$ \\
$\kappa^\lambda$ & depende de $\mathrm{cf}(\kappa)$ & $2^\lambda < \kappa$ \\
$\kappa + \kappa$ & $\kappa$ & $\kappa$ infinito \\
$\kappa \cdot \kappa$ & $\kappa$ & $\kappa$ infinito \\
$\kappa^{\mathrm{cf}(\kappa)}$ & $> \kappa$ (König) & $\kappa$ infinito \\
$\sum_i \kappa_i$ & $< \prod_i \lambda_i$ si $\kappa_i < \lambda_i$ & (König general) \\
$\mathrm{cf}(2^\kappa)$ & $> \kappa$ & $\kappa$ infinito \\
\hline
\end{tabular}
\end{center}

\bigskip

\begin{importante}
\textbf{Resumen conceptual de los tipos de cardinales infinitos:}
\begin{itemize}[leftmargin=2em]
  \item \textbf{Cardinales regulares}: $\mathrm{cf}(\kappa) = \kappa$. Ejemplos: $\aleph_0$ y todos los cardinales sucesores $\aleph_{\alpha+1}$. En ZFC no se puede probar la existencia de cardinales límite regulares mayores que $\aleph_0$.
  \item \textbf{Cardinales singulares}: $\mathrm{cf}(\kappa) < \kappa$. Ejemplos: $\aleph_\omega$, $\aleph_{\omega^2}$, $\aleph_{\omega_1}$, en general $\aleph_\lambda$ con $\lambda$ límite. La aritmética de cardinales singulares es el tema central de la investigación moderna en teoría de conjuntos.
  \item \textbf{Cardinales inaccesibles}: cardinales límite regulares. Su existencia es independiente de ZFC y su estudio pertenece a la teoría de grandes cardinales.
\end{itemize}
\end{importante}

% ===========================================================
\section{Problemas y tareas}
% ===========================================================

\begin{enumerate}[leftmargin=2em, label=\textbf{\arabic*.}]

\item \textbf{(Aritmética cardinal básica.)}
Calcule y justifique los siguientes cardinales:
\begin{enumerate}[label=(\alph*)]
  \item $\aleph_3 + \aleph_7$.
  \item $\aleph_0 \cdot \aleph_{\omega}$.
  \item $\aleph_5 \cdot \aleph_5$.
  \item $(\aleph_0 + \aleph_1) \cdot (\aleph_2 + \aleph_3)$.
  \item $\aleph_\omega \cdot \aleph_{\omega+1}$.
\end{enumerate}

\item \textbf{(Potenciación cardinal.)}
\begin{enumerate}[label=(\alph*)]
  \item Pruebe que $2^{\aleph_0} = \aleph_0^{\aleph_0}$.
  \item Pruebe que $\aleph_1^{\aleph_0} = 2^{\aleph_0}$ (sin ninguna hipótesis adicional).
  \item Pruebe que $\aleph_\alpha^{\aleph_\alpha} = 2^{\aleph_\alpha}$ para todo ordinal $\alpha$.
  \item Bajo GCH, calcule $\aleph_5^{\aleph_3}$, $\aleph_3^{\aleph_5}$ y $\aleph_\omega^{\aleph_0}$.
  \item Demuestre que $(2^{\aleph_0})^{\aleph_0} = 2^{\aleph_0}$ (el continuo elevado a $\aleph_0$ es el propio continuo).
\end{enumerate}

\item \textbf{(Teorema de König.)}
\begin{enumerate}[label=(\alph*)]
  \item Usando el Teorema de König, demuestre que $\aleph_\omega \neq 2^{\aleph_0}$.
  \item Pruebe que $\mathrm{cf}(2^{\aleph_0}) > \aleph_0$, es decir, el continuo no puede ser $\aleph_\omega$. (\emph{¿Puede ser $2^{\aleph_0} = \aleph_{\omega_1}$?})
  \item Demuestre que para cualquier cardinal $\kappa$, $\kappa^\kappa = 2^\kappa$ cuando $\mathrm{cf}(\kappa) = \kappa$ (es decir, cuando $\kappa$ es regular).
  \item Sea $\{A_n\}_{n<\omega}$ una familia de conjuntos con $|A_n| = \aleph_n$. Use el Teorema de König para demostrar que $\left|\prod_{n<\omega} A_n\right| > \aleph_\omega$.
  \item (Generalización) Enuncia y demuestra la versión del Teorema de König para una familia de dos cardinales: $\kappa < \lambda$ implica $\kappa + \kappa < \lambda \cdot \lambda$.
\end{enumerate}

\item \textbf{(Cofinalidad.)}
\begin{enumerate}[label=(\alph*)]
  \item Calcule $\mathrm{cf}(\aleph_{\omega \cdot 3})$, $\mathrm{cf}(\aleph_{\omega^3})$, $\mathrm{cf}(\aleph_{\omega^\omega})$, $\mathrm{cf}(\aleph_{\varepsilon_0})$.
  \item Demuestre que $\mathrm{cf}(\aleph_\lambda) = \mathrm{cf}(\lambda)$ para todo ordinal límite $\lambda$.
  \item Demuestre que para todo cardinal regular $\kappa > \aleph_0$ y toda familia $\{A_\alpha\}_{\alpha < \lambda}$ con $\lambda < \kappa$ y $|A_\alpha| < \kappa$, se tiene $\left|\bigcup_{\alpha < \lambda} A_\alpha\right| < \kappa$. (Esta es la propiedad de \emph{cierre por uniones pequeñas} de los cardinales regulares.)
  \item Pruebe que si $f \colon \kappa \to \kappa$ es estrictamente creciente y $\kappa$ es regular, entonces $f$ tiene un punto fijo, es decir, existe $\alpha < \kappa$ con $f(\alpha) = \alpha$. (\emph{Sugerencia}: considere $\sup_{n<\omega} f^n(0)$ si $\kappa$ es regular de cofinalidad $> \omega$; ajuste para $\kappa = \omega$.)
\end{enumerate}

\item \textbf{(Cardinales regulares y singulares.)}
\begin{enumerate}[label=(\alph*)]
  \item Demuestre que $\aleph_{\omega_1}$ es singular y determine su cofinalidad.
  \item Sea $\kappa$ un cardinal singular con $\mathrm{cf}(\kappa) = \delta$. Pruebe que $\kappa^\delta > \kappa$ pero $\kappa^\mu = \kappa$ para todo $\mu < \delta$ con $2^\mu < \kappa$.
  \item Demuestre que no existe ningún cardinal singular mínimo, es decir, para cada cardinal singular $\kappa$ existe un cardinal singular $\mu < \kappa$.
  \item Pruebe que la clase de cardinales regulares es cofinal en la clase de todos los cardinales.
\end{enumerate}

\item \textbf{(Cardinales inaccesibles.)}
\begin{enumerate}[label=(\alph*)]
  \item Un cardinal $\kappa$ es \textbf{débilmente inaccesible} si es un cardinal límite regular mayor que $\aleph_0$. Demuestre que si existe un cardinal débilmente inaccesible $\kappa$, entonces $\kappa = \aleph_\kappa$.
  \item Un cardinal $\kappa$ es \textbf{fuertemente inaccesible} si además es fuerte-límite: $2^\lambda < \kappa$ para todo $\lambda < \kappa$. Demuestre que si $\kappa$ es fuertemente inaccesible, entonces $V_\kappa \models \text{ZFC}$.
  \item (Investigación) Explique por qué la consistencia de ZFC no puede implicar la existencia de un cardinal inaccesible (use el Segundo Teorema de Incompletitud de Gödel).
\end{enumerate}

\item \textbf{(El Teorema de Silver: exploración.)}
Sea $\kappa$ un cardinal singular de cofinalidad $\aleph_1$.
\begin{enumerate}[label=(\alph*)]
  \item Escriba $\kappa = \sup_{\alpha < \omega_1} \kappa_\alpha$ con $\kappa_\alpha < \kappa$ estrictamente creciente. Suponga que $2^{\kappa_\alpha} = \kappa_\alpha^+$ para todo $\alpha < \omega_1$. Enuncie qué predice el Teorema de Silver sobre $2^\kappa$.
  \item Compare este resultado con el caso $\mathrm{cf}(\kappa) = \aleph_0$: ¿qué diferencias hay en lo que la aritmética cardinal singular puede o no puede controlar?
  \item Explique, informalmente, por qué la condición $\mathrm{cf}(\kappa) > \aleph_0$ en el Teorema de Silver es esencial, citando el ejemplo de $\aleph_\omega$.
\end{enumerate}

\item \textbf{(Problema de investigación.)}
La \textbf{función pcf} de Shelah (possible cofinalities) es el instrumento más potente conocido para estudiar la aritmética de cardinales singulares. Investiga los siguientes aspectos:
\begin{enumerate}[label=(\alph*)]
  \item ¿Qué es el conjunto $\mathrm{pcf}(A)$ para un conjunto progresivo de cardinales $A$?
  \item Enuncia el Teorema de Shelah: $\aleph_\omega^{\aleph_0} < \aleph_{\omega_4}$ (bajo la hipótesis $\aleph_\omega > 2^{\aleph_0}$). ¿Cuál es la cota actual conocida?
  \item ¿Es demostrable en ZFC que $\aleph_\omega^{\aleph_0} < \aleph_{\omega_1}$?
\end{enumerate}

\end{enumerate}

% ===========================================================
\section{Referencias de la Clase}
% ===========================================================

Los contenidos de esta clase se encuentran desarrollados con todo detalle en las siguientes obras. Para la aritmética cardinal elemental (suma, producto, potenciación) y el Teorema de Hessenberg, las referencias fundamentales son \cite{jech2003} (capítulos 5 y 6), \cite{enderton1977} (capítulo 6) y \cite{hrbacek1999} (capítulo 4). El artículo original del Teorema de König es \cite{konig1905}; una exposición moderna y completa se encuentra en \cite{jech2003} y \cite{kunen2011}. Para la cofinalidad y los cardinales regulares y singulares, la presentación estándar está en \cite{jech2003} (capítulo 3), \cite{kunen2011} y \cite{levy1979}. La teoría de grandes cardinales y los cardinales inaccesibles se desarrollan en \cite{kanamori2003}. El Teorema de Easton sobre la libertad de $\kappa \mapsto 2^\kappa$ para regulares es \cite{easton1970}. El Teorema de Silver se encuentra en \cite{silver1975}; véanse también las exposiciones en \cite{jech2003} y \cite{devlin1993}. Para la aritmética cardinal singular avanzada (teoría pcf de Shelah), la referencia principal es \cite{shelah1994}; una exposición introductoria accesible se encuentra en \cite{jech2003} (capítulo 24) y en \cite{kanamori2003}.

\nocite{sierpinski1958, suppes1972, halmos1960}

% ===========================================================
\chapter{Clase 9: La Hipótesis del Continuo: Enunciado, Historia e Independencia}
% ===========================================================

El \textbf{problema del continuo} es, sin duda, uno de los más profundos y revolucionarios de toda la historia de las matemáticas. Formulado por Georg Cantor en la década de 1870, pregunta con precisión cuántos puntos hay en la recta real comparado con la menor infinitud incontable. Su respuesta, o más bien la \emph{imposibilidad} de responderla dentro del sistema axiomático estándar, transformó para siempre nuestra comprensión de los fundamentos de las matemáticas.

En esta clase enunciamos con rigor la \textbf{Hipótesis del Continuo} (HC) y su generalización (HCG), narramos el recorrido histórico del problema desde Cantor hasta Paul Cohen, y exponemos los dos grandes resultados del siglo~XX que juntos establecen la \emph{independencia} de HC respecto de los axiomas de Zermelo--Fraenkel con Elección (ZFC):

\begin{enumerate}[leftmargin=2em, label=\textbf{\arabic*.}]
  \item \textbf{Kurt Gödel (1940):} HC es \emph{consistente} con ZFC; es decir, si ZFC no tiene contradicción, tampoco la tiene ZFC $+$ HC. Esto se establece mediante el \textbf{universo constructible} $L$.
  \item \textbf{Paul Cohen (1963):} $\neg$HC es \emph{consistente} con ZFC; es decir, si ZFC no tiene contradicción, tampoco la tiene ZFC $+$ $\neg$HC. Esto se establece mediante la técnica del \textbf{forcing}.
\end{enumerate}

Juntos, estos dos resultados prueban que HC es \textbf{independiente} de ZFC: ni se puede demostrar ni se puede refutar dentro de ese sistema. El trabajo de Gödel y Cohen constituye, posiblemente, el logro más celebrado de la lógica matemática del siglo~XX.

La exposición sigue los tratamientos clásicos de \cite{jech2003, kunen2011, cohen1966, godel1940}, con referencias adicionales a \cite{enderton1977, devlin1993, hrbacek1999}.

% ===========================================================
\section{El problema del continuo: contexto histórico}
% ===========================================================

\subsection{Cantor y la aritmética de lo infinito}

Para comprender la profundidad de la Hipótesis del Continuo es imprescindible remontarse a los trabajos de Georg Cantor de la década de 1870. Cantor demostró en 1873 que los números reales $\mathbb{R}$ no son enumerables: no existe ninguna función biyectiva de $\mathbb{N}$ en $\mathbb{R}$. Este resultado, publicado en 1874, inauguró la teoría de conjuntos y estableció que existen \emph{distintos tamaños de infinito}.

La prueba diagonal de Cantor (1891) generaliza este resultado: para cualquier conjunto $A$, el conjunto de todas sus partes $\mathcal{P}(A)$ tiene cardinalidad estrictamente mayor que $A$. En símbolos: $|A| < |\mathcal{P}(A)|$, o equivalentemente $\kappa < 2^\kappa$ para todo cardinal $\kappa$.

Aplicado a los números naturales: $|\mathbb{N}| = \aleph_0 < 2^{\aleph_0} = |\mathbb{R}|$. El cardinal del continuo $\mathfrak{c} = 2^{\aleph_0}$ es estrictamente mayor que $\aleph_0$. La pregunta inmediata que formuló Cantor fue:


¿Existe algún conjunto de números reales cuya cardinalidad sea estrictamente mayor que $\aleph_0$ (la de los naturales) y estrictamente menor que $\mathfrak{c} = 2^{\aleph_0}$ (la de los reales)?

\begin{importante}
Cantor conjeturó que la respuesta es \textbf{no}: no existe ningún cardinal $\kappa$ con $\aleph_0 < \kappa < 2^{\aleph_0}$. Esta conjetura es la \textbf{Hipótesis del Continuo}.
\end{importante}

Cantor intentó demostrarlo durante décadas sin éxito. Su fracaso no fue falta de ingenio: el problema resultó ser, en cierto sentido, más difícil que cualquier pregunta matemática ordinaria.

\subsection{El primer problema de Hilbert (1900)}

En el Congreso Internacional de Matemáticos de París de 1900, David Hilbert presentó su célebre lista de 23 problemas para el nuevo siglo. El \textbf{primer problema} de la lista era precisamente la Hipótesis del Continuo. Hilbert creía que toda pregunta matemática bien planteada podía ser resuelta dentro del sistema axiomático de la matemática. La historia demostró que estaba equivocado, al menos en este caso.

\subsection{Los esfuerzos previos a Gödel}

Durante las primeras décadas del siglo~XX, la Hipótesis del Continuo permaneció indemostrable. En 1905, König demostró que $2^{\aleph_0}$ no puede ser $\aleph_\omega$ (un cardinal de cofinalidad numerable), lo cual descartó algunos posibles valores del continuo pero no resolvió el problema. En 1925, Luzin y Sierpi\'nski estudiaron los subconjuntos del plano, los conjuntos proyectivos y la hipótesis de Suslin, todos relacionados con HC.

El panorama cambió radicalmente en 1940 con el trabajo de Gödel.

% ===========================================================
\section{Enunciado de la Hipótesis del Continuo}
% ===========================================================

\subsection{Recordatorio: la jerarquía aleph y la potenciación cardinal}

Recordamos las nociones esenciales de las clases anteriores. Los \textbf{números cardinales} se identifican con los ordinales iniciales: el menor ordinal de cada cardinalidad. La jerarquía $\aleph$ está definida por inducción transfinita:
\[
  \aleph_0 = \omega, \qquad
  \aleph_{\alpha+1} = \aleph_\alpha^+, \qquad
  \aleph_\lambda = \sup_{\beta < \lambda} \aleph_\beta \;\text{ para } \lambda \text{ límite}.
\]

El \textbf{cardinal del continuo} es $\mathfrak{c} = 2^{\aleph_0} = |\mathbb{R}| = |\mathcal{P}(\mathbb{N})|$. Por el teorema de Cantor, $\mathfrak{c} > \aleph_0$; por lo tanto $\mathfrak{c} \geq \aleph_1$. La Hipótesis del Continuo afirma que esta desigualdad es en realidad una igualdad.

\begin{definicion}{Hipótesis del Continuo (HC)}{hc}
La \textbf{Hipótesis del Continuo} es el enunciado:
\[
  \mathrm{HC} \;:\; 2^{\aleph_0} = \aleph_1.
\]
Equivalentemente, no existe ningún cardinal $\kappa$ tal que $\aleph_0 < \kappa < 2^{\aleph_0}$; es decir, no existe ningún conjunto $A \subseteq \mathbb{R}$ tal que $\aleph_0 < |A| < |\mathbb{R}|$.
\end{definicion}

\begin{observacion}{Formulaciones equivalentes de HC}{hcequiv}
Las siguientes afirmaciones son equivalentes a HC:
\begin{enumerate}[leftmargin=2em, label=(\alph*)]
  \item $|\mathbb{R}| = \aleph_1$.
  \item Todo subconjunto infinito de $\mathbb{R}$ es equipotente con $\mathbb{N}$ o con $\mathbb{R}$.
  \item $|\mathcal{P}(\mathbb{N})| = \aleph_1$.
  \item Todo subconjunto de la recta real que no sea numerable tiene la cardinalidad de la recta real.
  \item $|\{0,1\}^{\mathbb{N}}| = \aleph_1$ (el espacio de Cantor tiene cardinalidad $\aleph_1$).
\end{enumerate}
\end{observacion}

\begin{ejemplo}{Cardinalidades relevantes para HC}{}
Para poner HC en contexto, listamos algunas cardinalidades conocidas:
\begin{itemize}[leftmargin=2em]
  \item $|\mathbb{N}| = \aleph_0$: los naturales son el prototipo de conjunto infinito numerable.
  \item $|\mathbb{Z}| = \aleph_0$, $|\mathbb{Q}| = \aleph_0$: los enteros y los racionales son numerables.
  \item $|\mathbb{R}| = |\mathbb{C}| = |(0,1)| = |[0,1]^2| = 2^{\aleph_0}$: los reales, los complejos, cualquier intervalo abierto, el cuadrado unitario, todos tienen cardinalidad $\mathfrak{c}$.
  \item $|\mathcal{P}(\mathbb{R})| = 2^{\mathfrak{c}} = 2^{2^{\aleph_0}}$: las partes de $\mathbb{R}$ tienen cardinalidad estrictamente mayor que $\mathfrak{c}$.
\end{itemize}
HC afirma que en la secuencia $\aleph_0 < \aleph_1 < \aleph_2 < \cdots$, el primer eslabón por encima de $\aleph_0$ es exactamente $\mathfrak{c} = 2^{\aleph_0}$.
\end{ejemplo}

\subsection{La Hipótesis Generalizada del Continuo}

La Hipótesis Generalizada del Continuo extiende HC a todos los cardinales infinitos.

\begin{definicion}{Hipótesis Generalizada del Continuo (HCG)}{hcg}
La \textbf{Hipótesis Generalizada del Continuo} es el enunciado:
\[
  \mathrm{HCG} \;:\; \forall \alpha \;(\, 2^{\aleph_\alpha} = \aleph_{\alpha+1} \,).
\]
Es decir, para cada número ordinal $\alpha$, el cardinal $2^{\aleph_\alpha}$ es el inmediato sucesor de $\aleph_\alpha$ en la jerarquía aleph.
\end{definicion}

\begin{observacion}{HC como caso particular de HCG}{hchcg}
HC es el caso $\alpha = 0$ de HCG:
\[
  2^{\aleph_0} = \aleph_1.
\]
Los siguientes casos de HCG son:
\[
  2^{\aleph_1} = \aleph_2, \qquad 2^{\aleph_2} = \aleph_3, \qquad 2^{\aleph_\omega} = \aleph_{\omega+1}, \quad \ldots
\]
Cada uno de ellos afirma que la potenciación $2^{\aleph_\alpha}$ no produce un ``salto'' mayor de lo esperado: siempre da el primer cardinal disponible encima de $\aleph_\alpha$.
\end{observacion}

\begin{teorema}{HCG implica HC}{hcghc}
Si vale HCG, entonces vale HC.
\end{teorema}

\begin{proof}
Tomar $\alpha = 0$ en la definición de HCG da directamente $2^{\aleph_0} = \aleph_1$, que es HC.
\end{proof}

\begin{ejemplo}{Consecuencias de HCG: aritmética cardinal simplificada}{}
Bajo HCG, la potenciación cardinal queda completamente determinada para cardinales alephs. Calculamos algunos valores:
\begin{enumerate}[leftmargin=2em, label=(\alph*)]
  \item $2^{\aleph_0} = \aleph_1$ (caso $\alpha = 0$).
  \item $\aleph_1^{\aleph_0}$: como $\aleph_1^{\aleph_0} = (2^{\aleph_0})^{\aleph_0} = 2^{\aleph_0 \cdot \aleph_0} = 2^{\aleph_0} = \aleph_1$ (usando HC y que $\aleph_0 \cdot \aleph_0 = \aleph_0$).
  \item $2^{\aleph_1} = \aleph_2$ (caso $\alpha = 1$ de HCG).
  \item $\aleph_2^{\aleph_1} = (2^{\aleph_1})^{\aleph_1} = 2^{\aleph_1 \cdot \aleph_1} = 2^{\aleph_1} = \aleph_2$ (bajo HCG).
  \item En general, bajo HCG: $\aleph_\beta^{\aleph_\alpha} = \aleph_{\beta+1}$ siempre que $\aleph_\alpha \geq \mathrm{cf}(\aleph_\beta)$ y $\beta \geq \alpha$.
\end{enumerate}
Así, HCG no solo determina la función $\alpha \mapsto 2^{\aleph_\alpha}$, sino que simplifica toda la aritmética de potencias de alephs.
\end{ejemplo}

\begin{ejemplo}{Lo que ZFC sí puede demostrar sobre $2^{\aleph_0}$}{}
Sin asumir HC, ZFC solo puede establecer restricciones parciales sobre $\mathfrak{c} = 2^{\aleph_0}$:
\begin{enumerate}[leftmargin=2em, label=(\alph*)]
  \item $\mathfrak{c} \geq \aleph_1$ (por el teorema de Cantor).
  \item $\mathrm{cf}(\mathfrak{c}) > \aleph_0$, es decir, $\mathfrak{c}$ no puede tener cofinalidad numerable (corolario del teorema de König). En particular, $\mathfrak{c} \neq \aleph_\omega, \aleph_{\omega^2}, \aleph_{\omega^\omega}, \ldots$
  \item $\mathfrak{c}$ puede ser $\aleph_1, \aleph_2, \aleph_3, \aleph_{\omega_1}, \aleph_{\omega_1 + 1}, \ldots$ --- cualquier cardinal de cofinalidad no numerable y $\geq \aleph_1$ es posible, en el sentido de que hay modelos de ZFC donde $\mathfrak{c}$ toma ese valor.
\end{enumerate}
En particular, los valores $\mathfrak{c} = \aleph_1$ (HC), $\mathfrak{c} = \aleph_2$, $\mathfrak{c} = \aleph_{17}$, o $\mathfrak{c} = \aleph_{\omega_1}$ son todos consistentes con ZFC.
\end{ejemplo}

% ===========================================================
\section{El universo constructible de Gödel y la consistencia de HC}
% ===========================================================

\subsection{Motivación: la estrategia de Gödel}

La idea de Gödel para demostrar la consistencia de HC con ZFC fue construir un \emph{modelo} de ZFC dentro del cual HC también fuera verdadera. Si pudiera exhibirse tal modelo, quedaría demostrado que ZFC $+$ HC no puede ser contradictorio (pues ese modelo sería un testigo de la coherencia del sistema extendido).

La construcción que Gödel ideó —el \textbf{universo constructible} $L$— es un subcollection del universo $V$ de todos los conjuntos, construido de manera análoga a la jerarquía acumulativa $V_\alpha$, pero en cada etapa solo admitiendo los subconjuntos que son ``definibles'' mediante fórmulas del lenguaje de la teoría de conjuntos con parámetros de etapas anteriores.

\subsection{Preliminares: la jerarquía acumulativa}

Recordamos que el universo de todos los conjuntos $V$ está estratificado en la \textbf{jerarquía acumulativa de von Neumann}:
\[
  V_0 = \emptyset, \qquad
  V_{\alpha+1} = \mathcal{P}(V_\alpha), \qquad
  V_\lambda = \bigcup_{\beta < \lambda} V_\beta \;\text{ (para } \lambda \text{ límite)}.
\]
Todo conjunto pertenece a algún $V_\alpha$. El \textbf{rango} de un conjunto $x$ es el menor $\alpha$ tal que $x \in V_{\alpha+1}$.

\subsection{Definibilidad y la operación \texorpdfstring{$\mathrm{Def}$}{Def}}

La idea central de la construcción de $L$ es reemplazar la operación ``tomar todas las partes'' ($\mathcal{P}$) por una operación más restrictiva que solo produce las partes \emph{definibles}.

\begin{definicion}{Subconjuntos definibles sobre un conjunto}{definibles}
Sea $M$ un conjunto. Un subconjunto $A \subseteq M$ es \textbf{definible sobre $M$} (con parámetros en $M$) si existe una fórmula $\varphi(x, y_1, \ldots, y_n)$ del lenguaje $\{\in\}$ y parámetros $p_1, \ldots, p_n \in M$ tales que
\[
  A = \{x \in M : (M, \in) \models \varphi(x, p_1, \ldots, p_n)\}.
\]
El conjunto de todos los subconjuntos definibles sobre $M$ se denota
\[
  \mathrm{Def}(M) = \bigl\{\, \{x \in M : (M,\in) \models \varphi(x, p_1, \ldots, p_n)\} \;\bigm|\; \varphi \text{ fórmula},\; p_i \in M \,\bigr\}.
\]
\end{definicion}

\begin{observacion}{$\mathrm{Def}(M)$ versus $\mathcal{P}(M)$}{defvspow}
Siempre $\mathrm{Def}(M) \subseteq \mathcal{P}(M)$, pero la inclusión puede ser estricta. $\mathcal{P}(M)$ contiene \emph{todos} los subconjuntos de $M$, incluyendo los que no son describibles por ninguna fórmula. $\mathrm{Def}(M)$ solo retiene los que pueden ser ``caracterizados'' con una fórmula y parámetros de $M$.

Como $M$ es un conjunto, hay a lo sumo $|M| + \aleph_0$ fórmulas con parámetros en $M$ (puesto que las fórmulas forman un conjunto numerable y los parámetros son elementos de $M$); luego $|\mathrm{Def}(M)| \leq \max(|M|, \aleph_0)$. En cambio, $|\mathcal{P}(M)| = 2^{|M|}$, que puede ser mucho mayor.
\end{observacion}

\subsection{La jerarquía constructible y el universo \texorpdfstring{$L$}{L}}

\begin{definicion}{El universo constructible $L$ de Gödel}{universoL}
La \textbf{jerarquía constructible} se define por inducción transfinita en los ordinales:
\[
  L_0 = \emptyset, \qquad
  L_{\alpha+1} = \mathrm{Def}(L_\alpha), \qquad
  L_\lambda = \bigcup_{\beta < \lambda} L_\beta \;\text{ (para } \lambda \text{ límite)}.
\]
El \textbf{universo constructible} de Gödel es
\[
  L = \bigcup_{\alpha \in \mathrm{Ord}} L_\alpha.
\]
Un conjunto $x$ se dice \textbf{constructible} si $x \in L$. El \textbf{axioma de constructibilidad} es el enunciado $V = L$: ``todo conjunto es constructible''.
\end{definicion}

\begin{ejemplo}{Primeros niveles de la jerarquía $L$}{}
Calculamos los primeros niveles de $L$:
\begin{enumerate}[leftmargin=2em, label=(\alph*)]
  \item $L_0 = \emptyset$.
  \item $L_1 = \mathrm{Def}(L_0) = \mathrm{Def}(\emptyset)$. Los únicos subconjuntos definibles de $\emptyset$ son los definibles en la estructura $(\emptyset, \in)$, que es solo $\emptyset$ mismo. Así $L_1 = \{\emptyset\}$.
  \item $L_2 = \mathrm{Def}(L_1) = \mathrm{Def}(\{\emptyset\})$. Los subconjuntos definibles de $\{\emptyset\}$ son: $\emptyset$ (definido por la fórmula falsa) y $\{\emptyset\}$ (definido por $x = \emptyset$). Luego $L_2 = \{\emptyset, \{\emptyset\}\}$.
  \item $L_3 = \mathrm{Def}(L_2)$. Los subconjuntos definibles de $L_2 = \{\emptyset, \{\emptyset\}\}$ incluyen todos los subconjuntos de $L_2$ que pueden ser definidos, que en este caso son todos: $\emptyset, \{\emptyset\}, \{\{\emptyset\}\}, \{\emptyset, \{\emptyset\}\}$. Luego $L_3 = \mathcal{P}(L_2)$ y $|L_3| = 4$.
  \item Para $n < \omega$, $L_n$ es un conjunto finito y $L_n \subseteq V_n$.
  \item $L_\omega = \bigcup_{n < \omega} L_n$ coincide con $V_\omega$, el conjunto de todos los conjuntos hereditariamente finitos (todos los conjuntos finitos cuyo cierre transitivo es también finito).
\end{enumerate}
\end{ejemplo}

\begin{ejemplo}{Comparación entre $L_{\omega+1}$ y $V_{\omega+1}$}{}
Este ejemplo ilustra la diferencia crucial entre $L$ y $V$:
\begin{itemize}[leftmargin=2em]
  \item $V_{\omega+1} = \mathcal{P}(V_\omega) = \mathcal{P}(\omega)$ (identificando los hereditariamente finitos con $V_\omega$). Como $|\omega| = \aleph_0$, se tiene $|V_{\omega+1}| = 2^{\aleph_0} = \mathfrak{c}$. En $V_{\omega+1}$ están \emph{todos} los subconjuntos de $\omega$, incluyendo los no definibles.
  \item $L_{\omega+1} = \mathrm{Def}(L_\omega) = \mathrm{Def}(\omega)$. Solo contiene los subconjuntos de $\omega$ definibles con parámetros en $\omega$. Como $\omega$ es numerable y hay $\aleph_0$ fórmulas, $|L_{\omega+1}| = \aleph_0$.
\end{itemize}
Así, $L_{\omega+1}$ es \emph{numerable}, mientras que $V_{\omega+1}$ tiene cardinalidad $2^{\aleph_0}$. Dentro de $L$, el conjunto de todos los subconjuntos de $\omega$ es $L_{\omega+1} \cap \mathcal{P}(\omega)$, que es numerable bajo la perspectiva del universo $V$ pero que, dentro de $L$, ``se comporta como si tuviera cardinalidad $\aleph_1$''. Este fenómeno es el núcleo técnico de por qué HC vale en $L$.
\end{ejemplo}

\subsection{Propiedades fundamentales de \texorpdfstring{$L$}{L}}

\begin{teorema}{Propiedades básicas de $L$}{propsL}
La jerarquía constructible satisface:
\begin{enumerate}[leftmargin=2em, label=(\alph*)]
  \item \textbf{Transitividad:} Para todo ordinal $\alpha$, $L_\alpha$ es un conjunto transitivo y $L_\alpha \subseteq L_{\alpha+1} \subseteq L$.
  \item \textbf{Ordinales en $L$:} Todo ordinal es constructible: $\mathrm{Ord} \subseteq L$.
  \item \textbf{Absolutez de $L$:} La definición de ``$x \in L_\alpha$'' es absoluta entre modelos transitivos de ZF que contienen a $\alpha$. En particular, si $M$ es un modelo transitivo de ZF, entonces $(L)^M = L^M$ (el $L$ calculado dentro de $M$ es el mismo que el ``$L$ real'').
  \item \textbf{Cardinalidades en $L$:} Para todo ordinal infinito $\alpha$, $|L_\alpha| = |\alpha|$.
  \item \textbf{$L$ satisface ZFC:} $(L, \in)$ es un modelo de todos los axiomas de ZFC.
  \item \textbf{$L$ satisface $V = L$:} Dentro de $L$, todo conjunto es constructible.
\end{enumerate}
\end{teorema}

El resultado (d) es clave para la consistencia de HC. Lo enunciamos más precisamente:

\begin{proposicion}{Cardinalidad de los niveles de $L$}{cardL}
Para todo cardinal infinito $\kappa$, $|L_\kappa| = \kappa$. En particular:
\[
  |L_{\omega}| = \aleph_0, \qquad |L_{\omega_1}| = \aleph_1, \qquad |L_{\omega_2}| = \aleph_2, \qquad \ldots
\]
\end{proposicion}

\begin{proof}
Por inducción en $\kappa$. Para $\kappa = \omega$: $L_\omega = V_\omega$ es numerable. Para el paso sucesor: si $|L_\kappa| = \kappa$, entonces $|L_{\kappa+1}| = |\mathrm{Def}(L_\kappa)| \leq |\kappa|^{<\omega} \cdot \aleph_0 = \kappa$ (hay $\kappa$ muchas fórmulas con parámetros en $L_\kappa$ y $\kappa$ muchos parámetros). Por otro lado $L_\kappa \subseteq L_{\kappa+1}$, luego $|L_{\kappa+1}| = \kappa$. Para el caso límite: $|L_\lambda| = |\bigcup_{\alpha < \lambda} L_\alpha| \leq \sum_{\alpha < \lambda} |\alpha| = \lambda$.
\end{proof}

\subsection{El teorema de condensación de Gödel}

El resultado técnico central de la teoría del universo constructible es el \textbf{Lema de Condensación}, que es la herramienta clave para demostrar que HC vale en $L$.

\begin{teorema}{Lema de Condensación de Gödel}{condensacion}
Sea $X \prec L_\alpha$ una subestructura elemental de $(L_\alpha, \in)$ para algún ordinal límite $\alpha$. Entonces el colapsamiento de Mostowski de $X$ (la única estructura transitiva isomorfa a $X$) es un $L_\beta$ para algún $\beta \leq \alpha$. En particular, $X \cong (L_\beta, \in)$ para algún ordinal $\beta \leq \alpha$.
\end{teorema}

\begin{observacion}{Por qué el Lema de Condensación implica HC en $L$}{condensacionHC}
El Lema de Condensación puede usarse para controlar el tamaño de $L_{\omega_1}$ dentro de $L$. Aquí presentamos la idea intuitiva:

Todo subconjunto $A$ de $\omega$ que está en $L$ aparece en $L_\alpha$ para algún ordinal contable $\alpha < \omega_1$ (por un argumento de Löwenheim--Skolem: tomamos la clausura de Skolem de $\{A, \omega\}$ en $L_{\omega_1}$, que es numerable; por el Lema de Condensación, esta clausura es isomorfa a algún $L_\beta$ con $\beta < \omega_1$; entonces $A \in L_\beta \subseteq L_{\omega_1}$, pero en realidad $\beta < \omega_1$).

Por lo tanto, cada subconjunto de $\omega$ en $L$ aparece en el nivel $L_\beta$ para algún $\beta < \omega_1$. Como hay $\aleph_1$ muchos ordinales contables y $|L_\beta| = |\beta| \leq \aleph_0$ para cada $\beta < \omega_1$, hay a lo sumo $\aleph_1 \cdot \aleph_0 = \aleph_1$ subconjuntos de $\omega$ en $L$. Pero hay al menos $\aleph_1$ (pues los conjuntos $\{0\}, \{1\}, \{2\}, \ldots$ son numerables y distintos, y $\omega$ mismo es constructible, y $\aleph_1 = \omega_1$ muchos ordinales constructibles son subconjuntos de $\omega$ via la numeración). En conclusión, $L \models |\mathcal{P}(\omega)| = \aleph_1$, que es precisamente HC.
\end{observacion}

\subsection{El teorema principal: consistencia de HC con ZFC}

\begin{teorema}{Gödel, 1940: Consistencia de HC}{godelHC}
Si ZFC es consistente, entonces ZFC $+$ HC es consistente. Más precisamente:
\[
  \mathrm{Con}(\mathrm{ZFC}) \;\Rightarrow\; \mathrm{Con}(\mathrm{ZFC} + \mathrm{HC}).
\]
En particular, HC no puede ser \emph{refutada} en ZFC (si ZFC es consistente).

El mismo resultado vale para HCG: si ZFC es consistente, entonces ZFC $+$ HCG es consistente.
\end{teorema}

\begin{proof}[Esquema de la demostración]
El argumento es el siguiente. Dentro de ZFC, se construye la clase $L$ y se verifica que:
\begin{enumerate}[leftmargin=2em, label=(\roman*)]
  \item $(L, \in)$ satisface todos los axiomas de ZFC. Esto requiere verificar cada axioma de ZFC: extensionalidad, fundación, esquema de separación, potencia, unión, infinito, esquema de reemplazo y elección. La verificación es técnica pero completa. El axioma de elección vale en $L$ porque $L$ tiene un buen orden definible (el orden lexicográfico en las fórmulas y los parámetros).
  \item $(L, \in)$ satisface $V = L$ (axioma de constructibilidad).
  \item A partir de $V = L$, se deduce HC dentro de $L$ usando el Lema de Condensación como se bosquejó arriba.
\end{enumerate}
Por lo tanto, $L$ es un modelo de ZFC $+$ HC. Si ZFC fuera inconsistente, no podría existir ningún modelo de ZFC (por el teorema de completitud de Gödel), contradiciendo nuestra hipótesis. Luego ZFC $+$ HC es consistente.
\end{proof}

\begin{corolario}{HC no es refutable en ZFC}{hcnorefutable}
No existe ninguna demostración en ZFC de $\neg\mathrm{HC}$ (a menos que ZFC sea inconsistente). Formalmente: ZFC $\not\vdash \neg\mathrm{HC}$.
\end{corolario}

% ===========================================================
\section{El forcing de Cohen y la consistencia de \texorpdfstring{$\neg$HC}{negHC}}
% ===========================================================

\subsection{El problema que quedaba abierto}

El resultado de Gödel de 1940 demostró que HC no puede ser refutada en ZFC. Pero quedaba abierta la pregunta opuesta: ¿puede HC ser \emph{demostrada} en ZFC? ¿O es posible construir un modelo de ZFC donde $\neg$HC sea verdadera?

Esta pregunta permaneció abierta durante 23 años. En 1963, Paul Cohen inventó una técnica completamente nueva —el \textbf{forcing}— y la usó para construir un modelo de ZFC donde $2^{\aleph_0} = \aleph_2$, demostrando que $\neg$HC es consistente con ZFC.

\subsection{La idea intuitiva del forcing}

El forcing es una técnica para \emph{extender} modelos de ZFC. Comenzamos con un modelo $M$ de ZFC (llamado el \textbf{modelo base} o \textbf{modelo del suelo}, \emph{ground model} en inglés) y construimos un modelo más grande $M[G]$ (el \textbf{modelo genérico}) añadiendo un nuevo conjunto $G$ que no estaba en $M$ pero que satisface ciertas propiedades compatibles con ZFC.

La analogía más conocida es la de \textbf{adjuntar una raíz}: en álgebra, podemos extender $\mathbb{Q}$ a $\mathbb{Q}(\sqrt{2})$ adjuntando un nuevo elemento $\sqrt{2}$ que satisface $x^2 = 2$. De manera análoga, en forcing se adjunta al modelo base un nuevo conjunto $G$ (llamado el \textbf{genérico} o \textbf{filtro genérico}) que satisface ciertas propiedades controladas.

El mecanismo técnico del forcing involucra los siguientes ingredientes:

\begin{enumerate}[leftmargin=2em, label=\textbf{(\arabic*)}]
  \item \textbf{Una noción de forcing $\mathbb{P} = (P, \leq)$:} Un orden parcial en el modelo base $M$. Los elementos de $P$ se llaman \textbf{condiciones}. La relación $p \leq q$ se lee ``$p$ es más fuerte que $q$'' o ``$p$ extiende a $q$'': una condición más fuerte impone más restricciones al nuevo conjunto que se está añadiendo.

  \item \textbf{Filtros genéricos $G$:} Un subconjunto $G \subseteq P$ es un \textbf{filtro} si (i) es hacia arriba cerrado, (ii) es dirigido hacia abajo (cualesquiera $p, q \in G$ tienen una extensión común $r \in G$ con $r \leq p, q$). $G$ se llama \textbf{$M$-genérico} si intersecta a cada \textbf{conjunto denso} $D \subseteq P$ que pertenece a $M$ (un conjunto $D$ es denso si cada $p \in P$ tiene una extensión $q \leq p$ con $q \in D$). En general, si $M$ es numerable, los filtros $M$-genéricos existen (por el teorema de Rasiowa--Sikorski).

  \item \textbf{El modelo genérico $M[G]$:} La extensión mínima de $M$ que contiene a $G$ y satisface ZFC. Todo elemento de $M[G]$ se puede ``nombrar'' desde $M$ usando los llamados \textbf{nombres de forcing} (o \emph{$\mathbb{P}$-nombres}), que son objetos en $M$.

  \item \textbf{La relación de forcing $\Vdash$:} La relación $p \Vdash \varphi$ se lee ``la condición $p$ \emph{fuerza} el enunciado $\varphi$''. Esta relación es definible en $M$ y permite controlar qué enunciados son verdaderos en $M[G]$: $M[G] \models \varphi$ si y solo si existe $p \in G$ tal que $p \Vdash \varphi$.
\end{enumerate}

\begin{importante}
La magia del forcing consiste en que, a pesar de que $G$ no pertenece al modelo base $M$, la relación de forcing $\Vdash$ es \emph{definible en $M$}. Esto permite que los matemáticos que trabajan \emph{dentro de $M$} puedan razonar sobre propiedades del modelo extendido $M[G]$ \emph{antes} de construirlo.
\end{importante}

\subsection{La noción de forcing de Cohen para refutar HC}

Cohen diseñó una noción de forcing específica para producir ``muchos'' subconjuntos nuevos de $\omega$ en el modelo extendido, forzando así que $2^{\aleph_0}$ sea grande.

\begin{definicion}{Forcing de Cohen para $2^{\aleph_0} = \aleph_2$}{forcingCohen}
Sea $M$ un modelo numerable y transitivo de ZFC (con $M \models \mathrm{HC}$, aunque esto no es necesario). El \textbf{forcing de Cohen} para añadir $\aleph_2^M$ muchos subconjuntos reales usa la noción de forcing:
\[
  \mathbb{P} = \mathrm{Fn}(\omega_2^M \times \omega,\, 2) = \bigl\{\, p : s \to \{0,1\} \;\bigm|\; s \subseteq \omega_2^M \times \omega,\; |s| < \omega \,\bigr\},
\]
es decir, $P$ consiste en \textbf{funciones parciales finitas} de $\omega_2^M \times \omega$ en $\{0,1\}$, ordenadas por la extensión como función:
\[
  p \leq q \;\iff\; p \supseteq q \;\text{ (como funciones)}.
\]
(La condición más fuerte es la que tiene más información, es decir, el dominio más grande.)
\end{definicion}

\begin{observacion}{Intuición de las condiciones de Cohen}{intuicionCohen}
Cada condición $p \in \mathbb{P}$ puede verse como \emph{información parcial} sobre $\aleph_2$ muchos subconjuntos de $\omega$. Para cada índice $\xi < \omega_2^M$, la condición $p$ especifica, para ciertos $n \in \omega$, si $n$ pertenece o no al $\xi$-ésimo subconjunto real que se está añadiendo.

Cuando se forma el genérico $G = \bigcup \{p : p \in G_{\text{filtro}}\}$, se obtiene una función total $G : \omega_2^M \times \omega \to \{0,1\}$, que codifica $\aleph_2^M$ funciones $G_\xi : \omega \to \{0,1\}$ (es decir, $\aleph_2^M$ subconjuntos de $\omega$).
\end{observacion}

\subsection{Propiedades del forcing de Cohen}

Los dos hechos clave que Cohen necesitó verificar son:

\begin{teorema}{ZFC en el modelo genérico}{zcfgenerico}
El modelo genérico $M[G]$ satisface todos los axiomas de ZFC.
\end{teorema}

La demostración requiere verificar cada axioma. El más delicado es el \textbf{axioma de potencia}: en $M[G]$ el conjunto $\mathcal{P}(\omega)^{M[G]}$ debe ser un conjunto. Esto se sigue del hecho de que cada subconjunto de $\omega$ en $M[G]$ tiene un nombre en $M$, y los nombres forman un conjunto en $M$.

\begin{teorema}{Preservación de cardinales: la antichain condition}{cardpreserv}
Si $\mathbb{P} = \mathrm{Fn}(\omega_2 \times \omega, 2)$, entonces $\mathbb{P}$ satisface la \textbf{condición de cadena numerable} (c.c.n.): toda familia de condiciones mutuamente incompatibles es numerable. Formalmente: si $A \subseteq P$ es tal que $p \not\leq q$ y $q \not\leq p$ (son incompatibles) para cualesquiera $p \neq q$ en $A$, entonces $|A| \leq \aleph_0$.

Como consecuencia, si $\kappa$ es un cardinal en $M$, entonces $\kappa$ sigue siendo un cardinal en $M[G]$. En particular, $\aleph_1^M = \aleph_1^{M[G]}$ y $\aleph_2^M = \aleph_2^{M[G]}$.
\end{teorema}

\begin{proof}[Idea de la demostración de la c.c.n.]
Dos condiciones $p, q \in \mathrm{Fn}(\omega_2 \times \omega, 2)$ son compatibles si y solo si su unión $p \cup q$ es función (es decir, no difieren en ningún par $(\xi, n)$ que esté en el dominio de ambas). Ahora, cada $p$ tiene dominio finito $\mathrm{dom}(p) \subseteq \omega_2 \times \omega$. Para una familia antichain (mutuamente incompatible), cada dos condiciones deben ``contradecirse'' en al menos un punto. Un argumento de deltas-sistemas (\emph{delta-system lemma}) muestra que no puede haber $\aleph_1$ muchas condiciones mutuamente incompatibles.
\end{proof}

\begin{teorema}{Cohen, 1963: Consistencia de $\neg$HC}{cohenNHC}
Si ZFC es consistente, entonces ZFC $+$ $\neg$HC es consistente. Más precisamente:
\[
  \mathrm{Con}(\mathrm{ZFC}) \;\Rightarrow\; \mathrm{Con}(\mathrm{ZFC} + 2^{\aleph_0} = \aleph_2).
\]
En el modelo genérico $M[G]$ construido con el forcing de Cohen $\mathrm{Fn}(\omega_2 \times \omega, 2)$:
\[
  M[G] \models 2^{\aleph_0} = \aleph_2^M = \aleph_2^{M[G]}.
\]
Por tanto $M[G] \models \neg\mathrm{HC}$, y en consecuencia HC no puede ser demostrada en ZFC.
\end{teorema}

\begin{proof}[Esquema de la demostración]
\begin{enumerate}[leftmargin=2em, label=(\roman*)]
  \item \textbf{$M[G]$ satisface ZFC:} ya mencionado.
  \item \textbf{Los cardinales se preservan:} por la c.c.n. del forcing, $\aleph_1^{M[G]} = \aleph_1^M$ y $\aleph_2^{M[G]} = \aleph_2^M$.
  \item \textbf{$M[G]$ tiene $\aleph_2^M$ subconjuntos de $\omega$:} el genérico $G$ codifica, para cada $\xi < \omega_2^M$, la función característica del subconjunto $G_\xi = \{n \in \omega : G(\xi, n) = 1\}$. Estos $\omega_2^M$ subconjuntos son todos distintos (por la genericidad de $G$, cualquier par de índices distintos ``diverge'' en algún natural).
  \item \textbf{No hay más subconjuntos de $\omega$ en $M[G]$:} cada subconjunto de $\omega$ en $M[G]$ es ``nombrado'' por algún $\mathbb{P}$-nombre en $M$. La cardinalidad del conjunto de nombres en $M$ es $\aleph_2^M$ (pues cada nombre es una función de $\omega_2^M \times \omega$ a un conjunto en $M$, y hay $\aleph_2^M$ muchos tales), de modo que $|\mathcal{P}(\omega)^{M[G]}| = \aleph_2^{M[G]}$.
  \item \textbf{Conclusión:} $M[G] \models 2^{\aleph_0} = \aleph_2$, luego $M[G] \models \neg\mathrm{HC}$.
\end{enumerate}
\end{proof}

\subsection{El forcing puede producir cualquier valor de \texorpdfstring{$2^{\aleph_0}$}{el continuo} (con restricciones)}

El resultado de Cohen no se limita a $2^{\aleph_0} = \aleph_2$. Usando $\mathrm{Fn}(\omega_\alpha \times \omega, 2)$ con distintos valores de $\alpha$, se puede forzar $2^{\aleph_0} = \aleph_\alpha$ para cualquier cardinal $\aleph_\alpha$ con $\mathrm{cf}(\aleph_\alpha) > \aleph_0$.

\begin{teorema}{Easton, 1970: Arbitrariedad de la potenciación cardinal}{eastonforcing}
Bajo los axiomas de ZFC, la función $\alpha \mapsto 2^{\aleph_\alpha}$ puede tomar casi cualquier valor: para cualquier función $F : \mathrm{Card} \to \mathrm{Card}$ definida en los cardinales infinitos regulares tal que
\begin{enumerate}[leftmargin=2em, label=(\roman*)]
  \item $\alpha < \beta \Rightarrow F(\alpha) \leq F(\beta)$ (monotonicidad),
  \item $\mathrm{cf}(F(\alpha)) > \alpha$ para todo cardinal regular $\alpha$ (restricción de König),
\end{enumerate}
existe un modelo de ZFC donde $2^{\aleph_\alpha} = F(\aleph_\alpha)$ para todos los cardinales regulares $\aleph_\alpha$.
\end{teorema}

\begin{observacion}{Lo que Easton no cubre}{eastonnocovers}
El teorema de Easton controla la potenciación en cardinales \emph{regulares}, pero el comportamiento de $2^\kappa$ para cardinales singulares $\kappa$ no está completamente determinado. Saharon Shelah desarrolló la \textbf{aritmética cardinal pcf} (``possible cofinalities'') para entender mejor las restricciones en el caso singular. Esto pertenece a la investigación avanzada del siglo~XX y XXI en teoría de conjuntos.
\end{observacion}

% ===========================================================
\section{Independencia de HC respecto de ZFC}
% ===========================================================

\subsection{El teorema de independencia}

Los dos resultados de Gödel y Cohen juntos establecen la independencia completa de HC:

\begin{teorema}{Independencia de HC respecto de ZFC}{independenciaHC}
La Hipótesis del Continuo es \textbf{independiente} de ZFC:
\begin{enumerate}[leftmargin=2em, label=(\roman*)]
  \item ZFC $\not\vdash$ HC \quad (HC no es demostrable en ZFC).
  \item ZFC $\not\vdash$ $\neg$HC \quad ($\neg$HC no es demostrable en ZFC).
\end{enumerate}
Equivalentemente: si ZFC es consistente, entonces ZFC $+$ HC y ZFC $+$ $\neg$HC son ambos consistentes.
\end{teorema}

\begin{proof}
(i) sigue del teorema de Cohen (si ZFC $\vdash$ HC, entonces ZFC $+$ $\neg$HC sería inconsistente, contradiciendo la consistencia de ZFC $+$ $\neg$HC que Cohen demostró). (ii) sigue del teorema de Gödel (análogamente).
\end{proof}

\begin{observacion}{¿Qué significa la independencia?}{significaindependencia}
La independencia de HC respecto de ZFC tiene las siguientes consecuencias filosóficas y matemáticas:
\begin{enumerate}[leftmargin=2em, label=(\alph*)]
  \item \textbf{ZFC no determina $\mathfrak{c}$:} Los axiomas estándar de la teoría de conjuntos son insuficientes para decidir cuántos puntos tiene la recta real. Hay modelos de ZFC donde $\mathfrak{c} = \aleph_1$, modelos donde $\mathfrak{c} = \aleph_2$, modelos donde $\mathfrak{c} = \aleph_{\omega_1}$, etc.
  \item \textbf{No es un defecto de ZFC:} La independencia no significa que la pregunta no tenga sentido o que ZFC sea ``incompleto en un sentido patológico''. Por los teoremas de incompletitud de Gödel, cualquier sistema axiomático suficientemente expresivo tendrá enunciados independientes.
  \item \textbf{¿Hay axiomas adicionales naturales?} Algunos matemáticos han propuesto enriquecer ZFC con axiomas adicionales (axiomas de grandes cardinales, el axioma $\mathbf{MM}$ de Martin's Maximum, el axioma $\mathbf{PFA}$ de la propiedad de Proper Forcing Axiom) que sí determinan el valor de $\mathfrak{c}$. Varios de estos axiomas implican $\mathfrak{c} = \aleph_2$.
  \item \textbf{El punto de vista pluralista:} Algunos filósofos de la matemática (como Joel Hamkins) defienden un enfoque ``pluralista'' o ``multiverso'': no hay \emph{un} universo matemático, sino muchos, y la pregunta ``¿cuánto vale $\mathfrak{c}$?'' es tan sensata como la pregunta ``¿vale el quinto postulado?'' en geometría.
\end{enumerate}
\end{observacion}

\subsection{Independencia de HCG}

\begin{corolario}{Independencia de HCG}{independenciaHCG}
La Hipótesis Generalizada del Continuo es también independiente de ZFC: si ZFC es consistente, entonces ZFC $+$ HCG y ZFC $+$ $\neg$HCG son ambos consistentes.
\end{corolario}

\begin{proof}
La consistencia de ZFC $+$ HCG sigue del hecho de que $L \models \mathrm{HCG}$ (el universo constructible satisface HCG; la demostración generaliza el argumento del Lema de Condensación a todos los cardinales). La consistencia de ZFC $+$ $\neg$HCG sigue del forcing de Cohen (pues en $M[G]$ vale $2^{\aleph_0} = \aleph_2 \neq \aleph_1$, luego $\neg$HCG).
\end{proof}

% ===========================================================
\section{El legado histórico y matemático}
% ===========================================================

\subsection{La medalla Fields de Paul Cohen}

Paul Cohen recibió la Medalla Fields en 1966 por la invención del forcing y la prueba de la independencia de HC. Su trabajo se considera uno de los más originales y profundos de las matemáticas del siglo~XX: no solo resolvió el primer problema de Hilbert, sino que creó una herramienta (el forcing) que transformó toda la teoría de conjuntos.

\subsection{La proliferación del forcing}

Desde 1963, el forcing ha sido usado para demostrar la independencia de decenas (o cientos) de problemas matemáticos respecto de ZFC:
\begin{itemize}[leftmargin=2em]
  \item La \textbf{hipótesis de Suslin}: ¿toda recta lineal completa sin elementos máximo ni mínimo, separable y tal que toda familia de intervalos disjuntos es numerable, es isomorfa a $\mathbb{R}$? Independiente de ZFC.
  \item El \textbf{problema de Whitehead}: ¿todo grupo abeliano $G$ con $\mathrm{Ext}^1(G, \mathbb{Z}) = 0$ es libre? Independiente de ZFC (resultado de Shelah, 1974).
  \item El \textbf{axioma de Martin} ($\mathrm{MA}$): consistente con ZFC $+$ $\neg$HC. Implica muchos resultados de combinatoria infinita.
\end{itemize}

\subsection{¿Puede HC resolverse con nuevos axiomas?}

Esta es una de las grandes preguntas de la lógica matemática contemporánea. Algunos avances recientes son:

\begin{itemize}[leftmargin=2em]
  \item \textbf{Woodin y el programa $\Omega$:} Hugh Woodin ha desarrollado, desde la década de 1990, un programa para encontrar el ``valor correcto'' de $\mathfrak{c}$. Su programa sugiere que $\mathfrak{c} = \aleph_2$ es el valor más natural consistente con las consecuencias de los grandes cardinales.
  \item \textbf{Axiomas de forcing como PFA y MM:} El \emph{Proper Forcing Axiom} (PFA) y el \emph{Martin's Maximum} (MM) implican $\mathfrak{c} = \aleph_2$. Estos axiomas son consistentes con ZFC (relativa a grandes cardinales) y tienen muchas consecuencias matemáticas atractivas.
  \item \textbf{El programa de los grandes cardinales:} Los axiomas de grandes cardinales (measurable, supercompact, huge, etc.) no determinan por sí solos el valor de $\mathfrak{c}$, pero sí imponen restricciones indirectas a través de los resultados de absolutez.
\end{itemize}

% ===========================================================
\section{Esquema comparativo: $V$, $L$ y los modelos de forcing}
% ===========================================================

Para ayudar a la comprensión global, presentamos un esquema de los distintos ``universos'' de conjuntos relevantes para la independencia de HC:

\begin{ejemplo}{Tres universos, tres valores de $\mathfrak{c}$}{}
\begin{enumerate}[leftmargin=2em, label=(\alph*)]
  \item \textbf{El universo constructible $L$:} $L \models \mathrm{ZFC} + V = L + \mathrm{HCG}$. En $L$, $2^{\aleph_0} = \aleph_1$, $2^{\aleph_1} = \aleph_2$, etc. Todo es ``lo más pequeño posible''. Los subconjuntos de $\omega$ son exactamente los constructibles: hay $\aleph_1$ de ellos.
  \item \textbf{El modelo de Cohen $M[G]$ con $\mathrm{Fn}(\omega_2 \times \omega, 2)$:} $M[G] \models \mathrm{ZFC} + 2^{\aleph_0} = \aleph_2$. El continuo da un salto de tamaño $\aleph_2$: hay $\aleph_2$ subconjuntos de $\omega$. Los cardinales $\aleph_1$ y $\aleph_2$ son los mismos que en $M$.
  \item \textbf{Un modelo con $\mathfrak{c}$ grande:} Usando $\mathrm{Fn}(\omega_{\omega_1} \times \omega, 2)$ se puede forzar $2^{\aleph_0} = \aleph_{\omega_1}$. Aquí el continuo es un cardinal de cofinalidad $\omega_1 > \omega$, lo cual es consistente con el teorema de König ($\mathrm{cf}(\aleph_{\omega_1}) = \omega_1 > \omega$).
\end{enumerate}
\end{ejemplo}

% ===========================================================
\section{Resumen de los resultados principales}
% ===========================================================

Recapitulamos los resultados principales de esta clase en un cuadro de referencia:

\begin{importante}
\textbf{Resumen: la independencia de HC}

\begin{itemize}[leftmargin=2em]
  \item \textbf{HC}: $2^{\aleph_0} = \aleph_1$. \quad \textbf{HCG}: $\forall \alpha\; (2^{\aleph_\alpha} = \aleph_{\alpha+1})$.
  \item \textbf{Gödel, 1940:} El universo constructible $L$ satisface ZFC $+$ HCG. Luego ZFC $\not\vdash \neg$HC.
  \item \textbf{Cohen, 1963:} Mediante forcing, existe un modelo de ZFC donde $2^{\aleph_0} = \aleph_2$. Luego ZFC $\not\vdash$ HC.
  \item \textbf{Conclusión:} HC es \emph{independiente} de ZFC.
  \item \textbf{Easton, 1970:} La función $\kappa \mapsto 2^\kappa$ en cardinales regulares puede ser cualquier función monótona con $\mathrm{cf}(2^\kappa) > \kappa$.
  \item \textbf{König (restricción):} $\mathrm{cf}(2^{\aleph_0}) > \aleph_0$, es decir $2^{\aleph_0} \neq \aleph_\omega$, $\aleph_{\omega^2}$, etc.
\end{itemize}
\end{importante}

% ===========================================================
\section{Problemas y tareas}
% ===========================================================

\subsection*{Problemas de comprensión y cálculo}

\begin{enumerate}[leftmargin=2em, label=\textbf{P\arabic*.}]

  \item \textbf{(Formulaciones equivalentes de HC.)} Demuestre que las siguientes afirmaciones son equivalentes en ZFC:
  \begin{enumerate}[label=(\alph*)]
    \item $2^{\aleph_0} = \aleph_1$.
    \item Cada subconjunto infinito de $\mathbb{R}$ es numerable o tiene la cardinalidad de $\mathbb{R}$.
    \item $|\mathcal{P}(\mathbb{N})| = \aleph_1$.
    \item Todo conjunto $A$ con $\aleph_0 \leq |A| \leq 2^{\aleph_0}$ satisface $|A| = \aleph_0$ o $|A| = 2^{\aleph_0}$.
  \end{enumerate}

  \item \textbf{(Consecuencias de HCG.)} Asumiendo HCG, calcule los siguientes cardinales:
  \begin{enumerate}[label=(\alph*)]
    \item $2^{\aleph_3}$.
    \item $\aleph_2^{\aleph_1}$.
    \item $\aleph_\omega^{\aleph_0}$.
    \item $\aleph_{\omega+1}^{\aleph_\omega}$.
  \end{enumerate}
  \textit{(Sugerencia: use la ley de absorción $\kappa^\lambda = 2^\kappa$ cuando $\lambda < \kappa$ y $\mathrm{cf}(\kappa) > \lambda$, combinada con HCG.)}

  \item \textbf{(Restricciones sobre $\mathfrak{c}$ en ZFC.)} Demuestre en ZFC (sin asumir HC) que:
  \begin{enumerate}[label=(\alph*)]
    \item $\mathfrak{c} \geq \aleph_1$.
    \item $\mathrm{cf}(\mathfrak{c}) > \aleph_0$.
    \item $\mathfrak{c}$ no puede ser $\aleph_\omega$.
    \item $\mathfrak{c}$ no puede ser $\aleph_{\omega_1}$ si $\mathrm{cf}(\aleph_{\omega_1}) = \omega_1 > \omega$ (¿es esto una restricción real?).
  \end{enumerate}

  \item \textbf{(La jerarquía $L$ en los primeros niveles.)} Describa explícitamente los conjuntos que pertenecen a $L_4$ pero no a $L_3$. ¿Cuántos elementos tiene $L_4$? ¿Qué elementos de $V_4$ \emph{no} pertenecen a $L_4$? (Compare $|L_4|$ con $|V_4|$.)

  \item \textbf{(Cardinalidad de $L_\alpha$ para ordinales límite.)} Demuestre que para todo cardinal infinito regular $\kappa$, se tiene $|L_\kappa| = \kappa$. \textit{(Sugerencia: use inducción transfinita y el hecho de que $|\mathrm{Def}(M)| \leq \max(|M|, \aleph_0)$ para conjuntos infinitos $M$.)}

  \item \textbf{(Subconjuntos no constructibles.)} Dado que $|L_{\omega+1}| = \aleph_0$ mientras que $|\mathcal{P}(\omega)| = 2^{\aleph_0}$:
  \begin{enumerate}[label=(\alph*)]
    \item Argumente que en ZFC $+$ $\neg$HC (con $\mathfrak{c} = \aleph_2$, por ejemplo), la mayoría de los subconjuntos de $\omega$ no son constructibles.
    \item ¿Es posible, sin asumir HC, que \emph{todos} los subconjuntos de $\omega$ sean constructibles? ¿Qué axioma afirmaría eso?
    \item Bajo $V = L$, ¿cuántos subconjuntos de $\omega$ existen? Justifique su respuesta.
  \end{enumerate}

  \item \textbf{(Forcing elemental.)} Considere la noción de forcing $\mathbb{P} = \mathrm{Fn}(\omega, 2)$, es decir, las funciones parciales finitas de $\omega$ en $\{0,1\}$, con $p \leq q$ si $p \supseteq q$.
  \begin{enumerate}[label=(\alph*)]
    \item Describa qué tipo de objeto ``añade'' un filtro genérico $G$ sobre $\mathbb{P}$: ¿qué es $f_G = \bigcup G$?
    \item Muestre que $\mathrm{Fn}(\omega, 2)$ satisface la c.c.n.\ (condición de cadena numerable).
    \item Explique por qué este forcing añade exactamente \emph{un} subconjunto nuevo de $\omega$, y no cambia los cardinales.
    \item ¿Por qué este forcing, a diferencia del forcing de Cohen con $\mathrm{Fn}(\omega_2 \times \omega, 2)$, no puede refutar HC?
  \end{enumerate}

  \item \textbf{(El teorema de König y $\mathfrak{c}$.)} El teorema de König establece que $\mathrm{cf}(2^{\aleph_0}) > \aleph_0$.
  \begin{enumerate}[label=(\alph*)]
    \item Use el teorema de König para demostrar que $2^{\aleph_0} \neq \aleph_\omega$.
    \item Demuestre que $2^{\aleph_0} \neq \aleph_{\omega \cdot 2}$.
    \item En general, $2^{\aleph_0}$ no puede ser ningún cardinal de cofinalidad $\omega$. Liste cinco cardinales que por este motivo no pueden ser el valor de $\mathfrak{c}$.
    \item ¿Puede ser $\mathfrak{c} = \aleph_{\omega_1}$? Justifique su respuesta.
  \end{enumerate}

  \item \textbf{(Independencia y lógica.)} Explique, en sus propias palabras, la diferencia entre los siguientes conceptos:
  \begin{enumerate}[label=(\alph*)]
    \item ``HC es verdadera'' vs. ``HC es demostrable en ZFC''.
    \item ``HC es independiente de ZFC'' vs. ``HC no tiene valor de verdad''.
    \item ``ZFC es consistente'' vs. ``ZFC es completo''.
  \end{enumerate}
  Conecte su explicación con los teoremas de incompletitud de Gödel (clase~1).

  \item \textbf{(HCG y la aritmética cardinal completa.)} Suponiendo HCG, demuestre que para todo ordinal $\alpha$ y todo cardinal $\lambda < \aleph_\alpha$:
  \[
    \aleph_\alpha^\lambda = \aleph_\alpha.
  \]
  \textit{(Sugerencia: use que $\aleph_\alpha^\lambda \leq \aleph_\alpha^{|\alpha|} = (2^{\aleph_{|\alpha|}})^{|\alpha|} = 2^{\aleph_{|\alpha|} \cdot |\alpha|} = 2^{\aleph_{|\alpha|}}$ y aplique HCG.)}

\end{enumerate}

\subsection*{Problemas de profundización}

\begin{enumerate}[leftmargin=2em, label=\textbf{P\arabic*$'$.}]

  \item \textbf{(El lema de $\Delta$-sistemas.)} Un \textbf{sistema $\Delta$} (o \textbf{girasol}) es una familia $\{A_i : i \in I\}$ de conjuntos finitos tal que $A_i \cap A_j = r$ para todo $i \neq j$ (el ``núcleo'' $r$ es el mismo para todos los pares). Enuncie y demuestre el \textbf{lema del sistema $\Delta$}: toda familia no numerable de conjuntos finitos contiene una subfamilia no numerable que forma un sistema $\Delta$. Use este lema para demostrar la c.c.n.\ del forcing de Cohen $\mathrm{Fn}(\kappa \times \omega, 2)$ para cualquier cardinal $\kappa$.

  \item \textbf{(El axioma de Martin.)} El axioma de Martin ($\mathrm{MA}$) afirma que para todo orden parcial $\mathbb{P}$ con la c.c.n.\ y toda familia $\{D_i : i < \kappa\}$ de conjuntos densos en $\mathbb{P}$ con $\kappa < \mathfrak{c}$, existe un filtro $G$ que intersecta a todos los $D_i$. Investigue:
  \begin{enumerate}[label=(\alph*)]
    \item $\mathrm{MA}$ es consistente con ZFC $+$ $\neg$HC (de hecho, $\mathrm{MA}$ implica $\mathrm{cf}(\mathfrak{c}) > \aleph_0$ pero no determina el valor de $\mathfrak{c}$).
    \item $\mathrm{MA} + \neg\mathrm{HC}$ implica que toda unión de $\aleph_1$ muchos conjuntos de medida de Lebesgue cero tiene medida cero.
    \item ¿$\mathrm{MA}$ implica o es implicado por HC?
  \end{enumerate}

  \item \textbf{(El programa de Woodin.)} Investigar el enunciado: ``el axioma $(\ast)$ de Woodin implica que el valor de $2^{\aleph_0}$ es $\aleph_2$''. Relacione este resultado con el programa de grandes cardinales. ¿Cuáles son los grandes cardinales más usados en la investigación de HC?

  \item \textbf{(Independencia de la hipótesis de Suslin.)} La hipótesis de Suslin ($\mathrm{HS}$) afirma: toda recta lineal completa, sin extremos, separable (en sentido de Suslin: toda familia de intervalos disjuntos abiertos es numerable) es isomorfa a $\mathbb{R}$. Investigue: (a) ¿$\mathrm{HS}$ implica HC? (b) ¿HC implica $\mathrm{HS}$? (c) ¿$\mathrm{HS}$ es independiente de ZFC?

\end{enumerate}

% ===========================================================
\section*{Referencias de la Clase 9}
% ===========================================================

Los resultados de esta clase provienen de las siguientes fuentes fundamentales:

\begin{itemize}[leftmargin=2em]
  \item El resultado de Gödel sobre la consistencia de HC con ZFC fue publicado en \cite{godel1940}. Una presentación moderna detallada se encuentra en \cite{jech2003} (capítulos~13 y 14) y en \cite{kunen2011} (capítulo~VI).
  \item El forcing de Cohen fue publicado en \cite{cohen1963a, cohen1963b} y desarrollado en el libro \cite{cohen1966}. Una presentación moderna accesible está en \cite{kunen2011} (capítulo~VII) y en \cite{jech2003} (capítulo~14).
  \item El teorema de Easton sobre la independencia de la función $\kappa \mapsto 2^\kappa$ en cardinales regulares se encuentra en \cite{easton1970}.
  \item Exposiciones introductorias del forcing se encuentran en \cite{devlin1993} y en \cite{hrbacek1999}. Una introducción muy accesible en español es \cite{cori2001}.
  \item Para el programa de grandes cardinales y las perspectivas actuales sobre HC, ver \cite{kanamori2003} y \cite{woodin2001}.
\end{itemize}

% ===========================================================
\chapter{Clase 10: Horizontes --- Cardinales Grandes, Determinismo y Límites de ZFC}
% ===========================================================

Esta sesión de cierre sitúa los números transfinitos en el panorama más amplio de la investigación matemática contemporánea. Las clases anteriores construyeron el edificio axiomático de ZFC, desarrollaron la aritmética ordinal y cardinal, y culminaron en la independencia de la Hipótesis del Continuo. Ahora ascendemos hasta los horizontes actuales de la teoría de conjuntos.

Tres grandes arcos temáticos estructuran esta clase. En primer lugar, los \textbf{cardinales grandes}: extensiones axiomáticas de ZFC que postulan la existencia de cardinales de una grandeza tan extrema que su mera existencia es indemostrable dentro del sistema base. En segundo lugar, la \textbf{jerarquía de consistencia}: una escala lineal que ordena las extensiones de ZFC según su fuerza deductiva. En tercer lugar, el \textbf{Axioma de Determinismo} ($\mathrm{AD}$) y su radical incompatibilidad con el Axioma de Elección, junto con las consecuencias sorprendentes que $\mathrm{AD}$ tiene sobre los subconjuntos de la recta real. Concluimos con una reflexión sobre los límites epistemológicos del formalismo conjuntista y con un panorama de las líneas de investigación abiertas.

La exposición sigue los tratamientos canónicos de \cite{kanamori2003, jech2003, kunen2011}, y para los aspectos relativos al determinismo, \cite{moschovakis2009, woodin2010det}.

% ===========================================================
\section{Cardinales inaccesibles}
% ===========================================================

\subsection{Motivación: ¿puede ZFC hablar de sus propios modelos?}

La jerarquía de von Neumann $V = \bigcup_{\alpha \in \mathrm{Ord}} V_\alpha$ constituye el universo estándar de la teoría de conjuntos. En la Clase 5 vimos que $V_\omega \models \mathrm{ZFC} \setminus \{\mathrm{Inf}\}$: el estrato finito satisface todos los axiomas de ZFC salvo el de Infinitud. Es natural preguntarse: ¿existe algún estrato $V_\kappa$ que satisfaga \emph{todos} los axiomas de ZFC, incluidos los esquemas de separación y reemplazamiento?

La respuesta conduce a la noción de cardinal inaccesible.

\subsection{Cardinales límite regulares y la función beth}

\begin{definicion}{Cardinal regular}{regular}
Un cardinal infinito $\kappa$ es \textbf{regular} si $\mathrm{cf}(\kappa) = \kappa$; es decir, si $\kappa$ no puede escribirse como unión de menos de $\kappa$ conjuntos de cardinalidad menor que $\kappa$. En caso contrario, $\kappa$ es \textbf{singular}.
\end{definicion}

\begin{ejemplo}{Cardinales regulares y singulares}{}
\begin{enumerate}[leftmargin=2em, label=(\alph*)]
  \item Todo sucesor cardinal $\aleph_{\alpha+1}$ es regular: si $\aleph_{\alpha+1} = \bigcup_{i < \lambda} A_i$ con $\lambda < \aleph_{\alpha+1}$ y $|A_i| < \aleph_{\alpha+1}$, entonces $\lambda \leq \aleph_\alpha$ y $|A_i| \leq \aleph_\alpha$, así que $|\bigcup A_i| \leq \aleph_\alpha \cdot \aleph_\alpha = \aleph_\alpha < \aleph_{\alpha+1}$, contradicción.
  \item $\aleph_0 = \omega$ es regular: ningún conjunto finito de conjuntos finitos cubre a $\omega$.
  \item $\aleph_\omega = \sup_{n < \omega} \aleph_n$ es singular de cofinalidad $\omega$: es la reunión de los $\aleph_n$ con $n \in \omega$.
  \item En general, $\aleph_\lambda$ para $\lambda$ límite de cofinalidad numerable es singular.
\end{enumerate}
\end{ejemplo}

\begin{definicion}{Función beth}{beth}
La \textbf{función beth} $\beth$ se define por inducción transfinita:
\[
  \beth_0 = \aleph_0, \qquad \beth_{\alpha+1} = 2^{\beth_\alpha}, \qquad \beth_\lambda = \sup_{\beta < \lambda} \beth_\beta \;\text{ para } \lambda \text{ límite.}
\]
Un cardinal $\kappa$ es \textbf{fuerte límite} si $2^\lambda < \kappa$ para todo $\lambda < \kappa$, es decir, si $\kappa = \beth_\kappa$ (punto fijo de la función beth).
\end{definicion}

\begin{observacion}{Diferencia entre límite y fuerte límite}{lim_slim}
Un cardinal $\kappa$ es \textbf{límite} si $\kappa = \sup_{\lambda < \kappa} \lambda$, es decir, no es sucesor. Ser límite es más débil que ser fuerte límite: $\aleph_\omega$ es límite pero no fuerte límite, pues si la HCG falla puede ocurrir $2^{\aleph_0} \geq \aleph_\omega$.
\end{observacion}

\subsection{Cardinales débilmente y fuertemente inaccesibles}

\begin{definicion}{Cardinales inaccesibles}{inacc}
Sea $\kappa$ un cardinal no numerable.
\begin{enumerate}[leftmargin=2em, label=(\alph*)]
  \item $\kappa$ es \textbf{débilmente inaccesible} si es regular y límite (es decir, $\kappa = \aleph_\kappa$).
  \item $\kappa$ es \textbf{fuertemente inaccesible} (o simplemente \textbf{inaccesible}) si es regular y fuerte límite; equivalentemente, si $\kappa$ es regular y $2^\lambda < \kappa$ para todo $\lambda < \kappa$.
\end{enumerate}
\end{definicion}

\begin{importante}
Bajo la Hipótesis Generalizada del Continuo (HCG), débilmente inaccesible equivale a fuertemente inaccesible, pues HCG implica $\beth_\alpha = \aleph_\alpha$ para todo $\alpha$. Sin HCG, las dos nociones pueden diferir.
\end{importante}

\begin{teorema}{$V_\kappa$ modela ZFC}{vkappa_zfc}
Si $\kappa$ es un cardinal fuertemente inaccesible, entonces $V_\kappa \models \mathrm{ZFC}$.
\end{teorema}

\begin{proof}
Verificamos cada axioma de ZFC en $V_\kappa$:
\begin{itemize}[leftmargin=2em]
  \item \textbf{Extensionalidad, Fundación, Par, Unión, Partes, Infinitud:} Se verifican directamente en cualquier $V_\alpha$ con $\alpha > \omega$.
  \item \textbf{Separación:} Si $a \in V_\kappa$ y $\varphi$ es una fórmula, entonces $\{x \in a : V_\kappa \models \varphi(x)\} \subseteq a \in V_\kappa$, así que pertenece a $V_\kappa$ (pues $V_\kappa$ es transitivo).
  \item \textbf{Reemplazamiento:} Supongamos que $F: a \to V_\kappa$ es una función definible en $V_\kappa$ con $a \in V_\kappa$. Sea $\alpha = \mathrm{rango}(a)$, de modo que $a \subseteq V_\alpha$ con $|\alpha| < \kappa$. Cada $F(x)$ pertenece a algún $V_{\beta_x}$ con $\beta_x < \kappa$. La colección de los $\beta_x$ es un conjunto de cardinalidad $|a| \leq |V_\alpha| < \kappa$. Como $\kappa$ es regular, $\beta = \sup_{x \in a} \beta_x < \kappa$, y por tanto $F[a] \subseteq V_\beta \in V_\kappa$.
  \item \textbf{Elección:} Se hereda porque $V_\kappa \subseteq V$ y en $V$ existe un buen orden del universo.
\end{itemize}
\end{proof}

\begin{corolario}{Inaccesibles y consistencia de ZFC}{inacc_cons}
Si existe un cardinal fuertemente inaccesible $\kappa$, entonces $\mathrm{Con}(\mathrm{ZFC})$. En particular, la existencia de un cardinal inaccesible no es demostrable en ZFC (asumiendo que ZFC es consistente), pues por el segundo teorema de incompletitud de Gödel, ZFC no puede demostrar $\mathrm{Con}(\mathrm{ZFC})$.
\end{corolario}

\begin{ejemplo}{Puntos fijos y la jerarquía de inaccesibles}{}
Sea $I$ la clase de todos los cardinales fuertemente inaccesibles. Si $\kappa$ es inaccesible, entonces $\kappa = \aleph_\kappa = \beth_\kappa$: es un punto fijo de la función aleph \emph{y} de la función beth simultáneamente. 

Ilustración con los primeros puntos fijos de la función aleph: los cardinales de la forma $\aleph_\kappa = \kappa$ son exactamente los cardinales límite regulares. La primera solución de $\aleph_\kappa = \kappa$ se denota $\aleph_{\aleph_{\aleph_{\cdots}}}$ (límite de la torre $\aleph_0, \aleph_{\aleph_0}, \aleph_{\aleph_{\aleph_0}}, \ldots$), que es de cofinalidad $\omega$, luego \emph{singular}. El primer cardinal débilmente inaccesible, si existe, requiere un ``salto'' que ZFC no puede garantizar.
\end{ejemplo}

% ===========================================================
\section{Cardinales de Mahlo}
% ===========================================================

Los cardinales de Mahlo son, por así decir, inaccesibles de segundo orden: no solo son ellos mismos inaccesibles, sino que poseen una abundancia ``densa'' de cardinales inaccesibles por debajo de ellos.

\subsection{Conjuntos estacionarios y clubs}

Para definir los cardinales de Mahlo necesitamos el concepto de conjunto estacionario en un cardinal regular.

\begin{definicion}{Club y conjunto estacionario}{club_stat}
Sea $\kappa$ un cardinal regular no numerable.
\begin{enumerate}[leftmargin=2em, label=(\alph*)]
  \item Un conjunto $C \subseteq \kappa$ es \textbf{cerrado sin límite} (\textbf{club}) si:
    \begin{itemize}
      \item Es \emph{no acotado}: $\sup(C) = \kappa$.
      \item Es \emph{cerrado}: para toda sucesión creciente $(\alpha_i)_{i < \lambda}$ de elementos de $C$ con $\lambda < \kappa$ límite, $\sup_i \alpha_i \in C$.
    \end{itemize}
  \item Un conjunto $S \subseteq \kappa$ es \textbf{estacionario} en $\kappa$ si $S \cap C \neq \emptyset$ para todo club $C \subseteq \kappa$.
\end{enumerate}
\end{definicion}

\begin{ejemplo}{Clubs y estacionarios en $\omega_1$}{}
En $\kappa = \omega_1$:
\begin{enumerate}[leftmargin=2em, label=(\alph*)]
  \item $C = \{\alpha < \omega_1 : \alpha \text{ es un ordinal límite}\}$ es un club en $\omega_1$: es cerrado (el supremo de ordinales límite es límite) y no acotado.
  \item $S = \{\alpha < \omega_1 : \mathrm{cf}(\alpha) = \omega\}$ (ordinales de cofinalidad numerable por debajo de $\omega_1$) es estacionario: cualquier club en $\omega_1$ contiene puntos de cofinalidad $\omega$.
  \item El conjunto de ordinales sucesores $\{\alpha + 1 : \alpha < \omega_1\}$ es estacionario pero no es club (no es cerrado: el supremo de una sucesión de sucesores puede ser un límite).
\end{enumerate}
\end{ejemplo}

\begin{definicion}{Cardinal de Mahlo}{mahlo}
Un cardinal fuertemente inaccesible $\kappa$ es un \textbf{cardinal de Mahlo} si el conjunto
\[
  I_\kappa = \{\lambda < \kappa : \lambda \text{ es fuertemente inaccesible}\}
\]
es \textbf{estacionario} en $\kappa$.
\end{definicion}

\begin{observacion}{Interpretación intuitiva}{mahlo_intuit}
Decir que $I_\kappa$ es estacionario significa que los cardinales inaccesibles debajo de $\kappa$ son tan abundantes que no pueden ser evitados por ningún club: están ``diseminados densamente'' en $\kappa$ en el sentido del filtro de clubs. En particular, $\kappa$ no es solo inaccesible: por debajo de $\kappa$ hay una clase propia de cardinales inaccesibles (vistos desde dentro de $V_\kappa$).
\end{observacion}

\begin{teorema}{Fuerza de consistencia de Mahlo}{mahlo_cons}
\begin{enumerate}[leftmargin=2em, label=(\roman*)]
  \item Si $\kappa$ es de Mahlo, entonces $V_\kappa \models \text{``existen clase propias de cardinales inaccesibles''}$.
  \item $\mathrm{Con}(\mathrm{ZFC} + \text{``existe un cardinal de Mahlo''}) \Rightarrow \mathrm{Con}(\mathrm{ZFC} + \text{``existen infinitos cardinales inaccesibles''})$.
  \item La recíproca falla: la existencia de infinitos cardinales inaccesibles es consistente con la no existencia de ningún cardinal de Mahlo.
\end{enumerate}
\end{teorema}

\begin{ejemplo}{La jerarquía de cardinales de Mahlo}{}
Se puede iterar la noción de Mahlo:
\begin{itemize}[leftmargin=2em]
  \item $\kappa$ es \textbf{Mahlo de orden 1} (o simplemente Mahlo) si $I_\kappa$ es estacionario en $\kappa$.
  \item $\kappa$ es \textbf{Mahlo de orden 2} si el conjunto de cardinales Mahlo de orden 1 por debajo de $\kappa$ es estacionario en $\kappa$.
  \item $\kappa$ es \textbf{Mahlo de orden $\alpha$} por inducción sobre $\alpha$.
  \item $\kappa$ es \textbf{fuertemente Mahlo} si es Mahlo de orden $\alpha$ para todo $\alpha < \kappa$.
\end{itemize}
Cada nivel de esta jerarquía es estrictamente más fuerte en consistencia que el anterior.
\end{ejemplo}

% ===========================================================
\section{Cardinales medibles}
% ===========================================================

Los cardinales medibles, introducidos por Stanisław Ulam en 1930 y estudiados a fondo por Scott, Dana y Solovay en las décadas de 1960, representan un salto radical en la jerarquía de cardinales grandes.

\subsection{Ultrafiltros y medidas}

\begin{definicion}{Ultrafiltro $\kappa$-completo}{uf_kappa}
Sea $\kappa$ un cardinal no numerable. Una familia $\mathcal{U} \subseteq \mathcal{P}(\kappa)$ es un \textbf{ultrafiltro $\kappa$-completo no principal} sobre $\kappa$ si:
\begin{enumerate}[leftmargin=2em, label=(\roman*)]
  \item $\emptyset \notin \mathcal{U}$ y $\kappa \in \mathcal{U}$ (no trivial, contiene el total).
  \item $A \in \mathcal{U}$ y $A \subseteq B \subseteq \kappa$ implica $B \in \mathcal{U}$ (upward closed).
  \item $A, B \in \mathcal{U}$ implica $A \cap B \in \mathcal{U}$ (cerrado bajo intersecciones finitas).
  \item Para todo $A \subseteq \kappa$: $A \in \mathcal{U}$ o $\kappa \setminus A \in \mathcal{U}$, pero no ambos (ultrapropiedad).
  \item $\{x\} \notin \mathcal{U}$ para todo $x \in \kappa$ (\textbf{no principal}).
  \item Si $\lambda < \kappa$ y $\{A_i : i < \lambda\} \subseteq \mathcal{U}$, entonces $\bigcap_{i < \lambda} A_i \in \mathcal{U}$ (\textbf{$\kappa$-completitud}).
\end{enumerate}
\end{definicion}

\begin{definicion}{Cardinal medible}{medible}
Un cardinal no numerable $\kappa$ es \textbf{medible} si existe un ultrafiltro $\kappa$-completo no principal sobre $\kappa$.
\end{definicion}

\begin{observacion}{Origen del nombre}{nombre_med}
La terminología proviene de la \emph{teoría de la medida}: un cardinal $\kappa$ es medible si y solo si existe una medida $\mu: \mathcal{P}(\kappa) \to \{0,1\}$ (con valores en $\{0,1\}$, es decir, una medida de dos valores) que es $\kappa$-aditiva, no es idénticamente $0$, y da medida $0$ a todos los singletons. El ultrafiltro correspondiente es $\mathcal{U} = \{A \subseteq \kappa : \mu(A) = 1\}$. Ulam preguntó en 1930 si tal medida podía existir para $\kappa = \omega_1$; la respuesta (que no puede, en ZFC) llevó al estudio de los cardinales medibles.
\end{observacion}

\subsection{Cardinales medibles y ultrapotencias}

La herramienta principal para estudiar cardinales medibles es la construcción de \textbf{ultrapotencias} del universo.

\begin{definicion}{Ultrapotencia de $V$}{ultrapower}
Sea $\kappa$ un cardinal medible con ultrafiltro $\kappa$-completo no principal $\mathcal{U}$ sobre $\kappa$. La \textbf{ultrapotencia} $\mathrm{Ult}(V, \mathcal{U})$ es la estructura definida por:
\begin{itemize}[leftmargin=2em]
  \item El universo es el conjunto de clases de equivalencia de funciones $f: \kappa \to V$ bajo la relación: $f \sim g \iff \{\alpha < \kappa : f(\alpha) = g(\alpha)\} \in \mathcal{U}$.
  \item La relación de pertenencia está dada por: $[f] \in_{\mathcal{U}} [g] \iff \{\alpha < \kappa : f(\alpha) \in g(\alpha)\} \in \mathcal{U}$.
\end{itemize}
El \textbf{\textit{Teorema de Łoś}} garantiza que $\mathrm{Ult}(V, \mathcal{U}) \equiv V$ (son elementalmente equivalentes), y la función $j: V \to \mathrm{Ult}(V, \mathcal{U})$ dada por $j(x) = [c_x]$ (donde $c_x$ es la función constante con valor $x$) es una \textbf{inmersión elemental} no trivial.
\end{definicion}

\begin{teorema}{Cardinales medibles e inmersiones elementales}{med_embed}
Un cardinal $\kappa$ es medible si y solo si existe una inmersión elemental no trivial $j: V \to M$ para algún modelo transitivo $M$ de ZFC tal que $\kappa = \mathrm{crit}(j)$ (el \textbf{punto crítico} de $j$, el menor ordinal movido por $j$: $j(\kappa) > \kappa$ y $j(\alpha) = \alpha$ para todo $\alpha < \kappa$).
\end{teorema}

\begin{ejemplo}{Propiedades del punto crítico}{}
Sea $j: V \to M$ una inmersión elemental con punto crítico $\kappa$. Entonces:
\begin{enumerate}[leftmargin=2em, label=(\alph*)]
  \item $j(\kappa) > \kappa$: el punto crítico se ``mueve hacia arriba''.
  \item $\kappa < j(\kappa) < j(j(\kappa)) < \cdots$: podemos iterar la inmersión para obtener una cadena creciente de cardinales $\kappa_0 = \kappa < \kappa_1 = j(\kappa) < \kappa_2 = j(\kappa_1) < \cdots$
  \item La secuencia $\langle \kappa_n \rangle_{n < \omega}$ es una sucesión de cardinales medibles en $M$.
  \item Toda esta estructura es ``interna'' a $j$ e invisible desde $V$: no contradice la existencia de un único cardinal medible en $V$.
\end{enumerate}
\end{ejemplo}

\subsection{El teorema de Scott: $V \neq L$ si hay un cardinal medible}

El resultado más célebre relacionado con los cardinales medibles es el teorema de Scott de 1961, que establece que los cardinales medibles son incompatibles con el axioma $V = L$ (el universo constructible de Gödel).

\begin{teorema}{Teorema de Scott (1961)}{scott_thm}
Si existe un cardinal medible, entonces $V \neq L$.
\end{teorema}

\begin{proof}[Esquema de la demostración]
Sea $\kappa$ un cardinal medible con inmersión elemental $j: V \to M$. En $L$, todo conjunto es constructible. Si $V = L$, entonces $M = L$ también (pues $M$ es transitivo y contiene todos los conjuntos constructibles). Pero entonces $j: L \to L$ es una inmersión elemental no trivial de $L$ en sí mismo con punto crítico $\kappa$.

Ahora bien, en $L$ el orden $<_L$ (el buen-orden canónico del universo constructible) satisface $j(<_L) = <_L$ (porque $j$ es elemental). Pero $j(\kappa) > \kappa$ y el $\kappa$-ésimo elemento de $L$ en el orden $<_L$ es $\kappa$ mismo, por lo que $j$ desplaza $\kappa$ a $j(\kappa)$: el $j(\kappa)$-ésimo elemento de $j(L) = L$ es $j(\kappa)$. Esta contradicción con la rigidez de $L$ (no existe ninguna inmersión elemental de $L$ en sí mismo con punto crítico por debajo del primer cardinal medible de $L$) muestra que $V \neq L$. La demostración completa usa el hecho de que $L$ admite un parámetro de Skolem definible, y que cualquier inmersión elemental de $L$ en sí mismo debe ser la identidad.
\end{proof}

\begin{corolario}{Implicaciones de la medibilidad}{med_impl}
Si $\kappa$ es un cardinal medible, entonces:
\begin{enumerate}[leftmargin=2em, label=(\roman*)]
  \item $\kappa$ es fuertemente inaccesible.
  \item $\kappa$ es de Mahlo.
  \item $\kappa$ es Mahlo de todo orden $\alpha < \kappa$.
  \item Por debajo de $\kappa$ existen $\kappa$ muchos cardinales de Mahlo.
  \item $V \neq L$.
\end{enumerate}
En particular, $\mathrm{Con}(\mathrm{ZFC} + \text{``existe un cardinal medible''}) \Rightarrow \mathrm{Con}(\mathrm{ZFC} + \text{``existen infinitos cardinales de Mahlo''})$.
\end{corolario}

% ===========================================================
\section{La jerarquía de cardinales grandes}
% ===========================================================

Los cardinales grandes forman una cadena (lineal, en lo que concierne a la jerarquía de consistencia) de hipótesis cada vez más fuertes que extienden ZFC. Presentamos aquí una selección de los más importantes, en orden creciente de fuerza.

\subsection{Cardinales débilmente compactos}

\begin{definicion}{Cardinal débilmente compacto}{weakly_compact}
Un cardinal inaccesible $\kappa$ es \textbf{débilmente compacto} si satisface la \textbf{propiedad del árbol}: todo árbol de altura $\kappa$ en el que cada nivel tiene menos de $\kappa$ nodos tiene una rama de longitud $\kappa$.
\end{definicion}

\begin{observacion}{Equivalencias de débilmente compacto}{wc_equiv}
Las siguientes afirmaciones son equivalentes para un cardinal inaccesible $\kappa$:
\begin{enumerate}[leftmargin=2em, label=(\roman*)]
  \item $\kappa$ es débilmente compacto.
  \item $\kappa \to (\kappa)^2_2$ (propiedad de Ramsey para pares con 2 colores).
  \item $\kappa$ es $\Pi^1_1$-indescriptible: para toda fórmula $\Pi^1_1$ $\varphi$ y todo $R \subseteq V_\kappa$ con $\langle V_\kappa, \in, R \rangle \models \varphi$, existe $\alpha < \kappa$ con $\langle V_\alpha, \in, R \cap V_\alpha \rangle \models \varphi$.
  \item Existe un ultrafiltro $M$-completo no principal sobre $\kappa$ para algún modelo transitivo $M$ con $M^\kappa \subseteq M$.
\end{enumerate}
\end{observacion}

\subsection{Cardinales de Ramsey}

\begin{definicion}{Cardinal de Ramsey}{ramsey}
Un cardinal $\kappa$ es de \textbf{Ramsey} si para toda función de coloración $f: [\kappa]^{<\omega} \to 2$ (es decir, para todo subconjunto finito de $\kappa$ le asignamos uno de dos colores), existe un conjunto homogéneo $H \subseteq \kappa$ con $|H| = \kappa$: $f$ es constante sobre $[H]^n$ para todo $n < \omega$.

Formalmente: $\kappa \to (\kappa)^{<\omega}_2$.
\end{definicion}

\subsection{Cardinales supercompactos}

Los cardinales supercompactos, introducidos por Reinhardt y estudiados a fondo por Solovay y Magidor, son los más utilizados en las aplicaciones modernas de la teoría de conjuntos.

\begin{definicion}{Cardinal supercompacto}{supercompact}
Un cardinal $\kappa$ es \textbf{$\lambda$-supercompacto} (para $\lambda \geq \kappa$) si existe una inmersión elemental $j: V \to M$ con punto crítico $\kappa$, $j(\kappa) > \lambda$, y $M^\lambda \subseteq M$ (el modelo $M$ es cerrado bajo secuencias de longitud $\lambda$).

$\kappa$ es \textbf{supercompacto} si es $\lambda$-supercompacto para todo $\lambda \geq \kappa$.
\end{definicion}

\begin{importante}
La condición $M^\lambda \subseteq M$ es crucial: garantiza que el modelo $M$ ``ve'' todas las sucesiones de longitud $\lambda$ de sus propios elementos. Esto hace que los cardinales supercompactos sean mucho más fuertes que los cardinales medibles (donde solo se requiere $M^\omega \subseteq M$, es decir, que $M$ sea cerrado bajo sucesiones numerables).
\end{importante}

\subsection{Cardinales de Woodin}

Los cardinales de Woodin, descubiertos en la década de 1980, ocupan un lugar especial: son el umbral en el que la teoría de los conjuntos proyectivos (subconjuntos de $\mathbb{R}$ definibles) pasa a ser completamente determinada.

\begin{definicion}{Cardinal de Woodin}{woodin_card}
Un cardinal $\delta$ es un \textbf{cardinal de Woodin} si para toda función $f: \delta \to \delta$ existe un cardinal $\kappa < \delta$ tal que $\kappa$ es ``$f$-fuerte'': para todo $\lambda$ con $f``\kappa \subseteq \lambda$ existe una inmersión elemental $j: V \to M$ con $\mathrm{crit}(j) = \kappa$, $j(\kappa) \geq \lambda$, y $V_{j(f)(\kappa)} \subseteq M$.
\end{definicion}

\begin{teorema}{Woodin y determinismo proyectivo}{woodin_pd}
Si existen $n$ cardinales de Woodin con un cardinal medible por encima de todos ellos, entonces todos los enunciados $\Sigma^1_{n+1}$ sobre los reales son determinados (\textit{véase la sección sobre el Axioma de Determinismo}).
\end{teorema}

\subsection{La cadena de la jerarquía}

En la Tabla siguiente se resume la jerarquía de cardinales grandes, ordenados de menor a mayor fuerza de consistencia:

\begin{center}
\begin{tabular}{ll}
\hline
\textbf{Cardinal grande} & \textbf{Implicación} \\
\hline
Fuertemente inaccesible & $\Rightarrow$ $V_\kappa \models \mathrm{ZFC}$ \\
Mahlo & $\Rightarrow$ inaccesible \\
Débilmente compacto & $\Rightarrow$ Mahlo \\
Ramsey & $\Rightarrow$ débilmente compacto \\
Medible & $\Rightarrow$ Ramsey, $V \neq L$ \\
Fuerte & $\Rightarrow$ medible \\
Woodin & $\Rightarrow$ muchos Mahlo por debajo \\
Supercompacto & $\Rightarrow$ Woodin \\
Extendible & $\Rightarrow$ supercompacto \\
Casi enorme & $\Rightarrow$ extendible \\
Enorme & $\Rightarrow$ casi enorme \\
$n$-enorme & jerarquía de enormes \\
\hline
\end{tabular}
\end{center}

Cada flecha $\Rightarrow$ indica una implicación de consistencia estricta, no reversible.

% ===========================================================
\section{Jerarquía de consistencia entre extensiones de ZFC}
% ===========================================================

\subsection{El concepto de fuerza de consistencia}

\begin{definicion}{Fuerza de consistencia}{cons_strength}
Dado un sistema axiomático $T$ (que extiende ZFC), la \textbf{fuerza de consistencia} de $T$ es la ``posición'' de $T$ en la jerarquía de consistencia:
\[
  T_1 \leq_{\mathrm{Con}} T_2 \iff \mathrm{Con}(T_2) \Rightarrow \mathrm{Con}(T_1).
\]
Decimos que $T_1$ y $T_2$ tienen la \textbf{misma fuerza de consistencia} si $T_1 \equiv_{\mathrm{Con}} T_2$, es decir, $\mathrm{Con}(T_1) \iff \mathrm{Con}(T_2)$.
\end{definicion}

\begin{importante}
Uno de los resultados más notables de la teoría de conjuntos contemporánea es que prácticamente todos los enunciados independientes estudiados se pueden ubicar en esta jerarquía de consistencia. En particular, la jerarquía parece ser \textbf{lineal} (aunque esto no está demostrado en general): dados dos sistemas $T_1$ y $T_2$ razonables, casi siempre ocurre que $T_1 \leq_{\mathrm{Con}} T_2$ o $T_2 \leq_{\mathrm{Con}} T_1$.
\end{importante}

\subsection{La escala de consistencia}

Presentamos la jerarquía principal, de menor a mayor fuerza:

\begin{enumerate}[leftmargin=2em, label=\textbf{N\arabic*.}]
  \item \textbf{ZFC}: El sistema base. $\mathrm{Con}(\mathrm{ZFC})$ no puede ser demostrado en ZFC (Gödel, 1931).

  \item \textbf{ZFC + $\exists\kappa$ inaccesible}: Implica $\mathrm{Con}(\mathrm{ZFC})$. Equivalente a $\mathrm{ZFC}_2$ (ZFC de segundo orden, según un resultado de Zermelo).

  \item \textbf{ZFC + $\exists\kappa$ de Mahlo}: Estrictamente más fuerte que el anterior.

  \item \textbf{ZFC + $\exists\kappa$ débilmente compacto}: Implica la consistencia de ``infinitos Mahlo''.

  \item \textbf{ZFC + $\exists\kappa$ medible}: Un salto notable; implica $V \neq L$ y la consistencia de todos los anteriores.

  \item \textbf{ZFC + $\exists\kappa$ de Woodin}: Implica el determinismo de todos los conjuntos proyectivos.

  \item \textbf{ZFC + AD$^{L(\mathbb{R})}$}: Equivalente en consistencia a infinitos cardinales de Woodin.

  \item \textbf{ZFC + $\exists\kappa$ supercompacto}: Los cardinales supercompactos implican la consistencia de los cardinales de Woodin y mucho más.

  \item \textbf{ZFC + $\exists\kappa$ enorme}: Una de las hipótesis más fuertes conocidas que son ``seguras'' (aún no refutadas).

  \item \textbf{ZFC + $I_0$}: El cardinalidad $I_0$ (hay un $j: L(V_{\lambda+1}) \to L(V_{\lambda+1})$ elemental con punto crítico por debajo de $\lambda$) está cerca del límite de lo que la teoría de conjuntos puede estudiar sin caer en contradicción con AC.
\end{enumerate}

\subsection{Inconsistencia: el límite superior}

\begin{teorema}{Teorema de Kunen (1971)}{kunen_incons}
No existe ninguna inmersión elemental no trivial $j: V \to V$.
\end{teorema}

\begin{proof}[Esquema de la demostración]
Supongamos que $j: V \to V$ es elemental no trivial con punto crítico $\kappa$. Definimos la \emph{secuencia de Kunen}: $\kappa_0 = \kappa$, $\kappa_{n+1} = j(\kappa_n)$, y $\lambda = \sup_n \kappa_n$. Entonces $j(\lambda) = \lambda$ (porque $j$ envía $\kappa_n$ a $\kappa_{n+1}$ y el supremo se preserva).

Consideremos ahora el conjunto $A = \{\kappa_n : n < \omega\}$. La función $f: \lambda \to V$ definida por $f(\alpha) = j\restriction V_\alpha$ es un elemento de $V_{\lambda+2}$. Pero $j(f)(\lambda) = j \circ f \circ j^{-1}(\lambda) = j \circ f(\lambda)$... La derivación formal conduce a que el conjunto $\{A \subseteq \lambda : \kappa_0 \in j(A)\}$ es un ultrafiltro normal sobre $\lambda$ que es $\lambda^+$-completo, lo cual (junto con el teorema de Erdős-Hajnal) produce una contradicción con el Axioma de Elección.
\end{proof}

\begin{observacion}{Cardinales Reinhardt}{reinhardt}
Un \textbf{cardinal Reinhardt} sería el punto crítico de una inmersión elemental $j: V \to V$. El teorema de Kunen muestra que, bajo ZFC, no existen cardinales Reinhardt. Sin embargo, si se abandona el Axioma de Elección, Woodin ha propuesto la hipótesis de que cardinales Reinhardt son consistentes con ZF (sin AC). Este es un área de investigación activa.
\end{observacion}

% ===========================================================
\section{El Axioma de Determinismo}
% ===========================================================

\subsection{Juegos de longitud $\omega$ y estrategias}

El Axioma de Determinismo formaliza la intuición de que en cualquier juego de información perfecta de longitud infinita, uno de los dos jugadores debe tener una estrategia ganadora. Su incompatibilidad con el Axioma de Elección es uno de los resultados más sorprendentes de la teoría de conjuntos.

\begin{definicion}{Juego $G(A)$}{juego}
Sea $A \subseteq \omega^\omega$ (un subconjunto del espacio de Baire). El \textbf{juego de Gale--Stewart} $G(A)$ se desarrolla como sigue:
\begin{itemize}[leftmargin=2em]
  \item El \textbf{Jugador I} escoge $a_0 \in \omega$.
  \item El \textbf{Jugador II} escoge $b_0 \in \omega$.
  \item El \textbf{Jugador I} escoge $a_1 \in \omega$.
  \item El \textbf{Jugador II} escoge $b_1 \in \omega$.
  \item $\cdots$ (se continúa $\omega$ rondas).
\end{itemize}
Al cabo de $\omega$ rondas se produce una sucesión $x = (a_0, b_0, a_1, b_1, \ldots) \in \omega^\omega$. El \textbf{Jugador I gana} si $x \in A$; el \textbf{Jugador II gana} si $x \notin A$.
\end{definicion}

\begin{definicion}{Estrategia y determinismo}{estrategia}
\begin{enumerate}[leftmargin=2em, label=(\alph*)]
  \item Una \textbf{estrategia} para el Jugador I es una función $\sigma: \omega^{<\omega} \to \omega$ que asigna a cada secuencia finita de jugadas (la historia de la partida hasta ese momento) el siguiente movimiento del Jugador I. Una estrategia para el Jugador II se define análogamente.
  \item Una estrategia $\sigma$ es \textbf{ganadora} para el Jugador I si toda partida jugada siguiendo $\sigma$ produce una secuencia $x \in A$, independientemente de cómo juegue el Jugador II.
  \item El juego $G(A)$ es \textbf{determinado} si alguno de los dos jugadores tiene una estrategia ganadora.
\end{enumerate}
\end{definicion}

\begin{ejemplo}{Ejemplo concreto: juego con condición sobre paridades}{}
Sea $A = \{(a_0, b_0, a_1, b_1, \ldots) \in \omega^\omega : \text{la primera componente } a_0 \text{ es par}\}$.

\medskip
\textbf{Estrategia ganadora del Jugador I:} En el primer turno, el Jugador I elige $a_0 = 0$ (un número par). Después de esto, no importa qué elija el Jugador II (ni qué elija el propio Jugador I en los turnos posteriores), la secuencia resultante tendrá $a_0 = 0$, que es par, así que $x \in A$. El Jugador I gana.

\medskip
En este ejemplo, el juego es determinado y la estrategia ganadora del Jugador I es extremadamente simple.
\end{ejemplo}

\begin{ejemplo}{Juego con condición sobre límites}{}
Sea $A = \{(a_0, b_0, a_1, b_1, \ldots) \in \omega^\omega : \lim_{n \to \infty} a_n = \infty\}$ (la secuencia de movimientos del Jugador I diverge a infinito).

\medskip
\textbf{Estrategia ganadora del Jugador I:} En el turno $n$-ésimo, el Jugador I elige $a_n = n$. Entonces $a_n = n \to \infty$, luego $x \in A$ con certeza.

\medskip
Así, el Jugador I gana con una estrategia sencilla. Nótese que si la condición fuera $\limsup_{n \to \infty} a_n < \infty$, entonces la estrategia ganadora sería del Jugador II (eligiendo en cada turno $b_n$ de manera que ``fuerce'' a $a_n$ a ser acotado... pero como el Jugador I controla sus propias elecciones, en realidad el Jugador I puede siempre hacer $a_n = n$). El juego $G(\{x : \limsup a_n < \infty\})$ es determinado por el Jugador II, que no tiene control sobre $a_n$; de hecho, el Jugador I tiene la estrategia ganadora trivialmente al elegir $a_n = n$. El análisis depende de qué jugador controla qué componentes.
\end{ejemplo}

\subsection{El Axioma de Determinismo}

\begin{definicion}{Axioma de Determinismo (AD)}{AD}
El \textbf{Axioma de Determinismo} es el enunciado:
\[
  \mathrm{AD}: \quad \forall A \subseteq \omega^\omega, \;\; G(A) \text{ es determinado.}
\]
\end{definicion}

\begin{teorema}{Incompatibilidad de AD con AC}{AD_AC}
$\mathrm{AD}$ es incompatible con el Axioma de Elección $\mathrm{AC}$ en presencia de $\mathrm{ZF}$:
\[
  \mathrm{ZF} \vdash \mathrm{AD} \Rightarrow \neg\mathrm{AC}.
\]
\end{teorema}

\begin{proof}
Suponemos $\mathrm{ZF} + \mathrm{AC}$ y construimos un conjunto $A \subseteq \omega^\omega$ sin estrategia ganadora para ninguno de los dos jugadores (un juego no determinado).

Bajo AC, el conjunto $\omega^\omega$ tiene cardinalidad $2^{\aleph_0}$, y el conjunto de todas las estrategias para el Jugador I (resp. Jugador II) también tiene cardinalidad $2^{\aleph_0}$. Enumeremos todas las estrategias para el Jugador I como $\{\sigma_\alpha : \alpha < 2^{\aleph_0}\}$ y todas las estrategias para el Jugador II como $\{\tau_\alpha : \alpha < 2^{\aleph_0}\}$.

Construimos $A$ por inducción transfinita de longitud $2^{\aleph_0}$. En el paso $\alpha$, escogemos (usando AC) dos partidas:
\begin{itemize}[leftmargin=2em]
  \item $x_\alpha$: la partida que resulta de que el Jugador I siga $\sigma_\alpha$ y el Jugador II siga alguna estrategia fija (por ejemplo $\tau_0$). Ponemos $x_\alpha \notin A$ (para que $\sigma_\alpha$ no sea ganadora para I).
  \item $y_\alpha$: la partida que resulta de que el Jugador II siga $\tau_\alpha$ y el Jugador I siga $\sigma_0$. Ponemos $y_\alpha \in A$ (para que $\tau_\alpha$ no sea ganadora para II).
\end{itemize}
El conjunto $A$ así construido no tiene estrategia ganadora para ningún jugador: $\sigma_\alpha$ no es ganadora para I (porque $x_\alpha \notin A$) y $\tau_\alpha$ no es ganadora para II (porque $y_\alpha \in A$). Por lo tanto, $G(A)$ no es determinado, refutando AD.
\end{proof}

\begin{importante}
La construcción del contraejemplo usa de manera esencial la enumeración transfinita de $2^{\aleph_0}$ estrategias, que requiere el Axioma de Elección. Si renunciamos a AC, la construcción falla.
\end{importante}

\subsection{Determinismo de Borel: lo que ZFC sí puede demostrar}

\begin{teorema}{Determinismo de Borel (Martin, 1975)}{borel_det}
Dentro de $\mathrm{ZFC}$, todo juego $G(A)$ con $A$ un conjunto de Borel en $\omega^\omega$ es determinado. Formalmente:
\[
  \mathrm{ZFC} \vdash \forall A \in \mathcal{B}(\omega^\omega),\; G(A) \text{ es determinado.}
\]
\end{teorema}

\begin{proof}[Comentario sobre la demostración]
La demostración de Martin es una inducción transfinita sobre el rango de Borel de $A$. Para conjuntos abiertos, el resultado es elemental (conocido como el teorema de Gale--Stewart, 1953): el Jugador I tiene una estrategia ganadora si y solo si puede ``forzar'' que la partida entre en $A$, y en caso contrario el Jugador II tiene una estrategia ganadora cerrando afuera. Para conjuntos de rango de Borel $\alpha$, se reduce a juegos de rango $\beta < \alpha$ mediante una construcción de ``desenrollamiento'' (\textit{unraveling}). La demostración usa $\mathrm{ZFC}$ completo, incluyendo el axioma de reemplazamiento en rango alto.
\end{proof}

\begin{observacion}{Fuerza del Determinismo de Borel}{bd_strength}
El Determinismo de Borel requiere utilizar el axioma de reemplazamiento para sets de rango arbitrariamente alto; en particular, no es demostrable en $\mathrm{ZFC}$ sin reemplazamiento (Friedman, 1971). Este es uno de los ejemplos más concretos de la necesidad del esquema de Reemplazamiento en matemáticas ``convencionales''.
\end{observacion}

\subsection{AD bajo $L(\mathbb{R})$ y cardinales de Woodin}

Si bien AD es incompatible con ZFC (pues ZFC incluye AC), puede ser consistente con ZF. Más aún, puede ser verdadero en el modelo $L(\mathbb{R})$ (el universo constructible relativo a los reales).

\begin{teorema}{Woodin--Martin--Steel}{woodin_martin_steel}
Las siguientes afirmaciones son equiconsistentes:
\begin{enumerate}[leftmargin=2em, label=(\roman*)]
  \item Existen infinitos cardinales de Woodin.
  \item $L(\mathbb{R}) \models \mathrm{AD}$.
  \item Todos los juegos proyectivos son determinados ($\mathrm{PD}$: \textit{Projective Determinacy}).
\end{enumerate}
\end{teorema}

\begin{proof}[Comentario sobre la demostración]
La implicación (i) $\Rightarrow$ (ii) es el teorema de Woodin: si existen infinitos cardinales de Woodin, entonces en $L(\mathbb{R})$ vale el Axioma de Determinismo. La equivalencia con (iii) es el teorema de Martin--Steel: el Determinismo Proyectivo es equivalente en consistencia a la existencia de infinitos cardinales de Woodin. Estos resultados son de los más profundos de la teoría de conjuntos de las últimas décadas.
\end{proof}

\subsection{Consecuencias de AD}

Bajo $\mathrm{ZF} + \mathrm{AD}$, la teoría de los subconjuntos de $\mathbb{R}$ cambia radicalmente:

\begin{teorema}{Consecuencias de AD}{AD_conseq}
Bajo $\mathrm{ZF} + \mathrm{AD}$:
\begin{enumerate}[leftmargin=2em, label=(\roman*)]
  \item \textbf{(Mycielski--Steinhaus)} Todos los subconjuntos de $\mathbb{R}$ son Lebesgue medibles.
  \item Todos los subconjuntos de $\mathbb{R}$ tienen la propiedad de Baire.
  \item Todos los subconjuntos de $\mathbb{R}$ tienen la propiedad del conjunto perfecto (cualquier conjunto no numerable contiene un conjunto cerrado perfecto, luego tiene cardinalidad $\mathfrak{c}$).
  \item \textbf{(Solovay)} $\omega_1$ es un cardinal medible.
  \item \textbf{(Steel)} $L(\mathbb{R}) \models \mathrm{AD}$ implica que los cardinales de $L(\mathbb{R})$ hasta $\Theta$ (el mayor cardinal de $L(\mathbb{R})$) son todos Woodin en un cierto sentido generalizado.
  \item La Hipótesis del Continuo falla de la manera más extrema posible: $|\mathbb{R}|$ es incomparable con $\aleph_1$ en el sentido de que no existe ninguna inyección de $\mathbb{R}$ en $\omega_1$ ni viceversa (ya que AC falla).
\end{enumerate}
\end{teorema}

\begin{ejemplo}{Comparación entre $\mathrm{ZFC}$ y $\mathrm{ZF} + \mathrm{AD}$}{}
\begin{center}
\begin{tabular}{lll}
\hline
\textbf{Enunciado} & \textbf{Bajo ZFC} & \textbf{Bajo ZF + AD} \\
\hline
Todo subconjunto de $\mathbb{R}$ es medible & Falso (conjunto de Vitali) & Verdadero \\
$\omega_1$ es medible & Falso & Verdadero \\
Existen conjuntos de Vitali & Verdadero & Falso \\
Existe un buen orden de $\mathbb{R}$ & Verdadero (bajo AC) & Falso \\
Determinismo de Borel & Verdadero & Verdadero \\
Determinismo proyectivo & Independiente & Verdadero \\
Hipótesis del continuo & Independiente & No tiene sentido directo \\
\hline
\end{tabular}
\end{center}
La tabla ilustra cómo $\mathrm{AD}$ y $\mathrm{AC}$ dan lugar a universos matemáticos radicalmente distintos.
\end{ejemplo}

% ===========================================================
\section{Límites epistemológicos de ZFC y preguntas abiertas}
% ===========================================================

\subsection{El segundo teorema de incompletitud y sus consecuencias}

El segundo teorema de incompletitud de Gödel (1931) establece que ningún sistema consistente suficientemente expresivo puede demostrar su propia consistencia. En el contexto de ZFC:

\begin{teorema}{Segundo Teorema de Incompletitud (Gödel, 1931)}{godel2}
Si $\mathrm{ZFC}$ es consistente, entonces $\mathrm{ZFC} \not\vdash \mathrm{Con}(\mathrm{ZFC})$.
\end{teorema}

Esto significa que la jerarquía de cardinales grandes es \emph{inevitable}: para demostrar $\mathrm{Con}(\mathrm{ZFC})$, necesitamos un sistema más fuerte (como $\mathrm{ZFC} +$ ``existe un cardinal inaccesible''). Y para demostrar la consistencia de ese sistema, necesitamos uno aún más fuerte. No hay un ``techo'' alcanzable desde dentro.

\subsection{El problema del valor de $2^{\aleph_0}$}

El teorema de Easton (1970) establece que la función $\kappa \mapsto 2^\kappa$ en los cardinales regulares puede tomar casi cualquier valor compatible con el Teorema de König (que exige $\mathrm{cf}(2^\kappa) > \kappa$). Sin más axiomas, ZFC no determina el valor de $2^{\aleph_0}$.

Las propuestas actuales más estudiadas son:
\begin{enumerate}[leftmargin=2em, label=(\alph*)]
  \item \textbf{$2^{\aleph_0} = \aleph_1$ (HC):} La propuesta de Gödel y, más recientemente, del programa de Woodin ($\Omega$-lógica y el axioma $(\ast)$ llevan a $2^{\aleph_0} = \aleph_2$; en versiones más recientes, Woodin ha explorado $2^{\aleph_0} = \aleph_2$).
  \item \textbf{$2^{\aleph_0} = \aleph_2$:} Sugerida por el Axioma de Martin + $\neg\mathrm{HC}$, y por el axioma $(\ast)$ de Woodin.
  \item \textbf{Valores más grandes:} Ningún argumento filosófico o matemático convincente resuelve la cuestión. La elección depende de qué axiomas se prefieran añadir a ZFC.
\end{enumerate}

\subsection{La Hipótesis de los Cardinales Singulares}

\begin{definicion}{Hipótesis de los Cardinales Singulares (SCH)}{sch}
La \textbf{Hipótesis de los Cardinales Singulares} afirma: para todo cardinal singular $\kappa$ con $2^{\mathrm{cf}(\kappa)} < \kappa$,
\[
  \kappa^{\mathrm{cf}(\kappa)} = \kappa^+.
\]
\end{definicion}

\begin{teorema}{SCH y cardinales grandes}{sch_large}
\begin{enumerate}[leftmargin=2em, label=(\roman*)]
  \item \textbf{(Silver, 1975)} Si SCH falla para algún cardinal singular de cofinalidad no numerable, entonces hay muchos cardinales singulares donde falla.
  \item \textbf{(Shelah)} SCH puede fallar en $\aleph_\omega$ si existen cardinales medibles (Magidor construyó un modelo donde $2^{\aleph_\omega} = \aleph_{\omega+2}$).
  \item \textbf{(Solovay)} Si existe un cardinal supercompacto, entonces SCH holds por encima de dicho cardinal.
\end{enumerate}
\end{teorema}

\subsection{El programa de modelos interiores}

El \textbf{programa de modelos interiores} busca construir para cada cardinal grande $\kappa$ un modelo interno $L[\mu]$ (generalización del universo constructible $L$) que contenga $\kappa$ y sea tan ``delgado'' como sea posible, permitiendo analizar la aritmética cardinal en detalle.

\begin{itemize}[leftmargin=2em]
  \item Para cardinales medibles: $L[\mu]$ (Dodd-Jensen, 1981).
  \item Para cardinales fuertes y de Woodin: los modelos $\mathcal{M}_n$ de Steel.
  \item Para cardinales supercompactos: el programa de modelos interiores canónicos está incompleto; es uno de los mayores problemas abiertos.
\end{itemize}

\subsection{Preguntas abiertas en teoría de conjuntos contemporánea}

Concluimos con un panorama de los problemas más importantes de la teoría de conjuntos actual.

\begin{enumerate}[leftmargin=2em, label=\textbf{PA\arabic*.}]

  \item \textbf{(El programa de modelos interiores para supercompactos)} ¿Puede construirse un modelo interior canónico para un cardinal supercompacto? Esto requeriría, en particular, resolver la pregunta sobre la $\mathbf{\Sigma}^2_1$-absoluteness (absolutez de segundo nivel). Es posiblemente el problema técnico más difícil de la teoría de conjuntos actual.

  \item \textbf{(Determinismo y el universo $V$)} ¿Bajo qué hipótesis de cardinales grandes es $\mathrm{AD}^{V}$ consistente? (AD ``en el universo completo'', no solo en $L(\mathbb{R})$.) ¿Son los cardinales de Reinhardt (sin AC) consistentes con ZF?

  \item \textbf{(La Hipótesis del Continuo y la $\Omega$-lógica)} El programa de Woodin busca una noción de ``verdad absoluta'' para enunciados sobre $H(\omega_2)$ (los conjuntos de rango hereditario $\leq \aleph_1$) usando la $\Omega$-lógica. ¿Es el axioma $(\ast)$ de Woodin (que implica $2^{\aleph_0} = \aleph_2$) el ``axioma correcto'' para $H(\omega_2)$?

  \item \textbf{(La jerarquía de los grandes ordinales)} ¿Cuál es el cardinal Woodin más grande posible antes de llegar a una contradicción? ¿Cuáles son los límites superior e inferior de los cardinales ``seguros''?

  \item \textbf{(SCH en $\aleph_\omega$)} ¿Es consistente, relativo a ZFC solo (sin cardinales grandes), que $2^{\aleph_\omega} > \aleph_{\omega+1}$? El resultado de Shelah muestra que sí bajo cardinales medibles; ¿es necesaria esa hipótesis?

  \item \textbf{(El axioma de Forcing máximo)} El \textbf{Axioma de Forcing Máximo} ($\mathrm{MFP}$) o el \textbf{Axioma de Martin Máximo} ($\mathrm{MM}$) son principios potentes que implican $2^{\aleph_0} = \aleph_2$ y resuelven muchos problemas. ¿Cuál es la fuerza de consistencia exacta de MM? ¿Es consistente con un cardinal supercompacto?

  \item \textbf{(Problema $\mathfrak{p} = \mathfrak{t}$)} Este problema de la teoría de cardinales combinatorios fue resuelto por Malliaris y Shelah en 2017: $\mathfrak{p} = \mathfrak{t}$ en ZFC. Esto ilustra que incluso problemas de larga data sobre el continuo pueden resolverse sin cardinales grandes. ¿Hay otros problemas similares pendientes?

  \item \textbf{(Teoría de conjuntos finitos)} ¿Puede la teoría de conjuntos aportar métodos nuevos a la teoría de Ramsey finita (hipergrafos, combinatoria extremal) mediante herramientas ultrafilter? Esta dirección, impulsada por Baumgartner y Shelah, tiene resultados concretos pero deja muchas preguntas abiertas.

\end{enumerate}

\subsection{Reflexión epistemológica final}

Los resultados de esta clase —y del curso en su conjunto— permiten una reflexión sobre la naturaleza misma de las matemáticas:

\begin{enumerate}[leftmargin=2em, label=\textbf{R\arabic*.}]
  \item \textbf{Los fundamentos no son un ``suelo firme'':} ZFC es un sistema poderoso y preciso, pero no puede demostrar su propia consistencia, ni resolver preguntas como la Hipótesis del Continuo. Las matemáticas no tienen un fundamento absoluto e indubitable: cualquier sistema suficientemente fuerte es, en cierto sentido, incompleto.

  \item \textbf{Los cardinales grandes como extensión natural:} La jerarquía de cardinales grandes no es un capricho: cada nivel responde a preguntas concretas (el valor de $2^{\aleph_0}$, el determinismo de conjuntos proyectivos, etc.) y produce consecuencias verificables en matemáticas ``ordinarias'' (análisis, topología, combinatoria).

  \item \textbf{La multivers de la teoría de conjuntos:} Joel David Hamkins ha propuesto la imagen del \emph{multiverso conjuntista}: no existe un único universo de conjuntos $V$, sino una pluralidad de universos, cada uno con diferentes propiedades. Bajo esta visión, preguntar si HC es ``verdadera'' carece de sentido absoluto; lo que importa es qué universo (qué axiomas) estamos usando.

  \item \textbf{El platonismo de Gödel:} Por el contrario, Gödel mantuvo una visión platónica: existe un único universo verdadero de conjuntos, y nuestra tarea es descubrir los axiomas correctos que lo describen. Para Gödel, los cardinales grandes son herramientas para acercarse a esa verdad matemática objetiva.

  \item \textbf{La teoría de conjuntos como ciencia empírica de los axiomas:} Una postura intermedia, defendida por muchos logistas actuales, es tratar la teoría de conjuntos como una disciplina que estudia las consecuencias de distintos sistemas axiomáticos, sin comprometerse con una ``verdad absoluta''. Bajo esta luz, preguntas como ``¿vale HC?'' son cuestiones sobre qué universo de conjuntos es más útil o más natural para la matemática.
\end{enumerate}

\begin{importante}
El recorrido desde los números transfinitos de Cantor hasta los horizontes actuales de la teoría de conjuntos muestra que las matemáticas son una disciplina viva, en constante expansión. Cada respuesta genera nuevas preguntas, y los límites del formalismo son, a la vez, invitaciones a ir más allá.
\end{importante}

% ===========================================================
\section*{Problemas y tareas}
% ===========================================================

\subsection*{Problemas básicos}

\begin{enumerate}[leftmargin=2em, label=\textbf{P\arabic*.}]

  \item \textbf{(Cardinales regulares y singulares.)}
  \begin{enumerate}[label=(\alph*)]
    \item Demuestre que todo sucesor cardinal $\aleph_{\alpha+1}$ es regular.
    \item Demuestre que $\aleph_\omega$ es singular de cofinalidad $\omega$.
    \item ¿Es posible que exista un cardinal límite regular de cofinalidad $\omega_1$? Justifique.
    \item Demuestre que $\mathrm{cf}(\aleph_{\omega_1}) = \omega_1$.
  \end{enumerate}

  \item \textbf{(Cardinales inaccesibles y la jerarquía $V_\alpha$.)}
  \begin{enumerate}[label=(\alph*)]
    \item Demuestre que si $\kappa$ es fuertemente inaccesible, entonces $|V_\kappa| = \kappa$.
    \item Demuestre que $V_\kappa$ satisface el axioma de Partes: si $a \in V_\kappa$, entonces $\mathcal{P}(a) \in V_\kappa$.
    \item Use el Teorema~\ref{teo:vkappa_zfc} para concluir que si $\kappa$ es inaccesible y $T \subseteq V_\kappa$ es un enunciado de ZFC, entonces $V_\kappa \models T$.
    \item Explique por qué la existencia de un cardinal inaccesible no puede demostrarse en ZFC (argumento del segundo teorema de Gödel).
  \end{enumerate}

  \item \textbf{(Clubs y estacionarios.)}
  \begin{enumerate}[label=(\alph*)]
    \item Demuestre que la intersección de dos clubs en $\kappa$ (cardinal regular no numerable) es un club.
    \item Demuestre que la intersección de $\lambda < \kappa$ clubs en $\kappa$ es un club.
    \item Demuestre que todo club en $\kappa$ es estacionario, pero no todo conjunto estacionario es un club.
    \item \textit{(Presión del club.)} Demuestre que si $S \subseteq \kappa$ es estacionario y $f: S \to \kappa$ es una función de regresión (es decir, $f(\alpha) < \alpha$ para todo $\alpha \in S$), entonces existe $\gamma < \kappa$ tal que $f^{-1}(\gamma)$ es estacionario (Lema de Fodor o Presión del Club).
  \end{enumerate}

  \item \textbf{(Cardinales de Mahlo.)}
  \begin{enumerate}[label=(\alph*)]
    \item Demuestre que si $\kappa$ es de Mahlo, entonces el conjunto $\{\lambda < \kappa : \lambda \text{ es débilmente inaccesible}\}$ también es estacionario en $\kappa$.
    \item Use el Lema de Fodor para mostrar que si $\kappa$ es de Mahlo, entonces el conjunto de puntos fijos de la función aleph debajo de $\kappa$ es estacionario.
    \item Enuncie (sin demostrar) en qué sentido los cardinales de Mahlo ``ven'' a los cardinales inaccesibles como ``pequeños''.
  \end{enumerate}

  \item \textbf{(Ultrafiltros y completitud.)}
  \begin{enumerate}[label=(\alph*)]
    \item Demuestre que para $\kappa = \omega$, cualquier ultrafiltro no principal sobre $\omega$ es $\sigma$-completo si y solo si es $\omega_1$-completo.
    \item Demuestre que si $\mathcal{U}$ es un ultrafiltro $\kappa$-completo no principal sobre $\kappa$, entonces para toda partición $\{A_i : i < \lambda\}$ de $\kappa$ con $\lambda < \kappa$, exactamente uno de los $A_i$ pertenece a $\mathcal{U}$.
    \item Utilice (b) para demostrar que si $\kappa$ es medible, entonces $\kappa$ es cardinal límite (no sucesor).
  \end{enumerate}

  \item \textbf{(Juegos de Gale--Stewart.)}
  \begin{enumerate}[label=(\alph*)]
    \item Determine quién tiene estrategia ganadora en el juego $G(A)$ donde $A = \{x \in \omega^\omega : x(0) > x(1)\}$.
    \item Determine quién tiene estrategia ganadora en el juego $G(A)$ donde $A = \{x \in \omega^\omega : \exists n, x(n) = 0\}$.
    \item Demuestre que si $A$ es abierto en $\omega^\omega$ (topología producto), entonces $G(A)$ es determinado (caso base del Determinismo de Borel).
    \item \textit{(Juego de Banach--Mazur.)} Defina el juego $BM(A)$ para $A \subseteq [0,1]$ (los jugadores alternan eligiendo intervalos anidados de longitud decreciente, y el Jugador I gana si el punto límite está en $A$). ¿Qué relación tiene con el Axioma de Determinismo?
  \end{enumerate}

  \item \textbf{(Jerarquía de consistencia.)}
  \begin{enumerate}[label=(\alph*)]
    \item Ordene los siguientes sistemas por fuerza de consistencia, de menor a mayor: ZFC, ZFC + Mahlo, ZFC + $V = L$, ZFC + medible, ZFC + supercompacto, ZFC + Determinismo de Borel.
    \item ¿En qué posición de la jerarquía se sitúa ZFC + HC?
    \item ¿Es ZFC + $\neg$HC más fuerte en consistencia que ZFC? Justifique.
    \item Explique por qué ``ZFC + $V = L$'' tiene la misma fuerza de consistencia que ZFC.
  \end{enumerate}

\end{enumerate}

\subsection*{Problemas de profundización}

\begin{enumerate}[leftmargin=2em, label=\textbf{P\arabic*$'$.}]

  \item \textbf{(El teorema de Scott: detalles.)} Desarrolle los detalles de la demostración del Teorema de Scott (\ref{teo:scott_thm}):
  \begin{enumerate}[label=(\alph*)]
    \item Demuestre que si $j: V \to M$ es elemental con $V = L$ y $M = L$, entonces $j\restriction \mathrm{Ord}$ es una función estrictamente creciente de los ordinales en sí mismos, y que $j(\alpha) \geq \alpha$ para todo $\alpha$.
    \item Use el hecho de que $L$ tiene un parámetro de Skolem uniforme $h_L$ (una función definible $h_L: \mathrm{Ord} \times \omega \to L$ con $L = \{h_L(\alpha, n) : \alpha \in \mathrm{Ord}, n \in \omega\}$) para demostrar que cualquier inmersión elemental $j: L \to L$ debe ser la identidad.
    \item Concluya que si $\kappa$ es medible, entonces $V \neq L$.
  \end{enumerate}

  \item \textbf{(Determinismo proyectivo y cardinales de Woodin.)} El \textbf{Determinismo Proyectivo} (PD) afirma que todo juego $G(A)$ con $A$ un conjunto proyectivo (en la jerarquía proyectiva $\Sigma^1_n$, $\Pi^1_n$) es determinado. Investigue y exponga:
  \begin{enumerate}[label=(\alph*)]
    \item La demostración de Martin de que $\Sigma^1_1$ (analítico) implica determinismo bajo ZFC (sin cardinales grandes).
    \item La demostración de Martin de que la existencia de un cardinal medible implica el determinismo de todos los juegos $\Sigma^1_2$.
    \item El enunciado del teorema de Martin--Steel: la existencia de $n$ cardinales de Woodin con un cardinal medible por encima implica el determinismo de todos los juegos $\Sigma^1_{n+1}$.
    \item Por qué la existencia de infinitos cardinales de Woodin implica PD.
  \end{enumerate}

  \item \textbf{(El multiverso conjuntista.)} Lea el artículo seminal de Joel David Hamkins ``The set-theoretic multiverse'' (\textit{Rev. Symb. Log.}, 2012) y elabore un ensayo (2--3 páginas) que:
  \begin{enumerate}[label=(\alph*)]
    \item Explique la tesis del multiverso: los distintos modelos de ZFC son todos igualmente reales.
    \item Contraste esta postura con el platonismo conjuntista de Gödel.
    \item Discuta cómo el multiverso afecta la interpretación de resultados de independencia como HC.
    \item Presente al menos un argumento en favor y uno en contra de la tesis del multiverso.
  \end{enumerate}

  \item \textbf{(El axioma de Martin.)} El \textbf{Axioma de Martin} ($\mathrm{MA}$) afirma: para todo ccc (condición de cadena numerable) orden parcial $\mathbb{P}$ y toda familia $\{D_i : i < \kappa\}$ con $\kappa < \mathfrak{c}$ de abiertos densos en $\mathbb{P}$, existe un filtro que intersecta todos los $D_i$.
  \begin{enumerate}[label=(\alph*)]
    \item Demuestre que MA es un teorema cuando $\kappa < \aleph_0$ (Baire, el espacio de Baire).
    \item Demuestre que MA $+$ HC es consistente con ZFC.
    \item Demuestre que MA $+$ $\neg$HC es consistente con ZFC.
    \item Use MA para demostrar que bajo MA $+$ $\neg$HC, toda unión de $\omega_1$ muchos conjuntos nulos (Lebesgue) tiene medida cero.
  \end{enumerate}

  \item \textbf{(Independencia de SCH.)} La falla de la Hipótesis de los Cardinales Singulares en $\aleph_\omega$ es un resultado profundo de Magidor (1977). Investigue:
  \begin{enumerate}[label=(\alph*)]
    \item ¿Qué cardinalidad grande se usa en la construcción de Magidor?
    \item Enuncie el resultado de Magidor: es consistente (relativo a un cardinal supercompacto) que $2^{\aleph_\omega} = \aleph_{\omega+2}$.
    \item ¿Por qué el Teorema de König no prohíbe esta situación?
    \item ¿Qué dice el teorema de Shelah (\textit{pcf theory}) sobre el supremo de las potencias $2^{\aleph_n}$ para $n < \omega$?
  \end{enumerate}

\end{enumerate}

% ===========================================================
\section*{Referencias de la Clase 10}
% ===========================================================

Los resultados y conceptos de esta clase provienen de las siguientes fuentes fundamentales:

\begin{itemize}[leftmargin=2em]
  \item La referencia enciclopédica para la teoría de cardinales grandes es \cite{kanamori2003}. Cubre cardinales inaccesibles, Mahlo, débilmente compactos, Ramsey, medibles, fuertes, supercompactos, enormes y los límites de la jerarquía. Es la fuente primaria para todos los resultados de las Secciones 1--4 de esta clase.

  \item Para los cardinales inaccesibles y la verificación de $V_\kappa \models \mathrm{ZFC}$, véanse los capítulos introductorios de \cite{jech2003} (capítulo~12) y \cite{kunen2011} (capítulo~V). La definición y propiedades de los cardinales de Mahlo se tratan en \cite{jech2003} (§12.4) y \cite{kanamori2003} (§1.4).

  \item El teorema de Scott (si $\kappa$ es medible entonces $V \neq L$) fue publicado en \cite{scott1961}. Una demostración completa se encuentra en \cite{jech2003} (§17.1) y en \cite{kanamori2003} (§5.1).

  \item El Determinismo de Borel fue demostrado por Martin en \cite{martin1975}. La necesidad del Reemplazamiento para el Determinismo de Borel se debe a Friedman (1971); véase el análisis en \cite{kanamori2003} (§27.3).

  \item El teorema de Martin--Steel (la existencia de $n$ cardinales de Woodin con un cardinal medible por encima implica el determinismo $\Sigma^1_{n+2}$) se publicó en \cite{martinSteel1988}. La equivalencia entre infinitos cardinales de Woodin y $\mathrm{AD}^{L(\mathbb{R})}$ es el gran teorema de Woodin, expuesto en \cite{woodin2010det}.

  \item Para el Axioma de Determinismo y sus consecuencias (medibilidad de todos los subconjuntos de $\mathbb{R}$, propiedad de Baire, propiedad del conjunto perfecto), la referencia principal es \cite{moschovakis2009}, especialmente los capítulos 6 y 7. El libro de \cite{kanamori2003} (capítulo~28) también presenta estos resultados en el contexto de la jerarquía de cardinales grandes.

  \item Para la jerarquía de consistencia y el Axioma de Martin, véase \cite{kunen2011} (capítulo~III y~IV) y \cite{jech2003} (capítulo~16). El Axioma de Martin Máximo fue introducido por Foreman, Magidor y Shelah; véase \cite{foreman1988}.

  \item La tesis del multiverso conjuntista se expone en \cite{hamkins2012}. Una perspectiva alternativa (absolutismo de $\Omega$-lógica) se encuentra en \cite{woodin2001} y en los artículos ulteriores de Woodin sobre el Programa Final.

  \item Para el programa de modelos interiores, la referencia introductoria es \cite{steel2010}. Las aplicaciones a la Hipótesis de los Cardinales Singulares se exponen en \cite{shelah1994}.

  \item El teorema de Kunen (no existen inmersiones elementales $j: V \to V$ bajo ZFC) se publicó en \cite{kunen1971}. Una exposición moderna se encuentra en \cite{kanamori2003} (§23.2).
\end{itemize}


% ---------- FINAL ----------
\backmatter
\printbibliography[heading=bibintoc, title={Referencias}]

\end{document}
